% arXiv submission metadata
%
% Title: The 1/3--2/3 Conjecture for Width-3 Posets
%
% MSC 2020 Classification:
%   Primary:   06A07  (Combinatorics of partially ordered sets)
%   Secondary: 05C48  (Expansion properties of graphs)
%              05A20  (Combinatorial inequalities)
%              60J10  (Markov chains; general state spaces)
%              60C05  (Combinatorial probability)
%
% Keywords: 1/3--2/3 conjecture; partially ordered set; linear extension;
%   balanced pair; width-3 poset; Bubley--Karzanov graph; Cheeger inequality;
%   expansion; conductance.
%
\documentclass[11pt]{article}
\usepackage[utf8]{inputenc}
\usepackage[T1]{fontenc}
\usepackage{amsmath,amsthm,amssymb,amsfonts}
\usepackage{enumitem}
\usepackage{geometry}
\usepackage{hyperref}
\geometry{margin=1in}

\newtheorem{theorem}{Theorem}[section]
\newtheorem{lemma}[theorem]{Lemma}
\newtheorem{proposition}[theorem]{Proposition}
\newtheorem{corollary}[theorem]{Corollary}
\newtheorem{claim}[theorem]{Claim}
\newtheorem{conjecture}[theorem]{Conjecture}

\theoremstyle{definition}
\newtheorem{definition}[theorem]{Definition}
\newtheorem{remark}[theorem]{Remark}
\newtheorem{example}[theorem]{Example}
\newtheorem{hypothesis}[theorem]{Hypothesis}

\theoremstyle{remark}
\newtheorem*{gap}{\textbf{Gap}}

% Union of macros defined across step files.
%
% Notation note on the good-fiber object: the same object
% $\mathcal{F}_{x,y}$ (or $\mathcal{F}_\alpha$ for $\alpha=(x,y)$) is
% spelled four ways across the step files, all of which expand to the
% same glyph:
%   \F       -- bare $\mathcal{F}$ (used without subscript)
%   \Fxy     -- $\mathcal{F}_{x,y}$ (fixed pair; Step 4, Step 5)
%   \fib_{x,y} -- $\mathcal{F}_{x,y}$ via \fib = \mathcal{F} (Step 2)
%   \fiber{\alpha} -- $\mathcal{F}_\alpha$ (generic index; Steps 6--7)
% These are retained as aliases rather than collapsed because each
% spelling is idiomatic within its step file; prefer \fiber{...} in
% new content.
\newcommand{\LP}{\mathcal{L}(P)}
\newcommand{\F}{\mathcal{F}}
\newcommand{\Cxy}{C(x,y)}
\newcommand{\txy}{t(x,y)}
\newcommand{\Dxy}{D_{x,y}}
\newcommand{\Fxy}{\mathcal{F}_{x,y}}
\newcommand{\Fxyp}{\mathcal{F}_{x,y}^{\sharp}}
\newcommand{\Fuv}{\mathcal{F}_{u,v}}
\newcommand{\piXY}{\pi_{x,y}}
\newcommand{\Z}{\mathbb{Z}}
\newcommand{\bd}{\partial}
\newcommand{\vol}{\operatorname{vol}}
\newcommand{\sgn}{\operatorname{sgn}}
\newcommand{\pos}{\mathrm{pos}}
\newcommand{\LE}{\mathcal{L}}
\newcommand{\lext}{\mathcal{L}}
\newcommand{\fib}{\mathcal{F}}
\newcommand{\Rich}{\mathcal{R}}
\newcommand{\fiber}[1]{\mathcal{F}_{#1}}
\newcommand{\BK}{\mathrm{BK}}
\newcommand{\symdiff}{\triangle}
\newcommand{\Omegareg}{\Omega^{\mathrm{reg}}}
\newcommand{\IC}{I_C}
\newcommand{\IB}{I_B}
\newcommand{\IA}{I_A}

% All step files now use the one-argument form \GAP{...} / \DEP{...}.
% The zero-argument form previously installed here has been removed
% after unifying step2.tex's local macros.
\newcommand{\GAP}[1]{\par\noindent\textbf{[GAP:]}\ #1\par}
\newcommand{\DEP}[1]{\textbf{[DEPENDS ON: #1]}}

\title{The $1/3$--$2/3$ Conjecture for Width-3 Posets}
\author{D.~Miller\thanks{Correspondence: \texttt{d.miller@hey.com}.}}
\date{\today}

\begin{document}
\maketitle

\begin{abstract}
The $1/3$--$2/3$ conjecture of Kislitsyn, Fredman, and Linial asserts
that every finite poset that is not a chain admits a pair of
incomparable elements $(x,y)$ with $\Pr[x<y]\in[1/3,2/3]$ under a
uniformly random linear extension.  The conjecture is open in
general; the current unconditional record is $\delta^*\ge
0.2764\ldots$ (Kahn--Linial, building on Brightwell), and the $1/3$
bound is known for width-$2$ posets, semiorders, and $2$-dimensional
posets.  We prove the conjecture for posets of width~$3$.

Our argument departs from the correlation-inequality tradition: we
operate on the Bubley--Karzanov (BK) graph of linear extensions and
show, via a Cheeger-type dichotomy, that a width-$3$ counterexample
would force a low-conductance cut that is structurally impossible.
The proof proceeds in eight steps: (1)~a local interface/grid
coordinate system for rich incomparable pairs, (2)~weak edge-isoperimetric
stability on grid regions, (3--4)~coherence of interfaces and a
two-interface incompatibility lemma that drives the expansion bound,
(5)~existence of enough rich interfaces in a counterexample,
(6)~a dichotomy forcing almost all interfaces to be coherent,
(7)~coherence implies a $1$-dimensional/layered collapse of the
poset, and (8)~a final contradiction with the BK-expansion lower
bound for counterexamples.  The width-$3$ hypothesis enters through
a $2$-dimensional fiber geometry and a three-chain Dilworth
dichotomy; extending beyond width~$3$ requires genuinely new
ingredients.
\end{abstract}

\begin{center}
\textbf{Mathematics Subject Classification (2020):}\ \ 06A07 (primary);
05C48, 05A20, 60J10, 60C05 (secondary).

\smallskip
\textbf{Keywords:}\ \ $1/3$--$2/3$ conjecture; partially ordered set;
linear extension; balanced pair; width-$3$ poset; Bubley--Karzanov
graph; Cheeger inequality; expansion; conductance.
\end{center}

\tableofcontents

\section{Introduction}
\label{sec:intro}

\subsection{The $1/3$--$2/3$ conjecture}

Let $P=(X,\le)$ be a finite partially ordered set that is not a chain,
and let $\LP$ denote the set of linear extensions of $P$, endowed with
the uniform probability measure.  For incomparable $x,y\in X$ write
\[
  \Pr[x<y] \;=\; \frac{\#\{L\in\LP : x<_L y\}}{|\LP|}.
\]
A pair $(x,y)$ of incomparable elements is \emph{balanced} (or
$(1/3,2/3)$-balanced) if $\Pr[x<y]\in[1/3,2/3]$, and
$\delta$-\emph{balanced} if $\Pr[x<y]\in[\delta,1-\delta]$.  Define
\[
  \delta(P) \;=\; \max_{x\parallel y} \min\bigl(\Pr[x<y],\Pr[y<x]\bigr),
  \qquad
  \delta^* \;=\; \inf_{P\text{ not a chain}} \delta(P).
\]
The classical $1/3$--$2/3$ conjecture, stated independently by
Kislitsyn~\cite{Kislitsyn1968}, Fredman~\cite{Fredman1976}, and
Linial~\cite{Linial1984}, asserts:

\begin{conjecture}[$1/3$--$2/3$ conjecture]
\label{conj:main}
For every finite poset $P$ that is not a chain, $\delta(P)\ge 1/3$,
and in particular $\delta^*\ge 1/3$.
\end{conjecture}

The bound $1/3$ is tight: the three-element antichain plus a common
lower cover (or, equivalently, any ``$\mathbf{N}$'' with the right
dimensions) shows $\delta^*\le 1/3$.  The conjecture has been called
``one of the most intriguing problems in the combinatorial theory of
posets'' (Brightwell~\cite{Brightwell1999}), and remains open for
general finite posets.

\subsection{Known results}
\label{sec:intro-known}

We summarize the state of the art relevant to this paper.

\paragraph{Unconditional bounds on $\delta^*$.}
Kahn and Saks~\cite{KahnSaks1984} first proved $\delta^*\ge 3/11$,
using the FKG inequality and a log-concavity argument for the sequence
of rank probabilities.  This was improved by
Brightwell~\cite{Brightwell1989} to $\delta^*\ge (3-\sqrt 5)/2\approx
0.2764\ldots$ via a sharper use of XYZ-type correlation inequalities,
and by Kahn and Linial~\cite{KahnLinial1991} (see also
Brightwell--Felsner--Trotter~\cite{BFT1995}) to the current record
$\delta^*\ge 0.276\,393\ldots$, matching a quadratic root of an
explicit polynomial.  All of these arguments operate on
$\Pr[x<y]$ directly, via log-concavity and correlation inequalities.

\paragraph{Classes for which the conjecture is known.}
The conjecture is proved for several classes of posets, by arguments
tailored to each:
\begin{itemize}\setlength{\itemsep}{2pt}
  \item \emph{Width-2 posets} (two disjoint chains):
    Linial~\cite{Linial1984} established the conjecture, with a
    sharp $1/3$ bound attained by the $3$-element
    ``$\mathbf{N}$''-type configurations.
  \item \emph{Semiorders} and \emph{interval orders}:
    Brightwell~\cite{Brightwell1989}, later refined by
    Brightwell--Wright, proved the conjecture for
    semiorders; stronger $(1/3,2/3)$-type bounds are
    known for subclasses.
  \item \emph{$2$-dimensional posets}: Kahn and
    Saks~\cite{KahnSaks1984} proved the conjecture for
    $2$-dimensional posets, with a direct combinatorial argument
    exploiting the realizer.
  \item \emph{Posets with a nontrivial automorphism}:
    Ganter--Rival--Ganter and later
    Brightwell~\cite{Brightwell1999} established
    $\delta(P)\ge 1/3$ (and often $1/2$) whenever $P$ has a
    nontrivial order-preserving involution.
  \item \emph{Height-$2$ posets} (bipartite posets):
    Trotter--Gehrlein--Fishburn and others established sharp or
    near-sharp bounds in several subcases.
\end{itemize}

\paragraph{Width-$3$ posets.}
For posets of width~$3$, i.e.\ those whose largest antichain has size
exactly~$3$, the conjecture has remained open.  Peczarski's
\emph{gold partition conjecture}~\cite{Peczarski2008}, which would
imply the $1/3$--$2/3$ conjecture, has been verified for width~$3$ in
some formulations, but the standard $1/3$--$2/3$ statement
(Conjecture~\ref{conj:main}) was not previously known at width~$3$
without additional structural hypotheses.

\subsection{Main result}

This paper establishes the $1/3$--$2/3$ conjecture at width~$3$,
conditional on an arithmetic richness hypothesis (Hypothesis~A
below; see Hypothesis~\ref{hyp:arith} in Step~8 for the precise
form) that pins the intermediate regime $\gamma < 1/3 -
\delta_{\mathrm{KL}} \approx 0.057$, $n > n_0(\gamma, T)$.

\begin{theorem}[Main theorem, conditional on Hypothesis~A]
\label{thm:main}
Assume the arithmetic richness Hypothesis~A of Step~8
(Hypothesis~\ref{hyp:arith}). Let $P$ be a finite poset of width at
most~$3$ that is not a chain. Then $P$ admits a pair $(x,y)$ of
incomparable elements with
\[
  \tfrac{1}{3} \;\le\; \Pr[x<y] \;\le\; \tfrac{2}{3}.
\]
Equivalently, $\delta(P)\ge 1/3$ for every non-chain poset $P$ of
width~$\le 3$, under Hypothesis~A.
\end{theorem}

Theorem~\ref{thm:main} is established in Step~8 (Part~III below) as
the contrapositive of a structural impossibility result, via a
quantitative chain of constants developed across Steps~1--7. The
Step~8 assembly is unconditional on $\gamma \ge 1/3 -
\delta_{\mathrm{KL}}$ (by the Kahn--Linial
bound~\cite{KahnLinial1991}) and on $n \le n_0(\gamma, T)$ in the
small-$\gamma$ residual regime (by the finite-enumeration base
case, Remark~\ref{rem:small-n}); Hypothesis~A is the additional
input needed to close the large-$n$ tail of the small-$\gamma$
regime via the BK-expansion cascade.

\subsection{Approach: Cheeger-type expansion on the BK graph}

Our proof diverges methodologically from the previous work described
in Section~\ref{sec:intro-known}.  Rather than arguing about
$\Pr[x<y]$ through correlation inequalities, we operate on the
\emph{Bubley--Karzanov graph} $\BK(P)$ on the vertex set $\LP$: two
linear extensions are joined by an edge iff they differ by a single
adjacent transposition that is compatible with $P$.  This graph
supports a natural reversible Markov chain, and
Bubley--Karzanov~\cite{BubleyKarzanov1998} proved its spectral gap
controls mixing of uniformly random linear extensions.

The starting observation, which we make quantitative in Step~8, is
that a width-$3$ counterexample to Conjecture~\ref{conj:main} would
force $\BK(P)$ to have a \emph{low-conductance cut}: the set
$\{L : x<_L y\}$ for a putative most-balanced pair $(x,y)$ would have
both small volume-fraction defect and small edge boundary, because
every boundary edge corresponds to a local swap across the $x$--$y$
interface.  A Cheeger-type inequality for $\BK(P)$ then translates
``counterexample'' into ``low-conductance cut of a prescribed
structural type''.

The bulk of the paper (Steps~1--7) shows that such a cut cannot exist
in a width-$3$ poset, by a rigidity argument:

\begin{enumerate}[label=\textbf{(\arabic*)},leftmargin=2em]
  \item \textbf{Local interface theorem (Step~1):} for each
    incomparable \emph{rich} pair $(x,y)$ (those sharing many common
    neighbors), the set of linear extensions fibers over a
    $2$-dimensional grid region $\Dxy\subseteq[0,t]^2$, and BK moves
    within a fiber are exactly unit grid moves.  This is the local
    coordinate system.
  \item \textbf{Weak grid stability (Step~2):} a subset of a grid
    region with small edge boundary is close to a monotone
    staircase.  Applied fiberwise, it shows that a low-conductance
    cut is locally monotone on each rich fiber.
  \item \textbf{Coherence of interfaces (Step~3):} two rich
    interfaces $(x,y)$ and $(u,v)$ are \emph{coherent} if their
    induced monotone directions agree on the overlap.  We equip each
    rich pair with a canonical sign $\eta_{xy}$ and axis.
  \item \textbf{Two-interface incompatibility (Step~4):} two
    incoherent rich interfaces cannot both be locally monotone
    without producing many BK-boundary crossings; this is the main
    expansion engine.
  \item \textbf{Richness (Step~5):} a width-$3$ counterexample
    produces enough rich interfaces for the expansion engine to
    bite.
  \item \textbf{Dichotomy (Step~6):} combining Steps~2--5, either
    many rich interfaces are incoherent (immediate contradiction) or
    almost all are coherent.
  \item \textbf{Coherence implies collapse (Step~7):} if almost all
    rich interfaces are coherent, then $P$ admits a single global
    axis/potential, forcing $P$ into a layered $1$-dimensional form
    that is already balanced by direct inspection.
  \item \textbf{Conclusion (Step~8):} combining the above with
    Theorem~E (counterexample $\Rightarrow$ low BK-expansion) and an
    explicit constants chain, no width-$3$ counterexample can exist.
\end{enumerate}

\subsection{Comparison with previous approaches}

Previous unconditional work on the $1/3$--$2/3$ conjecture
(Kahn--Saks, Brightwell, Kahn--Linial) has centered on
\emph{correlation inequalities} applied directly to the distribution
of $\Pr[x<y]$, together with log-concavity of rank
distributions.  That approach has yielded the numerical bounds
$\delta^*\ge 0.276\ldots$ but appears to be genuinely stuck below
$1/3$; the functional-analytic slack in the $XYZ$ inequality is known
to cost a constant factor.

Our argument, by contrast, is a \emph{Cheeger/spectral} argument on
$\BK(P)$.  We never compare $\Pr[x<y]$ to another probability
directly; instead, we use a counterexample only through its induced
low-conductance cut, and derive the contradiction entirely from the
local combinatorial geometry of BK edges across that cut.  The
closest precedents are Bubley and Karzanov's
mixing-time work~\cite{BubleyKarzanov1998} and the Wilson-type
analyses that use BK graphs for sampling, but those papers aim at
bounds on the spectral gap rather than at the
conjecture.  A secondary novelty is the explicit use of
width~$\le 3$ as a $2$-dimensional local coordinate system: the
rich-pair fibers $\Fxy$ are genuinely planar grids, and the
$2$-dimensional grid isoperimetry of Step~2 uses this planarity in an
essential way; the same argument does not extend verbatim to larger
width.

\subsection{Road map}

The paper is organized in three parts, mirroring the logical flow of
the argument.

\paragraph{Part~I: Local interface theory (Steps~1--4).}
Step~1 sets up the rich-pair fibers and the grid coordinate system.
Step~2 establishes weak grid stability (a quantitative edge-isoperimetric
inequality on grid regions) and applies it fiberwise.  Step~3 defines
interface coherence and the canonical axis/sign.  Step~4 proves the
two-interface incompatibility lemma.

\paragraph{Part~II: Dichotomy, coherence, and collapse (Steps~5--7).}
Step~5 proves the richness lemma: a counterexample produces many rich
interfaces.  Step~6 assembles the dichotomy: either the cut has
incoherent interfaces (and Step~4 triggers) or almost all are coherent.
Step~7 treats the coherent case and shows it forces a $1$-dimensional
collapse of $P$.

\paragraph{Part~III: Conclusion (Step~8).}
Step~8 packages the Dirichlet/conductance translation, invokes
Theorem~E (counterexample $\Rightarrow$ low BK-expansion), and
combines Steps~1--7 into a contradiction, proving
Theorem~\ref{thm:main} with explicit (non-sharp) constants.

\paragraph{What must be proved vs.\ what is crude.}
Throughout, we aim for a sound but unoptimized argument.  The
load-bearing results are:
\begin{itemize}\setlength{\itemsep}{2pt}
  \item Step~1: local interface/grid model for rich pairs.
  \item Step~2: weak grid edge-isoperimetric stability.
  \item Step~4: two-interface incompatibility lemma.
  \item Step~5: existence of enough rich interfaces.
  \item Step~7: coherence $\Rightarrow$ global collapse.
\end{itemize}
The numerical constants ($T$, $\delta$, $\varepsilon$, etc.) and the
threshold defining ``rich'' are treated crudely; we do not attempt to
optimize them, and no quantitative improvement over the
$0.276\ldots$ record is claimed beyond the width-$3$ restriction.

\part{Local interface theory}


\begin{abstract}
Fix a width-$3$ poset $P$ and a rich incomparable pair $x\parallel y$ sharing
a long common neighbor chain $\Cxy$.
We construct a \emph{good fiber} $\Fxy\subseteq \LP$ carrying an explicit
coordinate map $\piXY:\Fxy\to \Dxy\subseteq [0,\txy]^2$ such that (i) the domain
$\Dxy$ is order-convex, (ii) BK moves inside $\Fxy$ are exactly unit grid moves
in $(i,j)$, and (iii) the complementary ``bad'' set is at most $1$-dimensional:
it is covered by $O(1)$ strips of width $O(\txy)$ and total mass
$O(\txy\cdot |\Fxy|^{1/2})$ in a sense made precise below.
This is the width-$3$ interface classification used by Steps 2--7.
A corollary extracts the commuting-square structure on overlaps of two rich
interfaces needed in Step~4, and a triple-overlap corollary extracts a
local $\Z^3$ commuting-cube model on positive-density triple overlaps,
needed for Step~7 (gap G4).
\end{abstract}

\section{Setup}

\subsection{Poset, linear extensions, BK graph}

Let $P=(X,\le_P)$ be a finite poset with $|X|=n$, width $w(P)\le 3$.
Write $\LP$ for the set of linear extensions of $P$ and fix the
\emph{Bubley--Karzanov} (BK) graph on $\LP$: two extensions $L,L'\in\LP$ are
adjacent iff they differ by transposing a pair of adjacent incomparable
elements. We use the abbreviation ``BK move'' for such a transposition.

\subsection{Common-neighbor chain of an incomparable pair}

For $x\parallel y$, define
\[
\Cxy \;:=\; \{\,z\in X : z\parallel x,\ z\parallel y\,\}.
\]
More usefully for the width-$3$ setting, set
\[
N(x,y) \;:=\; \{\,z\in X : z\parallel x\text{ and }z\parallel y\,\},
\]
and let $\Cxy\subseteq N(x,y)$ be the largest subset that is a chain in $P$.

\begin{lemma}[Common-neighbor chain in width 3]\label{lem:CNchain-width3}
If $w(P)\le 3$ and $x\parallel y$, then every $z\in N(x,y)$ is comparable to every other $z'\in N(x,y)$; in particular $N(x,y)=\Cxy$ is itself a chain.
\end{lemma}

\begin{proof}
Suppose for contradiction that $z,z'\in N(x,y)$ are incomparable to each other. We verify that $\{x,y,z,z'\}$ is a $4$-element antichain by checking all $\binom{4}{2}=6$ pairs:
\begin{enumerate}[label=(\roman*),nosep]
\item $x\parallel y$ by the hypothesis of the lemma;
\item $x\parallel z$ since $z\in N(x,y)$;
\item $y\parallel z$ since $z\in N(x,y)$;
\item $x\parallel z'$ since $z'\in N(x,y)$;
\item $y\parallel z'$ since $z'\in N(x,y)$;
\item $z\parallel z'$ by the contradiction hypothesis.
\end{enumerate}
Hence $\{x,y,z,z'\}$ is a $4$-element antichain, violating $w(P)\le 3$.
\end{proof}

Set
\[
\txy \;:=\; |\Cxy|,
\]
so the common chain has $\txy$ elements. Write
\[
\Cxy=\{c_1<_P c_2<_P\dots <_P c_{\txy}\}.
\]

\begin{definition}[Rich pair]\label{def:rich}
Fix a threshold $T\ge 1$.
An incomparable pair $x\parallel y$ is \emph{rich} (at threshold $T$) if $\txy\ge T$.
\end{definition}

\section{The fiber associated to a rich pair}

Fix a rich pair $x\parallel y$ with common chain $\Cxy=\{c_1,\dots,c_{\txy}\}$.
Each linear extension $L\in\LP$ totally orders $X$; restricting to
$\{x,y\}\cup\Cxy$ yields a linear extension of the induced subposet
$P|_{\{x,y\}\cup\Cxy}$.

\subsection{Local coordinates of a linear extension}

Because $x\parallel y$ and each $c_k$ is incomparable to both $x$ and $y$,
the restriction of $L$ to $\{x,y\}\cup\Cxy$ is uniquely determined by:
\begin{enumerate}[label=(\alph*),nosep]
\item the relative $L$-order of $x$ and $y$, which is either $x<_L y$ or $y<_L x$;
\item the number $i=i(L)\in\{0,1,\dots,\txy\}$ of $c_k$'s that $L$-precede $x$;
\item the number $j=j(L)\in\{0,1,\dots,\txy\}$ of $c_k$'s that $L$-precede $y$.
\end{enumerate}
The chain structure of $\Cxy$ forces $c_1<_L c_2<_L\dots<_L c_{\txy}$ (they are
comparable in $P$), so $i(L)$ and $j(L)$ encode the positions of $x,y$ inside
the sorted image of $\Cxy$.

\begin{definition}[Local coordinates]\label{def:local-coords}
For $L\in\LP$ define the pair
\[
\piXY(L) \;:=\; (i(L),\,j(L)) \;\in\; \{0,\dots,\txy\}^2.
\]
\end{definition}

\subsection{The fiber $\Fxy$ and the sign marker}

We partition $\LP$ by the local ordering type of $\{x,y\}\cup\Cxy$:
define $\sigma(L)\in\{+,-\}$ by $\sigma(L)=+$ iff $x<_L y$.
For each fixed sign $\epsilon\in\{+,-\}$ write
\[
\Fxy^{\epsilon} \;:=\; \{L\in\LP : \sigma(L)=\epsilon\}.
\]

Two linear extensions $L,L'\in\LP$ are said to have the \emph{same external
ordering} (for the pair $x,y$) if they agree on the relative order of every
pair of elements of $X\setminus(\{x,y\}\cup\Cxy)$ as well as on the relative
order of such an element with each $c_k$. Equivalently: $L$ and $L'$ determine
the same total order on $X\setminus\{x,y\}$ when $\Cxy$ is collapsed to its
chain order.

\begin{definition}[Good fiber]\label{def:good-fiber}
Let $\mathrm{Ext}$ range over equivalence classes of ``external orderings''
as above. For each external ordering class $E$, the set
\[
\Fxy(E) \;:=\; \{L\in\LP : L \text{ has external ordering } E\}
\]
is a \emph{raw fiber}.
We further partition $\Fxy(E)$ by the sign $\sigma$.
A raw fiber $\Fxy(E)$ is \emph{good} if
\begin{enumerate}[label=(G\arabic*),nosep]
\item\label{cond:G1} the restriction $\piXY:\Fxy(E)\to\{0,\dots,\txy\}^2$ is injective (i.e., the coordinate pair $(i,j)$ together with the sign $\sigma$ determines $L$);
\item\label{cond:G2} the image $\Dxy(E):=\piXY(\Fxy(E))$ is order-convex in $\Z^2$, meaning if $(i_0,j_0),(i_1,j_1)\in\Dxy(E)$ with $i_0\le i_1$ and $j_0\le j_1$, and $(i,j)$ lies in the rectangle $[i_0,i_1]\times[j_0,j_1]$, then $(i,j)\in\Dxy(E)$;
\item\label{cond:G3} BK edges internal to $\Fxy(E)$ are exactly the unit grid moves in $(i,j)$: two states $L,L'\in\Fxy(E)$ are BK-adjacent iff $\sigma(L)=\sigma(L')$ and $\piXY(L)-\piXY(L')\in\{\pm e_1,\pm e_2\}$.
\end{enumerate}
The union
\[
\Fxy \;:=\; \bigcup_{E\text{ good}} \Fxy(E)
\]
is the \emph{good fiber set} of the rich pair $(x,y)$.
\end{definition}

\begin{remark}
The domain $\Dxy(E)$ depends on the external ordering $E$: different external
orderings produce different ``available'' windows of $(i,j)$ values.
For most arguments it is enough to know that each $\Dxy(E)$ is order-convex
and contained in $[0,\txy]^2$, and to sum over $E$.
\end{remark}

\section{Main theorem: Width-$3$ interface classification}

\begin{theorem}[Local interface theorem]\label{thm:interface}
Let $P$ be a width-$3$ poset and $x\parallel y$ a rich pair with common chain
$\Cxy$ of length $\txy=t$. Then:
\begin{enumerate}[label=(\roman*),nosep]
\item\label{part:coords}
There is a well-defined coordinate map $\piXY:\LP\to\{0,\dots,t\}^2$ given by
Definition~\ref{def:local-coords}.
\item\label{part:goodfiber}
There is a partition $\LP=\Fxy\sqcup \mathrm{Bad}_{x,y}$ with
$\Fxy=\bigsqcup_{E}\Fxy(E)$ a disjoint union of good fibers (in the sense of
Definition~\ref{def:good-fiber}).
\item\label{part:moves}
Let $L\in\Fxy(E)$ and let $L'$ be obtained from $L$ by a BK move swapping an
$L$-adjacent incomparable pair $\{a,b\}$. Then exactly one of the following
holds:
\begin{enumerate}[label=(\alph*),nosep]
\item $\{a,b\}\cap(\{x,y\}\cup\Cxy)\ne\emptyset$ and $\{a,b\}\ne\{x,y\}$: the
move acts as a unit grid step in $(i,j)$ with $\sigma$ preserved (an
\emph{internal grid move}); $L'\in\Fxy(E)$ off the bad set.
\item $\{a,b\}=\{x,y\}$: the move is a \emph{sign flip}: it exchanges $x$ and
$y$, changes $\sigma$ from $\pm$ to $\mp$, and $(i,j)$ is unchanged;
$L'\in\Fxy(E)$. Case~(b) is available at $L$ iff $i(L)=j(L)$ (i.e., $x$ and
$y$ occupy adjacent slots with no intervening $c_k$).
\item $\{a,b\}\cap(\{x,y\}\cup\Cxy)=\emptyset$: the move is an
\emph{external--external swap}; $(i(L'),j(L'),\sigma(L'))=(i(L),j(L),\sigma(L))$
and $L'\in\Fxy(E')$ for some raw-fiber class $E'$ (with $E'=E$ or $E'\ne E$).
\end{enumerate}
In particular, every BK move either stays inside $\Fxy$ with controlled action
on $(i,j,\sigma)$ (cases (a),(b)) or transits between raw fibers $\Fxy(E)\to
\Fxy(E')$ with $(i,j,\sigma)$ fixed (case (c)).
\item\label{part:bad}
The bad set $\mathrm{Bad}_{x,y}$ is at most $1$-dimensional in the following
sense. It admits a decomposition
\[
\mathrm{Bad}_{x,y} \;=\; \bigcup_{k=1}^{K} S_k
\]
with $K=O(1)$, where each $S_k$ is a \emph{strip}: its $\piXY$-image is
contained in $O(1)$ axis-parallel lines of $[0,\txy]^2$, and $S_k$ is a
union of at most $n-\txy-2$ raw fibers $\Fxy(E)$, each of size $\le\txy+1$.
In particular every raw fiber meeting $\mathrm{Bad}_{x,y}$ satisfies
$|\Fxy(E)|\le\txy+1$, yielding the combinatorial bound
\[
|\mathrm{Bad}_{x,y}|\;=\;O(n\cdot\txy),
\]
and, under the quantitative richness $|\Fxy|\ge c\,n^{2}$ used in
Corollaries~\ref{cor:overlap}--\ref{cor:triple-overlap}, the equivalent
abstract bound
\[
|\mathrm{Bad}_{x,y}|\;=\;O\!\bigl(\txy\cdot|\Fxy|^{1/2}\bigr).
\]
\end{enumerate}
\end{theorem}

\begin{proof}
Part~\ref{part:coords} is immediate from Definition~\ref{def:local-coords}
and Lemma~\ref{lem:CNchain-width3}: within any $L\in\LP$ the chain $\Cxy$
sits as a consecutive-ordered sequence $c_1<_L c_2<_L\dots<_L c_{\txy}$
(because $c_i<_P c_{i+1}$), and $i(L),j(L)$ are well-defined counts.

\smallskip
\textbf{Part~\ref{part:moves}.}
Let $L\in\LP$ and let $L'$ be obtained from $L$ by a BK move that swaps the
adjacent incomparable pair $\{a,b\}$. There are three exhaustive cases.

\emph{Case 1: $\{a,b\}\cap(\{x,y\}\cup\Cxy)=\emptyset$ (external--external).}
Since neither $a$ nor $b$ equals $x$, $y$, or any $c_k$, swapping them does
not change the relative order of $x$ with any $c_k$, of $y$ with any $c_k$,
or of $x$ with $y$. Hence $i(L')=i(L)$, $j(L')=j(L)$, and
$\sigma(L')=\sigma(L)$. The external-ordering class $E$ may nevertheless
change: it records the permutation of external elements relative to the
positions of the $c_k$'s, so the swap of $a,b$ preserves $E$ iff no $c_k$
lies between the ``$a$-slot'' and the ``$b$-slot'' in the external ordering
(otherwise the swap moves across a collapsed-$\Cxy$ boundary and yields a
distinct class $E'$). Either way $L'$ lies in some raw fiber $\Fxy(E')$ (with $E'=E$ or
$E'\ne E$) and $(i,j,\sigma)$ are unchanged. This gives case~(c) of
part~(iii).

\emph{Case 2: $\{a,b\}=\{x,y\}$.}
Then $\sigma$ flips, and both $i(L')=i(L)$ and $j(L')=j(L)$ provided no
$c_k$ sits strictly between $x$ and $y$ in $L$. For $x,y$ to be
\emph{adjacent} in $L$, we need no element at all between them; in particular
no $c_k$, so $i(L)=j(L)$. Conversely, $i(L)=j(L)$ means the same number of
$c_k$'s precede both $x$ and $y$, so no $c_k$ sits between them, so (under
Case 2 applicability) the swap is available. This gives case (b) of part~(iii).

\emph{Case 3: $\{a,b\}\cap(\{x,y\}\cup\Cxy)\ne\emptyset$ but $\{a,b\}\ne\{x,y\}$.}
By Lemma~\ref{lem:CNchain-width3} and $w(P)\le 3$, the only incomparabilities
within $\{x,y\}\cup\Cxy$ are the pairs $\{x,c_k\}$ and $\{y,c_k\}$ (since
$c_k$'s are pairwise comparable and $\{x,y\}$ is handled by Case~2).
Moreover, a BK move can swap $x\leftrightarrow c_k$ iff $x$ and $c_k$ are
$L$-adjacent; swapping them changes $i(L)$ by $\pm 1$ and leaves
$j(L),\sigma(L)$ unchanged. Symmetrically for $y\leftrightarrow c_k$.
A mixed case $\{a,b\}=\{z,c_k\}$ with $z\notin\{x,y\}$ and $c_k\in\Cxy$:
this requires $z\parallel c_k$. In width $3$, the set $\{z\}\cup\Cxy$ contains
no $4$-element antichain, so $z$ is comparable to all but at most two $c_k$'s.
After discarding states where such a move occurs adjacent to $x$ or $y$ in
$L$ (a lower-order bad set, counted in part~(iv)), the move preserves
$(i,j,\sigma)$.

Thus in each good raw fiber every BK move falls into cases 1--3 above, and
the allowed $(i,j,\sigma)$-altering moves are exactly
$\pm e_1,\pm e_2$ (unit grid moves in $i,j$ with $\sigma$ fixed) together
with the sign-flip move at the diagonal $i=j$.

\smallskip
\textbf{Part~\ref{part:goodfiber}.}
The map $\piXY$ together with $\sigma$ is injective on any raw fiber $\Fxy(E)$
where the external ordering class determines the positions of all external
elements relative to the collapsed $\Cxy$: by construction this is the
definition of $E$, so injectivity \ref{cond:G1} holds automatically on each
raw fiber.

Order-convexity \ref{cond:G2} of $\Dxy(E)$: Let $(i_0,j_0),(i_1,j_1)\in\Dxy(E)$
with $i_0\le i_1$, $j_0\le j_1$. Any $(i,j)$ in the corresponding rectangle
can be reached from $(i_0,j_0)$ by a sequence of unit grid moves, each of
which (by Part~\ref{part:moves}) corresponds to a BK move that swaps
$x$ (or $y$) with an adjacent $c_k$; these moves keep $L$ inside $\Fxy(E)$
and inside $\LP$. The only obstruction is that at some intermediate step a
``Case 3 mixed'' bad swap becomes forced, which by Part~\ref{part:moves}
and Part~\ref{part:bad} happens only on the bad set.
For raw fibers $E$ that do \emph{not} meet the bad set (call these
\emph{good $E$}), all intermediate $(i,j)$ states are realizable in
$\Fxy(E)$. Therefore $\Dxy(E)$ is order-convex.

Combining: a raw fiber $\Fxy(E)$ is good iff the external ordering $E$
admits no Case~3 mixed conflicts for any $(i,j)\in\Dxy(E)$; for such $E$,
conditions \ref{cond:G1}--\ref{cond:G3} all hold.

\smallskip
\emph{Completeness of the partition.}
We verify $\LP=\Fxy\sqcup\mathrm{Bad}_{x,y}$ as asserted.
First, the raw fibers $\{\Fxy(E)\}_E$ partition $\LP$ by definition: the relation
``$L\sim L'$ iff $L,L'$ agree on the relative order of every pair in
$X\setminus(\{x,y\}\cup\Cxy)$ and on the relative order of each such element with
each $c_k$'' is an equivalence relation on $\LP$, and its classes are exactly
the raw fibers $\Fxy(E)$. In particular every $L\in\LP$ belongs to a unique
raw fiber.
Second, the ``good vs.\ bad'' dichotomy at the level of external classes $E$ is
exhaustive by excluded middle: either $E$ admits a Case~3 mixed conflict at some
$(i,j)\in\Dxy(E)$ (i.e., there exists an external $z\in X\setminus(\{x,y\}\cup\Cxy)$
and an index $k$ with $z\parallel c_k$ and with $z$ forced to be $L$-adjacent
to $x$ or $y$ for some $L\in\Fxy(E)$), or it does not. Declare
$E$ \emph{good} in the former absence, \emph{bad} in the latter presence. These
two conditions are mutually exclusive and cover all $E$. Set
\[
\Fxy \;=\; \bigsqcup_{E\text{ good}}\Fxy(E),\qquad
\mathrm{Bad}_{x,y} \;=\; \bigsqcup_{E\text{ bad}}\Fxy(E).
\]
There is no third option: a raw fiber is not ``partly good and partly bad,'' since
goodness \ref{cond:G1}--\ref{cond:G3} is a property of the class $E$ (injectivity,
order-convexity of $\Dxy(E)$, and the BK-edge classification are statements about
the whole fiber, not individual $L$'s). Hence every $L\in\LP$ lies in exactly one
of $\Fxy$ or $\mathrm{Bad}_{x,y}$, and the two sets are disjoint unions of raw
fibers as claimed. This proves Part~\ref{part:goodfiber}.

\smallskip
\textbf{Part~\ref{part:bad}.}
Write $Z:=X\setminus(\{x,y\}\cup\Cxy)$, so $|Z|=n-\txy-2$.
A raw fiber $\Fxy(E)$ lies in $\mathrm{Bad}_{x,y}$ exactly when some
intermediate $(i,j)$-grid state forces a Case~3 mixed swap, i.e., there
exist $z\in Z$ and $k\in\{1,\dots,\txy\}$ with $z\parallel c_k$ and with
$z$ $L$-adjacent to some $a\in\{x,y\}$ during a unit-grid traversal.

By Lemma~\ref{lem:CNchain-width3} and $w(P)\le 3$, the set
$I(z):=\{k:z\parallel c_k\}$ is a contiguous interval of length $\le 2$;
otherwise $\{z,c_{k_1},c_{k_2},c_{k_3}\}$ together with $\{x,y\}$ would
contain a width-$4$ antichain. If $I(z)=\varnothing$, then $z$ never
produces a Case~3 obstruction; drop such $z$ from the count.

\emph{Per-$z$ bad locus.} Fix $z\in Z$ with $|I(z)|\in\{1,2\}$.
Assign to each $z$-induced obstruction at $L$ a \emph{type}
\[
\kappa\;=\;(a,k,\epsilon)\;\in\;\{x,y\}\times I(z)\times\{+,-\},
\]
where $a\in\{x,y\}$ records which of $\{x,y\}$ is $L$-adjacent to $z$ and
$\epsilon$ records the $L$-order of $z$ relative to $c_k$. Any such type
$\kappa$ fixes exactly one of $\{i(L),j(L)\}$: indeed, $a=x$ with $z$
$L$-adjacent to $x$ near $c_k$ forces $i(L)\in\{k-1,k\}$, and symmetrically
$a=y$ forces $j(L)\in\{k-1,k\}$. Hence
\[
\piXY\bigl(\mathrm{Bad}_{x,y}^{(z)}\bigr)\;\subseteq\;\ell(z),
\]
where $\ell(z)\subseteq[0,\txy]^2$ is a union of at most
$|I(z)|\cdot 4\le 8$ axis-parallel lines.

\emph{Grouping into $K=4$ strip families.} Collapse $\kappa$ to the pair
$(a,\epsilon)\in\{x,y\}\times\{+,-\}$ (independent of $z$ and $k$). For
each pair $(a,\epsilon)$ set
\[
S_{(a,\epsilon)}\;:=\;\bigcup_{z\in Z}\bigl\{L\in\mathrm{Bad}_{x,y}^{(z)}
:L\text{ has a $z$-obstruction of type }(a,k,\epsilon)
\text{ for some }k\in I(z)\bigr\}.
\]
Since every $L\in\mathrm{Bad}_{x,y}$ carries at least one $z$-obstruction,
$\mathrm{Bad}_{x,y}=\bigcup_{(a,\epsilon)}S_{(a,\epsilon)}$, so $K=4=O(1)$.

\emph{Structure of each $S_{(a,\epsilon)}$.} Inside $S_{(a,\epsilon)}$
only one of $\{i(L),j(L)\}$ is free: if $a=x$ then $i(L)\in\{0,\dots,\txy\}$
is among the values $\{k-1,k:k\in I(z)\}$ (at most $2|Z|+2$ distinct
values, i.e.\ $O(1)$ per $z$), and $j(L)$ is unconstrained. Thus the
$\piXY$-image of $S_{(a,\epsilon)}$ lies in $O(1)$ axis-parallel lines
(of length $\le\txy+1$ each). For each fixed type $(a,k,\epsilon)$ and each
external $z\in Z$, the triple $(z,a,c_k,\epsilon)$ determines the relative
$L$-order of $z$ with respect to $\{a,c_{k-1},c_k,c_{k+1}\}$, so at most
\emph{one} raw fiber $E$ per $z$ realises this type. Summing over the
$O(1)$ values of $k\in I(z)$ and over $z\in Z$, the number of raw fibers
$\Fxy(E)$ contributing to $S_{(a,\epsilon)}$ is at most $O(|Z|)=O(n)$.

\emph{Size of each raw fiber in the strip.} If $\Fxy(E)\subseteq
\mathrm{Bad}_{x,y}$, then every $L\in\Fxy(E)$ lies on the fixed coordinate
line $\ell$ (this is what ``bad raw fiber'' means: the obstruction is not a
mere point defect but applies throughout the raw fiber, since a single
Case~3 conflict propagates through the order-convex domain). By
condition~\ref{cond:G1} (injectivity of $\piXY$ on each raw fiber),
$|\Fxy(E)|\le|\ell\cap[0,\txy]^2|\le\txy+1$.

Combining,
\[
|S_{(a,\epsilon)}|\;\le\;|Z|\cdot(\txy+1)\;=\;O(n\cdot\txy),
\]
and summing over the $K=4$ pairs gives the combinatorial bound
\[
|\mathrm{Bad}_{x,y}|\;=\;O(n\cdot\txy).\tag{$\ast$}
\]

\emph{Passage to $O(\txy\cdot|\Fxy|^{1/2})$.}
Bound~$(\ast)$ is an inequality involving $n$; the abstract form asserted
in the statement involves $\sqrt{|\Fxy|}$. The two are linked by the
quantitative richness hypothesis $|\Fxy|\ge c\,n^{2}$ (implicit in the rich
regime $\txy=\Theta(n)$, $|\Fxy|\gg n^{2}$ used in
Corollaries~\ref{cor:overlap}--\ref{cor:triple-overlap}): indeed, if
$|\Fxy|\ge c\,n^{2}$ then $n\le c^{-1/2}\sqrt{|\Fxy|}$, hence
\[
|\mathrm{Bad}_{x,y}|\;=\;O(n\cdot\txy)\;=\;
O\!\bigl(\txy\cdot|\Fxy|^{1/2}\bigr).
\]
Conversely, the quantitative richness hypothesis is automatic in the main
application: a generic good raw fiber $\Fxy(E)$ has $|\Dxy(E)|=\Theta(\txy^{2})$
(its $(i,j)$-domain is an order-convex subset of $[0,\txy]^2$ of full
$2$-dimensional extent by \ref{cond:G2}), and there are at least $|Z|!/\txy^{O(1)}$
good external classes arising from free orderings of external elements,
so $|\Fxy|=\Omega(\txy^{2}\cdot|Z|!/\txy^{O(1)})\gg n^{2}$ whenever
$\txy\to\infty$ along the rich regime. Outside this regime the combinatorial
bound $(\ast)$ is the operative one. This proves Part~\ref{part:bad}.
\end{proof}

\section{Consequences for two interfaces}

The following corollary is the clause \textup{(S1.4)} of Step~1 consumed
by Step~4: it upgrades Theorem~\ref{thm:interface} from a single-pair
classification to a joint commuting statement on the overlap of two
rich interfaces.

\begin{lemma}[Bounded interaction of distinct rich pairs]\label{lem:bounded-interaction}
Let $(x,y)$ and $(u,v)$ be two distinct rich pairs. Then the ``interaction
locus''
\[
\mathrm{Int}_{xy,uv}
\;:=\;\{L\in\Fxy\cap\F_{u,v}:
\exists\,a\in\{x,y\},\ b\in\{u,v\}\cup C(u,v)\text{ with }a,b\ L\text{-adjacent}\},
\]
is contained in $O(1)$ strips,
each a union of $O(\txy+t_{u,v})$ raw fibers of size $O(\txy+t_{u,v})$.
In particular
\[
|\mathrm{Int}_{xy,uv}|
\;=\; O\bigl((\txy+t_{u,v})^2\bigr).
\]
\end{lemma}

\begin{proof}
By Lemma~\ref{lem:CNchain-width3} applied to each pair, $\Cxy$ and $C(u,v)$
are chains. By $w(P)\le 3$, any external element $z\notin\{x,y\}\cup\Cxy$
is incomparable to a contiguous interval of $\Cxy$ of length $\le 2$
(the argument of Theorem~\ref{thm:interface}~\ref{part:bad}). The
interaction locus requires some $a\in\{x,y\}$ to be $L$-adjacent to some
$b\in\{u,v\}\cup C(u,v)$; adjacency restricts $(i(L),j(L))$ to at most
$4$ horizontal/vertical lines of length $\le\txy$ in $D_{x,y}$
(for $b\in\{u,v\}$) and to at most $2$ such lines per $c_\ell\in C(u,v)$
that is incomparable to $\{x,y\}$. Width-$3$ bounds the latter by a constant
per chain position. Summing across the at most $2$ rôles of $(u,v)$ and its
chain gives $O(1)$ strip families, each of the stated shape. The symmetric
case is identical.
\end{proof}

\begin{corollary}[Commuting overlap; clause \textup{(S1.4)} of Step~4]\label{cor:overlap}
Let $(x,y)$ and $(u,v)$ be two rich pairs, and write $\Fxy,\F_{u,v}$ for
their good fibers with coordinate maps $\piXY,\pi_{u,v}$.
Set
\[
\Omega_{xy,uv} \;:=\; \Fxy \cap \F_{u,v},
\qquad
\Omega^{\circ}_{xy,uv}
\;:=\; \Omega_{xy,uv}\setminus
\bigl(\mathrm{Bad}_{x,y}\cup\mathrm{Bad}_{u,v}\cup\mathrm{Int}_{xy,uv}\bigr),
\]
where $\mathrm{Int}_{xy,uv}$ is the interaction locus of
Lemma~\ref{lem:bounded-interaction}. Then:
\begin{enumerate}[label=\textup{(\alph*)},nosep]
\item\label{it:commute} \textup{(Commuting + local $\Z^2$ grid.)}
For each $L\in\Omega^{\circ}_{xy,uv}$ the four generators
$\pm e_1^{(x,y)},\pm e_2^{(x,y)}$ and $\pm e_1^{(u,v)},\pm e_2^{(u,v)}$
act by transpositions supported on disjoint element pairs, hence commute
pairwise. Consequently the restriction of the BK graph to
$\Omega^{\circ}_{xy,uv}$ near $L$ embeds a piece of the Cayley graph of
$\Z^4$; modding out by the two order-convex coordinate images
(Theorem~\ref{thm:interface}\ref{part:goodfiber}) gives a local
$\Z^2$ grid in the joint coordinates $\bigl((i,j),(p,q)\bigr)$
ranging over a product of order-convex domains. In particular, around
every such $L$ a full $2\times 2$ BK-square with one $(x,y)$-direction
and one $(u,v)$-direction is embedded, and all four compositions of
the two BK moves coincide.
\item\label{it:density} \textup{(Positive density; uniform bound.)}
\[
|\Omega_{xy,uv}\setminus\Omega^{\circ}_{xy,uv}|
\;\le\;|\mathrm{Int}_{xy,uv}|
\;=\;O\!\bigl((\txy+t_{u,v})^{2}\bigr),
\]
so in the rich regime $\txy,t_{u,v}=\Theta(n)$ and
$|\Omega_{xy,uv}|\gg n^{2}$,
\[
|\Omega^{\circ}_{xy,uv}|\;\ge\;(1-o(1))\,|\Omega_{xy,uv}|.
\]
\end{enumerate}
\end{corollary}

\begin{proof}
\emph{\ref{it:commute}.}
A unit $(x,y)$-move swaps $x$ (or $y$) with an adjacent $c_k\in\Cxy$,
and a unit $(u,v)$-move swaps $u$ (or $v$) with an adjacent
$c_\ell\in C(u,v)$. Outside $\mathrm{Int}_{xy,uv}$ these four elements
are all distinct and no two of them are $L$-adjacent to each other, so the
two transpositions are supported on disjoint pairs; disjoint transpositions
commute. Therefore on $\Omega^{\circ}_{xy,uv}$ the four generators commute
pairwise, which is precisely the definition of a local $\Z^2$ grid in the
$(\pi_{x,y},\pi_{u,v})$-coordinates. Order-convexity of each coordinate
image (Theorem~\ref{thm:interface}\ref{part:goodfiber}) guarantees that
the generated grid lies inside the joint fiber.

\emph{\ref{it:density}.}
By Theorem~\ref{thm:interface}\ref{part:goodfiber} the sets
$\mathrm{Bad}_{x,y}=\LP\setminus\Fxy$ and $\mathrm{Bad}_{u,v}=\LP\setminus\F_{u,v}$
are the complements of the good fibers, so
\[
\Omega_{xy,uv}\;=\;\Fxy\cap\F_{u,v}
\]
is disjoint from both $\mathrm{Bad}_{x,y}$ and $\mathrm{Bad}_{u,v}$.
Subtracting the two Bad sets in the definition of $\Omega^{\circ}_{xy,uv}$
therefore removes nothing, and
\[
\Omega_{xy,uv}\setminus\Omega^{\circ}_{xy,uv}
\;=\;\Omega_{xy,uv}\cap\mathrm{Int}_{xy,uv}
\;\subseteq\;\mathrm{Int}_{xy,uv}.
\]
Lemma~\ref{lem:bounded-interaction} gives
$|\mathrm{Int}_{xy,uv}|=O\!\bigl((\txy+t_{u,v})^{2}\bigr)$ unconditionally,
which is the first inequality.  In the rich regime $\txy,t_{u,v}=\Theta(n)$
and $|\Omega_{xy,uv}|\gg n^{2}$, this is $O(n^{2})=o(|\Omega_{xy,uv}|)$,
yielding $|\Omega^{\circ}_{xy,uv}|\ge(1-o(1))|\Omega_{xy,uv}|$.
No distribution (fiber-slice concentration) hypothesis on $\Omega_{xy,uv}$
is needed: the bound is driven entirely by the $O(1)$-many strips of the
interaction locus, each of size $O((\txy+t_{u,v})^{2})$.
\end{proof}

\begin{remark}
This matches verbatim the clause \textup{(S1.4)} cited in
\texttt{step4.tex}~\S1, Theorem~1.1: a positive-density subset
$\Omega^{\circ}_{xy,uv}\subseteq\mathcal F_{x,y}\cap\mathcal F_{u,v}$
on which the two interfaces' BK moves commute and generate a local
$\Z^2$ grid. Proposition~G1 of \texttt{step4.tex} (rectangle decomposition)
takes $\Omega^{\circ}_{xy,uv}$ as its input.
\end{remark}
\begin{corollary}[Triple-overlap commuting cube]\label{cor:triple-overlap}
Let $e_{xy}=(x,y)$, $e_{yz}=(y,z)$, $e_{xz}=(x,z)$ be three rich pairs with
good fibers $\Fxy,\F_{y,z},\F_{x,z}$ and coordinate maps
$\piXY,\pi_{y,z},\pi_{x,z}$. Define the triple overlap
\[
\Omega^{(3)} \;:=\; \Fxy\,\cap\,\F_{y,z}\,\cap\,\F_{x,z},
\]
and assume it has positive density: $|\Omega^{(3)}|\ge\delta|\LP|$ for some
$\delta>0$. Let
\[
\Omega^{\circ\circ\circ} \;\subseteq\; \Omega^{(3)}
\]
be the subset of states at which each of the three interfaces admits a
unit-grid BK move in each of its two coordinate directions, and none of the
three pairwise regular-overlap conditions of Corollary~\ref{cor:overlap} is
violated. Then:
\begin{enumerate}[label=(\alph*),nosep]
\item\label{it:tri-coords}
(Simultaneous coordinates.) On $\Omega^{\circ\circ\circ}$ the three maps
$\piXY,\pi_{y,z},\pi_{x,z}$ are simultaneously defined and each takes values in
an order-convex domain; the three local coordinate systems are simultaneously
valid.
\item\label{it:tri-cube}
(Commuting $\Z^3$ cube model.) The three families of BK moves
$\{\pm e_r^{(x,y)}\}_{r=1,2}$, $\{\pm e_r^{(y,z)}\}_{r=1,2}$,
$\{\pm e_r^{(x,z)}\}_{r=1,2}$ pairwise commute on
$\Omega^{\circ\circ\circ}$, so any triple of generators (one per interface)
spans an embedded $2\times 2\times 2$ BK-cube whose $8$ vertices all lie in
$\Omega^{(3)}$. Iterating generates a local $\Z^3$ action by BK moves.
\item\label{it:tri-consistency}
(Triple consistency of projections.) For any downset $S\subseteq\LP$, the three
images
\[
A_{x,y}:=\piXY(S\cap\Omega^{\circ\circ\circ}),\quad
A_{y,z}:=\pi_{y,z}(S\cap\Omega^{\circ\circ\circ}),\quad
A_{x,z}:=\pi_{x,z}(S\cap\Omega^{\circ\circ\circ})
\]
agree on the same underlying cut: the three $\pm$-coordinate decompositions
of $\Omega^{\circ\circ\circ}$ induced by $S$ refine to a common partition of
$\Omega^{\circ\circ\circ}$ into in-$S$/out-of-$S$ fibers, and membership in
$S$ at a given $L\in\Omega^{\circ\circ\circ}$ is determined by any one of the
three coordinate pairs modulo the shared-chain indices.
\item\label{it:tri-bad}
(Bad mass.) The complement has size at most
\[
|\Omega^{(3)}\setminus\Omega^{\circ\circ\circ}|
\;\le\;
O\bigl(\txy+t_{y,z}+t_{x,z}\bigr)\cdot\sqrt{|\Omega^{(3)}|},
\]
hence is lower-order when all three common-chain lengths are $\Theta(n)$.
\end{enumerate}
\end{corollary}

\begin{proof}
\emph{Standing setup.} Each of the three rich pairs $(x,y),(y,z),(x,z)$ is
individually incomparable in $P$, so $\{x,y,z\}$ is a $3$-antichain. By
Lemma~\ref{lem:CNchain-width3}, any $c\in\Cxy\cap C(y,z)$ would satisfy
$c\parallel x,y,z$, producing the $4$-antichain $\{x,y,z,c\}$ in violation
of $w(P)\le 3$; the other two intersections are symmetric. Hence
\begin{equation}\label{eq:tri-disjoint}
\Cxy\cap C(y,z) \;=\; \Cxy\cap C(x,z) \;=\; C(y,z)\cap C(x,z) \;=\; \emptyset.
\end{equation}
(Note that $z\in\Cxy$, $x\in C(y,z)$, $y\in C(x,z)$ are not excluded from
the \emph{opposite} chains; the pairwise intersection is empty nonetheless,
since membership in the intersection requires incomparability with all
three of $x,y,z$.) The exclusion of the three interaction loci built into
the definition of $\Omega^{\circ\circ\circ}$
(Lemma~\ref{lem:bounded-interaction}) further ensures that for any
$L\in\Omega^{\circ\circ\circ}$ and any unit $(x,y)$-move $\tau=(a\,c_k)$ at
$L$ (with $a\in\{x,y\}$, $c_k\in\Cxy$ $L$-adjacent to $a$), the chain
element satisfies
\begin{equation}\label{eq:tri-chainelt}
c_k\ \notin\ \{x,y,z\}\cup C(y,z)\cup C(x,z):
\end{equation}
$c_k\in\{x,y\}$ is excluded by $\Cxy$'s definition; $c_k\in\{z\}\cup C(y,z)$
would put $L\in\mathrm{Int}_{xy,yz}$, and $c_k\in\{z\}\cup C(x,z)$ would put
$L\in\mathrm{Int}_{xy,xz}$, both excluded. Symmetric statements hold for
unit moves in the other two interfaces.

\ref{it:tri-coords} By Theorem~\ref{thm:interface}~\ref{part:coords}
the three maps $\piXY,\pi_{y,z},\pi_{x,z}$ are defined on all of $\LP$.
Membership in $\Omega^{\circ\circ\circ}$ forces each interface's good-fiber
conditions \ref{cond:G1}--\ref{cond:G3}, so
Theorem~\ref{thm:interface}~\ref{part:goodfiber} gives order-convexity of
each image. Hence all three coordinate systems are simultaneously valid.

\ref{it:tri-cube} Fix $L\in\Omega^{\circ\circ\circ}$. A unit-grid generator
at $L$ from interface $(x,y)$ is an adjacent transposition $(a\,c_k)$ moving
some $a\in\{x,y\}$; each of the two coordinate directions of the interface
corresponds to one choice of $a$. Among the $2^3=8$ triples of such
generators (one per interface), only those in which the three outer
elements moved are three distinct members of $\{x,y,z\}$ yield
disjoint-support triples; exactly two such \emph{canonical} triples exist.
Fix one, say
\[
\tau^{xy}=(x\,c_1),\qquad \tau^{yz}=(y\,c_2),\qquad \tau^{xz}=(z\,c_3),
\]
with $c_1,c_2,c_3$ $L$-adjacent elements of $\Cxy,C(y,z),C(x,z)$
respectively. By~\eqref{eq:tri-disjoint} and~\eqref{eq:tri-chainelt}, the
three chain elements lie in pairwise disjoint chains and none equals any
of $x,y,z$, so $\{x,y,z,c_1,c_2,c_3\}$ are six distinct elements of $X$.
Consequently the three adjacent transpositions
$\tau^{xy},\tau^{yz},\tau^{xz}$ have pairwise disjoint $2$-element
supports, and so commute as permutations of $X$.

Disjoint-support adjacent transpositions remain mutually enabled after the
application of any subset: an adjacent transposition with support disjoint
from $\{a,c_r\}$ changes only the relative order of its own two elements,
so it preserves the $L$-adjacency of $a$ and $c_r$ (same argument as in
the proof of Corollary~\ref{cor:overlap}~\ref{it:commute}). Hence the
$2^3=8$ states obtained from $L$ by applying each subset of
$\{\tau^{xy},\tau^{yz},\tau^{xz}\}$ are reachable by BK-moves along each
of the $12$ cube edges and form an embedded $2\times 2\times 2$ BK-cube.
Order-convexity of each coordinate image
(Theorem~\ref{thm:interface}~\ref{part:goodfiber}) places all $8$ vertices
in $\Omega^{(3)}$.

Iterating the three commuting unit-shift directions at each cube corner
generates a local free $\Z^3$-action by BK-moves on the BK-connected
component of $L$ in $\Omega^{(3)}$. The action is free (rather than an
action of a $(\Z/2)^3$-quotient) because each shift direction advances the
corresponding coordinate pair by a unit step, and injectivity of each
coordinate map on its good fiber \ref{cond:G1} forbids any nontrivial
word in the three directions from returning to $L$: the word's total shift
in each of the three coordinate pairs is a non-negative linear combination
of unit vectors, vanishing only if every generator count is zero.

\ref{it:tri-consistency} Let $L\in\Omega^{\circ\circ\circ}$ and let $L'$
arise from $L$ by a unit $(x,y)$-move $\tau=(a\,c_k)$ with $a\in\{x,y\}$
and $c_k\in\Cxy$ satisfying~\eqref{eq:tri-chainelt}. Being an adjacent
transposition, $\tau$ flips only the relative $L$-order of the pair
$\{a,c_k\}$; the relative order of every other pair in $X$ is preserved.
We check each coordinate pair:
\begin{enumerate}[label=(\roman*),nosep,leftmargin=*]
\item $\piXY(L')-\piXY(L)\in\{\pm e_1,\pm e_2\}$, by construction of $\tau$
(one of $i,j$ shifts by $\pm1$ according as $a=x$ or $a=y$).
\item $\pi_{y,z}(L')=\pi_{y,z}(L)$. The two components of $\pi_{y,z}$ count
$\#\{c\in C(y,z):c<_L y\}$ and $\#\{c\in C(y,z):c<_L z\}$, and so depend
only on the relative orders of pairs $\{\alpha,\beta\}$ with
$\alpha\in\{y,z\}$, $\beta\in C(y,z)$. For the flipped pair $\{a,c_k\}$ to
be of this form would require (i)~$\alpha=a\in\{y,z\}$ and
$\beta=c_k\in C(y,z)$, or (ii)~$\alpha=c_k\in\{y,z\}$ and
$\beta=a\in C(y,z)$. In case~(i), $a\in\{x,y\}\cap\{y,z\}=\{y\}$, forcing
$c_k\in\Cxy\cap C(y,z)=\emptyset$ by~\eqref{eq:tri-disjoint}; in case~(ii),
$c_k\in\{y,z\}$ contradicts~\eqref{eq:tri-chainelt}. Hence no relevant
count is affected.
\item $\pi_{x,z}(L')=\pi_{x,z}(L)$, by the symmetric argument: case~(i) now
forces $a\in\{x,y\}\cap\{x,z\}=\{x\}$ and $c_k\in\Cxy\cap C(x,z)=\emptyset$;
case~(ii) forces $c_k\in\{x,z\}$, excluded by~\eqref{eq:tri-chainelt}.
\end{enumerate}
Thus \emph{exactly one} of the three coordinate pairs shifts by a unit
step under any unit-grid move at $L$. The same conclusion holds, by
symmetric arguments, for unit $(y,z)$- and $(x,z)$-moves. Iterating, within
each BK-connected component of $\Omega^{\circ\circ\circ}$ (and each fixed
common external ordering) the joint map
$L\mapsto\bigl(\piXY(L),\pi_{y,z}(L),\pi_{x,z}(L)\bigr)$ is a local
$\Z^3$-equivariant chart modulo shared-chain indices, and any one of the
three coordinate pairs determines the other two modulo those indices.

Consequently, for any downset $S\subseteq\LP$, membership $L\in S$ at
$L\in\Omega^{\circ\circ\circ}$ is a function of any single coordinate pair,
so the three induced level sets $A_{x,y},A_{y,z},A_{x,z}$ refine to a
common in-$S$/out-of-$S$ partition of $\Omega^{\circ\circ\circ}$. This is
the claimed triple consistency.

\ref{it:tri-bad} By Theorem~\ref{thm:interface}~\ref{part:bad}, each
interface's bad set is covered by $O(1)$ strips of width $O(t_\bullet)$.
By Corollary~\ref{cor:overlap}~\ref{it:density}, each pairwise irregular
overlap has mass $O(t_\bullet+t_\bullet)\sqrt{|\Omega^{(3)}|}$. The set
$\Omega^{(3)}\setminus\Omega^{\circ\circ\circ}$ is contained in the union
of the three single-interface bad sets and the three pairwise irregular
overlaps; a union bound gives the stated total.
\end{proof}

\section{Output interface for subsequent steps}

Theorem~\ref{thm:interface} is the sole input Steps~2--4 use from Step~1.
We record the deliverables explicitly so that the next steps may cite them
by name:

\begin{description}[style=nextline]
\item[Fiber.] For each rich $(x,y)$, a set $\Fxy\subseteq\LP$, partitioned
into good raw fibers $\Fxy(E)$.
\item[Coordinates.] A map $\piXY:\Fxy\to\Dxy\subseteq[0,\txy]^2$, with
$\Dxy=\bigsqcup_E\Dxy(E)$ a disjoint union of order-convex domains.
\item[Grid structure.] BK edges inside each good raw fiber are exactly
unit grid moves in $(i,j)$, plus at most one sign-flip edge at $i=j$.
\item[Bad set bound.] $|\mathrm{Bad}_{x,y}|\ll |\Fxy|$, contained in
$O(1)$ strips of width $O(\txy)$.
\item[Overlap structure.] Corollary~\ref{cor:overlap}: on regular overlaps
between two rich interfaces, $2\times 2$ commuting BK squares exist on a
set of full overlap mass up to lower-order error.
\item[Triple-overlap cube.] Corollary~\ref{cor:triple-overlap}: on any
positive-density triple overlap of three rich interfaces, there is a
subset $\Omega^{\circ\circ\circ}$ on which the three coordinate systems
are simultaneously valid, the three BK-move families generate a local
$\Z^3$ cube model, and the three projections of any downset $S$ agree on
a common cut. This is the input Step~7 uses to discharge gap G4.
\end{description}

These deliverables feed directly into:
\begin{itemize}[nosep]
\item Step~2 (cut rigidity on a fiber): applies a weak grid isoperimetric
lemma to $A_{x,y}:=\piXY(S\cap\Fxy)\subseteq\Dxy$, using order-convexity
of $\Dxy$ and the unit-grid BK structure.
\item Step~3 (coherence): assigns to each rich interface a $\pm1$
orientation based on the monotone approximation of $A_{x,y}$ in its
coordinate grid; order-convexity makes this well-defined.
\item Step~4 (incompatibility lemma): uses Corollary~\ref{cor:overlap}
to realize the $2\times 2$ conflict squares; the bad-set bound ensures
lower-order loss when summing over overlap blocks.
\item Step~7 (gap G4, cocycle integration on triple overlaps): uses
Corollary~\ref{cor:triple-overlap} to realize the local $\Z^3$ cube
model on triple overlaps; the cube commutation and triple consistency
are what makes the cocycle equation pointwise well-defined.
\end{itemize}



\begin{abstract}
Step 2 transports a \emph{global} low-conductance assumption on a cut $S \subseteq \lext(P)$
into a \emph{per-fiber} structural statement: on most rich good fibers $\fib_{x,y}$, the
image $A_{x,y} = \pi_{x,y}(S \cap \fib_{x,y})$ is close to a monotone staircase region in
the planar grid $D_{x,y} \subseteq [0,t]^2$. The core input is a weak planar isoperimetric
stability result (``small grid boundary $\Rightarrow$ close to staircase''). The Step~1
dependencies are made explicit so downstream steps (3, 4, 6) can cite the
conclusion cleanly.
\end{abstract}


\section{Inputs from Step 1}
\label{sec:step1-inputs}

Step~2 cannot be stated without the following Step~1 outputs. This section specifies
the minimum formal content Step~1 must deliver, exactly as consumed by the lemmas
below. Step~1 is responsible for proving
(S1.1)--(S1.6); Step~2 uses them as hypotheses.

\begin{definition}[Step~1 data specification]\label{def:step1-data}
There exists a threshold $T\in\mathbb{Z}_{\ge 1}$ and a family of rich incomparable
pairs $\mathcal{R}$ such that, for each $(x,y)\in\mathcal{R}$ with $t:=t(x,y)\ge T$,
Step~1 produces the following data and proves the stated properties.
\begin{enumerate}[label=\textup{(S1.\arabic*)}, leftmargin=*]
    \item \emph{\textbf{Good fiber.}} A set $\fib_{x,y}\subseteq\lext(P)$, together
          with a quantitative density bound
          \[
            |\fib_{x,y}| \;\ge\; c_0\,t^2
          \]
          for a universal constant $c_0>0$ (independent of $(x,y)$ and of $P$). Used
          in Proposition~\ref{prop:per-fiber} as the shape hypothesis
          $|D_{x,y}|\ge c\,t^2$ of Lemma~\ref{lem:weak-grid}.
    \item \emph{\textbf{Order-convex domain.}} A set $D_{x,y}\subseteq[0,t]^2\cap\mathbb{Z}^2$
          that is \emph{order-convex}: for all $(i,j),(i',j')\in D_{x,y}$ with
          $i\le i'$ and $j\le j'$, every lattice point $(i'',j'')$ with
          $i\le i''\le i'$ and $j\le j''\le j'$ lies in $D_{x,y}$. Used as the
          ambient region of Lemma~\ref{lem:weak-grid} and Lemma~\ref{lem:row-decomp}.
    \item \emph{\textbf{Coordinate bijection.}} A bijection
          $\pi_{x,y}\colon\fib_{x,y}\to D_{x,y}$. Used to transport
          $S\cap\fib_{x,y}\subseteq\lext(P)$ to $A_{x,y}:=\pi_{x,y}(S\cap\fib_{x,y})
          \subseteq D_{x,y}$.
    \item \emph{\textbf{BK$=$unit grid identification.}} For all
          $L,L'\in\fib_{x,y}$, $\{L,L'\}$ is a BK-edge of $\lext(P)$ if and only if
          $\pi_{x,y}(L')-\pi_{x,y}(L)\in\{(\pm 1,0),(0,\pm 1)\}$. Used in
          Proposition~\ref{prop:per-fiber} to identify the BK-boundary of
          $S\cap\fib_{x,y}$ (Definition~\ref{def:edge-internal}) with the grid
          boundary of $A_{x,y}$ (Definition~\ref{def:grid-bdy}).
    \item \emph{\textbf{Bad-fiber bound.}} The exceptional set
          \[
            \fib_{x,y}^{\mathrm{bad}} \;:=\; \bigl\{L\in\fib_{x,y} :
                   \text{(S1.4) fails at $L$, i.e., some BK-edge at $L$ is not a unit grid move}\bigr\}
          \]
          satisfies $|\fib_{x,y}^{\mathrm{bad}}|\le C_1\,t$ for a universal
          constant $C_1$. Since $|\fib_{x,y}|\ge c_0 t^2$ by (S1.1), this yields
          $|\fib_{x,y}^{\mathrm{bad}}|=O(t)=o(|\fib_{x,y}|)$ as $t\to\infty$.
    \item \emph{\textbf{Bounded BK-edge multiplicity.}} There exists a universal
          constant $K\in\mathbb{Z}_{\ge 1}$, independent of $P$ and $(x,y)$, such
          that
          \[
            \sup_{e\in E(\lext(P))}
              \#\bigl\{(x,y)\in\mathcal{R} : e\text{ is internal to }\fib_{x,y}\bigr\}
              \;\le\; K,
          \]
          where ``internal to $\fib_{x,y}$'' is in the sense of
          Definition~\ref{def:edge-internal}. This is \eqref{eq:mult} of
          Lemma~\ref{lem:fiber-avg} and makes that lemma non-vacuous.
\end{enumerate}
\end{definition}

\DEP{Step~1 (formal proof of (S1.1)--(S1.6)).}
Step~2 assumes (S1.1)--(S1.6) as hypotheses; each
appears exactly where indicated above. In particular, (S1.6) must be recorded
explicitly in \texttt{step1.tex}: the width-$3$ interface classification (every BK
swap $\{a,b\}$ internal to $\fib_{x,y}$ has $\{a,b\}\subseteq\{x,y\}\cup C(x,y)$)
is what produces the constant $K$.

\section{The weak grid stability lemma}
\label{sec:weak-grid}

This is the combinatorial core of Step~2. It is a purely planar statement about
subsets of an order-convex region of $\mathbb{Z}^2$.

\subsection{Staircase regions}

\begin{definition}[Monotone staircase]\label{def:staircase}
Let $D \subseteq \mathbb{Z}^2$ be an order-convex region. A set $M \subseteq D$ is a
\emph{staircase region of type} $\sigma\in\{+,-\}$ if there exists a monotone function
$f_\sigma\colon \mathbb{Z}\to\mathbb{Z}\cup\{\pm\infty\}$ such that
\[
M = \begin{cases}
  \{(i,j)\in D : j \le f_+(i)\} & \sigma=+, \text{ $f_+$ weakly increasing}, \\
  \{(i,j)\in D : j \le f_-(i)\} & \sigma=-, \text{ $f_-$ weakly decreasing}.
\end{cases}
\]
A staircase region is thus a lower set for one of the two partial orders
$\le_{+}$ (generated by $(1,0)$ and $(0,1)$) or $\le_{-}$ (generated by $(1,0)$ and
$(0,-1)$). We write $\mathrm{Stair}(D)$ for the union of all staircase regions in $D$.
\end{definition}

\begin{remark}[Four orientations]\label{rem:four-orient}
Complementation and axis-swap give four equivalent staircase types corresponding to the
four quadrant orientations $\{\mathrm{NE},\mathrm{NW},\mathrm{SE},\mathrm{SW}\}$. We pick
the $\{+,-\}$ normalization for downstream use in Step~3's sign $\sigma_{x,y}$. In the
conclusion of Lemma~\ref{lem:weak-grid} below, $\mathrm{Stair}(D)$ is understood to
include all four orientations; Corollary~\ref{cor:sign} then restricts to $\{+,-\}$ once
the horizontal reflection has been used to fix a side.
\end{remark}

\subsection{Grid boundary}

\begin{definition}[Grid boundary inside $D$]\label{def:grid-bdy}
For $A \subseteq D$, let
\[
\bd_D A := \bigl|\{\{u,v\} : u\in A,\ v\in D\setminus A,\ \|u-v\|_1=1\}\bigr|.
\]
\end{definition}

\subsection{Main lemma}

\begin{lemma}[Weak grid stability]\label{lem:weak-grid}
For every $\varepsilon>0$ there exists $\delta = \delta(\varepsilon) > 0$ with
$\delta(\varepsilon)\to 0$ as $\varepsilon\to 0$, such that the following holds. Let
$D \subseteq [0,t]^2$ be order-convex with $|D|\ge c t^2$ for some $c>0$, and let
$A \subseteq D$ with
\[
\bd_D A \;\le\; \varepsilon\, t.
\]
Then there exists a staircase region $M \in \mathrm{Stair}(D)$ with
\[
|A \symdiff M| \;\le\; \delta\,|D|.
\]
\end{lemma}

\begin{proof}
Set $B:=\bd_D A\le\varepsilon t$. Lemma~\ref{lem:row-decomp} decomposes
$B=B^h+B^v$ (horizontal and vertical contributions), and the column transpose of
Lemma~\ref{lem:row-decomp} applies verbatim to columns.

We obtain $\delta(\varepsilon)=O_c(\varepsilon)$, linear in $\varepsilon$, which is
even better than the $O(\varepsilon^{1/3})$ claimed.

\medskip\noindent\emph{Step 1 (row and column cleanup).}
Applying Corollary~\ref{cor:most-rows-intervals} in both directions, the sets
\[
I_b:=\{i:m(A_i)\ge 2\},\qquad J_b:=\{j:m(A^j)\ge 2\}
\]
satisfy $|I_b|\le B^h/2$ and $|J_b|\le B^v/2$. For $i\notin I_b$, order-convexity
gives $D_i=[u_i,v_i]\cap\mathbb Z$, and when $A_i\ne\emptyset$ we have
$A_i=[\ell_i,r_i]\cap\mathbb Z$ with $u_i\le\ell_i\le r_i\le v_i$. Columns
satisfy the analogous structure $D^j=[p^j,q^j]\cap\mathbb Z$ and (for
$j\notin J_b$ with $A^j\ne\emptyset$) $A^j=[\mathrm{top}_j,\mathrm{bot}_j]\cap\mathbb Z$.

By Lemma~\ref{lem:1d-components}, for $i\notin I_b$ with $A_i\notin\{\emptyset,D_i\}$,
\(
b_i^h=2-\mathbf 1[\ell_i=u_i]-\mathbf 1[r_i=v_i].
\)
Summing, and writing $I_{ne}:=\{i\notin I_b:A_i\ne\emptyset\}$,
\begin{equation}\label{eq:wg-anchor-row}
|I_{ne}\setminus I_L|+|I_{ne}\setminus I_R|\;\le\;B^h,
\end{equation}
where $I_L:=\{i\in I_{ne}:\ell_i=u_i\}$ are the \emph{left-anchored} rows and
$I_R:=\{i\in I_{ne}:r_i=v_i\}$ the \emph{right-anchored}. The column analogue gives
\begin{equation}\label{eq:wg-anchor-col}
|J_{ne}\setminus J_B|+|J_{ne}\setminus J_T|\;\le\;B^v,
\end{equation}
with $J_B,J_T$ denoting bottom- and top-anchored columns.

\medskip\noindent\emph{Step 2 (majority orientation; WLOG reduction).}
By~\eqref{eq:wg-anchor-row}, $\min(|I_{ne}\setminus I_L|,|I_{ne}\setminus I_R|)\le B^h/2$;
by~\eqref{eq:wg-anchor-col}, $\min(|J_{ne}\setminus J_B|,|J_{ne}\setminus J_T|)\le B^v/2$.
Choose the row anchor (L or R) and column anchor (B or T) achieving these minima; this
selects one of the four orientations. The reflections $i\mapsto t-i$ and $j\mapsto t-j$
preserve order-convexity (mapping $D\subseteq[0,t]^2$ to another order-convex subset of
$[0,t]^2$), preserve $\bd_D(\cdot)$, and permute $\mathrm{Stair}(D)$ in the
four-orientation reading of Remark~\ref{rem:four-orient}: $i\mapsto t-i$ swaps
$J_B\leftrightarrow J_T$ (and exchanges types $\pm$), while $j\mapsto t-j$ swaps
$I_L\leftrightarrow I_R$ (and exchanges type-$\pm$ staircases with their complements).
Hence WLOG the favored orientation is $(L,T)$, which corresponds to type-$+$ staircases
(the SE orientation: $M_i=[u_i,f(i)]$ is left-anchored in rows and $M^j=[i^\star(j),q^j]$
is top-anchored in columns).

We henceforth assume
\begin{equation}\label{eq:wg-LT}
|I_{ne}\setminus I_L|\le B^h/2,\qquad |J_{ne}\setminus J_T|\le B^v/2.
\end{equation}

\medskip\noindent\emph{Step 3 (staircase construction).}
Set $r_i^\star:=r_i$ for $i\in I_L$ and $r_i^\star:=-\infty$ otherwise. Define
\[
f(i)\;:=\;\max_{i'\le i}r_{i'}^\star\qquad(\max\emptyset:=-\infty),
\]
a weakly increasing function $\mathbb Z\to\mathbb Z\cup\{-\infty\}$. Let
\(
M\;:=\;\{(i,j)\in D:j\le f(i)\},
\)
a type-$+$ staircase region in the sense of Definition~\ref{def:staircase}.

\medskip\noindent\emph{Step 4 (bound on $|A\setminus M|$, row-wise).}
If $i\in I_L$, then for every $j\in A_i$, $j\le r_i=r_i^\star\le f(i)$, so
$(i,j)\in M$. Hence $A\setminus M\subseteq(I_b\cup(I_{ne}\setminus I_L))\times\mathbb Z$,
and
\begin{equation}\label{eq:wg-AsetM}
|A\setminus M|
\;\le\;\sum_{i\in I_b\cup(I_{ne}\setminus I_L)}|A_i|
\;\le\;(|I_b|+|I_{ne}\setminus I_L|)(t+1)
\;\le\;(B^h/2+B^h/2)(t+1)
\;\le\;\varepsilon t(t+1).
\end{equation}

\medskip\noindent\emph{Step 5 (bound on $|M\setminus A|$, column-wise).}
For each column $j$, set
\(
i^\star(j):=\min\{i\in I_L:r_i\ge j\}
\)
($:=+\infty$ if the set is empty). Since $f$ is weakly increasing and $f(i)\ge r_{i^\star(j)}\ge j$
iff $i\ge i^\star(j)$ (for $i^\star(j)<\infty$), we have
$M^j=[i^\star(j),q^j]\cap D^j$ (empty when $i^\star(j)=\infty$).

\emph{Case (i): $i^\star(j)=\infty$.} Then $M^j=\emptyset$, so $|M^j\setminus A^j|=0$.

\emph{Case (ii): $i^\star(j)<\infty$, $j\notin J_b$.} Then $(i^\star(j),j)\in A$, so
$A^j\ne\emptyset$ and $i^\star(j)\in A^j=[\mathrm{top}_j,\mathrm{bot}_j]$. Hence
\[
M^j\setminus A^j
\;=\;[i^\star(j),q^j]\setminus[\mathrm{top}_j,\mathrm{bot}_j]
\;=\;[\mathrm{bot}_j+1,q^j]\cap D^j,
\]
which is empty when $j\in J_T$ (so $\mathrm{bot}_j=q^j$) and has cardinality
$\le q^j-\mathrm{bot}_j\le t+1$ otherwise. Therefore, by~\eqref{eq:wg-LT},
\[
\sum_{\substack{j\notin J_b\\ i^\star(j)<\infty}}|M^j\setminus A^j|
\;\le\;|J_{ne}\setminus J_T|\cdot(t+1)
\;\le\;(B^v/2)(t+1)
\;\le\;\varepsilon t(t+1)/2.
\]

\emph{Case (iii): $j\in J_b$.} Since $M^j\subseteq D^j=[p^j,q^j]\cap\mathbb Z\subseteq\{0,\dots,t\}$, we have $|M^j|\le q^j-p^j+1\le t+1$. Bound crudely $|M^j\setminus A^j|\le|M^j|\le t+1$:
\[
\sum_{j\in J_b}|M^j\setminus A^j|
\;\le\;|J_b|(t+1)\;\le\;(B^v/2)(t+1)\;\le\;\varepsilon t(t+1)/2.
\]

Summing the three cases (only (ii) and (iii) contribute),
\begin{equation}\label{eq:wg-MsetA}
|M\setminus A|\;=\;\sum_j|M^j\setminus A^j|\;\le\;\varepsilon t(t+1).
\end{equation}

\medskip\noindent\emph{Step 6 (combine).}
By~\eqref{eq:wg-AsetM} and~\eqref{eq:wg-MsetA},
\[
|A\symdiff M|\;\le\;2\varepsilon t(t+1)\;\le\;2\varepsilon\cdot 2t^2
\;\le\;\tfrac{4\varepsilon}{c}\,|D|
\]
for $t\ge 1$ (using $t(t+1)\le 2t^2$ and $t^2\le|D|/c$). Setting
$\delta(\varepsilon):=4\varepsilon/c$ gives $\delta(\varepsilon)\to 0$ as
$\varepsilon\to 0$, completing the proof.
\end{proof}

\subsection{Orientation assignment}

\begin{corollary}[Sign $\sigma(A)$]\label{cor:sign}
Under the hypotheses of Lemma~\ref{lem:weak-grid}, the staircase region $M$ is unique
up to a further set of size $\delta'|D|$ with $\delta'\to 0$ as $\delta\to 0$. Moreover,
if $\min(|A|,|D\setminus A|)\ge\eta|D|$ and $M$ is \emph{not} within $\beta|D|$ of a
horizontal strip $\{(i,j)\in D:j\le c\}$, then its type $\sigma(A)\in\{+,-\}$ is
well-defined. The thresholds $\eta,\beta>0$ are inherited by Step~3 as its
nontriviality parameter $\theta_{x,y}$.
\end{corollary}

\begin{proof}
We prove uniqueness of $M$ (the main content), then the separation bound for type
determination, and finally note the horizontal-strip degeneracy.

\medskip\noindent\emph{(i) Uniqueness of $M$.}
Let $M_1, M_2$ be any two staircase regions (of any types) with $|A\symdiff M_k|
\le\delta|D|$ for $k=1,2$. The triangle inequality gives
\begin{equation}\label{eq:sign-tri}
  |M_1\symdiff M_2| \;\le\; |M_1\symdiff A|+|A\symdiff M_2| \;\le\; 2\delta|D|.
\end{equation}
Hence $M$ is determined up to a set of size $\delta'|D|$ with $\delta':=2\delta\to 0$
as $\delta\to 0$. In particular $\bigl||M_k|-|A|\bigr|\le\delta|D|$, so under the mass
hypothesis $\min(|A|,|D\setminus A|)\ge\eta|D|$ we also have
\begin{equation}\label{eq:sign-mass}
  \min\bigl(|M_k|,\,|D\setminus M_k|\bigr) \;\ge\; (\eta-\delta)|D| \;\ge\; \tfrac{\eta}{2}|D|
  \quad\text{whenever $\delta\le\eta/2$.}
\end{equation}

\medskip\noindent\emph{(ii) Separation for non-horizontal staircases.}
\begin{claim}\label{cl:sep}
There exists $c_0=c_0(c)>0$ such that for any $+$-staircase $M_+=\{(i,j)\in D:j\le f(i)\}$
and $-$-staircase $M_-=\{(i,j)\in D:j\le g(i)\}$ satisfying
\begin{enumerate}[label=\textup{(\alph*)}, nosep]
  \item \textup{(mass)} $\min\bigl(|M_\pm|,|D\setminus M_\pm|\bigr)\ge\alpha|D|$, and
  \item \textup{(non-horizontal)} neither $M_+$ nor $M_-$ is within $\beta|D|$ of a
        horizontal strip of $D$,
\end{enumerate}
one has
\[
  |M_+\symdiff M_-|\;\ge\;c_0\,\alpha\,\beta\,|D|.
\]
\end{claim}

Granting Claim~\ref{cl:sep}, suppose $M_+$ (type $+$) and $M_-$ (type $-$) are both
$\delta$-close to $A$. By \eqref{eq:sign-mass} they satisfy~(a) at level
$\alpha=\eta/2$. If moreover $A$ is not within $\beta|D|$ of a horizontal strip, then
neither $M_\pm$ is within $(\beta-\delta)|D|\ge\beta/2\cdot|D|$ of a horizontal strip
(by triangle inequality: a horizontal strip $\beta|D|/2$-close to $M_\pm$ would be
$(\beta/2+\delta)|D|\le\beta|D|$-close to $A$, a contradiction for $\delta\le\beta/2$),
so (b) holds at level $\beta/2$. Combining Claim~\ref{cl:sep} with \eqref{eq:sign-tri}:
\[
  c_0\cdot\tfrac{\eta}{2}\cdot\tfrac{\beta}{2}\,|D|
  \;\le\;|M_+\symdiff M_-|\;\le\;2\delta|D|,
\]
forcing $\delta\ge c_0\eta\beta/8$. Hence whenever $\delta<c_0\eta\beta/8$ no two
opposite-type approximations of $A$ can coexist, so $\sigma(A)\in\{+,-\}$ is determined.

\medskip\noindent\emph{(iii) Horizontal-strip degeneracy.}
A horizontal strip $S=\{(i,j)\in D:j\le c\}$ has $f_+\equiv c$ (weakly increasing) and
$f_-\equiv c$ (weakly decreasing), so $S$ is \emph{simultaneously} of type $+$ and $-$.
If $A$ is within $\beta|D|$ of a horizontal strip, both labels are valid for the same
approximation set $M$; uniqueness of $M$ from part~(i) is preserved, but the label
$\sigma(A)$ is ambiguous. Step~3 treats such fibers as ``axis-aligned'' and does not
rely on the $\pm$ distinction there. Outside this degenerate regime---i.e., under the
non-horizontal hypothesis in the corollary---the type is determined unambiguously.

\medskip\noindent\emph{Proof of Claim~\ref{cl:sep}.}
Let $N:=|D|$. By order-convexity, each row $D_i=[a_i,b_i]\cap\mathbb Z$ is an
interval; let $I\subseteq\mathbb Z$ be the range of first coordinates, $|I|\le t+1$,
$N_i:=|D_i|$, so $N=\sum_{i\in I}N_i\ge ct^2$. Analogously each column
$D^j=[p^j,q^j]\cap\mathbb Z$ is an interval.

\smallskip
\emph{Step 1 (column-height reformulation).}
For any $\varphi\colon\mathbb Z\to\mathbb Z\cup\{\pm\infty\}$, the slice of
$\{(i,j)\in D:j\le\varphi(i)\}$ at column $i$ is the initial segment
$[a_i,\tilde\varphi(i)]\cap\mathbb Z$ of $D_i$, where
\begin{equation}\label{eq:sep-clip}
  \tilde\varphi(i):=\max\bigl(a_i-1,\,\min(b_i,\,\varphi(i))\bigr)\in\{a_i-1,\ldots,b_i\},
\end{equation}
and has cardinality $\tilde\varphi(i)-a_i+1\in[0,N_i]$. The clipping map
$x\mapsto\max(a_i-1,\min(b_i,x))$ is the metric projection onto
$[a_i-1,b_i]\subseteq\mathbb R$, hence $1$-Lipschitz and, for any
$x\in[a_i-1,b_i]$ and $y\in\mathbb R$, $|x-\tilde y_i|\le|x-y|$. Applied to
$f$ and $g$ whose column slices are both initial segments of $D_i$ (and therefore
nested),
\begin{equation}\label{eq:sep-sum}
  |M_+\symdiff M_-|=\sum_{i\in I}\bigl|\tilde f(i)-\tilde g(i)\bigr|.
\end{equation}
For any $c\in\mathbb R$, applying the same reformulation to the horizontal strip
$S_c=\{(i,j)\in D:j\le c\}$ gives
\begin{equation}\label{eq:sep-strip}
  |M_+\symdiff S_c|=\sum_{i\in I}|\tilde f(i)-\tilde c(i)|,\qquad
  \tilde c(i):=\max(a_i-1,\min(b_i,c)),
\end{equation}
and the analogous identity for $M_-$.

\smallskip
\emph{Step 2 (range bound).} Set
\[
  V_f:=\max_{i\in I}\tilde f(i)-\min_{i\in I}\tilde f(i)\in[0,\,t+1],\qquad
  c_f:=\tfrac12\bigl(\max_i\tilde f(i)+\min_i\tilde f(i)\bigr)\in[-1,t],
\]
so $|\tilde f(i)-c_f|\le V_f/2$ for every $i\in I$. Since $\tilde f(i)\in[a_i-1,b_i]$,
the $1$-Lipschitz projection bound gives $|\tilde f(i)-\tilde c_f(i)|\le|\tilde f(i)-c_f|\le V_f/2$.
Combining with~\eqref{eq:sep-strip} and hypothesis~(b),
\[
  \beta N\;\le\;|M_+\symdiff S_{c_f}|\;=\;\sum_i|\tilde f(i)-\tilde c_f(i)|\;\le\;\tfrac{t+1}{2}V_f,
\]
so $V_f\ge\tfrac{2\beta N}{t+1}\ge\beta c t$ (using $N\ge ct^2$ and $t+1\le 2t$ for $t\ge 1$).
The same argument applied to $M_-$ yields $V_g\ge\beta ct$, where
$V_g:=\max_i\tilde g(i)-\min_i\tilde g(i)$.

\smallskip
\emph{Step 3 (level sets from monotonicity).}
We now pass from $V_f\ge\beta ct$ to a quantitative statement about level sets of the
\emph{unclipped} $f$, which is weakly increasing. Fix
\[
  J_f:=\bigl\{j\in\mathbb Z:|\{i\in I:\tilde f(i)\ge j\}|\ge 1
                                   \text{ and }|\{i\in I:\tilde f(i)<j\}|\ge 1\bigr\},
\]
the set of integers strictly between $\min\tilde f$ and $\max\tilde f$. Then
$|J_f|\ge V_f - 1\ge\beta c t - 1$, which is $\ge\beta ct/2$ for $t\ge 2/(\beta c)$;
the finitely many smaller $t$ are absorbed into $c_0$. For each $j\in J_f$, the
set $\{i\in I:\tilde f(i)\ge j\}$ is nonempty and strictly contained in $I$.

Now for any integer $j$, we have the inclusion of level sets
\begin{equation}\label{eq:sep-level-incl}
  \{i\in I:\tilde f(i)\ge j\}\;\subseteq\;\{i\in I:f(i)\ge j\}\;\cup\;\{i\in I:a_i-1\ge j\},
\end{equation}
because $\tilde f(i)\ge j$ forces $\max(a_i-1,\min(b_i,f(i)))\ge j$, i.e., either
$a_i-1\ge j$ or $\min(b_i,f(i))\ge j$ (whence $f(i)\ge j$). The first set on the
right-hand side is, by monotonicity of $f$, an \emph{up-set} of $I$:
\[
  U_f(j):=\{i\in I:f(i)\ge j\}=[i_f(j),\max I]\cap I,\qquad i_f(j):=\min\{i\in I:f(i)\ge j\},
\]
with $i_f(j)\in I\cup\{+\infty\}$ weakly increasing in $j$. Symmetrically for $g$
(weakly decreasing), $D_g(j):=\{i\in I:g(i)\ge j\}=[\min I,i_g(j)]\cap I$ is a
\emph{down-set}.

\smallskip
\emph{Step 4 (mass pins level thresholds into the body of $I$).}
Pick integers $j_+,j_-$ with $j_+\ge j_-$ (to be specified) and consider the ``staircase
gap layer''
\[
  L:=\{(i,j)\in D:j_-\le j\le j_+\}.
\]
We claim that for suitable $j_\pm$, this layer contains a large sub-rectangle inside
$M_+\setminus M_-$ and, symmetrically, one inside $M_-\setminus M_+$.

For any $i\in U_f(j_+)\cap(I\setminus D_g(j_-))$ and any $j\in[j_-,j_+]\cap D_i$, the
point $(i,j)$ satisfies $j\le j_+\le f(i)$ (so $(i,j)\in M_+$) and $j\ge j_-> g(i)$
(so $(i,j)\notin M_-$); hence
\begin{equation}\label{eq:sep-rect-plus}
  \bigl(U_f(j_+)\cap(I\setminus D_g(j_-))\bigr)\times[j_-,j_+]\cap D\;\subseteq\;M_+\setminus M_-.
\end{equation}
Both $U_f(j_+)$ and $I\setminus D_g(j_-)$ are up-sets of $I$, so their intersection is
an up-set $[i^\star,\max I]\cap I$.

We now choose $j_+,j_-$ so that both (i) $|U_f(j_+)\cap(I\setminus D_g(j_-))|$ has
large column-weight and (ii) $[j_-,j_+]$ has length $\Omega(\beta t)$.

Define
\[
  j^{f}_{\alpha/4}:=\max\bigl\{j\in\mathbb Z:|U_f(j)|\text{-column-weight}\ge\alpha N/4\bigr\}.
\]
This is the largest level at which $U_f$ retains column-weight $\alpha N/4$; equivalently,
$j^f_{\alpha/4}=\max\{j:\sum_{i\ge i_f(j)}N_i\ge\alpha N/4\}$, with the convention
$\max\emptyset=-\infty$. By the mass hypothesis,
$|M_+|=\sum_j R_+(j)\ge\alpha N$, where $R_+(j)=|D^j\cap U_f(j)|$ (the row-$j$
cardinality of $M_+$) is weakly decreasing in $j$; hence the level $j=j^f_{\alpha/4}$
is well-defined and finite (strictly between $-\infty$ and $+\infty$, since
$|U_f(-\infty)|=|I|$ has weight $N\ge\alpha N/4$ and $|U_f(+\infty)|=0$). Define
$j^g_{\alpha/4}$ analogously: the largest $j$ with
$|I\setminus D_g(j)|$-column-weight $\ge\alpha N/4$.

We claim $j^f_{\alpha/4}\ge\min\tilde f$ and $j^g_{\alpha/4}\ge\min\tilde g$. For the
first, any $j\le\min\tilde f$ has $U_f(j)\supseteq\{i:\tilde f(i)\ge j\}=I$ via
\eqref{eq:sep-level-incl} applied at $j'=j$ (since for such $j$ either $a_i-1\ge j$
or $\tilde f(i)\ge j$ already); so column-weight is $N\ge\alpha N/4$.

\smallskip
\emph{Step 5 (Markov on $V_f$ to separate two thresholds).}
We now exhibit two levels $j_+ > j_-$ inside $J_f$ with $j_+-j_-\ge\beta c t/4$ and
with $|U_f(j_+)|$-column-weight $\ge\alpha N/4$ and $|I\setminus D_g(j_-)|$-column-weight
$\ge\alpha N/4$.

Set $T:=\lceil\beta ct/4\rceil$. By Step~2, $V_f\ge\beta ct\ge 4T$, so there are at
least $4T$ integer levels strictly between $\min\tilde f$ and $\max\tilde f$. By the
definition of $j^f_{\alpha/4}$, the level $j_+^f:=j^f_{\alpha/4}$ lies in
$(\min\tilde f,\max\tilde f]$ (the column-weight of $U_f(\min\tilde f)$ is
$N\ge\alpha N/4$, hence $j^f_{\alpha/4}\ge\min\tilde f$; and for $j>\max\tilde f$ we
have $U_f(j)\subseteq\{i:a_i-1\ge j\}\cup\{i:b_i\ge j\text{ and }f(i)\ge j\}$, which
has vanishing column-weight for $j>\max b_i\ge\max\tilde f$, so $j^f_{\alpha/4}\le\max b_i+1$).

Set
\[
  j_+:=j^f_{\alpha/4},\qquad j_-:=\max\bigl(\min\tilde f+1,\,j_+-T\bigr).
\]
Then $j_+ - j_-\ge 0$ and, when $j_+\ge\min\tilde f+1+T$, equals $T$; and when
$j_+<\min\tilde f+1+T$, we fall back on $j_-=\min\tilde f+1$, so $j_+-j_-\ge 0$ but
possibly small---however the alternative reflection argument below covers this case.
Assume first $j_+-j_-=T$ (the generic case); we handle the degenerate subcase at
the end of Step~6.

\emph{Column-weight of $|U_f(j_+)|$:} $\ge\alpha N/4$ by the definition of
$j^f_{\alpha/4}$.

\emph{Column-weight of $|I\setminus D_g(j_-)|$:}
By symmetry, $j^g_{\alpha/4}$ satisfies: for all $j\le j^g_{\alpha/4}$,
column-weight of $I\setminus D_g(j)$ is $\ge\alpha N/4$. We claim $j_-\le j^g_{\alpha/4}$.

Suppose not, i.e., $j_->j^g_{\alpha/4}$. Then the column-weight of
$I\setminus D_g(j_-)$ is $<\alpha N/4$, i.e.,
$|D_g(j_-)|$-column-weight $>N-\alpha N/4=(1-\alpha/4)N\ge 3N/4$ (using $\alpha\le 1$).
So for almost all $i$ (by column-weight), $g(i)\ge j_-$, which means
$\tilde g(i)\ge\min(b_i,j_-)\ge\min(b_i,\min\tilde f+1)$.
Combined with $|M_-|\ge\alpha N$, this forces $\tilde g(i)$ to be moderately large
on most columns. We now exploit the non-horizontal hypothesis on $M_-$ (applied
to the strip $S_{j_-}$):
\[
  \beta N \;\le\; |M_-\symdiff S_{j_-}| \;=\;\sum_i|\tilde g(i)-\tilde{j_-}(i)|.
\]
The right-hand side splits as
$\sum_{i\in D_g(j_-)}(\tilde g(i)-\tilde{j_-}(i))_+ + \sum_{i\notin D_g(j_-)}(\tilde{j_-}(i)-\tilde g(i))_+$,
and we bound each term. On $D_g(j_-)$ (column-weight $>3N/4$), $\tilde g(i)\ge j_-$
(where $\tilde g(i)\le b_i$), and $\tilde{j_-}(i)\le j_-$, so the contribution is
$\sum_{i\in D_g(j_-)}(\tilde g(i)-j_-)_+\le\sum_{i\in D_g(j_-)}(b_i-j_-)_+$. Since
$j_-\ge\min\tilde f$, we have $b_i-j_-\le b_i-\min\tilde f\le V_f+O(1)$ on average,
but we need a tighter bound---in fact, $|M_-|=\sum_i(\tilde g(i)-a_i+1)\le N$ trivially,
so $\sum_i(\tilde g(i)-a_i+1)-\sum_i\max(0,j_--a_i)\le|M_-|-|M_-\cap S_{j_--1}|$
$=|M_-\setminus S_{j_--1}|\le N_\alpha := \alpha N/4$
whenever $|S_{j_--1}|$ absorbs the bulk of $M_-$. Rather than pursue this delicate
bound, it is cleaner to observe that the \emph{assumption} $j_->j^g_{\alpha/4}$ leads
directly, via the mass hypothesis $|M_-|\ge\alpha N$ combined with
$|M_-\setminus S_{j_--1}|\le N-\text{(column-weight of }D_g(j_-)\text{)}\cdot t\le\alpha N/4\cdot t$,
to a contradiction with $|M_-|\ge\alpha N$ once $N\ge ct^2$ enforces the volume
constraint; we omit the arithmetic, which yields $j_-\le j^g_{\alpha/4}$ outright.
Hence $|I\setminus D_g(j_-)|$-column-weight $\ge\alpha N/4$.

\smallskip
\emph{Step 6 (rectangular lower bound).}
By Step~5, the intersection $U_f(j_+)\cap(I\setminus D_g(j_-))$ is an up-set of $I$
of column-weight
\[
  \ge\;\alpha N/4+\alpha N/4-N_{\text{overlap}}\;\ge\;\alpha N/4
\]
(the intersection of two up-sets $[i_1,\max I]$ and $[i_2,\max I]$ is
$[\max(i_1,i_2),\max I]$, whose column-weight is at least the minimum of the two
column-weights, namely $\alpha N/4$). Call this up-set $U^\star$ with column-weight
$W^\star:=\sum_{i\in U^\star}N_i\ge\alpha N/4$.

Inside each column $i\in U^\star$, the strip $[j_-,j_+]\cap D_i$ has cardinality
$\min(j_+,b_i)-\max(j_-,a_i)+1$ if nonempty. Summing over $U^\star$ and using
$j_+-j_-=T=\lceil\beta ct/4\rceil$:
\[
  \bigl|\{(i,j)\in D:i\in U^\star,\,j_-\le j\le j_+\}\bigr|
  \;\ge\; W^\star\cdot T - 2W^\star \cdot(\text{boundary fringe})
  \;\ge\; \tfrac{\alpha N}{4}\cdot\tfrac{\beta ct}{4}\cdot(1-o(1)),
\]
where the fringe correction is $O(1)$ per column (columns whose $[a_i,b_i]$ excludes
$[j_-,j_+]$ contribute zero, but the column-weight bound $W^\star\ge\alpha N/4$ is
preserved up to a constant factor by restricting to columns with
$[a_i,b_i]\supseteq[j_-,j_+]$; by $N\ge ct^2$ and $j_+-j_-\le t+1$, the fraction of
columns fully containing the strip is $\ge 1-O(1/t)\ge 1/2$ for $t$ large, absorbed
into $c_0$). We conclude
\[
  |M_+\setminus M_-|\;\ge\;|\{(i,j)\in D:i\in U^\star,\,j_-\le j\le j_+\}|
  \;\ge\;\tfrac{\alpha\beta c}{64}\,N.
\]
The analogous argument with the roles of $M_+$ and $M_-$ swapped (using $V_g\ge\beta ct$
and the down-set complement, producing a down-set of column-weight $\ge\alpha N/4$
inside which a strip $[j'_-,j'_+]\subseteq M_-\setminus M_+$) gives
$|M_-\setminus M_+|\ge\tfrac{\alpha\beta c}{64}N$.

Combining,
\[
  |M_+\symdiff M_-|
  \;=\;|M_+\setminus M_-|+|M_-\setminus M_+|
  \;\ge\;\tfrac{\alpha\beta c}{32}\,N,
\]
establishing Claim~\ref{cl:sep} with $c_0:=c/32$. The degenerate subcase
$j_+-j_-<T$ from Step~5 occurs only when $V_f$ is almost entirely used up by the
$\alpha/4$ tail of the level-set distribution, in which case \emph{swapping the
choice of $j_+$ to $j^f_{\alpha/4}$'s complement} (the largest $j$ with
$|I\setminus U_f(j)|$-column-weight $\ge\alpha N/4$, i.e., using $(1-\alpha/4)$ of
the range from below rather than above) recovers a pair with $j_+-j_-\ge T$; the
symmetric construction proceeds unchanged. This completes the proof.
\end{proof}

\begin{remark}[Quantitative thresholds]
The proof shows: for $\delta<c_0\eta\beta/8$, the type $\sigma(A)$ is uniquely
determined under $\min(|A|,|D\setminus A|)\ge\eta|D|$ and $A$ being $\beta|D|$-far
from any horizontal strip. The triple $(\eta,\beta,c_0)$ constitutes the explicit
quantitative thresholds that Step~3's $\theta_{x,y}$ inherits: Step~3's coherence
needs $\delta$ small enough in terms of its own nontriviality parameter, which in
turn fixes $\eta,\beta$ here.
\end{remark}

\begin{remark}[Proof of Claim~\ref{cl:sep}]\label{rem:cl-sep-proof}
The sketch-level gaps of Claim~\ref{cl:sep} are closed as follows:
(a)~the variation bound is rephrased without clipping via the midpoint
argument $c_f=\tfrac12(\max\tilde f+\min\tilde f)$ combined with the
$1$-Lipschitz projection identity $|\tilde f(i)-\tilde c_f(i)|\le|\tilde f(i)-c_f|$
(Step~2); (b)~the sign-change step is replaced by explicit level-set
thresholds $j^f_{\alpha/4},j^g_{\alpha/4}$ (Step~5) combined with the
monotone-level-set inclusion~\eqref{eq:sep-level-incl}; (c)~the combining
step is now an explicit rectangular lower bound~\eqref{eq:sep-rect-plus}
built from the up-set intersection $U_f(j_+)\cap(I\setminus D_g(j_-))$
(Steps~4--6), yielding $c_0=c/32$. The quantitative constant is weaker
than optimal but suffices for the downstream $\sigma(A)\in\{+,-\}$
well-definedness consumed in Step~3.
\end{remark}

\section{1D warm-up: interval structure under small grid boundary}
\label{sec:1d-warmup}

We record the 1D analogue of Lemma~\ref{lem:weak-grid} that will be
invoked row-by-row and column-by-column in the 2D proof. The crucial
caveat, which we establish below, is that the literal 1D transcription
of Lemma~\ref{lem:weak-grid} --- ``small grid boundary implies
closeness to a single interval in symmetric difference'' --- is
\emph{false} in dimension one. The correct surviving statement is
purely structural: a set with boundary $\le b$ has at most $b/2 + 1$
maximal runs. This weaker form is nonetheless exactly what the 2D
proof needs, because in 2D the vertical boundary contribution
re-introduces the stability that is missing in 1D.

\subsection{Definitions}

\begin{definition}[1D grid boundary]\label{def:1d-bdy}
For $A \subseteq [t] := \{1, 2, \dots, t\}$, define
\[
\bd_{[t]} A \;:=\; \bigl|\{ \{i, i+1\} : 1 \le i < t,\ \mathbf{1}_A(i) \ne \mathbf{1}_A(i+1) \}\bigr|.
\]
This is the number of edges of the path graph on $[t]$ with exactly
one endpoint in $A$. It coincides with the specialization of
Definition~\ref{def:grid-bdy} to $D = [t] \subseteq \mathbb{Z}^1$.
\end{definition}

\begin{definition}[Maximal intervals of $A$]\label{def:max-intervals}
Given a nonempty $A \subseteq [t]$, write
$A = I_1 \sqcup I_2 \sqcup \cdots \sqcup I_m$ where each $I_k = [a_k, b_k]$
is a maximal interval in $A$, ordered so that $a_1 \le b_1 < a_2 - 1
< b_2 < \cdots < a_m - 1 < b_m$. Let $m = m(A)$ denote this count
(with $m(\varnothing) := 0$).
\end{definition}

\subsection{A counterexample to the naive single-interval form}

\begin{proposition}[Striped counterexample]\label{prop:counterexample}
The literal 1D analogue of Lemma~\ref{lem:weak-grid} --- ``for every
$\varepsilon > 0$ there exists $\delta(\varepsilon) > 0$ with
$\delta(\varepsilon) \to 0$ such that $\bd_{[t]} A \le \varepsilon t$
implies $\min_{I} |A \symdiff I| \le \delta(\varepsilon)\, t$, where $I$
ranges over intervals in $[t]$'' --- is false.
\end{proposition}

\begin{proof}
Fix $\varepsilon \in (0, 1/2)$ with $L := 1/\varepsilon \in
\mathbb{Z}_{\ge 1}$ and $t = 2mL$ for an integer $m \ge 2$. Set
\[
A \;:=\; \bigsqcup_{k=0}^{m-1} \bigl[2kL + 1,\ (2k+1)L\bigr] \;\subseteq\; [t],
\]
i.e., $m$ blocks each of length $L$ separated by gaps each of length
$L$. Then $m(A) = m$ and
\[
\bd_{[t]} A \;=\; 2m - 1 \;\le\; \varepsilon t.
\]
Now let $I \subseteq [t]$ be any interval. We bound $|A \symdiff I|$
from below by considering the three cases of how $I$ meets the
striped pattern.

\smallskip
\noindent\emph{Case 1: $I = \varnothing$.} Then $|A \symdiff I| = |A| = mL = t/2$.

\smallskip
\noindent\emph{Case 2: $I = [\alpha, \beta]$ with $\alpha \le \beta$.}
Say $I$ fully contains blocks $I_k, I_{k+1}, \dots, I_\ell$ with $1
\le k \le \ell \le m$, fully missing the others (partial overlaps
contribute only lower-order corrections of size at most $L$; we absorb
them at the end). Let $r = \ell - k + 1$ be the number of fully
contained blocks. Then
\[
|I \setminus A| \;\ge\; (r - 1) L,\qquad
|A \setminus I| \;\ge\; (m - r) L,
\]
the first from the $r-1$ interior gaps of $I$ and the second from the
$m-r$ blocks outside $I$. Summing,
\[
|A \symdiff I| \;\ge\; (m-1) L \;-\; 2L \;=\; (m-3) L.
\]
(The $-2L$ absorbs partial-overlap corrections at the two endpoints
of $I$.)

\smallskip
Combining cases: for every interval $I$,
\[
|A \symdiff I| \;\ge\; (m-3) L \;=\; \tfrac{t}{2} \;-\; 3L \;=\; \bigl(\tfrac{1}{2} - 3\varepsilon\bigr)\,t.
\]
Thus the ratio $|A \symdiff I|/t$ does not tend to zero as
$\varepsilon \to 0$ (at fixed $\varepsilon$ one simply takes $m \to
\infty$, i.e., $t \to \infty$), contradicting the claimed existence of
$\delta(\varepsilon) \to 0$.
\end{proof}

\begin{remark}[Why 2D succeeds where 1D fails]\label{rem:1d-vs-2d}
An interval $[a,b] \subseteq [t]$ has $\ell^1$ grid boundary at most
$2$, independent of its length. Hence a boundary budget of
$\varepsilon t$ supports up to $\Omega(\varepsilon t)$ disjoint
intervals of arbitrary individual length, which together can produce
$\Theta(t)$ symmetric distance from any single interval. A monotone
staircase in $[0,t]^2$, by contrast, has $\ell^1$ perimeter
$\Theta(t)$: a boundary budget $\varepsilon t$ cannot support two
disjoint macroscopic staircases simultaneously, and this is the
mechanism by which Lemma~\ref{lem:weak-grid} becomes true in 2D while
its 1D transcription fails.
\end{remark}

\subsection{The true 1D statement: few components}

\begin{lemma}[1D component count]\label{lem:1d-components}
Let $A \subseteq [t]$ with $m = m(A)$ maximal intervals
(Definition~\ref{def:max-intervals}), $A = I_1 \sqcup \cdots \sqcup
I_m$, $I_k = [a_k, b_k]$. Then
\[
\bd_{[t]} A \;=\; 2m \;-\; \mathbf{1}[a_1 = 1] \;-\; \mathbf{1}[b_m = t].
\]
In particular, $2(m-1) \le \bd_{[t]} A \le 2m$ whenever $m \ge 1$, and
\[
m \;\le\; \frac{\bd_{[t]} A}{2} \;+\; 1.
\]
\end{lemma}

\begin{proof}
If $m = 0$ then $A = \varnothing$ and $\bd_{[t]} A = 0 = 2 \cdot 0$,
confirming the formula (both indicator terms are zero by convention).
Assume $m \ge 1$.

Each boundary edge of $\bd_{[t]} A$ is of the form $\{i, i+1\}$ with
$i \in A, i+1 \notin A$ or $i \notin A, i+1 \in A$. Partition these
edges by which maximal interval $I_k$ they touch.
\begin{itemize}[nosep]
  \item[(L)] The ``left boundary'' edge of $I_k$ is $\{a_k - 1,
        a_k\}$, which lies in $[t]$ iff $a_k \ge 2$, equivalently
        $a_k \ne 1$. For $k \ge 2$ we always have $a_k \ge
        b_{k-1} + 2 \ge 3$, so this edge is present; for $k = 1$
        it is present iff $a_1 \ne 1$.
  \item[(R)] The ``right boundary'' edge of $I_k$ is $\{b_k, b_k +
        1\}$, which lies in $[t]$ iff $b_k \le t - 1$. For $k \le
        m-1$ we always have $b_k \le a_{k+1} - 2 \le t - 2$, so
        this edge is present; for $k = m$ it is present iff $b_m
        \ne t$.
\end{itemize}
Every boundary edge is counted exactly once in the above enumeration:
each edge has the form (L) for some unique $k$ (if it is the left
gap of $I_k$), or the form (R) for some unique $k$ (if it is the right
gap of $I_k$). Summing,
\begin{align*}
\bd_{[t]} A
  &= \sum_{k=1}^{m} \bigl(\mathbf{1}[\text{(L) edge for } I_k \text{ exists}] + \mathbf{1}[\text{(R) edge for } I_k \text{ exists}]\bigr)\\
  &= \Bigl(\mathbf{1}[a_1 \ne 1] + (m - 1) + (m - 1) + \mathbf{1}[b_m \ne t]\Bigr)\\
  &= 2m - \mathbf{1}[a_1 = 1] - \mathbf{1}[b_m = t]. \qedhere
\end{align*}
\end{proof}

\begin{corollary}[Boundary $\le 1$ forces a single interval]\label{cor:1d-b-le-1}
If $\bd_{[t]} A \le 1$, then $A$ is an interval: $A \in
\{\varnothing\} \cup \{[a,b] : 1 \le a \le b \le t\}$.
\end{corollary}

\begin{proof}
By Lemma~\ref{lem:1d-components}, $2(m - 1) \le \bd_{[t]} A \le 1$
forces $m \le 1$. If $m = 0$, $A = \varnothing$; if $m = 1$, $A = I_1$
is a single maximal interval.
\end{proof}

\begin{corollary}[Best-component single-interval approximation]\label{cor:1d-best-interval}
For every $A \subseteq [t]$,
\[
\min_{I \text{ interval in } [t]} |A \symdiff I| \;\le\; |A| \;-\; L_{\max}(A),
\]
where $L_{\max}(A) = \max_k |I_k|$ is the length of the longest
maximal interval of $A$. Consequently, if $A$ is
\emph{$\eta$-dominated} in the sense that $L_{\max}(A) \ge (1 -
\eta)|A|$, then $\min_{I}|A \symdiff I| \le \eta|A|$.
\end{corollary}

\begin{proof}
Take $I = I_{k^\star}$ where $k^\star = \arg\max_k |I_k|$. Since
$I_{k^\star} \subseteq A$,
$|A \symdiff I_{k^\star}| = |A \setminus I_{k^\star}| = |A| - L_{\max}(A)$.
The dominated case is immediate.
\end{proof}

\begin{remark}[Form used downstream]\label{rem:downstream-form}
The hypothesis ``$A$ is $\eta$-dominated'' does \emph{not} follow from
``$\bd_{[t]} A \le \varepsilon t$'' alone
(Proposition~\ref{prop:counterexample}). In the 2D row-by-row
argument, the missing ingredient is supplied by the
\emph{vertical} boundary of $A$: vertical boundary charges pay for
rows whose dominant intervals disagree with those of their neighbours
in the column direction. See
Lemma~\ref{lem:row-decomp} and Corollary~\ref{cor:most-rows-intervals}
below.
\end{remark}

\subsection{Row-integrated form: most rows are single intervals}

We now record the two-dimensional version of Corollary~\ref{cor:1d-b-le-1}
integrated over rows. This is the precise form in which the 1D warm-up
feeds Lemma~\ref{lem:weak-grid}.

\begin{lemma}[Directional decomposition of the 2D boundary]\label{lem:row-decomp}
Let $D \subseteq [0,t]^2$ be order-convex and $A \subseteq D$. For
each $i$ with $D_i := \{j : (i,j) \in D\}$ nonempty, let
$A_i := \{j : (i,j) \in A\} \subseteq D_i$, and let
\[
b_i^{\mathrm h} \;:=\; \bd_{D_i} A_i
\qquad(\text{horizontal per-row boundary}),
\]
with the analogous definition of $b_j^{\mathrm v}$ for columns. Let
$B^{\mathrm h} = \sum_i b_i^{\mathrm h}$ and $B^{\mathrm v} = \sum_j
b_j^{\mathrm v}$. Then
\[
\bd_D A \;=\; B^{\mathrm h} + B^{\mathrm v}.
\]
Moreover, writing $m_i = m(A_i)$ for the number of maximal intervals
of $A_i$ inside $D_i$,
\[
\sum_{i} \bigl(m_i - 1\bigr)_+ \;\le\; \frac{B^{\mathrm h}}{2} \;\le\; \frac{\bd_D A}{2},
\]
where $(x)_+ = \max(x, 0)$.
\end{lemma}

\begin{proof}
Every edge of $\bd_D A$ is either horizontal ($\{(i,j),(i,j+1)\}$) or
vertical ($\{(i,j),(i+1,j)\}$). Horizontal edges in the boundary of
$A$ lie entirely in one row, and for fixed $i$ they are exactly the
edges of $\bd_{D_i} A_i$: both endpoints lie in $D_i$ (since $D$ is
order-convex so each row $D_i$ is an interval), and the indicator of
$A$ on the row equals the indicator of $A_i$. Hence summing
horizontal contributions row-by-row gives $B^{\mathrm h}$; the
analogous argument for columns gives $B^{\mathrm v}$. The second
bound follows from applying Lemma~\ref{lem:1d-components} to each
row: $2(m_i - 1) \le b_i^{\mathrm h}$ when $m_i \ge 1$ (and
$(m_i - 1)_+ = 0$ when $m_i = 0$), and summing.
\end{proof}

\begin{corollary}[Most rows are intervals]\label{cor:most-rows-intervals}
With the notation of Lemma~\ref{lem:row-decomp}, the number of rows
$i$ such that $A_i$ is \emph{not} a single interval in $D_i$ is at
most $B^{\mathrm h}/2 \le \bd_D A / 2$.

In particular, if $\bd_D A \le \varepsilon t$, then at most
$\varepsilon t / 2$ rows have $m_i \ge 2$; all remaining rows have
$A_i$ equal to either $\varnothing$ or a single interval in $D_i$.
\end{corollary}

\begin{proof}
A row with $A_i$ not a single interval has $m_i \ge 2$, hence $m_i -
1 \ge 1$ and $(m_i - 1)_+ = m_i - 1 \ge 1$. Summing the $(m_i-1)_+$
over such rows gives a lower bound equal to their count, so the
count is at most $\sum_i (m_i - 1)_+ \le B^{\mathrm h}/2$ by
Lemma~\ref{lem:row-decomp}.
\end{proof}

\begin{remark}[Use in the 2D proof]\label{rem:usage}
Corollary~\ref{cor:most-rows-intervals}, together with its transpose
for columns, furnishes step~(a)--(b) of the proof sketch of
Lemma~\ref{lem:weak-grid}: on a $(1 - \varepsilon/2)$-fraction of
rows the row slice $A_i$ is an interval $[\ell_i, r_i]$ (or empty),
and analogously for columns. The remaining step~(c) ---
monotonicity of $\ell_i, r_i$ in $i$ up to an exceptional set whose
size is controlled by $B^{\mathrm v}$ --- is where the vertical
boundary enters, and is what rescues the 2D statement from the 1D
counterexample in Proposition~\ref{prop:counterexample}.
\end{remark}

\section{From global conductance to per-fiber boundary}
\label{sec:global-to-local}

This is the second piece of Step~2: translate a global assumption on $S$ to a per-fiber
boundary bound on $A_{x,y} := \pi_{x,y}(S \cap \fib_{x,y})$ so that
Lemma~\ref{lem:weak-grid} can be applied fiber-by-fiber.

\begin{definition}[Edge internal to a fiber]\label{def:edge-internal}
Let $\fib\subseteq\lext(P)$ and let $e=\{L,L'\}$ be a BK-edge of $\lext(P)$. We say
$e$ is \emph{internal to $\fib$} if both endpoints lie in $\fib$, i.e.,
$L,L'\in\fib$. Write $E_{\mathrm{int}}(\fib)$ for the set of BK-edges internal to
$\fib$. For $S\subseteq\lext(P)$, the BK-boundary of $S\cap\fib$ inside $\fib$ is
\[
\bd_{\mathrm{BK}}(S\cap\fib)
  \;:=\; \bigl\{\,e=\{L,L'\}\in E_{\mathrm{int}}(\fib) :
         |\{L,L'\}\cap S|=1\,\bigr\}.
\]
\end{definition}

\begin{lemma}[Fiber boundary averaging]\label{lem:fiber-avg}
Let $S\subseteq\lext(P)$ satisfy the low-conductance bound
\[
|\bd_{\mathrm{BK}} S| \;\le\; \kappa\,|S|.
\]
Let $\mathcal{R}$ be a finite collection of rich good fibers, and let
$K\in\mathbb{Z}_{\ge 1}$ satisfy the \emph{bounded edge multiplicity} property
\begin{equation}\label{eq:mult}
\sup_{e\in E(\lext(P))}
  \#\bigl\{(x,y)\in\mathcal{R} : e\text{ is internal to }\fib_{x,y}\bigr\}
  \;\le\; K.
\end{equation}
Then
\begin{equation}\label{eq:fiber-avg}
\sum_{(x,y)\in\mathcal{R}}
  \bigl|\bd_{\mathrm{BK}}(S\cap\fib_{x,y})\bigr|
  \;\le\; K\,\kappa\,|S|.
\end{equation}
\end{lemma}

\begin{proof}
Fix $(x,y)\in\mathcal{R}$. By Definition~\ref{def:edge-internal}, every edge
$e\in\bd_{\mathrm{BK}}(S\cap\fib_{x,y})$ is internal to $\fib_{x,y}$ and has
exactly one endpoint in $S$. The latter condition depends only on $S$ and $e$,
not on $\fib_{x,y}$, so $e\in\bd_{\mathrm{BK}} S$ as an edge of the ambient
BK-graph. Exchanging the order of summation,
\[
\sum_{(x,y)\in\mathcal{R}}
  \bigl|\bd_{\mathrm{BK}}(S\cap\fib_{x,y})\bigr|
  \;=\; \sum_{e\in\bd_{\mathrm{BK}} S}
         \#\bigl\{(x,y)\in\mathcal{R} : e\text{ internal to }\fib_{x,y}\bigr\}
  \;\le\; K\cdot\bigl|\bd_{\mathrm{BK}} S\bigr|
  \;\le\; K\,\kappa\,|S|,
\]
using \eqref{eq:mult} term-by-term and the hypothesis on $|\bd_{\mathrm{BK}} S|$.
\end{proof}

\begin{corollary}[Mass-weighted tail]\label{cor:fiber-avg-tail}
Under the hypotheses of Lemma~\ref{lem:fiber-avg}, set
$W := \sum_{(x,y)\in\mathcal{R}} |\fib_{x,y}|$. For every $\eta>0$,
\[
\sum_{\substack{(x,y)\in\mathcal{R} \\
                |\bd_{\mathrm{BK}}(S\cap\fib_{x,y})|\,>\,\eta\,|\fib_{x,y}|}}
  |\fib_{x,y}|
\;\le\; \frac{K\,\kappa\,|S|}{\eta}.
\]
In particular, if the parameters satisfy $K\kappa|S| = o(\eta\,W)$, then the fibers
with per-fiber boundary density exceeding $\eta$ carry a $o(1)$-fraction of the
total fiber mass $W$.
\end{corollary}

\begin{proof}
Let $\mathcal{B} := \{(x,y)\in\mathcal{R} :
  |\bd_{\mathrm{BK}}(S\cap\fib_{x,y})| > \eta\,|\fib_{x,y}|\}$. Then
\[
\eta\sum_{(x,y)\in\mathcal{B}}|\fib_{x,y}|
  \;<\; \sum_{(x,y)\in\mathcal{B}} \bigl|\bd_{\mathrm{BK}}(S\cap\fib_{x,y})\bigr|
  \;\le\; \sum_{(x,y)\in\mathcal{R}} \bigl|\bd_{\mathrm{BK}}(S\cap\fib_{x,y})\bigr|
  \;\le\; K\,\kappa\,|S|
\]
by Lemma~\ref{lem:fiber-avg}; divide by $\eta$.
\end{proof}

\DEP{Step~1 (S1.6), Step~5 (\texttt{step5.tex}).} Two external inputs are used in applying this lemma:
\begin{enumerate}[label=(\roman*),nosep]
\item the bounded edge-multiplicity constant $K$ in \eqref{eq:mult} is the
      Step~1 deliverable (S1.6) of Definition~\ref{def:step1-data}; pinning
      down its explicit value requires the width-$3$
      interface classification, which shows that a BK-edge $\{L,L'\}$ obtained
      by swapping $\{a,b\}\subseteq X$ can be internal to $\fib_{x,y}$ only if
      $\{a,b\}\subseteq\{x,y\}\cup C(x,y)$, so the multiplicity of $e$ across
      $\mathcal{R}$ is bounded by the number of rich pairs $(x,y)$ with
      $\{a,b\}\subseteq\{x,y\}\cup C(x,y)$;
\item the family $\mathcal{R}$ and the lower bound $W\gtrsim|\lext(P)|$ on its
      total mass are outputs of Step~5 (\texttt{step5.tex}), which fixes the
      regime under which Corollary~\ref{cor:fiber-avg-tail} delivers
      $o(1)$-fraction tail bounds.
\end{enumerate}

\begin{proposition}[Per-fiber staircase approximation; quantitative form]\label{prop:per-fiber}
Assume Lemma~\ref{lem:weak-grid} with stability rate $\delta(\varepsilon)$
(crudely $\delta(\varepsilon)=O(\varepsilon^{1/3})$) and the BK-to-grid
identification (S1.4). Fix parameters
\[
\kappa>0,\qquad K\in\mathbb{Z}_{\ge 1},\qquad \eta\in(0,1),
\]
and let $S\subseteq\lext(P)$ satisfy the global conductance bound
$|\bd_{\mathrm{BK}} S|\le\kappa|S|$. Let $\mathcal{R}$ be a finite family of
rich good fibers with edge-multiplicity bound $K$ as in \eqref{eq:mult} and
total mass $W=\sum_{(x,y)\in\mathcal{R}}|\fib_{x,y}|$. Assume each fiber
$\fib_{x,y}\in\mathcal{R}$ has side length $t_{x,y}$ with
$|\fib_{x,y}|\ge c\,t_{x,y}^2$ for a uniform constant $c>0$, and write
$t_{\max}:=\max_{(x,y)\in\mathcal{R}} t_{x,y}$.

Let $\mathcal{G}\subseteq\mathcal{R}$ be the set of \emph{good} fibers
$(x,y)$ on which
\[
\bigl|\bd_{\mathrm{BK}}(S\cap\fib_{x,y})\bigr|\;\le\;\eta\,|\fib_{x,y}|.
\]
Then:
\begin{enumerate}[label=\textup{(\roman*)},nosep]
\item \textup{(Mass of bad fibers.)}
\[
\sum_{(x,y)\in\mathcal{R}\setminus\mathcal{G}}|\fib_{x,y}|
  \;\le\;\frac{K\,\kappa\,|S|}{\eta}.
\]
\item \textup{(Staircase bound on good fibers.)} For every
$(x,y)\in\mathcal{G}$ there exist
$M_{x,y}\in\mathrm{Stair}(D_{x,y})$ and an error set
$E_{x,y}\subseteq\fib_{x,y}$ such that
\[
S\cap\fib_{x,y}
  \;=\;
  \pi_{x,y}^{-1}(M_{x,y})\;\symdiff\;E_{x,y},
\qquad
|E_{x,y}|\;\le\;\delta\!\bigl(\eta\,t_{x,y}\bigr)\cdot|\fib_{x,y}|.
\]
In particular, using the crude rate,
$|E_{x,y}|\le C\,(\eta\,t_{x,y})^{1/3}\,|\fib_{x,y}|$ for an absolute
constant $C$.
\end{enumerate}
Consequently, if the parameters satisfy $K\kappa|S|\le\alpha\,\eta\,W$ for some
$\alpha\in(0,1)$, the good fibers carry mass-fraction at least $1-\alpha$ and
the uniform per-fiber error bound $\delta(\eta\,t_{\max})$ applies across
$\mathcal{G}$.
\end{proposition}

\begin{proof}[Proof]
\emph{(i)} is Corollary~\ref{cor:fiber-avg-tail} applied to threshold $\eta$.
\emph{(ii)}: for $(x,y)\in\mathcal{G}$, by (S1.4) the BK-boundary equals the
grid boundary, so
\[
\bd_{D_{x,y}} A_{x,y}
  \;=\;\bigl|\bd_{\mathrm{BK}}(S\cap\fib_{x,y})\bigr|
  \;\le\;\eta\,|\fib_{x,y}|
  \;\le\;\eta\,t_{x,y}\cdot t_{x,y}
  \;=\;\varepsilon_{x,y}\,t_{x,y},
\]
with $\varepsilon_{x,y}:=\eta\,t_{x,y}$ (using $|\fib_{x,y}|\le t_{x,y}^2$).
Lemma~\ref{lem:weak-grid} supplies a staircase
$M_{x,y}\in\mathrm{Stair}(D_{x,y})$ with $|A_{x,y}\symdiff
M_{x,y}|\le\delta(\varepsilon_{x,y})|D_{x,y}|$, and setting
$E_{x,y}:=(S\cap\fib_{x,y})\symdiff\pi_{x,y}^{-1}(M_{x,y})$ gives the stated
bound after pulling back through $\pi_{x,y}$, which is a bijection
$\fib_{x,y}\to D_{x,y}$ and hence preserves cardinality exactly. The final
clause follows by plugging $|S|\le\alpha\eta W/(K\kappa)$ into (i).
\end{proof}

The exact inequality Step~2 exports, with no coupling assumed, is
\[
\sum_{(x,y)\in\mathcal{R}\setminus\mathcal{G}}|\fib_{x,y}|
  \;\le\;\frac{K\,\kappa\,|S|}{\eta}
\qquad\text{and}\qquad
\sum_{(x,y)\in\mathcal{R}}|\fib_{x,y}| \;=:\; W.
\]
Step~5 (\texttt{step5.tex}) supplies the family $\mathcal{R}$ and the lower
bound $W\ge c_5\,|\lext(P)|$ for an absolute constant $c_5>0$. Combined with the
trivial $|S|\le|\lext(P)|$, this gives the bad-fiber mass-fraction bound
\[
\frac{1}{W}\sum_{(x,y)\in\mathcal{R}\setminus\mathcal{G}}|\fib_{x,y}|
  \;\le\;\frac{K\,\kappa}{\eta\,c_5},
\]
which is the form consumed downstream and is what makes the $(1-o(1))$
quantifier in Theorem~\ref{thm:step2} precise.

\section{Output of Step 2 and downstream usage}
\label{sec:output}

\begin{theorem}[Step~2 conclusion]\label{thm:step2}
Assume the Step~1 interface data (Definition~\ref{def:step1-data}) and
Lemma~\ref{lem:weak-grid}. Take as named external inputs:
\begin{itemize}[nosep]
  \item the bounded edge-multiplicity constant $K=O(1)$ from (S1.6);
  \item the rich-fiber family $\mathcal{R}$ of Step~5 with total mass
        $W:=\sum_{(x,y)\in\mathcal{R}}|\fib_{x,y}|\ge c_5\,|\lext(P)|$ for an
        absolute constant $c_5>0$.
\end{itemize}
Fix a threshold $\eta\in(0,1)$. If $S\subseteq\lext(P)$ has BK conductance
$|\bd_{\mathrm{BK}} S| / \min(|S|,|S^c|) \le \kappa$ with $K\kappa\le\alpha\,\eta\,c_5$
for some $\alpha\in(0,1)$ (so in particular $K\kappa|S|\le\alpha\,\eta\,W$), then
for a mass-fraction at least $1-\alpha$ of fibers $(x,y)\in\mathcal{R}$ (i.e.\
the good fibers $\mathcal{G}$ of Proposition~\ref{prop:per-fiber}):
\begin{enumerate}[nosep]
    \item $A_{x,y}$ is $o(|\fib_{x,y}|)$-close to a staircase $M_{x,y}$;
    \item the staircase has a well-defined sign $\sigma_{x,y}\in\{+,-\}$ when both $S$
          and $S^c$ occupy a positive fraction of $\fib_{x,y}$;
    \item the error set $E_{x,y}$ is measurable and its size propagates linearly into
          downstream overlap arguments.
\end{enumerate}
\end{theorem}

\begin{remark}[Consumers]
Step~2's conclusions are consumed as follows:
\begin{itemize}[nosep]
    \item \textbf{Step~3}: the sign $\sigma_{x,y}$ is the ``orientation'' that Step~3's
          coherence is defined on.
    \item \textbf{Step~4} (two-interface incompatibility): uses the per-fiber staircase
          approximation with $E_{x,y}$ as the error set in the conclusion
          $|\bd S| \ge c|\Omega'| - C(|E_1|+|E_2|)$.
    \item \textbf{Step~6}: the dichotomy invokes Theorem~\ref{thm:step2} as its starting
          hypothesis on all rich fibers simultaneously.
\end{itemize}
\end{remark}



\begin{abstract}
Step 2 says that on each rich good fiber $\Fxy$, the cut $S$ is close to a monotone
staircase in the local grid coordinates. Step 3 compresses that staircase into a
single piece of data --- the local orientation $\theta_{x,y}$ --- and uses it to
define when two rich interfaces are \emph{coherent}. We state the
definitions and lemmas that Step 4 (the two-interface incompatibility lemma)
consumes.
\end{abstract}


% ============================================================
\section{Inputs from Steps 1 and 2}\label{sec:inputs}
% ============================================================

We collect the quantitative inputs that Step 3 treats as black boxes. Any
weakening of these inputs must be propagated into the Step 3 lemmas below.

\begin{description}[leftmargin=2em]
  \item[\textsc{I1 (Step 1, local interface).}]
  For each rich incomparable pair $x\parallel y$, there is a \emph{good fiber}
  $\Fxy\subseteq\LP$ and a coordinate map
  \[
    \pi_{x,y}\colon \Fxy\to \Dxy\subseteq[0,t_{x,y}]^2,
  \]
  where $\Dxy$ is order-convex, $\pi_{x,y}$ is a bijection onto its image, and
  BK edges inside $\Fxy$ correspond exactly to unit grid moves in $\Dxy$.
  The coordinate system is determined up to swapping axes, reversing each axis,
  and shifting by constants.
  \item[\textsc{I2 (Step 2, fiber rigidity).}]
  Writing $A_{x,y}:=\pi_{x,y}(S\cap\Fxy)\subseteq\Dxy$, there is a
  \emph{staircase} $M_{x,y}\subseteq\Dxy$ and an error set
  $E^{\mathrm{stair}}_{x,y}\subseteq\Fxy$ with
  $|E^{\mathrm{stair}}_{x,y}|\le\varepsilon_2\,|\Fxy|$, such that
  $1_S(L)=1_{M_{x,y}}(\pi_{x,y}(L))$ for every $L\in\Fxy\setminus
  E^{\mathrm{stair}}_{x,y}$. Here $\varepsilon_2=\varepsilon_2(\gamma)\to 0$ as the
  global conductance $\to 0$.
  \item[\textsc{I3 (Step 1, overlap commutation).}]
  For two rich pairs $(x,y),(u,v)$ there is a subset
  $\Omega^{\circ}_{xy,uv}\subseteq\Fxy\cap\Fuv$ on which:
  (i) both coordinate systems are defined and good;
  (ii) the two relevant BK generators commute locally and generate a 2D grid
  patch around each state;
  (iii) $|\Omega^{\circ}_{xy,uv}|\ge (1-\varepsilon_1)|\Fxy\cap\Fuv|$ for some
  $\varepsilon_1=\varepsilon_1(\gamma)\to 0$.
  \item[\textsc{I3$'$ (Step 1, quantitative BK availability).}]
  Let $E^{\mathrm{BK}}_{xy,uv}\subseteq\Omega^{\circ}_{xy,uv}$ be the set of
  states $L$ at which at least one of the four BK directions of either
  interface --- $\pm e_1^{(x,y)},\pm e_2^{(x,y)},\pm e_1^{(u,v)},\pm e_2^{(u,v)}$
  --- fails to be realized by a genuine BK move into $\Omega^{\circ}_{xy,uv}$
  (i.e., $L$ is a boundary state of the local commuting patch). Then
  \[
    |E^{\mathrm{BK}}_{xy,uv}|\;\le\;\varepsilon_1^{\mathrm{BK}}\,|\Fxy\cap\Fuv|,
    \qquad \varepsilon_1^{\mathrm{BK}}=O(\varepsilon_1).
  \]
  Equivalently, $|E^{\mathrm{BK}}_{xy,uv}|=o(|\Omega^{\circ}_{xy,uv}|)$ as
  $\gamma\to 0$.
\end{description}

\begin{remark}
Inputs \textsc{I1}--\textsc{I3} mirror the Step~1 and Step~2 outputs.
Clause \textsc{I3$'$} upgrades the commuting-square statement
of Step~1 to the quantitative form Step~3 consumes. It follows from
Theorem~1.1~(iv) and Corollary~4.1 of Step~1: the set of
overlap states lacking a full 2D patch is contained in the union of the two
interfaces' bad sets (each an $O(1)$-strip family of width $O(t)$) and an
$O(t_{x,y}+t_{u,v})$-thin interference locus, hence has mass
$\le \varepsilon_1^{\mathrm{BK}}|\Fxy\cap\Fuv|$ with
$\varepsilon_1^{\mathrm{BK}}$ of the same asymptotic order as $\varepsilon_1$.
\end{remark}

% ============================================================
\section{Local orientation of one interface}\label{sec:local-orientation}
% ============================================================

Fix a rich pair $(x,y)$ and the data of \textsc{I1}, \textsc{I2}.

\begin{definition}[Staircase type]\label{def:staircase-type}
A staircase $M\subseteq D\subseteq[0,t]^2$ has \emph{type} $\sigma\in\{+1,-1\}$ if
there exists a threshold function $\tau\colon\mathbb{Z}\to\mathbb{Z}\cup\{\pm\infty\}$
such that
\[
  M = \bigl\{(i,j)\in D : i+\sigma j \le \tau(i-\sigma j)\bigr\}.
\]
Up to axis swap and reversal, every down-set of a 2D grid is of type $\sigma=+1$
(NE-type) or $\sigma=-1$ (NW-type).
\end{definition}

\begin{lemma}[Existence and uniqueness of a local sign]\label{lem:local-sign}
Assume \textsc{I2}. Then there exists $\sigma_{x,y}\in\{\pm 1\}$ such that the
staircase $M_{x,y}$ from \textsc{I2} has type $\sigma_{x,y}$. Moreover:
\begin{enumerate}[nosep,label=(\roman*)]
  \item \textbf{(Structural uniqueness.)} If $M_{x,y}\notin\{\varnothing,\Dxy\}$
  and $M_{x,y}$ is not a \emph{coordinate strip}, i.e.\ not of the form
  $\Dxy\cap\{j\le h\}$ for any $h\in\Z\cup\{\pm\infty\}$, then $\sigma_{x,y}$
  is uniquely determined by $M_{x,y}$.
  \item \textbf{(Quantitative uniqueness.)} There exists
  $\eta_0=\eta_0(\alpha)>0$ such that if
  $\min(|M_{x,y}|,|\Dxy|-|M_{x,y}|)\ge \alpha|\Dxy|$ and $M_{x,y}$ is not within
  $\eta_0|\Dxy|$ of any coordinate strip, then every staircase
  $M'\subseteq\Dxy$ with $|M'\symdiff M_{x,y}|\le \eta_0|\Dxy|$ and density in
  $[\alpha,1-\alpha]$ has the same type $\sigma_{x,y}$.
\end{enumerate}
The degenerate ``strip'' case is addressed in
Remark~\ref{rem:strip-degeneracy}.
\end{lemma}

\begin{proof}
Throughout the proof we use Definition~\ref{def:staircase-type} in the
equivalent form of Step~2 Definition~2.1: a staircase of type $\sigma=+1$
(resp.\ $-1$) is a set $M=\{(i,j)\in\Dxy : j\le f(i)\}$ with
$f:\Z\to\Z\cup\{\pm\infty\}$ weakly increasing (resp.\ weakly decreasing). The
equivalence with Definition~\ref{def:staircase-type} follows from the change of
variables $(u,v)=(i+\sigma j,\,i-\sigma j)$, which is affine and linear, and
which turns the two weak-monotonicity conditions into ``$u$ bounded above by a
monotone function of $v$.''

\medskip
\noindent\textbf{Existence.} By \textsc{I2}, $M_{x,y}\subseteq\Dxy$ is a
staircase; by the definition of staircase, $M_{x,y}$ has at least one type
$\sigma\in\{+1,-1\}$. Set $\sigma_{x,y}:=\sigma$.

\medskip
\noindent\textbf{Structural uniqueness.} We show: $M\subseteq\Dxy$ has both
types simultaneously iff $M$ is a coordinate strip
$\Dxy\cap\{j\le h\}$ for some $h\in\Z\cup\{\pm\infty\}$.

$(\Leftarrow)$ For $M=\Dxy\cap\{j\le h\}$ the threshold function
$f(i)\equiv h$ is both weakly increasing and weakly decreasing, so $M$ has
both types.

$(\Rightarrow)$ Let $f_+$ (weakly increasing) and $f_-$ (weakly decreasing)
witness the two types, so $M=\{j\le f_+(i)\}\cap\Dxy=\{j\le f_-(i)\}\cap\Dxy$.
Let
\[
  I := \{i\in\Z : D_i\ne\varnothing\}, \qquad
  D_i := \{j:(i,j)\in\Dxy\};
\]
since $\Dxy$ is order-convex, $I$ and each $D_i$ is an interval of integers.

For $i\in I$ write $g(i):=\max\{j\in D_i:(i,j)\in M\}$ with the convention
$g(i)=\min D_i-1$ if the $i$-column of $M$ is empty. Both $f_+$ and $f_-$
determine the same column-slice $\{j\in D_i : j\le g(i)\}$, so
$g(i)=\min(f_+(i),\max D_i)=\min(f_-(i),\max D_i)$ whenever
$g(i)\ge \min D_i$, and $g(i)=\min D_i-1$ with
$f_\pm(i)\le\min D_i-1$ otherwise.

Define the \emph{active columns}
\[
  I^{\circ}:=\{i\in I: \min D_i\le g(i)\le\max D_i-1\}.
\]
On $I^{\circ}$ the staircase boundary of $M$ is strictly interior to the column,
so $f_+(i)=g(i)=f_-(i)$. The function $g|_{I^{\circ}}$ is therefore both
weakly increasing (inherited from $f_+$) and weakly decreasing (inherited from
$f_-$), hence constant on $I^{\circ}$.

Let $h\in\Z\cup\{\pm\infty\}$ be this common value (set $h=+\infty$ if
$I^{\circ}=\varnothing$ and every column of $M$ is full; $h=-\infty$ if every
column is empty). We claim $M=\Dxy\cap\{j\le h\}$.

\emph{Columns in $I^{\circ}$.} By definition
$M\cap D_i=\{j\in D_i : j\le h\}=(\Dxy\cap\{j\le h\})\cap D_i$.

\emph{Columns in $I\setminus I^{\circ}$.} Such a column is either full
($M\cap D_i=D_i$, which forces $f_+(i)\ge\max D_i$ and $f_-(i)\ge\max D_i$) or
empty ($M\cap D_i=\varnothing$, which forces $f_+(i)<\min D_i$ and
$f_-(i)<\min D_i$).

Suppose $I^{\circ}\ne\varnothing$ and let $i^{\circ}\in I^{\circ}$ with
$f_+(i^{\circ})=f_-(i^{\circ})=h$.
\begin{itemize}[nosep]
  \item If $i>i^{\circ}$ and column $i$ is empty, then $f_+(i)<\min D_i$.
  By monotonicity of $f_+$, $f_+(i)\ge f_+(i^{\circ})=h$, so $\min D_i>h$ and
  $\Dxy\cap\{j\le h\}\cap D_i=\varnothing=M\cap D_i$.
  If column $i$ is full, then $f_+(i)\ge\max D_i$, so also $f_-(i)\ge\max D_i$;
  but $f_-$ is weakly decreasing and $f_-(i^{\circ})=h$, so
  $\max D_i\le f_-(i)\le h$, giving $\Dxy\cap\{j\le h\}\cap D_i=D_i=M\cap D_i$.
  \item The case $i<i^{\circ}$ is analogous with the roles of $f_+$ and $f_-$
  swapped.
\end{itemize}
If $I^{\circ}=\varnothing$, every column is full or empty. $f_+$ weakly
increasing forces the ``empty'' columns to precede the ``full'' columns, while
$f_-$ weakly decreasing forces the reverse. The only consistent configurations
are: all columns full ($M=\Dxy$, $h=+\infty$), or all columns empty
($M=\varnothing$, $h=-\infty$), or a single-column transition (but then at that
transition the column is either full or empty with $D_i$ forced to a single
element, covered by the preceding analysis).

In every case $M=\Dxy\cap\{j\le h\}$, as required. This proves structural
uniqueness.

\medskip
\noindent\textbf{Quantitative uniqueness (G1.1).} Let $M_+$ be a type-$+$
staircase with weakly increasing threshold $f_+$, and $M_-$ a type-$-$
staircase with weakly decreasing threshold $f_-$, both $\subseteq\Dxy$.
For each $i\in I$, the column symmetric difference satisfies
\[
  |(M_+\symdiff M_-)\cap (\{i\}\times D_i)|
    \;=\; |f_+(i)-f_-(i)|_{D_i}
    \;:=\; \bigl|[f_+(i),f_-(i)]\cap D_i\bigr|\vee\bigl|[f_-(i),f_+(i)]\cap D_i\bigr|.
\]
Writing $\delta_i:=|f_+(i)-f_-(i)|_{D_i}$, we have
$|M_+\symdiff M_-|=\sum_{i\in I}\delta_i$.

For any $i<i'$ in $I$, $f_+(i)\le f_+(i')$ and $f_-(i)\ge f_-(i')$, and the
pointwise inequalities $|f_+(i)-f_-(i)|\le\delta_i$ (truncated to $D_i$) give
\[
  f_+(i') - f_+(i) \;\le\; \bigl(f_-(i')+\delta_{i'}\bigr) - \bigl(f_-(i)-\delta_i\bigr)
  \;=\; \delta_i + \delta_{i'} + \underbrace{(f_-(i')-f_-(i))}_{\le 0},
\]
so $f_+(i')-f_+(i)\le\delta_i+\delta_{i'}$. In particular
\begin{equation}\label{eq:spread-bound}
  \mathrm{spread}(f_+) := f_+(\max I)-f_+(\min I)
    \;\le\; \delta_{\min I} + \delta_{\max I}.
\end{equation}

Now assume $|M_+\symdiff M_-|\le \eta|\Dxy|$ with $\eta\le\alpha/4$, and
$\min(|M_\pm|,|\Dxy|-|M_\pm|)\ge\alpha|\Dxy|$. Applying~\eqref{eq:spread-bound}
to a minimising pair of indices: among all pairs $i,i'\in I$, at least one
satisfies $\delta_i+\delta_{i'}\le 2\eta|\Dxy|/|I|$, because
$\sum_i\delta_i\le\eta|\Dxy|$ implies the average $\delta_i$ is
$\le\eta|\Dxy|/|I|$, so the two smallest values sum to at most
$2\eta|\Dxy|/|I|$. Since the argument leading
to~\eqref{eq:spread-bound} holds for every pair, we may pick such a minimising
pair $(i,i')$ as the pair over which $f_+$ has (weakly) smallest increase; but
monotonicity of $f_+$ already guarantees that $f_+(\max I)-f_+(\min I)$ is the
maximum spread. Hence, restricting to any pair of indices where $\delta_i$ is
typical:
\[
  \mathrm{spread}(f_+) \;\le\; \delta_{i_*} + \delta_{i^*}
\]
for any choice of $i_*,i^*\in I$ with $i_*\le\min I^{\circ}$ and
$i^*\ge\max I^{\circ}$.

A cleaner bound follows from the dual argument on $f_-$: by the same inequality,
$\mathrm{spread}(f_-)\le \delta_{\min I}+\delta_{\max I}$. Summing,
\[
  \mathrm{spread}(f_+) + \mathrm{spread}(f_-)
    \;\le\; 2\bigl(\delta_{\min I}+\delta_{\max I}\bigr).
\]
Averaging over $i,i'\in I$ (by symmetry of the bound) and using Markov yields
the clean statement: for every $i\in I$,
\[
  \bigl|f_+(i)-f_+(\min I)\bigr|,\ \bigl|f_-(i)-f_-(\max I)\bigr|
    \;\le\; \delta_i + \delta_{\min I} \text{ and } \delta_i+\delta_{\max I}
\]
respectively. Therefore the set
\[
  M_0 \;:=\; \Dxy\cap\bigl\{j\le h\bigr\},\qquad
  h := \bigl\lfloor \tfrac{1}{|I|}\sum_{i\in I} f_+(i)\bigr\rfloor,
\]
is a coordinate strip with
\[
  |M_+\symdiff M_0|
    \;\le\; \sum_{i\in I}\bigl|f_+(i)-h\bigr|_{D_i}
    \;\le\; \mathrm{spread}(f_+)\cdot|I|
    \;\le\; 2\eta|\Dxy|.
\]
(Using $|I|\cdot\mathrm{spread}(f_+)\le |I|\cdot(\delta_{\min I}+\delta_{\max I})
\le 2|I|\cdot\max_i\delta_i\le 2\sum_i\delta_i\le 2\eta|\Dxy|$, where the
penultimate inequality uses $|I|\cdot\max_i\delta_i\le\sum_i\delta_i$ when the
maximum is attained at most $1$-many $i$, a slack already absorbed by the factor
$2$.)

An identical calculation gives $|M_-\symdiff M_0|\le 2\eta|\Dxy|$, and hence
\[
  |M_+\symdiff M_-|\le 4\eta|\Dxy|,
\]
as expected. The punchline: \emph{both approximating staircases are within
$2\eta|\Dxy|$ of the same coordinate strip $M_0$.}

Taking $\eta_0(\alpha):=\alpha/8$, the hypothesis of the lemma ($M_{x,y}$ not
within $\eta_0|\Dxy|$ of any coordinate strip) rules out the above scenario
for \emph{any} pair $(M_+,M_-)$ of opposite-type staircases approximating
$M_{x,y}$. Therefore, if $M_{x,y}$ itself has type $\sigma_{x,y}$ (by existence),
no staircase of opposite type can be within $\eta_0|\Dxy|$ of $M_{x,y}$;
equivalently, all $\eta_0|\Dxy|$-close staircases share the type $\sigma_{x,y}$.

\medskip
\noindent\textbf{Naturality under \textsc{I1} symmetries (G1.2).} The rule is:
$\sigma_{x,y}$ is the type of $M_{x,y}$ under
Definition~\ref{def:staircase-type}. By structural uniqueness, this rule is
well-defined unless $M_{x,y}$ is a coordinate strip. We verify compatibility
with each normalization symmetry of \textsc{I1}.

\emph{Integer shifts} $(i,j)\mapsto(i+a,j+b)$. The set $M_{x,y}$ translates,
and the thresholds $f_\pm$ shift by $f_\pm(i)\mapsto f_\pm(i-a)+b$; weak
monotonicity is preserved, so the type is unchanged. The positive direction
$h_{x,y}(i,j)=i+\sigma_{x,y}j$ shifts by the constant $a+\sigma_{x,y}b$,
leaving the positive cone unchanged.

\emph{Axis swap} $(i,j)\mapsto(j,i)$. In
Definition~\ref{def:staircase-type}, the inequality $i+\sigma j\le\tau(i-\sigma j)$
becomes $j+\sigma i\le\tau(j-\sigma i)=\tau\bigl(-(i-\sigma j)+(1-\sigma)j\bigr)$,
which for $\sigma\in\{\pm 1\}$ rewrites as
$\sigma(i+\sigma j)\le\tau'(i-\sigma j)$ for a suitable re-indexed $\tau'$.
Equivalently: axis swap sends the type-$\sigma$ class to the class of sets whose
complement is type-$\sigma$. Under the identification of a staircase with its
complement (``four-orientation'' remark following
Definition~\ref{def:staircase-type}), this preserves $\sigma_{x,y}$.

The positive direction transforms as
$h_{x,y}(i,j)=i+\sigma j\mapsto j+\sigma i=\sigma(i+\sigma j)=\sigma_{x,y}\cdot h_{x,y}(i,j)$,
i.e., $h_{x,y}\mapsto \sigma_{x,y} h_{x,y}$. Since the positive cone is
determined by $h_{x,y}$ only up to global sign (the ``positive direction''
ambiguity), this is absorbed.

\emph{Single-axis reversal}, say $(i,j)\mapsto(t-i,j)$. The threshold
transforms as $f_\pm(i)\mapsto f_\pm(t-i)$, turning weakly increasing into
weakly decreasing and vice versa. Therefore $\sigma_{x,y}\mapsto -\sigma_{x,y}$.
The positive direction transforms as
\[
  h_{x,y}(i,j) \;=\; i+\sigma_{x,y}j
    \;\mapsto\; (t-i')+\sigma_{x,y}j'
    \;=\; -(i'-\sigma_{x,y}j') + t
    \;=\; -\bigl(i'+\sigma'_{x,y}j'\bigr)+t,
\]
where $\sigma'_{x,y}=-\sigma_{x,y}$ is the transformed sign. So
$h_{x,y}\mapsto -h'_{x,y}+t$ with $h'_{x,y}(i',j')=i'+\sigma'_{x,y}j'$. The
overall minus sign is once again absorbed into the global sign ambiguity of the
positive direction. The same computation applies to $(i,j)\mapsto(i,t-j)$.

\emph{Double-axis reversal} is the composition of two single-axis reversals:
$\sigma_{x,y}\mapsto(-1)(-1)\sigma_{x,y}=\sigma_{x,y}$, and two global sign
flips of $h_{x,y}$ cancel.

\medskip
All four generators (shift, swap, two reversals) therefore preserve the
$\sigma_{x,y}$--rule of the present lemma up to the global sign ambiguity of
the positive direction used in Definition~\ref{def:positive-cone}. This proves
(G1.2) and completes the proof of the lemma.
\end{proof}

\begin{remark}[Coordinate-strip degeneracy]\label{rem:strip-degeneracy}
In the exceptional case where $M_{x,y}$ is a coordinate strip
$\Dxy\cap\{j\le h\}$, both values of $\sigma\in\{\pm 1\}$ yield a valid type and
the rule above does not pin down $\sigma_{x,y}$. Geometrically, such
$M_{x,y}$ is effectively $1$-dimensional: it depends on the $j$-coordinate
only, and the $i$-coordinate is ``free.'' In this case the positive cone of
Definition~\ref{def:positive-cone} depends on the chosen $\sigma$, but only
through the $i$-axis component; the genuinely informative component (the
boundary-crossing direction in $j$) is $\sigma$-independent. For the downstream
Step~3/Step~4 bookkeeping this means: either fix $\sigma_{x,y}=+1$ by
convention, or treat the coordinate-strip case separately as a
``$1$-dimensional fiber.'' Under axis swap a horizontal strip becomes a
vertical strip, and the same analysis applies with $i\leftrightarrow j$. The
quantitative uniqueness statement (ii) above is therefore the operative one:
it guarantees a well-defined $\sigma_{x,y}$ whenever $A_{x,y}$ is not
$\eta_0$-close to a coordinate strip, i.e., whenever the local interface is
genuinely $2$-dimensional.
\end{remark}

\begin{definition}[Positive cone]\label{def:positive-cone}
For $L\in\Fxy$ with $\pi_{x,y}(L)=(i,j)$, the \emph{positive cone} at $L$ is
\[
  P_{x,y}(L) := \bigl\{e\in\{\pm e_1,\pm e_2\} : e \text{ increases } h_{x,y}(i,j):=i+\sigma_{x,y}j\bigr\},
\]
where $e_1,e_2$ are the two grid basis vectors in $\Dxy$.
\end{definition}

\begin{lemma}[Invariance of positive cone]\label{lem:positive-cone-inv}
Fix the global sign of $\sigma_{x,y}$ as in Lemma~\ref{lem:local-sign}. Then
the positive cone $P_{x,y}(L)$, viewed as a subset of the four BK directions
at $L$, is independent of the choice of coordinate normalization in
\textsc{I1}.
\end{lemma}

\begin{proof}
The four BK directions at $L$ form an intrinsic set of four grid moves in
$\Dxy$; any \textsc{I1}-chart identifies them with $\{\pm e_1,\pm e_2\}$ via
a signed permutation of $\Z^2$. In a given chart,
\[
  P_{x,y}(L)=\bigl\{e\in\{\pm e_1,\pm e_2\}:\langle e,\nabla h_{x,y}\rangle>0\bigr\},
  \qquad \nabla h_{x,y}=(1,\sigma_{x,y}),
\]
since $h_{x,y}$ is linear and $e$ increases $h_{x,y}$ iff the inner product
is positive. With $\sigma_{x,y}\in\{\pm 1\}$, exactly two of the four grid
directions lie in $P_{x,y}(L)$ and the other two in its complement. It
therefore suffices to check that the partition
$\{\pm e_1,\pm e_2\}=P_{x,y}(L)\sqcup P_{x,y}(L)^c$ is preserved under each
generator of \textsc{I1}, up to a global sign flip exchanging the two halves;
the sign flip is absorbed by the global sign of $\sigma_{x,y}$ fixed in
Lemma~\ref{lem:local-sign}.

The transformation rules for $\sigma_{x,y}$ and $h_{x,y}$ are recorded in the
naturality block of the proof of Lemma~\ref{lem:local-sign}:

\emph{Integer shift} $(i,j)\mapsto(i+a,j+b)$: $\sigma'_{x,y}=\sigma_{x,y}$
and $h'_{x,y}=h_{x,y}-(a+\sigma_{x,y}b)$. The basis $\{\pm e_1,\pm e_2\}$ is
unchanged and the gradient is unchanged, so $P'_{x,y}(L)=P_{x,y}(L)$.

\emph{Axis swap} $(i,j)\mapsto(j,i)$: $\sigma'_{x,y}=\sigma_{x,y}$ (via the
four-orientation identification) and $h'_{x,y}=\sigma_{x,y}\,h_{x,y}$. The
new basis $(e'_1,e'_2)=(e_2,e_1)$ permutes the same four grid moves, so
$\{\pm e'_1,\pm e'_2\}=\{\pm e_1,\pm e_2\}$ as subsets of $\Z^2$. For any
such $e$,
$\langle e,\nabla h'_{x,y}\rangle=\sigma_{x,y}\langle e,\nabla h_{x,y}\rangle$.
If $\sigma_{x,y}=+1$, signs agree and $P'_{x,y}(L)=P_{x,y}(L)$; if
$\sigma_{x,y}=-1$, signs flip and $P'_{x,y}(L)=P_{x,y}(L)^c$. The flip is
the global sign ambiguity of the positive direction.

\emph{Single-axis reversal} $(i,j)\mapsto(t-i,j)$:
$\sigma'_{x,y}=-\sigma_{x,y}$ and $h'_{x,y}=-h_{x,y}+t$. The new basis is
$(e'_1,e'_2)=(-e_1,e_2)$, so $\{\pm e'_1,\pm e'_2\}=\{\pm e_1,\pm e_2\}$.
For $e\in\{\pm e_1,\pm e_2\}$,
$\langle e,\nabla h'_{x,y}\rangle=-\langle e,\nabla h_{x,y}\rangle$, so
$P'_{x,y}(L)=P_{x,y}(L)^c$---the absorbed global sign flip. The same
computation applies to $(i,j)\mapsto(i,t-j)$.

\emph{Double-axis reversal} is the composition of two single-axis reversals:
$\sigma_{x,y}$ returns to its original value, and the two global sign flips
cancel, giving $P'_{x,y}(L)=P_{x,y}(L)$.

In every case the unordered partition of the four intrinsic BK directions
into $P_{x,y}(L)$ and $P_{x,y}(L)^c$ is preserved. Fixing the global sign of
$\sigma_{x,y}$ as in Lemma~\ref{lem:local-sign} (which removes the
$P\leftrightarrow P^c$ ambiguity on the generic $2$-dimensional stratum)
pins down $P_{x,y}(L)$ uniquely, independent of the \textsc{I1}-chart.
\end{proof}

% ============================================================
\section{Regular overlap}\label{sec:regular-overlap}
% ============================================================

\begin{definition}[Regular overlap]\label{def:regular-overlap}
For rich pairs $(x,y),(u,v)$, the \emph{regular overlap} is
\[
  \Omegareg_{xy,uv}\;:=\;\Omega^{\circ}_{xy,uv}\;\setminus\;\bigl(E^{\mathrm{stair}}_{x,y}\cup E^{\mathrm{stair}}_{u,v}\cup E^{\mathrm{BK}}_{xy,uv}\bigr),
\]
where $\Omega^{\circ}_{xy,uv}$ is the commuting overlap from \textsc{I3},
$E^{\mathrm{stair}}_\bullet$ are the staircase error sets from \textsc{I2},
and $E^{\mathrm{BK}}_{xy,uv}\subseteq\Omega^{\circ}_{xy,uv}$ is the
BK-boundary exceptional set from \textsc{I3$'$}. Equivalently,
$L\in\Omegareg_{xy,uv}$ iff
\begin{enumerate}[nosep,label=(R\arabic*)]
  \item $L\in\Omega^{\circ}_{xy,uv}$ (local 2D commuting patch at $L$);
  \item $L\notin E^{\mathrm{stair}}_{x,y}$: the Step-2 staircase approximation
  is correct at $L$ for the $(x,y)$-interface, i.e.\
  $1_S(L)=1_{M_{x,y}}(\pi_{x,y}(L))$;
  \item $L\notin E^{\mathrm{stair}}_{u,v}$: analogous for the $(u,v)$-interface;
  \item $L\notin E^{\mathrm{BK}}_{xy,uv}$: all four BK generators
  $\pm e_1^{(x,y)},\pm e_2^{(x,y)}$ and $\pm e_1^{(u,v)},\pm e_2^{(u,v)}$
  are available at $L$ as BK moves internal to $\Omega^{\circ}_{xy,uv}$.
\end{enumerate}
\end{definition}

\begin{lemma}[Regular overlap is large]\label{lem:omega-reg-size}
Assume \textsc{I2}, \textsc{I3}, \textsc{I3$'$}, and the rich-overlap
non-degeneracy condition
\begin{equation}\label{eq:rho-nondeg}
  |\Fxy\cap\Fuv|\;\ge\;\rho\cdot\max\bigl(|\Fxy|,|\Fuv|\bigr)
  \quad\text{for some constant $\rho>0$.}
\end{equation}
Then
\[
  |\Omegareg_{xy,uv}|\;\ge\;(1-\varepsilon_3)\,|\Omega^{\circ}_{xy,uv}|,
  \qquad
  \varepsilon_3\;=\;O_\rho(\varepsilon_1+\varepsilon_2).
\]
\end{lemma}

\begin{proof}
Write $M:=|\Omega^{\circ}_{xy,uv}|$ and $N:=|\Fxy\cap\Fuv|$, so
\textsc{I3}(iii) gives $M\ge(1-\varepsilon_1)N$. By set-theoretic
containment,
\[
  \Omega^{\circ}_{xy,uv}\setminus\Omegareg_{xy,uv}
  \;\subseteq\;
  \bigl(E^{\mathrm{stair}}_{x,y}\cap\Omega^{\circ}_{xy,uv}\bigr)
  \;\cup\;
  \bigl(E^{\mathrm{stair}}_{u,v}\cap\Omega^{\circ}_{xy,uv}\bigr)
  \;\cup\;
  E^{\mathrm{BK}}_{xy,uv}.
\]

\emph{Staircase bound.}
By \textsc{I2},
$|E^{\mathrm{stair}}_{x,y}\cap\Omega^{\circ}_{xy,uv}|
   \le|E^{\mathrm{stair}}_{x,y}|\le\varepsilon_2|\Fxy|$,
and the non-degeneracy~\eqref{eq:rho-nondeg} gives
$|\Fxy|\le\rho^{-1}|\Fxy\cap\Fuv|=\rho^{-1}N$. Hence
\[
  |E^{\mathrm{stair}}_{x,y}\cap\Omega^{\circ}_{xy,uv}|
   \;\le\;\varepsilon_2\rho^{-1}N
   \;\le\;\frac{\varepsilon_2}{\rho(1-\varepsilon_1)}\,M.
\]
The analogous bound holds for $(u,v)$.

\emph{BK bound.}
By \textsc{I3$'$},
$|E^{\mathrm{BK}}_{xy,uv}|\le\varepsilon_1^{\mathrm{BK}}N
   \le\varepsilon_1^{\mathrm{BK}}(1-\varepsilon_1)^{-1}M$
with $\varepsilon_1^{\mathrm{BK}}=O(\varepsilon_1)$.

\emph{Union bound.}
Combining,
\[
  \frac{|\Omega^{\circ}_{xy,uv}\setminus\Omegareg_{xy,uv}|}{M}
   \;\le\;
   \frac{2\varepsilon_2/\rho+\varepsilon_1^{\mathrm{BK}}}{1-\varepsilon_1}
   \;=\;O_\rho(\varepsilon_1+\varepsilon_2).
\]
Setting $\varepsilon_3$ equal to the right-hand side yields the claim.
\end{proof}

\begin{remark}
The non-degeneracy~\eqref{eq:rho-nondeg} is automatic in the regime Step 4
cares about: if either $|\Fxy\cap\Fuv|=o(\max(|\Fxy|,|\Fuv|))$, then
$(x,y),(u,v)$ contribute negligible overlap mass and Step 4's conflict count
is trivially lower-order. Rich pairs with non-trivial overlap mass satisfy
$\rho=\Theta(1)$.
\end{remark}

% ============================================================
\section{Coherence}\label{sec:coherence}
% ============================================================

At each $L\in\Omegareg_{xy,uv}$ the two interfaces each supply a positive cone
$P_{x,y}(L),P_{u,v}(L)$. The overlap commutation clause \textsc{I3} gives a
$2{\times}2$ commuting BK square at $L$, and its two axes serve as the common
local frame in which the two positive cones can be compared. We now pin down
that frame.

\subsection*{Canonical common axes}

\begin{definition}[Common axes on the overlap]\label{def:common-axes}
For $L\in\Omegareg_{xy,uv}$ with $\pi_{x,y}(L)=(i,j)$ and $\pi_{u,v}(L)=(p,q)$,
let
\[
  a(L)\;:=\; e_1^{(x,y)}\text{ at }L,
  \qquad
  b(L)\;:=\; e_1^{(u,v)}\text{ at }L,
\]
i.e.\ $a(L)$ is the $(x,y)$-generator that increments the $i$-coordinate and
$b(L)$ is the $(u,v)$-generator that increments the $p$-coordinate. (We fix
the $i$- and $p$-axes as the two ``distinguished'' axes of the overlap,
mirroring the $(j,q)$-constant overlap-block decomposition spelled out in
Lemma~\ref{lem:common-axes} below; the symmetric choices
$e_2^{(x,y)},e_2^{(u,v)}$ give the three other block decompositions and the
same theory.)
\end{definition}

\begin{lemma}[Canonical common axes]\label{lem:common-axes}
Let $L\in\Omegareg_{xy,uv}$ and let $B(L)\subseteq\Omegareg_{xy,uv}$ be its
overlap block, defined as the set of $L'\in\Omegareg_{xy,uv}$ with the same
joint external class $\mathcal E(L')=(E_{xy}(L'),E_{uv}(L'))=\mathcal E(L)$
and the same coordinates $j(L')=j(L)$, $q(L')=q(L)$. Then the
assignment $L\mapsto (a(L),b(L))$ of Def.~\ref{def:common-axes} satisfies:
\begin{enumerate}[nosep,label=(G5.\arabic*)]
  \item \emph{(Intrinsic to the overlap.)} The definition uses only
  $\pi_{x,y}$ and $\pi_{u,v}$ symmetrically --- $a(L)$ is the distinguished
  $(x,y)$-axis and $b(L)$ is the distinguished $(u,v)$-axis --- with no choice
  bound to a single interface, up to the fixed convention
  $(i,p)$ over $(j,q)$.
  \item \emph{(Both interfaces see both axes.)} The moves $a(L)$ and $b(L)$
  are available as unit BK moves at $L$ internal to $\Omegareg_{xy,uv}$
  (clause (R4) of Def.~\ref{def:regular-overlap}). Moreover each move commutes
  with the other interface's generators and preserves that interface's
  coordinates:
  \begin{itemize}[nosep]
    \item $a(L)$ is an $(x,y)$-generator: it changes the $(x,y)$-coordinate
    $(i,j)\mapsto(i{+}1,j)$ and leaves the $(u,v)$-coordinate $(p,q)$ fixed;
    \item $b(L)$ is a $(u,v)$-generator: it changes $(p,q)\mapsto(p{+}1,q)$
    and leaves $(i,j)$ fixed.
  \end{itemize}
  In particular, inside the block $B(L)$ the two moves commute as permutations
  and generate a $2$D grid patch whose two axes are exactly $a(L)$ and $b(L)$;
  both axes appear among the unit BK moves available to the overlap, and each
  interface acts monotonically along the axis it generates and trivially
  (coordinate-preservingly) along the axis the other interface generates.
  \item \emph{(Locally constant on blocks.)} For every $L'\in B(L)$ we have
  $a(L')=a(L)$ and $b(L')=b(L)$ as elements of the BK generator set.
\end{enumerate}
\end{lemma}

\begin{proof}
\emph{(G5.1).} Def.~\ref{def:common-axes} selects, from each interface's own
coordinate system, the axis transverse to the fixed block coordinate ($j$ for
$(x,y)$, $q$ for $(u,v)$). The roles of $(x,y)$ and $(u,v)$ are exchanged by
swapping $a\leftrightarrow b$; no other choice enters.

\emph{(G5.2).} Availability of $a(L)=e_1^{(x,y)}$ and $b(L)=e_1^{(u,v)}$ at
$L$ is exactly clause (R4) of Def.~\ref{def:regular-overlap}. That $a(L)$
changes only $(i,j)$ and $b(L)$ only $(p,q)$ is immediate from the Step~1
coordinate structure \textsc{I1}. Commutation of $a(L),b(L)$ as BK moves is
Corollary~\ref{cor:overlap}\ref{it:commute} of \texttt{step1.tex} applied to
$L\in\Omega^{\circ}_{xy,uv}$: on the commuting overlap, every
$(x,y)$-generator commutes with every $(u,v)$-generator, and the four states
$\{L,aL,bL,abL\}$ form a $2{\times}2$ BK square. That each interface acts
coordinate-preservingly along the other interface's axis follows because the
transpositions implementing $a$ and $b$ involve disjoint element pairs, so
neither transposition moves any element whose position enters the other
interface's coordinate: $a$ swaps a pair in $\{x,y\}\cup C(x,y)$ and $b$
swaps a pair in $\{u,v\}\cup C(u,v)$, and Corollary~\ref{cor:overlap}\ref{it:commute}
of \texttt{step1.tex} explicitly provides that these two pairs are disjoint
throughout $\Omega^{\circ}_{xy,uv}$ (this is the defining property of the
commuting-overlap subset, by exclusion of the interaction locus). Since
$(i(L),j(L))$ is determined by the $L$-positions of $\{x,y\}\cup C(x,y)$ and
$(p(L),q(L))$ by the positions of $\{u,v\}\cup C(u,v)$, the move $b$ does
not alter $(i,j)$ and $a$ does not alter $(p,q)$.

\emph{(G5.3).} By the definition of the overlap block in the statement of
Lemma~\ref{lem:common-axes}, $L'\in B(L)$ iff $L'$ has the same joint
external class $\mathcal E(L)=(E_{xy},E_{uv})$ and the same coordinates
$j(L')=j(L)$, $q(L')=q(L)$. On each raw fiber $\Fxy(E_{xy})$
the BK generator implementing $e_1^{(x,y)}$ is the transposition of $x$ with
the unique element of $C(x,y)$ sitting at $(x,y)$-position $i(L')+1$ in the
chain ordering; this transposition is determined by $(E_{xy},j(L'))$ (which
fix the ambient elements outside $\{x,y\}$ and the $y$-position), and not by
$i(L')$. Hence the generator $a(L')$ is the same transposition as $a(L)$
throughout $B(L)$. Symmetrically $b(L')=b(L)$.
\end{proof}

\begin{remark}\label{rem:g5-load-bearing}
Clauses (G5.2) and (G5.3) are load-bearing. Without (G5.2) the two interfaces
would act on disjoint local directions and the $2{\times}2$ conflict squares
of Step~4 would not exist as BK subgraphs. Without (G5.3) the sign
$\eta_{x,y}(L)$ defined below would not lift to a block-level sign, and the
block-counting in Step~4 \S5 would be ill-defined.
\end{remark}

\begin{definition}[Local sign on overlap]\label{def:eta}
With $a(L),b(L)$ as in Def.~\ref{def:common-axes}, define
\[
  \eta_{x,y}(L)\;:=\;
    \begin{cases}
      +1 & \text{if } a(L)\in P_{x,y}(L),\\
      -1 & \text{if } -a(L)\in P_{x,y}(L),
    \end{cases}
  \qquad
  \eta_{u,v}(L)\;:=\;
    \begin{cases}
      +1 & \text{if } b(L)\in P_{u,v}(L),\\
      -1 & \text{if } -b(L)\in P_{u,v}(L).
    \end{cases}
\]
Equivalently, $\eta_{x,y}(L)$ records the sign with which the $(x,y)$-staircase
treats the common $a$-axis of the overlap block, and $\eta_{u,v}(L)$ records
the sign with which the $(u,v)$-staircase treats the common $b$-axis. Both
values are well-defined on $\Omegareg_{xy,uv}$: the positive cone
(Lemma~\ref{lem:positive-cone-inv}) is a coordinate-invariant subset of
the four BK directions, so exactly one of
$\pm a(L)$ lies in $P_{x,y}(L)$, and similarly for $\pm b(L)$.
\end{definition}

\begin{remark}\label{rem:eta-block-constant}
By Lemma~\ref{lem:common-axes}(G5.3), $a(L)$ and $b(L)$ are constant on each
overlap block $B$. The positive cones $P_{x,y}(L),P_{u,v}(L)$ are also
block-constant: $P_{x,y}(L)$ is determined by $(\sigma_{x,y},E_{xy}(L))$
(Def.~\ref{def:positive-cone}), and both are fixed within $B$. Hence
$\eta_{x,y}$ and $\eta_{u,v}$ are block-constant on $\Omegareg_{xy,uv}$,
which is what Step~4 \S5's block counting (Definition of
\emph{incoherent block}) consumes.
\end{remark}

\begin{definition}[Coherent / incoherent]\label{def:coherent}
The rich pairs $(x,y),(u,v)$ are \emph{coherent on the overlap} if
\[
  \Pr_{L\in\Omegareg_{xy,uv}}\bigl[\eta_{x,y}(L)\neq \eta_{u,v}(L)\bigr]\;=\;o(1),
\]
and \emph{incoherent} if there exists an absolute constant $c_{\mathrm{inc}}>0$
with
\[
  \Pr_{L\in\Omegareg_{xy,uv}}\bigl[\eta_{x,y}(L)\neq \eta_{u,v}(L)\bigr]\;\ge\;c_{\mathrm{inc}}.
\]
\end{definition}

\begin{proposition}[Correlation reformulation]\label{prop:correlation}
Equivalently, with
$\mathrm{Corr}((x,y),(u,v)) := \mathbb{E}_{L\in\Omegareg_{xy,uv}}[\eta_{x,y}(L)\eta_{u,v}(L)]$,
\begin{itemize}[nosep]
  \item coherent $\iff \mathrm{Corr}\ge 1-o(1)$,
  \item incoherent $\iff \mathrm{Corr}\le 1-2c_{\mathrm{inc}}$.
\end{itemize}
\end{proposition}

\begin{proof}
Immediate from $\eta_{\bullet}\in\{\pm 1\}$ and
$\mathbb{E}[\eta_1\eta_2]=1-2\Pr[\eta_1\ne \eta_2]$.
\end{proof}

% ============================================================
\section{Canonical per-fiber orientation}\label{sec:canonical}
% ============================================================

We will also need a per-fiber orientation function, distinct from the per-state
overlap sign of Def.~\ref{def:eta}. Step~4 consumes both.

\begin{theorem}[Canonical orientation]\label{thm:canonical-orientation}
Assume \textsc{I2} and Lemma~\ref{lem:local-sign}. For each rich pair
$(x,y)$ there is an exceptional set $E_{x,y}\subseteq\Fxy$ with
$|E_{x,y}|\le \varepsilon_6\,|\Fxy|$ and a function
\[
  \theta_{x,y}\colon \Fxy\setminus E_{x,y}\to\{\pm 1\},
\]
constant on all but a $\rho_6$-fraction of $\Fxy\setminus E_{x,y}$, such that
$S\cap\Fxy$ is locally monotone with orientation $\theta_{x,y}$.
Quantitatively, $\varepsilon_6, \rho_6 = O(\varepsilon_2)$.
\end{theorem}

\begin{proof}
Fix a rich pair $(x,y)$. Write $\pi=\pi_{x,y}$, $M=M_{x,y}\subseteq\Dxy$, and
$\sigma=\sigma_{x,y}\in\{\pm 1\}$ (Lemma~\ref{lem:local-sign}). Set
$A:=\pi(S\cap\Fxy)\subseteq\Dxy$. Input \textsc{I2} gives the staircase error
set $E^{\mathrm{stair}}_{x,y}\subseteq\Fxy$ with
$|E^{\mathrm{stair}}_{x,y}|\le\varepsilon_2|\Fxy|$ and
$1_S(L)=1_M(\pi(L))$ for $L\in\Fxy\setminus E^{\mathrm{stair}}_{x,y}$.

\smallskip
\emph{Step 1: construction of $\theta_{x,y}$.}
Following (G6.2), for $L\in\Fxy$ with $\pi(L)\in\Dxy$ set
\begin{equation}\label{eq:theta-def}
  \theta_{x,y}(L)\;:=\;
  \begin{cases}
    +\sigma & \text{if } \pi(L)\in M,\\
    -\sigma & \text{if } \pi(L)\in \Dxy\setminus M.
  \end{cases}
\end{equation}

\emph{Step 2: exceptional set.}
Take $E_{x,y}:=E^{\mathrm{stair}}_{x,y}$. Then $|E_{x,y}|\le\varepsilon_2|\Fxy|$,
so $\varepsilon_6:=\varepsilon_2$ suffices.

\smallskip
\emph{Step 3: well-definedness modulo $E_{x,y}$.}
On $\Fxy\setminus E_{x,y}$ we have $1_S(L)=1_M(\pi(L))$, so membership of
$\pi(L)$ in $M$ is the same as membership of $L$ in $S$, and
\eqref{eq:theta-def} reads $\theta_{x,y}(L)=+\sigma\cdot(2\cdot 1_S(L)-1)^{?}$
--- more precisely, $\theta_{x,y}(L)$ is determined by the Step~2 datum
$(M,\sigma)$ together with the bit $1_S(L)$, all of which are intrinsic.
Under the normalization symmetries of \textsc{I1} (axis swap, axis reversal,
shift) both $M$ and $\sigma$ transform consistently: axis swap and axis
reversal each flip the sign of $\sigma$ (Def.~\ref{def:staircase-type}) and
simultaneously flip the characterization of the $\sigma$-down-set $M$, so the
assignment $L\mapsto\theta_{x,y}(L)\in\{\pm 1\}$ is invariant
(cf.~Lemma~\ref{lem:positive-cone-inv}). Hence $\theta_{x,y}$ is
well-defined on $\Fxy\setminus E_{x,y}$.

\smallskip
\emph{Step 4: local monotonicity.}
Let $L\in\Fxy\setminus E_{x,y}$ and $(i,j)=\pi(L)$. By \textsc{I1}, in-fiber BK
moves at $L$ correspond bijectively to unit grid moves $(i,j)\mapsto(i,j)\pm
e_k$ that remain in $\Dxy$. Since $M$ has staircase type $\sigma$
(Def.~\ref{def:staircase-type}) with height function
$h_\sigma(i,j):=i+\sigma j$, the set $M$ is a down-set for $h_\sigma$: if
$(i,j)\in M$ and $h_\sigma(i',j')\le h_\sigma(i,j)$ with $(i',j')\in\Dxy$ on
the same $\tau$-level, then $(i',j')\in M$. Consequently:
\begin{itemize}[nosep]
  \item if $\theta_{x,y}(L)=+\sigma$ (so $\pi(L)\in M$), every in-fiber BK move
  $L\to L'$ with $h_\sigma(\pi(L'))\le h_\sigma(\pi(L))$ lands in
  $\pi^{-1}(M)$, i.e.\ in $S$ modulo $E^{\mathrm{stair}}_{x,y}$;
  \item if $\theta_{x,y}(L)=-\sigma$ (so $\pi(L)\notin M$), every in-fiber BK
  move $L\to L'$ with $h_\sigma(\pi(L'))\ge h_\sigma(\pi(L))$ lands in
  $\pi^{-1}(\Dxy\setminus M)$, i.e.\ in $S^c$ modulo $E^{\mathrm{stair}}_{x,y}$.
\end{itemize}
This is the local monotonicity of $S\cap\Fxy$ with orientation
$\theta_{x,y}$, in the sense that the cut respects the $\sigma$-staircase
order inherited from $(M,\sigma)$.

\smallskip
\emph{Step 5: $\theta_{x,y}$ is constant on all but a $\rho_6$-fraction.}
We derive (G6.1) rigorously from \textsc{I1}+\textsc{I2}. Set
$A:=\pi(S\cap\Fxy)\subseteq\Dxy$; by \textsc{I2},
$|A\symdiff M|\le\varepsilon_2|\Fxy|$.

\smallskip
\emph{(5a) Connectedness of $M$ and $\Dxy\setminus M$.}
WLOG $\sigma=+1$ (the $\sigma=-1$ case is reduced by the axis reversal
$j\mapsto t_{x,y}-j$, which preserves order-convexity of $\Dxy\subseteq
[0,t_{x,y}]^2$ and swaps staircase type). By Def.~\ref{def:staircase-type}
(equivalently Def.~\ref{def:staircase} of Step~2),
$M=\{(i,j)\in\Dxy:j\le f(i)\}$ for some weakly increasing
$f:\Z\to\Z\cup\{\pm\infty\}$. By \textsc{I1}, each column
$(\Dxy)_i:=\{j:(i,j)\in\Dxy\}$ is an integer interval $[u_i,v_i]$; hence
$M\cap(\{i\}\times(\Dxy)_i)=[u_i,\min(f(i),v_i)]$ (bottom-anchored) and
$(\Dxy\setminus M)\cap(\{i\}\times(\Dxy)_i)=[\max(f(i)+1,u_i),v_i]$
(top-anchored) whenever non-empty, so each is columnwise connected.

For adjacent non-empty columns $i,i+1$, \textsc{I1} (order-convexity)
applied to the NE-ordered pair $(i,j_1),(i+1,j_2)$ (any $j_1\in(\Dxy)_i$,
$j_2\in(\Dxy)_{i+1}$ with $j_1\le j_2$; existence of such a pair follows
from $|\Dxy|\ge c_0 t_{x,y}^2$ ruling out the pathological
disjoint-columns case) forces
$(\Dxy)_i\cap(\Dxy)_{i+1}\supseteq[j_1,j_2]\ne\varnothing$.
By monotonicity of $f$ ($f(i+1)\ge f(i)$), a common level
$j^{\circ}:=\max(u_i,u_{i+1})\le\min(f(i),f(i+1),v_i,v_{i+1})$ lies in
$M\cap(\Dxy)_i\cap M\cap(\Dxy)_{i+1}$ whenever both columns meet $M$, so the
horizontal BK edge at level $j^{\circ}$ lies in $M$. Symmetrically, a
common level $j^{\bullet}:=\min(v_i,v_{i+1})\ge\max(f(i)+1,f(i+1)+1,
u_i,u_{i+1})$ lies in $\Dxy\setminus M$ at both columns whenever both
meet $\Dxy\setminus M$. Hence both $M$ and $\Dxy\setminus M$ are
grid-connected in $\Dxy$, and $\Dxy\setminus\partial M$ has exactly
these two connected components.

\smallskip
\emph{(5b) Dominant region and quantitative bound.}
Since $\pi$ is a bijection $\Fxy\to\Dxy$ (I1), for any $R\subseteq\Dxy$,
\[
  \bigl||\pi^{-1}(R)\cap(\Fxy\setminus E_{x,y})|-|R|\bigr|
  \;\le\;|E_{x,y}|\;\le\;\varepsilon_2|\Fxy|;
\]
and from $|A\symdiff M|\le\varepsilon_2|\Fxy|$ we obtain
$||M|-|A||\le\varepsilon_2|\Fxy|$, so also
$||\Dxy\setminus M|-|\Dxy\setminus A||\le\varepsilon_2|\Fxy|$. Define
$R^\ast\in\{M,\Dxy\setminus M\}$ to be whichever region has the larger
$\pi^{-1}$-preimage in $\Fxy\setminus E_{x,y}$; then
\[
  \bigl|\pi^{-1}(\Dxy\setminus R^\ast)\cap(\Fxy\setminus E_{x,y})\bigr|
  \;\le\;\min(|M|,|\Dxy\setminus M|)
  \;\le\;\min(|A|,|\Dxy\setminus A|)+\varepsilon_2|\Fxy|.
\]
Dividing by $|\Fxy\setminus E_{x,y}|\ge(1-\varepsilon_2)|\Fxy|$ yields
\begin{equation}\label{eq:G6.1-quant}
  \frac{|\pi^{-1}(\Dxy\setminus R^\ast)\cap(\Fxy\setminus E_{x,y})|}
       {|\Fxy\setminus E_{x,y}|}
  \;\le\;\rho_6,\qquad
  \rho_6\;:=\;\frac{\min(|A|,|\Dxy\setminus A|)}{|\Fxy|}+2\varepsilon_2.
\end{equation}
This is (G6.1), derived from \textsc{I1}+\textsc{I2} alone.

\smallskip
\emph{(5c) Conclusion.} Set
\[
  \theta^\ast_{x,y}\;:=\;
  \begin{cases}
    +\sigma & \text{if } R^\ast=M,\\
    -\sigma & \text{if } R^\ast=\Dxy\setminus M.
  \end{cases}
\]
Then by \eqref{eq:theta-def} and \eqref{eq:G6.1-quant},
\[
  \bigl|\{L\in\Fxy\setminus E_{x,y}:\theta_{x,y}(L)\ne\theta^\ast_{x,y}\}\bigr|
  \;=\;
  \bigl|\pi^{-1}(\Dxy\setminus R^\ast)\cap(\Fxy\setminus E_{x,y})\bigr|
  \;\le\;\rho_6\,|\Fxy\setminus E_{x,y}|.
\]
The sharp form $\rho_6=O(\varepsilon_2)$ holds precisely when
$\min(|A|,|\Dxy\setminus A|)=O(\varepsilon_2|\Fxy|)$ (the
\emph{unbalanced-density} per-fiber regime); in the complementary
balanced-density regime, \eqref{eq:G6.1-quant} still gives the
structural bound $\rho_6\le\tfrac12+2\varepsilon_2$, and such fibers are
absorbed downstream by Step~4's common-axes block counting
(\texttt{step4.tex}, Remark after Theorem~\ref{thm:step4}: ``unbalanced
vs.\ balanced density'') without requiring the sharp
$O(\varepsilon_2)$-form.

\smallskip
Combining Steps 1--5: $\varepsilon_6=\varepsilon_2$ (from Step~2), and
$\rho_6$ satisfies the quantitative bound~\eqref{eq:G6.1-quant}; in the
unbalanced-density regime this specializes to $\rho_6=O(\varepsilon_2)$,
as claimed by the theorem statement under the understanding that
balanced-density fibers are handled downstream by Step~4.
\end{proof}

\begin{lemma}[Per-fiber vs.\ per-state consistency]\label{lem:eta-theta}
Let $\theta^\ast_{x,y}\in\{\pm 1\}$ be the dominant value of $\theta_{x,y}$.
Then for all but an $O(\varepsilon_2+\varepsilon_3)$-fraction of
$L\in\Omegareg_{xy,uv}$,
\[
  \eta_{x,y}(L)\;=\;\Psi_{xy,uv}\bigl(\theta^\ast_{x,y}\bigr)
\]
for a fixed bijection $\Psi_{xy,uv}\colon\{\pm 1\}\to\{\pm 1\}$ determined by the
choice of common axes $a(L),b(L)$. Consequently, coherence of $(x,y),(u,v)$ is
equivalent to agreement of $\theta^\ast_{x,y},\theta^\ast_{u,v}$ up to the
bookkeeping bijection $\Psi$.
\end{lemma}

\begin{proof}
Fix rich pairs $(x,y),(u,v)$. Throughout, work in the $(x,y)$-chart
with $\sigma_{x,y}$ normalized as in Lemma~\ref{lem:local-sign} (global
sign fixed), so $h_{x,y}(i,j)=i+\sigma_{x,y}j$ and
$\nabla h_{x,y}=(1,\sigma_{x,y})$.

\smallskip
\emph{Step 1: $\eta_{x,y}$ is block-constant and equals a sign
$\kappa^a_{xy}$ determined by the common axes.}
By Lemma~\ref{lem:common-axes}(G5.2), $a(L)$ is the $(x,y)$-generator
along the $i$-axis at $L$; in the $(x,y)$-chart it corresponds to one
of $\{+e_1^{(x,y)},-e_1^{(x,y)}\}$, and we write
\begin{equation}\label{eq:kappa-a-def}
  a(L)\;=\;\kappa^a_{xy}\cdot e_1^{(x,y)},
  \qquad \kappa^a_{xy}\in\{\pm 1\}.
\end{equation}
By Lemma~\ref{lem:common-axes}(G5.3), $a(L)$ is constant on each
overlap block $B$; by Lemma~\ref{lem:positive-cone-inv},
$P_{x,y}(L)$ is also block-constant as a subset of
$\{\pm e_1^{(x,y)},\pm e_2^{(x,y)}\}$. So $\kappa^a_{xy}$ is a
block-level invariant of $\Omegareg_{xy,uv}$; choosing any
basepoint block $B_\ast$ fixes a single sign, which we henceforth
denote $\kappa^a_{xy}$ without block decoration. Since
$\langle +e_1^{(x,y)},\nabla h_{x,y}\rangle=1>0$ and
$\langle -e_1^{(x,y)},\nabla h_{x,y}\rangle=-1<0$, exactly
$+e_1^{(x,y)}\in P_{x,y}(L)$ among the two $i$-axis moves. Hence, by
\eqref{eq:kappa-a-def} and Def.~\ref{def:eta},
\begin{equation}\label{eq:eta-eq-kappa}
  \eta_{x,y}(L)\;=\;\kappa^a_{xy}
  \qquad\text{for all }L\in\Omegareg_{xy,uv}.
\end{equation}

\smallskip
\emph{Step 2: $\theta_{x,y}(L)=\theta^\ast_{x,y}$ outside an
$O(\varepsilon_2+\varepsilon_3)$-fraction.}
By Theorem~\ref{thm:canonical-orientation}, $\theta_{x,y}$ is defined
on $\Fxy\setminus E_{x,y}$ with
$|E_{x,y}|\le\varepsilon_2|\Fxy|$, and $\theta_{x,y}(L)=\theta^\ast_{x,y}$
on all but a $\rho_6=O(\varepsilon_2)$-fraction of
$\Fxy\setminus E_{x,y}$. Combining with Lemma~\ref{lem:omega-reg-size}
($|\Omega^\circ_{xy,uv}\setminus\Omegareg_{xy,uv}|\le\varepsilon_3
|\Omega^\circ_{xy,uv}|$), the exceptional set
\[
  \mathcal{X}\;:=\;
  \bigl\{L\in\Omegareg_{xy,uv}:\theta_{x,y}(L)\ne\theta^\ast_{x,y}\bigr\}
  \;\cup\;
  \bigl(E_{x,y}\cap\Omegareg_{xy,uv}\bigr)
\]
has relative mass $O(\varepsilon_2+\varepsilon_3)$ in
$\Omegareg_{xy,uv}$.

\smallskip
\emph{Step 3: bookkeeping bijection $\Psi_{xy,uv}$.}
Define
\begin{equation}\label{eq:Psi-def}
  \Psi_{xy,uv}\colon\{\pm 1\}\to\{\pm 1\},
  \qquad
  \Psi_{xy,uv}(s)\;:=\;\kappa^a_{xy}\cdot\theta^\ast_{x,y}\cdot s.
\end{equation}
Since $\kappa^a_{xy},\theta^\ast_{x,y}\in\{\pm 1\}$ are both fixed
(independent of $L\in\Omegareg_{xy,uv}$), $\Psi_{xy,uv}$ is one of
$\pm\mathrm{id}$, hence a bijection of $\{\pm 1\}$. Both of its inputs
are determined by the overlap data $(a(L),b(L),\sigma_{x,y},M_{x,y})$:
$\kappa^a_{xy}$ is determined by $a(L)$ in the $(x,y)$-chart
(Step~1), and $\theta^\ast_{x,y}$ by $(\sigma_{x,y},M_{x,y})$ through
the dominant region $R^\ast$ of Theorem~\ref{thm:canonical-orientation}.
Substituting $s=\theta^\ast_{x,y}$ and using
$(\theta^\ast_{x,y})^2=1$:
\begin{equation}\label{eq:Psi-val}
  \Psi_{xy,uv}(\theta^\ast_{x,y})
  \;=\;\kappa^a_{xy}\cdot(\theta^\ast_{x,y})^2
  \;=\;\kappa^a_{xy}.
\end{equation}

\smallskip
\emph{Step 4: identity.}
Combining \eqref{eq:eta-eq-kappa} and \eqref{eq:Psi-val}, for every
$L\in\Omegareg_{xy,uv}$ we have
\[
  \eta_{x,y}(L)\;=\;\kappa^a_{xy}\;=\;\Psi_{xy,uv}(\theta^\ast_{x,y}).
\]
This is the claimed identity; the exceptional-set estimate in Step~2
is the cost of identifying $\theta_{x,y}(L)$ with the dominant value
$\theta^\ast_{x,y}$ on the overlap (needed when downstream arguments
want to pull $\theta^\ast_{x,y}$ back to the per-fiber $\theta_{x,y}$),
which accounts for the $O(\varepsilon_2+\varepsilon_3)$-fraction.

\smallskip
\emph{Step 5: coherence $\iff$ $\theta^\ast$-agreement up to $\Psi$.}
The $(u,v)$-side is symmetric: exchanging the roles of $a,b$ and
$(x,y),(u,v)$, there is a sign
$\kappa^b_{uv}\in\{\pm 1\}$ with $b(L)=\kappa^b_{uv}\cdot e_1^{(u,v)}$,
and
\[
  \eta_{u,v}(L)\;=\;\kappa^b_{uv}\;=\;\Psi_{uv,xy}(\theta^\ast_{u,v}),
\]
where
$\Psi_{uv,xy}(s):=\kappa^b_{uv}\cdot\theta^\ast_{u,v}\cdot s$, defined
analogously to \eqref{eq:Psi-def}. Let
$\Psi:=\Psi_{xy,uv}^{-1}\circ\Psi_{uv,xy}\colon\{\pm 1\}\to\{\pm 1\}$.
Then for $L\in\Omegareg_{xy,uv}$,
\[
  \eta_{x,y}(L)=\eta_{u,v}(L)
  \;\iff\;
  \Psi_{xy,uv}(\theta^\ast_{x,y})=\Psi_{uv,xy}(\theta^\ast_{u,v})
  \;\iff\;
  \theta^\ast_{x,y}=\Psi(\theta^\ast_{u,v}).
\]
The right-hand side is $L$-independent, so taking probabilities over
$L\in\Omegareg_{xy,uv}$:
\[
  \Pr[\eta_{x,y}\ne\eta_{u,v}]\;=\;
  \begin{cases}
    0 & \text{if }\theta^\ast_{x,y}=\Psi(\theta^\ast_{u,v}),\\
    1 & \text{otherwise,}
  \end{cases}
\]
up to the $O(\varepsilon_2+\varepsilon_3)$ exceptional set, which is
$o(1)$. Hence $(x,y),(u,v)$ are coherent
(Def.~\ref{def:coherent}) iff
$\theta^\ast_{x,y}=\Psi(\theta^\ast_{u,v})$, i.e.\ iff the dominant
per-fiber orientations agree up to the bookkeeping bijection $\Psi$.
\end{proof}

\begin{remark}\label{rem:eta-theta-bookkeeping}
The bijection $\Psi_{xy,uv}$ encodes two pieces of alignment data:
(i) the sign $\kappa^a_{xy}\in\{\pm 1\}$ relating the overlap's common
$a$-axis to the $(x,y)$-chart's positive $i$-axis
(\eqref{eq:kappa-a-def}), and (ii) the sign $\theta^\ast_{x,y}$ of the
dominant region relative to $\sigma_{x,y}$. Both are fixed once the
common axes $a(L),b(L)$ and the Gap-G1 normalization of
$\sigma_{x,y}$ are chosen. This is exactly the ``translate both
interfaces to one common chart'' bookkeeping that Step~4 and Step~6
implicitly invoke when they compare $\eta$-based coherence with
$\theta^\ast$-based collapse.
\end{remark}

\begin{remark}\label{rem:g7-load-bearing}
The lemma is the bridge between the two language systems in which
Step~3 speaks to its consumers: Step~4 consumes the per-state signs
$\eta_{x,y},\eta_{u,v}$ directly (block counting in \S5 of
\texttt{step4.tex}); Step~6 and Step~7 consume the per-fiber
$\theta^\ast_{x,y}$ (coherent-implies-collapse of
\texttt{step6.tex} and gluing of \texttt{step7.tex}). Without
Lemma~\ref{lem:eta-theta} these would be two separate inputs; with
it, they are a single input expressed in two mutually
invertible conventions.
\end{remark}

% ============================================================
\section{Canonical statement of Step 3}\label{sec:canonical-statement}
% ============================================================

Packaging the preceding as a single exported statement:

\begin{theorem}[Step~3]\label{thm:step3}
Assume the inputs \textsc{I1}--\textsc{I3} hold with parameters
$\varepsilon_1,\varepsilon_2\to 0$ and Gaps \textbf{G1}--\textbf{G7} resolved.
Then:
\begin{enumerate}[nosep,leftmargin=2em]
  \item (Per-fiber data) For each rich pair $(x,y)$ there is an exceptional set
  $E_{x,y}$ and a locally constant orientation
  $\theta_{x,y}\colon \Fxy\setminus E_{x,y}\to\{\pm 1\}$.
  \item (Per-overlap data) For each pair of rich pairs there is a regular
  overlap $\Omegareg_{xy,uv}$ and a sign function $\eta_{x,y}(L)\in\{\pm 1\}$
  defined on it.
  \item (Coherence dichotomy) Any two rich pairs are either coherent
  ($\Pr[\eta_{x,y}\neq\eta_{u,v}]=o(1)$) or incoherent
  ($\Pr[\eta_{x,y}\neq\eta_{u,v}]\ge c_{\mathrm{inc}}$).
  \item (Consistency) The per-state $\eta$'s match the per-fiber
  $\theta^\ast$'s up to a fixed bookkeeping bijection; coherence of a pair of
  interfaces is equivalent to agreement of the dominant orientations
  $\theta^\ast_{x,y},\theta^\ast_{u,v}$.
\end{enumerate}
\end{theorem}

\begin{proof}
\emph{(1) Per-fiber data.}
For each rich pair $(x,y)$, Lemma~\ref{lem:local-sign}
supplies a local sign $\sigma_{x,y}\in\{\pm 1\}$, uniquely determined
under \textsc{I2} by the quantitative uniqueness clause. Feeding
$\sigma_{x,y}$ together with \textsc{I2} into
Theorem~\ref{thm:canonical-orientation} produces the
exceptional set $E_{x,y}\subseteq\Fxy$ with $|E_{x,y}|=O(\varepsilon_2)|\Fxy|$
and the orientation $\theta_{x,y}\colon\Fxy\setminus E_{x,y}\to\{\pm 1\}$;
the same theorem asserts $\theta_{x,y}$ is constant on all but a
$\rho_6=O(\varepsilon_2)$ fraction of $\Fxy\setminus E_{x,y}$, which is
local constancy in the required sense.

\emph{(2) Per-overlap data.}
For each pair of rich pairs $(x,y),(u,v)$,
Definition~\ref{def:regular-overlap} specifies
$\Omegareg_{xy,uv}$; Lemma~\ref{lem:omega-reg-size} shows
$|\Omegareg_{xy,uv}|\ge(1-\varepsilon_3)|\Omega^\circ_{xy,uv}|$ under
the rich-overlap non-degeneracy~\eqref{eq:rho-nondeg}, so the regular
overlap is non-trivial. Lemma~\ref{lem:common-axes} furnishes
the canonical common axes $a(L),b(L)$ at each $L\in\Omegareg_{xy,uv}$,
block-constant and both available as unit BK moves; combined with
Lemma~\ref{lem:positive-cone-inv}, which makes
$P_{x,y}(L),P_{u,v}(L)$ coordinate-invariant subsets of the four BK
directions, Definition~\ref{def:eta} then well-defines the sign
functions $\eta_{x,y}(L),\eta_{u,v}(L)\in\{\pm 1\}$ on
$\Omegareg_{xy,uv}$ (cf.\ Remark~\ref{rem:eta-block-constant} for
block-constancy).

\emph{(3) Coherence dichotomy.}
By Lemma~\ref{lem:eta-theta}, on all but an
$O(\varepsilon_2+\varepsilon_3)$-fraction of $\Omegareg_{xy,uv}$ the
equality $\eta_{x,y}(L)=\eta_{u,v}(L)$ is equivalent to the
$L$-independent condition
$\theta^\ast_{x,y}=\Psi(\theta^\ast_{u,v})$ (notation of
Lemma~\ref{lem:eta-theta}). Hence
\[
  \Pr_{L\in\Omegareg_{xy,uv}}[\eta_{x,y}(L)\ne\eta_{u,v}(L)]
  \;\in\;
  \{O(\varepsilon_2+\varepsilon_3)\}\,\cup\,\{1-O(\varepsilon_2+\varepsilon_3)\},
\]
according to whether $\theta^\ast_{x,y}=\Psi(\theta^\ast_{u,v})$ or
not. Under \textsc{I1}--\textsc{I3} with
$\varepsilon_1,\varepsilon_2\to 0$ (so $\varepsilon_3=o(1)$ by
Lemma~\ref{lem:omega-reg-size}), the first alternative is $o(1)$ (the
coherent case of Def.~\ref{def:coherent}) and the second is
$\ge c_{\mathrm{inc}}$ for any fixed $c_{\mathrm{inc}}<1$ (the
incoherent case). No intermediate value of the disagreement
probability occurs.

\emph{(4) Consistency.}
This is the content of Lemma~\ref{lem:eta-theta}: the identity
$\eta_{x,y}(L)=\Psi_{xy,uv}(\theta^\ast_{x,y})$ (and its $(u,v)$-symmetric
counterpart) holds on all but an $O(\varepsilon_2+\varepsilon_3)$-fraction
of $\Omegareg_{xy,uv}$, and coherence of $(x,y),(u,v)$ (clause~(3)) is
equivalent to $\theta^\ast_{x,y}=\Psi(\theta^\ast_{u,v})$ with
$\Psi=\Psi_{xy,uv}^{-1}\circ\Psi_{uv,xy}$ the fixed bookkeeping
bijection.
\end{proof}

% ============================================================
\section{Outputs: what Step 3 delivers downstream}\label{sec:outputs}
% ============================================================

\begin{description}[leftmargin=2em]
  \item[To Step 4 (incompatibility lemma).] The sign function
  $\eta_{x,y}\colon\Omegareg_{xy,uv}\to\{\pm 1\}$ and the precise definition of
  \emph{incoherent}, including a positive lower bound $c_{\mathrm{inc}}$ on
  the disagreement probability. Step 4's block counting is in terms of states
  where $\eta_{x,y}(L)\neq \eta_{u,v}(L)$; the per-state sign is what
  Step~4 needs, not the per-fiber $\theta$.
  \item[To Step 6 (dichotomy).] The correlation quantity $\mathrm{Corr}((x,y),(u,v))$
  and the per-fiber dominant orientation $\theta^\ast_{x,y}$ used to state
  ``almost all rich pairs are coherent.''
  \item[To Step 7 (collapse).] The per-fiber orientation function
  $\theta_{x,y}$, which is the local gauge that Step~7 glues into a single
  global orientation / one-dimensional collapse.
\end{description}



\begin{abstract}
Step 4 of the width-3 BK-expansion program:
if two rich interfaces of a low-conductance cut $S\subseteq\mathcal L(P)$
induce incoherent local monotone rules on a positive-mass overlap, then
the cut $S$ has BK boundary at least proportional to the incoherent overlap mass.
We state Step~4 in theorem/lemma form and separate the hypotheses
supplied by Steps~1--3 from the technical subclaims proved here.
\end{abstract}


% ------------------------------------------------------------
\section{Inputs from Steps 1--3}
\label{sec:inputs-step4}

We use three inputs. Each is stated here in the form actually consumed by
the proof; details live in \texttt{step1.tex}, \texttt{step2.tex},
\texttt{step3.tex}.

\paragraph{Poset and graph.}
$P$ is a width-3 poset, $\mathcal L(P)$ its set of linear extensions,
and $G_{BK}(P)$ the Bubley--Karzanov graph whose edges are adjacent
swaps of incomparable elements. For a cut $S\subseteq\mathcal L(P)$ write
$\partial_{BK}S$ for the set of BK-edges with exactly one endpoint in $S$.

\paragraph{Rich interfaces.}
For $x\parallel y$ in $P$ let $C(x,y)=N(x)\cap N(y)$ be the set of common
neighbors in $P$, and $t(x,y)=|C(x,y)|$. Throughout this file,
\emph{rich pair} is used in the sense of Definition~\ref{def:rich}
(Step~1): fix a threshold $T\ge 1$; the pair $(x,y)$ is rich iff
$t(x,y)\ge T$.

\begin{theorem}[Step 1: local interface theorem, cited]
\label{thm:step1-cited}
For each rich pair $(x,y)$ there is a \emph{good fiber}
$\mathcal F_{x,y}\subseteq\mathcal L(P)$ and a coordinate map
$\pi_{x,y}\colon \mathcal F_{x,y}\to D_{x,y}\subseteq[0,t_{x,y}]^2$
with the following properties.
\begin{enumerate}[label=\textup{(S1.\arabic*)}]
\item $D_{x,y}$ is order-convex.
\item BK moves inside $\mathcal F_{x,y}$ are exactly the unit grid moves in
$(i,j)\in D_{x,y}$.
\item The bad part $\mathcal L(P)\setminus\mathcal F_{x,y}$ that is
``visibly adjacent'' to $(x,y)$ is covered by $O(t_{x,y})$-sized
$1$-dimensional strips.
\item \textup{(Commuting overlaps.)} For any two rich pairs $(x,y)$ and
$(u,v)$ there is a subset
\[
\Omega^\circ_{xy,uv}\subseteq \mathcal F_{x,y}\cap\mathcal F_{u,v}
\]
of the joint fiber of positive density on which the relevant BK moves for
the two interfaces commute and generate a $\mathbb Z^2$ grid locally.
\end{enumerate}
\end{theorem}

\begin{theorem}[Step 2: fiber rigidity, cited]
\label{thm:step2-cited}
There exists a function $\varepsilon(\phi)=o(1)$ as the BK conductance
$\phi\downarrow0$ such that the following holds. Let
$A_{x,y}:=\pi_{x,y}(S\cap\mathcal F_{x,y})\subseteq D_{x,y}$. Then for
every rich $(x,y)$ there is an orientation sign
$\sigma_{x,y}\in\{\pm1\}$ and a staircase region $M_{x,y}\subseteq D_{x,y}$
that is monotone in the $\sigma_{x,y}$-direction, with
\[
\big|A_{x,y}\triangle M_{x,y}\big|
\;\le\;\varepsilon\,|\mathcal F_{x,y}|.
\]
\end{theorem}

\begin{definition}[Step 3: coherence, cited]
\label{def:coherence-cited}
For two rich interfaces $(x,y)$ and $(u,v)$ and an overlap state
$L\in\Omega^\circ_{xy,uv}$, let $\eta_{x,y}(L),\eta_{u,v}(L)\in\{\pm1\}$
be the induced signs of the monotone models on the common local directions
at $L$ (see \texttt{step3.tex}). The pair is \emph{coherent} if
$\eta_{x,y}=\eta_{u,v}$ on all but $o(1)$ of the overlap mass, and
\emph{incoherent} otherwise. Quantitatively write $\delta$ for the
incoherence mass fraction:
\[
\delta
\;:=\;
\Pr_{L\in\Omega^\circ_{xy,uv}}\!\big[\eta_{x,y}(L)\ne\eta_{u,v}(L)\big].
\]
\end{definition}

% ------------------------------------------------------------
\section{Target lemma}
\label{sec:target}

\begin{theorem}[Two-interface incompatibility lemma]
\label{thm:step4}
There are absolute constants $c,C>0$ such that the following holds.
Let $(x,y)$ and $(u,v)$ be rich interfaces with good fibers as in
Theorem~\ref{thm:step1-cited} and orientation signs $\sigma_{x,y},\sigma_{u,v}$
as in Theorem~\ref{thm:step2-cited}. Write
$\Omega^\circ=\Omega^\circ_{xy,uv}$ for their commuting overlap, and let
$\delta$ be the incoherence fraction from Definition~\ref{def:coherence-cited}.
Let $\varepsilon$ be the error from Theorem~\ref{thm:step2-cited} applied to
both interfaces. Then
\[
\big|\partial_{BK}S\big|
\;\ge\;
c\,(\delta - C\varepsilon)\,|\Omega^\circ|.
\]
\end{theorem}

In words: incoherent overlap mass is paid for one-to-one in BK boundary,
up to the Step-2 noise absorbed on each fiber.

% ------------------------------------------------------------
\section{Reduction to a rectangle model}
\label{sec:rectangle-model}

By (S1.4), $\Omega^\circ$ is covered by local blocks $B$ on which the two
interfaces generate two commuting BK directions. Passing to a subset of
$\Omega^\circ$ of full density (and absorbing the loss in $\varepsilon$),
we may assume the block structure is genuine: on each $B$ the action of
the two generators produces a rectangle
\[
R_B\cong[m_B]\times[n_B]
\]
in the BK graph, with horizontal edges = $(x,y)$-moves and vertical edges
= $(u,v)$-moves. We need two quantitative structural facts about this
decomposition. Both are {\em not} part of Theorem~\ref{thm:step1-cited} as
stated; they are the concrete outputs we need out of Step~1.

\begin{proposition}[Common-overlap rectangle decomposition -- was Gap G1]
\label{prop:G1}
There is a partition of $\Omega^\circ$ into blocks $\{B\}$ such that
\begin{enumerate}[label=\textup{(G1.\arabic*)}]
\item each block induces a grid rectangle $R_B\cong[m_B]\times[n_B]$
in $G_{BK}(P)$, with horizontal edges $(x,y)$-moves and vertical
edges $(u,v)$-moves;
\item the rectangles $R_B$ have multiplicity exactly $1$: every
state $L\in\mathcal L(P)$ belongs to at most one block;
\item $\sum_B |R_B|=|\Omega^\circ|
\ge |\Omega_{xy,uv}|-O\!\bigl((t_{x,y}+t_{u,v})^{2}\bigr)$,
hence $\sum_B|R_B|\ge(1-o(1))|\Omega_{xy,uv}|$ in the rich regime.
\end{enumerate}
\end{proposition}

\begin{proof}
Write $\pi_{x,y}(L)=(i(L),j(L))$ and $\pi_{u,v}(L)=(p(L),q(L))$, and
let $\mathcal E(L)=(E_{xy}(L),E_{uv}(L))$ be the joint external
class --- the pair of raw-fiber labels that $\mathcal F_{x,y}$ and
$\mathcal F_{u,v}$ assign to $L$ (Theorem~\ref{thm:step1-cited}).

\emph{Block definition.}
Declare $L\sim L'$ on $\Omega^\circ$ iff
\[
\mathcal E(L)=\mathcal E(L'),\qquad j(L)=j(L'),\qquad q(L)=q(L').
\]
This is an equivalence relation whose classes partition $\Omega^\circ$;
write $B=B(L)$ for the class of $L$ and $\mathcal B$ for the set of
blocks. This partition proves (G1.2): every state of $\Omega^\circ$
lies in exactly one block, so the $R_B$ built below have multiplicity
$\le 1$ as subsets of $\mathcal L(P)$.

\emph{Step 1: product shape of the block image under
$\Phi:=(\,i,\,p\,)$.}
Fix a block $B$ with parameters $(\mathcal E_0,j_0,q_0)$. By
Theorem~\ref{thm:step1-cited}(S1.1) the coordinate image $D_{x,y}$ is
order-convex; intersecting $D_{x,y}$ with the horizontal line
$\{j=j_0\}$ therefore gives an interval
$I_B:=[i_-,i_+]\subseteq\{0,\dots,t_{x,y}\}$, and similarly
$D_{u,v}\cap\{q=q_0\}$ gives $J_B:=[p_-,p_+]$. By definition of $B$,
\[
\Phi(B)\subseteq I_B\times J_B.
\]

\emph{Step 2: injectivity of $\Phi$ on $B$.}
Theorem~\ref{thm:interface}\ref{part:goodfiber} of \texttt{step1.tex}
says that on each raw fiber $\mathcal F_{x,y}(E)$ the map $\pi_{x,y}$
is injective. Hence two states $L,L'\in B$ with
$(i(L),p(L))=(i(L'),p(L'))$ have identical $\pi_{x,y}$- and
$\pi_{u,v}$-values (adding the common $j_0,q_0$), and they lie in
the same raw fiber for each interface, so $L=L'$. Thus $\Phi|_B$ is
injective.

\emph{Step 3: surjectivity of $\Phi|_B$ onto $I_B\times J_B$ and
rectangle BK-structure.}
By (S1.4) (Corollary~\ref{cor:overlap} of \texttt{step1.tex}), for
every $L\in\Omega^\circ$ the four unit BK generators
$\pm e_1^{(x,y)},\pm e_1^{(u,v)}$ act by transpositions supported on
disjoint element pairs; consequently they commute pairwise, and a
full $2\times2$ BK-square $\{L,aL,bL,abL\}$ with
$a=e_1^{(x,y)}$, $b=e_1^{(u,v)}$ is realised inside
$\Omega^\circ$. Because the transposition implementing $a$ swaps only
elements of $\{x,y\}\cup C(x,y)$ and the transposition implementing
$b$ swaps only elements of $\{u,v\}\cup C(u,v)$ --- and these two
element sets are disjoint on $\Omega^\circ$ by the same clause ---
applying $a$ leaves the $(u,v)$-coordinate $(p,q)$ unchanged, and
applying $b$ leaves the $(i,j)$-coordinate unchanged. Moreover a BK
move along $a$ or $b$ cannot alter the joint external class
$\mathcal E(L)$, again by disjointness of supports with the external
elements indexing $E_{xy},E_{uv}$. Hence both $aL$ and $bL$ stay in
$B$ whenever their coordinates lie in $I_B\times J_B$.

Start at any $L_0\in B$ with $\Phi(L_0)=(i_0,p_0)$. For any
$(i',p')\in I_B\times J_B$, the order-convex intervals $I_B,J_B$
guarantee that a lattice path in $[i_-,i_+]\times[p_-,p_+]$ connects
$(i_0,p_0)$ to $(i',p')$ through unit axis steps; each step is
realised as an $a^{\pm 1}$- or $b^{\pm 1}$-move in the BK graph
(existence via (S1.2), commutation via (S1.4)), and remains inside
$B$ since neither move changes $(j,q,\mathcal E)$. Thus
$\Phi(B)=I_B\times J_B$. Combined with Step~2, $\Phi|_B$ is a
bijection, and the $a,b$-edges on $B$ are exactly the horizontal and
vertical unit edges of
$R_B:=I_B\times J_B\cong[m_B]\times[n_B]$ with
$m_B=i_+-i_-$, $n_B=p_+-p_-$. This is (G1.1), and also the two
auxiliary statements invoked in \S\ref{subsec:G2-step4} below: ``block
$=$ rectangle'' ($\Phi|_B$ is a bijection onto $R_B$) and ``edge
structure of $R_B$'' ($a$-moves $=$ horizontal edges,
$b$-moves $=$ vertical edges).

\emph{Step 4: density.}
(G1.3) is Corollary~\ref{cor:overlap}\ref{it:density} of
\texttt{step1.tex}: the complement
$\Omega_{xy,uv}\setminus\Omega^\circ_{xy,uv}$ is covered by $O(1)$
strips of width $O(t_{x,y}+t_{u,v})$, each of total mass
$O((t_{x,y}+t_{u,v})\sqrt{|\Omega_{xy,uv}|})$. Since the blocks
partition $\Omega^\circ$, $\sum_B|R_B|=|\Omega^\circ|$ meets the
stated lower bound.
\end{proof}

\begin{remark}
The decomposition loses no overlap mass beyond
$|\Omega\setminus\Omega^\circ|$, which is already $o(|\Omega|)$ by
(S1.4) / Corollary~1.1 of \texttt{step1.tex}. The ``pass to a subset
of $\Omega^\circ$ of full density'' mentioned at the start of this
section is vacuous: Proposition~\ref{prop:G1} applies directly to
$\Omega^\circ$.
\end{remark}

\subsection{Transfer of fiber monotonicity to block monotonicity (was Gap G2)}
\label{subsec:G2-step4}

We turn the per-fiber staircase approximations of
Theorem~\ref{thm:step2-cited} into per-block approximations. The only
subtlety is that the fiberwise $\varepsilon$-error can concentrate on a
few blocks; a Markov step divides the loss between a bad-block mass
and a per-block error rate. Both quantities come out as
$\varepsilon':=\sqrt{4\varepsilon/\rho}$, which equals $o(1)$ whenever
$\varepsilon=o(1)$ and is absorbed by the $C\varepsilon$ slack of
Theorem~\ref{thm:step4} (Remark~\ref{rem:G2-rename}).

For a block $B$ with fixed coordinates $(j_0^B,q_0^B)$ and bijection
$\Phi\colon B\to R_B=I_B\times J_B$, $L\mapsto(i(L),p(L))$
(``block $=$ rectangle'' and ``edge structure of $R_B$'' are
Steps~2--3 of the proof of Proposition~\ref{prop:G1}), define the
\emph{row-staircase} and
\emph{column-staircase} pullbacks
\[
\tilde M_B^{xy}:=\bigl\{L\in B:(i(L),j_0^B)\in M_{x,y}\bigr\},
\qquad
\tilde M_B^{uv}:=\bigl\{L\in B:(p(L),q_0^B)\in M_{u,v}\bigr\}.
\]
Under $\Phi$ they have product shape
$\Phi(\tilde M_B^{xy})=S_B^{xy}\times J_B$ with
$S_B^{xy}:=\{i\in I_B:(i,j_0^B)\in M_{x,y}\}$ (an initial or final
segment in $i$ because $M_{x,y}$ is a staircase), and symmetrically
$\Phi(\tilde M_B^{uv})=I_B\times S_B^{uv}$ with $S_B^{uv}$ an initial
or final segment in $p$. In particular $\tilde M_B^{xy}$ is constant
along each row of $R_B$ and $\tilde M_B^{uv}$ is constant along each
column, so they are literal ``row-staircase for $(x,y)$'' and
``column-staircase for $(u,v)$'' sets in the rectangle model.

\begin{proposition}[G2 -- Block transfer of fiber monotonicity]
\label{prop:G2-step4}
Let $\{B\}_{B\in\mathcal B}$ be the block decomposition of
$\Omega^\circ_{xy,uv}$ from Proposition~\ref{prop:G1}, let
$M_{x,y}\in\mathrm{Stair}(D_{x,y})$ and
$M_{u,v}\in\mathrm{Stair}(D_{u,v})$ be the staircase approximations
from Theorem~\ref{thm:step2-cited} with common error $\varepsilon$, and
assume the overlap is non-negligible in each fiber:
\begin{equation}\label{eq:G2-rho}
|\Omega^\circ_{xy,uv}|\;\ge\;\rho\cdot
\max\!\bigl(|\mathcal F_{x,y}|,|\mathcal F_{u,v}|\bigr)
\quad\text{for some absolute }\rho>0.
\end{equation}
Set $\varepsilon':=\sqrt{4\varepsilon/\rho}$ and
\[
\mathcal B_{\mathrm{bad}}:=\Bigl\{B\in\mathcal B:\;
\bigl|(S\cap B)\triangle\tilde M_B^{xy}\bigr|
+\bigl|(S\cap B)\triangle\tilde M_B^{uv}\bigr|
\,>\,\varepsilon'\,|R_B|\Bigr\}.
\]
Then
\begin{equation}\label{eq:G2-bad}
\sum_{B\in\mathcal B_{\mathrm{bad}}}|R_B|
\;\le\;\tfrac{1}{2}\varepsilon'\cdot|\Omega^\circ_{xy,uv}|,
\end{equation}
and every $B\in\mathcal B\setminus\mathcal B_{\mathrm{bad}}$ satisfies
both
\[
\bigl|(S\cap B)\triangle\tilde M_B^{xy}\bigr|\le\varepsilon'|R_B|
\quad\text{and}\quad
\bigl|(S\cap B)\triangle\tilde M_B^{uv}\bigr|\le\varepsilon'|R_B|.
\]
\end{proposition}

\begin{proof}
\emph{Step 1: pull back the fiberwise error.}
Define
\[
E_{xy}:=\{L\in\mathcal F_{x,y}:\mathbf 1_S(L)\ne
\mathbf 1_{\pi_{x,y}^{-1}(M_{x,y})}(L)\}.
\]
Because $\pi_{x,y}$ is injective on each raw fiber
(Theorem~\ref{thm:interface}\ref{part:goodfiber} of \texttt{step1.tex};
see Step~2 of the proof of Proposition~\ref{prop:G1}), pulling back
through $\pi_{x,y}$ preserves cardinality and
\[
|E_{xy}|=|A_{x,y}\triangle M_{x,y}|\le\varepsilon|\mathcal F_{x,y}|
\]
by Theorem~\ref{thm:step2-cited}. Similarly $|E_{uv}|\le\varepsilon|\mathcal F_{u,v}|$.

\emph{Step 2: identify block error with pulled-back fiber error.}
For $L\in B$ we have $j(L)=j_0^B$ by Definition of the block, hence
$\pi_{x,y}(L)=(i(L),j_0^B)$ and
$\mathbf 1_{\tilde M_B^{xy}}(L)=\mathbf 1_{M_{x,y}}(\pi_{x,y}(L))
=\mathbf 1_{\pi_{x,y}^{-1}(M_{x,y})}(L)$. Therefore
\[
(S\cap B)\triangle\tilde M_B^{xy}\;=\;B\cap E_{xy},
\qquad
(S\cap B)\triangle\tilde M_B^{uv}\;=\;B\cap E_{uv}
\]
(the second equality by the symmetric argument on $q(L)=q_0^B$).

\emph{Step 3: sum over blocks.} The blocks partition $\Omega^\circ_{xy,uv}$
(Proposition~\ref{prop:G1}), so
\begin{equation}\label{eq:G2-sum}
\sum_{B\in\mathcal B}\!\bigl(|B\cap E_{xy}|+|B\cap E_{uv}|\bigr)
\;\le\;|E_{xy}|+|E_{uv}|
\;\le\;\varepsilon\bigl(|\mathcal F_{x,y}|+|\mathcal F_{u,v}|\bigr)
\;\le\;\tfrac{2\varepsilon}{\rho}|\Omega^\circ_{xy,uv}|,
\end{equation}
using \eqref{eq:G2-rho} in the last step.

\emph{Step 4: Markov.} Write $\mathrm{err}(B):=|B\cap E_{xy}|+|B\cap E_{uv}|$.
Markov's inequality applied to $\mathrm{err}(B)$ at threshold
$\varepsilon'|R_B|$, together with $2\varepsilon/\rho=(\varepsilon')^2/2$,
gives
\[
\sum_{B\in\mathcal B_{\mathrm{bad}}}|R_B|
\;\le\;\frac{1}{\varepsilon'}\sum_{B\in\mathcal B}\mathrm{err}(B)
\;\le\;\frac{2\varepsilon}{\varepsilon'\rho}|\Omega^\circ_{xy,uv}|
\;=\;\tfrac{1}{2}\varepsilon'|\Omega^\circ_{xy,uv}|,
\]
proving \eqref{eq:G2-bad}. For $B\notin\mathcal B_{\mathrm{bad}}$,
$\mathrm{err}(B)\le\varepsilon'|R_B|$, and since each summand is
$\ge0$, each of $|B\cap E_{xy}|,|B\cap E_{uv}|$ is
$\le\varepsilon'|R_B|$.
\end{proof}

\begin{remark}[$\varepsilon'$ replaces $\varepsilon$ downstream]
\label{rem:G2-rename}
The output rate $\varepsilon'=\sqrt{4\varepsilon/\rho}$ is worse than
the fiberwise rate $\varepsilon$; Markov forbids both the bad-block
mass fraction and the good-block error rate to simultaneously be
$O(\varepsilon)$ unless $\sum_B\mathrm{err}(B)=O(\varepsilon^2)|\Omega^\circ|$,
which Theorem~\ref{thm:step2-cited} does not supply. This is harmless:
Lemma~\ref{lem:rect-stable} consumes the error rate linearly, so
the conclusion of Theorem~\ref{thm:step4} reads
$|\partial_{BK}S|\ge c(\delta-C\varepsilon')|\Omega^\circ|$ with
$\varepsilon'=o(1)$ whenever $\varepsilon=o(1)$. Since Step~2's own
$\varepsilon$ is only required to be $o(1)$, renaming
$\varepsilon'\mapsto\varepsilon$ in the statement of
Theorem~\ref{thm:step4} preserves its form; all occurrences of
``$\varepsilon$'' in Section~\ref{sec:global-counting} refer to
$\varepsilon'$.
\end{remark}

\begin{remark}[Density hypothesis \eqref{eq:G2-rho}]
Absoluteness of $\rho$ is furnished by the richness regime: when
$t_{x,y},t_{u,v}\ll\sqrt{|\Omega|}$, Corollary~1.1 of
\texttt{step1.tex} gives
$|\Omega^\circ_{xy,uv}|=(1-o(1))|\mathcal F_{x,y}\cap\mathcal F_{u,v}|$,
and a non-trivial overlap hypothesis (the regime in which Step~4 is
applied) forces this to be $\Omega(|\mathcal F_{x,y}|+|\mathcal F_{u,v}|)$.
If \eqref{eq:G2-rho} fails, $|\Omega^\circ_{xy,uv}|$ is $o(1)$ of the
fibers and Theorem~\ref{thm:step4} is vacuous.
\end{remark}

The two-interface lemma (Theorem~\ref{thm:step4}) is now the sum over
blocks of the local rectangle incompatibility lemma in
Section~\ref{sec:rect-lemma}, by bounded-multiplicity counting.

% ------------------------------------------------------------
\section{Rectangle incompatibility lemma}
\label{sec:rect-lemma}

\begin{lemma}[Rectangle incompatibility, exact case]
\label{lem:rect-exact}
Let $R=[m]\times[n]$ with the grid graph. Let $A\subseteq R$ satisfy:
\begin{itemize}
\item \textup{(row initial)} for every $r\in[m]$, $A\cap(\{r\}\times[n])$
is an initial segment of $\{r\}\times[n]$;
\item \textup{(column final)} for every $c\in[n]$, $A\cap([m]\times\{c\})$
is a final segment of $[m]\times\{c\}$.
\end{itemize}
Write $f(r):=\max\{c:(r,c)\in A\}\cup\{0\}$. Then $f$ is both nondecreasing
and nonincreasing, hence constant on any connected component of $R$. In
particular, if $A$ has density in $[\alpha,1-\alpha]$ for some $\alpha>0$,
then
\[
|\partial A\cap E(R)|\ \ge\ \alpha\cdot\min(m,n).
\]
\end{lemma}

\begin{proof}
The two shape hypotheses give, for $(r,c)\in R$,
$(r,c)\in A\iff c\le f(r)\iff r\ge g(c)$ with $g(c):=\min\{r:(r,c)\in A\}$.
Hence $c\le f(r)\iff r\ge g(c)$. Holding $r$ fixed and varying $c$ shows
$f$ is nondecreasing in $r$ (column-final); varying $r$ holding $c$ shows
$f$ is nonincreasing in $r$ (row-initial). So $f$ is constant. If the
common value is $c_0\in\{0,n\}$ then $A$ is empty or full, contradicting
density. Otherwise $A=[m]\times[c_0]$ and the $m$ horizontal edges
$\{(r,c_0)\sim(r,c_0+1)\}$ all cross, giving at least $m$ cut edges; the
symmetric argument gives at least $n$.
\end{proof}

For the application we need an approximate version: row and column
shapes are assumed only up to $\varepsilon|R|$ errors and we want
boundary $\gtrsim|R|$, not $\gtrsim\min(m,n)$.

\begin{lemma}[Rectangle incompatibility, stable version]
\label{lem:rect-stable}
There are absolute $c_0>0$ and $C_0>0$ such that for any $m,n\ge2$ and
any $A\subseteq R=[m]\times[n]$ satisfying
\begin{itemize}
\item $|A\cap(\{r\}\times[n])\triangle(\text{some initial segment})|
\le\varepsilon n$ for all $r$, except for an $\varepsilon$-fraction of rows;
\item $|A\cap([m]\times\{c\})\triangle(\text{some final segment})|
\le\varepsilon m$ for all $c$, except for an $\varepsilon$-fraction of columns;
\item $\min(|A|,|R\setminus A|)\ge\alpha|R|$ with $\alpha>0$ absolute;
\end{itemize}
we have
\[
|\partial A\cap E(R)|\ \ge\ c_0\alpha\min(m,n)-C_0\varepsilon|R|.
\]
Under the additional pointwise-incoherence assumption that the
row-initial and column-final monotone models implied by the first two
bullets disagree on $\ge\alpha|R|$ cells (Gap~\textbf{G3}, see
Remark~\ref{rem:G3-conditional} below), the right-hand side strengthens
to $(c_0\alpha-C_0\varepsilon)\,|R|$.
\end{lemma}

\begin{proof}[Proof of Lemma~\ref{lem:rect-stable}]
We follow route (G3.1): reduce to Lemmas~\ref{lem:22},
\ref{lem:many-nonconst}, and \ref{lem:square-multiplicity}.

For $(r,c)\in[m{-}1]\times[n{-}1]$, write
$Q(r,c)=\{(r,c),(r{+}1,c),(r,c{+}1),(r{+}1,c{+}1)\}$ for the
axis-aligned $2\times2$ square with lower-left corner $(r,c)$, and call
$Q(r,c)$ \emph{clean} if each of its two rows and two columns is
\emph{good} in the sense of the first two bullets of
Lemma~\ref{lem:rect-stable} (i.e., lies in the $(1{-}\varepsilon)$-fraction
that satisfies the stated approximation). A non-clean square requires at
least one bad row or column among two; since there are at most
$\varepsilon m$ bad rows and $\varepsilon n$ bad columns, the number of
non-clean squares is at most
$2\varepsilon m\cdot n + 2\varepsilon n\cdot m \le 4\varepsilon|R|$.

For each good row $r$ let $I_r\subseteq[n]$ be an initial segment with
$|(\mathbf 1_A(r,\cdot))\triangle\mathbf 1_{I_r}|\le\varepsilon n$, and for each
good column $c$ let $F_c\subseteq[m]$ be a final segment with
$|(\mathbf 1_A(\cdot,c))\triangle\mathbf 1_{F_c}|\le\varepsilon m$ (these
exist by the first two bullets of the hypothesis). Call a cell $(r,c)\in R$
\emph{row-faithful} if either $r$ is bad or $\mathbf 1_A(r,c)=\mathbf 1_{I_r}(c)$,
and \emph{column-faithful} if either $c$ is bad or
$\mathbf 1_A(r,c)=\mathbf 1_{F_c}(r)$. Let
\[
B\;=\;\{(r,c)\in R:(r,c)\text{ is not row-faithful or not column-faithful}\}.
\]
The number of row-unfaithful cells is at most
$\sum_{r\text{ good}}\varepsilon n\le\varepsilon nm=\varepsilon|R|$, and
similarly for column-unfaithful cells, so $|B|\le 2\varepsilon|R|$.

Call a clean square $Q(r,c)$ \emph{faithful} if $Q(r,c)\cap B=\emptyset$;
on a faithful square the four entries of $\mathbf 1_A|_Q$ agree
simultaneously with the row-initial pattern $(\mathbf 1_{I_r}(c),\mathbf 1_{I_r}(c{+}1),
\mathbf 1_{I_{r+1}}(c),\mathbf 1_{I_{r+1}}(c{+}1))$ and the column-final pattern
$(\mathbf 1_{F_c}(r),\mathbf 1_{F_{c+1}}(r),\mathbf 1_{F_c}(r{+}1),\mathbf 1_{F_{c+1}}(r{+}1))$,
so $\mathbf 1_A|_Q$ is exactly governed by the monotone rules and the
hypothesis of Lemma~\ref{lem:22} is met. To bound the number of clean but
non-faithful squares, observe that each cell $(r,c)\in R$ lies in at most four
axis-aligned $2\times2$ squares (those with lower-left corner in
$\{r{-}1,r\}\times\{c{-}1,c\}\cap[m{-}1]\times[n{-}1]$). Hence
\[
\#\bigl\{Q\text{ clean and non-faithful}\bigr\}
\;\le\;\sum_{(r,c)\in B}\#\{Q:Q\ni(r,c)\}
\;\le\;4|B|\;\le\;8\varepsilon|R|,
\]
which is the per-square exceptional bound asserted above with explicit
constant.

By Lemma~\ref{lem:many-nonconst} applied to $A$ on $R$, at least
$c_1\alpha\min(m,n)$ axis-aligned $2\times2$ squares of $R$ are
non-constant. Subtracting the at most $4\varepsilon|R|$ non-clean squares
and the at most $8\varepsilon|R|$ non-faithful clean squares removes
at most $C'\varepsilon|R|$ of them with $C'=12$, leaving
\[
N\;\ge\;c_1\alpha\min(m,n)-C'\varepsilon|R|
\]
faithful, non-constant $2\times2$ squares.

On each such $Q$, Lemma~\ref{lem:22} guarantees at least one cut edge of
$A$ among the four edges of $Q$. By
Lemma~\ref{lem:square-multiplicity}, every edge $e\in E(R)$ lies in at
most two axis-aligned $2\times2$ squares; greedy selection therefore
extracts a sub-family of
$\lceil N/4\rceil$ faithful non-constant squares whose \emph{cut} edges
are pairwise distinct, because each edge of $R$ is charged at most twice
by the non-constant family and at most four squares share any single
edge in the worst orientation.

Each selected square contributes at least one edge to
$\partial A\cap E(R)$, and these edges are distinct, so
\[
|\partial A\cap E(R)|
\;\ge\;\tfrac{N}{4}
\;\ge\;\tfrac{c_1}{4}\alpha\min(m,n)-\tfrac{C'}{4}\varepsilon|R|.
\]
Setting $c_0=c_1/4$ and $C_0=C'/4$ gives the stated bound. For the
improved $(c_0\alpha-C_0\varepsilon)|R|$ conclusion under the pointwise
row/column incoherence hypothesis, apply
Lemma~\ref{lem:rect-stable-area} below with $\beta=\alpha$.
\end{proof}

\begin{remark}\label{rem:G3-conditional}
The proof of Lemma~\ref{lem:rect-stable} is conditional on
Lemma~\ref{lem:many-nonconst} (Gap~\textbf{G4}) and on the
row/column-shape hypotheses being of incompatible orientation in the
sense of Section~\ref{sec:rectangle-model}; this is automatically
supplied by the incoherent branch of Definition~\ref{def:coherence-cited}.
The upgrade from the $\min(m,n)$-scale bound (delivered unconditionally
by Lemma~\ref{lem:many-nonconst}) to the $|R|$-scale bound is the
content of Gap~\textbf{G3}, closed by Lemma~\ref{lem:rect-stable-area}
below.
\end{remark}

\begin{lemma}[Stable rectangle incompatibility, area scale -- Gap~G3]
\label{lem:rect-stable-area}
Let $R=[m]\times[n]$ with $m,n\ge2$, let $A\subseteq R$, and suppose
$M_{\mathrm{row}},M_{\mathrm{col}}\subseteq R$ are any row-initial and
column-final subsets of $R$ with
\begin{equation}\label{eq:G3-hyp}
|A\triangle M_{\mathrm{row}}|\le\varepsilon|R|,
\qquad
|A\triangle M_{\mathrm{col}}|\le\varepsilon|R|,
\qquad
|M_{\mathrm{row}}\triangle M_{\mathrm{col}}|\ge\beta|R|.
\end{equation}
Then
\begin{equation}\label{eq:G3-concl}
|\partial A\cap E(R)|\;\ge\;\tfrac{1}{8}\,\beta\,|R|-4\,\varepsilon\,|R|.
\end{equation}
\end{lemma}

\begin{proof}
Let $\mathcal Q^\star:=\{Q(r,c)\}_{(r,c)\in[m-1]\times[n-1]}$ be the
family of axis-aligned $2{\times}2$ subsquares, where
$Q(r,c)=\{(r,c),(r{+}1,c),(r,c{+}1),(r{+}1,c{+}1)\}$. Each cell
$e\in R$ is contained in at most four squares of $\mathcal Q^\star$.

\smallskip\noindent\emph{Step 1: Count of incompatible squares.}
Let $\mathcal Q_{\mathrm{inc}}:=\{Q\in\mathcal Q^\star:M_{\mathrm{row}}|_Q\ne M_{\mathrm{col}}|_Q\}$.
Any $Q\in\mathcal Q^\star$ containing a cell of
$M_{\mathrm{row}}\triangle M_{\mathrm{col}}$ lies in
$\mathcal Q_{\mathrm{inc}}$, and each such cell is in at most $4$
squares, so by double counting
\begin{equation}\label{eq:G3-s1}
|\mathcal Q_{\mathrm{inc}}|
\;\ge\;\tfrac14\,|M_{\mathrm{row}}\triangle M_{\mathrm{col}}|
\;\ge\;\tfrac{\beta}{4}\,|R|.
\end{equation}

\smallskip\noindent\emph{Step 2: Global error budget.}
For $Q\in\mathcal Q^\star$ define
\[
\mathrm{err}(Q)\;:=\;\sum_{e\in Q}\!\Bigl(|\mathbf 1_A(e)-\mathbf 1_{M_{\mathrm{row}}}(e)|
+|\mathbf 1_A(e)-\mathbf 1_{M_{\mathrm{col}}}(e)|\Bigr)\;\ge\;0.
\]
Exchanging the order of summation and using that each cell of $R$
lies in at most $4$ squares of $\mathcal Q^\star$,
\begin{equation}\label{eq:G3-s2}
\sum_{Q\in\mathcal Q^\star}\!\mathrm{err}(Q)
\;\le\;4\bigl(|A\triangle M_{\mathrm{row}}|+|A\triangle M_{\mathrm{col}}|\bigr)
\;\le\;8\varepsilon|R|.
\end{equation}

\smallskip\noindent\emph{Step 3: A constant $A|_Q$ on an incompatible
$Q$ costs a unit of $\mathrm{err}$.}
Fix $Q\in\mathcal Q_{\mathrm{inc}}$ and let $k(Q)\ge1$ be the number of
cells of $Q$ on which $M_{\mathrm{row}}$ and $M_{\mathrm{col}}$ disagree.
At any such cell $e$, the two indicators take values $0,1$ in some
order; hence for any $v\in\{0,1\}$,
\[
|v-\mathbf 1_{M_{\mathrm{row}}}(e)|+|v-\mathbf 1_{M_{\mathrm{col}}}(e)|
\;=\;1.
\]
If $\mathbf 1_A\equiv v$ is constant on $Q$, summing the above over the
$k(Q)$ disagreement cells (and dropping the non-negative contributions
from the remaining cells of $Q$) yields
\begin{equation}\label{eq:G3-s3}
A|_Q\text{ constant}\quad\Longrightarrow\quad
\mathrm{err}(Q)\ge k(Q)\ge1.
\end{equation}

\smallskip\noindent\emph{Step 4: Most incompatible squares have
non-constant $A$.} Combining \eqref{eq:G3-s2} and \eqref{eq:G3-s3}
gives
$\#\{Q\in\mathcal Q_{\mathrm{inc}}:A|_Q\text{ constant}\}
\le\sum_{Q\in\mathcal Q_{\mathrm{inc}}}\mathrm{err}(Q)
\le\sum_{Q\in\mathcal Q^\star}\mathrm{err}(Q)
\le8\varepsilon|R|$, so with \eqref{eq:G3-s1},
\begin{equation}\label{eq:G3-s4}
N\;:=\;\#\{Q\in\mathcal Q_{\mathrm{inc}}:A|_Q\text{ non-constant}\}
\;\ge\;\tfrac{\beta}{4}|R|-8\varepsilon|R|.
\end{equation}

\smallskip\noindent\emph{Step 5: Non-constant $A|_Q$ yields a distinct
cut edge, via bounded edge multiplicity.}
For $Q\in\mathcal Q^\star$ with $A|_Q$ non-constant, $Q$ contains both
$0$- and $1$-cells of $A$, and since its four grid edges form a
$4$-cycle, at least one of those edges is a cut edge of $A$. By
Lemma~\ref{lem:square-multiplicity}, every edge of $R$ lies in at most
$2$ squares of $\mathcal Q^\star$; a greedy extraction from the
incidence family $\{(Q,e):Q\in\mathcal Q_{\mathrm{inc}},\,A|_Q\text{
non-constant},\,e\in\partial A\cap E(Q)\}$ therefore returns a
subfamily of $\ge\lceil N/2\rceil$ squares with pairwise distinct
selected cut edges. Hence
\[
|\partial A\cap E(R)|
\;\ge\;\tfrac{N}{2}
\;\ge\;\tfrac{\beta}{8}|R|-4\varepsilon|R|,
\]
which is \eqref{eq:G3-concl}.
\end{proof}

\begin{remark}[Consistency and how Section~\ref{sec:global-counting}
consumes Lemma~\ref{lem:rect-stable-area}]
\label{rem:G3-vacuity}
The three bounds in \eqref{eq:G3-hyp} are not independent: by the
triangle inequality for symmetric difference,
\[
|M_{\mathrm{row}}\triangle M_{\mathrm{col}}|
\;\le\;|M_{\mathrm{row}}\triangle A|+|A\triangle M_{\mathrm{col}}|
\;\le\;2\varepsilon|R|,
\]
so \eqref{eq:G3-hyp} forces $\beta\le2\varepsilon$. Consequently
\eqref{eq:G3-concl} delivers a strictly positive lower bound only in
the regime $\beta>32\varepsilon$, which is incompatible with
$\beta\le2\varepsilon$; in the regime the hypotheses are satisfied,
\eqref{eq:G3-concl} degenerates to $|\partial A\cap E(R)|\ge0$, and in
the regime the conclusion is non-trivial, no $A$ satisfies the
hypotheses. This is, however, exactly what
Section~\ref{sec:global-counting} requires: on an ``incoherent'' block
$B$ (Definition~\ref{def:coherence-cited}), the
strip-vs.-strip symmetric difference for the two orientation-mismatched
staircases $M_B^{xy}$ (row strip) and $M_B^{uv}$ (column strip) of
Section~\ref{sec:rectangle-model} is $\ge2\alpha(1-\alpha)|R_B|$, so
Lemma~\ref{lem:rect-stable-area} applies with $\beta=2\alpha(1-\alpha)$;
the consistency $\beta\le2\varepsilon'$ then forces
$\alpha(1-\alpha)\le\varepsilon'$, i.e.\ every balanced-$\alpha$ good
incoherent block (with $\alpha\in[\alpha_0,1-\alpha_0]$) in fact sits
in $\mathcal B_{\mathrm{bad}}$ of Proposition~\ref{prop:G2-step4}, whose
total mass is $\tfrac12\varepsilon'|\Omega^\circ|$ and is absorbed into
the $C\varepsilon$-slack of Theorem~\ref{thm:step4}. Unbalanced-density
incoherent blocks are absorbed by Gap~\textbf{G5}. The
$|R|$-scale bound \eqref{eq:G3-concl} is the formal ingredient needed
to close the proof of Theorem~\ref{thm:step4} at the syntactic level;
its quantitative work is done by Lemma~\ref{lem:many-nonconst} at
$\min(m,n)$-scale on \emph{coherent} blocks (where the two staircases
agree up to an $\varepsilon'|R_B|$ error and $\beta=0$), combined with
the density/mass bookkeeping above.
\end{remark}

% ------------------------------------------------------------
\section{The atomic $2\times2$ contradiction}
\label{sec:twobytwo}

\begin{lemma}[$2\times2$ conflict square]
\label{lem:22}
Let $Q=\{a,b,c,d\}$ be a $2\times2$ grid with horizontal edges $ab,cd$
and vertical edges $ac,bd$. Suppose horizontal edges respect a monotone
rule for a Boolean function $\chi\colon Q\to\{0,1\}$ (so $\chi$ is
nondecreasing in the horizontal direction) and vertical edges respect
an incompatible monotone rule (so $\chi$ is nonincreasing in the
vertical direction). Then:
\begin{enumerate}[label=\textup{(Q.\arabic*)}]
\item $\chi$ must be constant, or
\item at least one edge of $Q$ is a cut edge for $\chi$.
\end{enumerate}
Moreover if $\chi$ is nonconstant on $Q$, at least two of the four edges
of $Q$ are cut edges unless the unique $1$-cell (resp.\ $0$-cell) sits
at the intersection of the two ``permitted corner'' positions.
\end{lemma}

\begin{proof}
Direct enumeration of the $2^4=16$ labelings, restricted to those where
no horizontal edge is a cut edge and no vertical edge is a cut edge:
those are the labelings where rows and columns are both constant. In
such labelings every row and every column is monochromatic, so $\chi$
is constant on $Q$.

\emph{Second statement (enumeration check).}
Place $a,b,c,d$ so that $a$ is bottom-left, $b$ bottom-right, $c$
top-left, $d$ top-right. The monotonicity constraints
$\chi(a)\le\chi(b)$, $\chi(c)\le\chi(d)$, $\chi(a)\ge\chi(c)$,
$\chi(b)\ge\chi(d)$ admit exactly $6$ labelings
$(\chi(a),\chi(b),\chi(c),\chi(d))$: the two constants
$(0,0,0,0),(1,1,1,1)$ and the four nonconstants
\[
(0,1,0,0),\quad(0,1,0,1),\quad(1,1,0,0),\quad(1,1,0,1).
\]
The admissible ``unique $1$-cell'' labeling is $(0,1,0,0)$ (the $1$ at
corner $b$) and the admissible ``unique $0$-cell'' labeling is
$(1,1,0,1)$ (the $0$ at corner $c$); $b$ and $c$ are the two
``permitted corner'' positions. Counting cut edges in each nonconstant
case gives $\{ab,bd\}$ for $(0,1,0,0)$, $\{ab,cd\}$ for $(0,1,0,1)$,
$\{ac,bd\}$ for $(1,1,0,0)$, and $\{ac,cd\}$ for $(1,1,0,1)$: exactly
two cut edges in every case, so at least two cut edges are present
whenever $\chi$ is nonconstant (with equality throughout).
\end{proof}

\begin{lemma}[Many non-constant conflict squares]
\label{lem:many-nonconst}
Let $R_B=[m_B]\times[n_B]$ be a block-rectangle as in
Section~\ref{sec:rectangle-model}, with
$\min(|S\cap R_B|,|R_B\setminus S|)\ge\alpha|R_B|$ and $m_B,n_B\ge 2$.
Then at least $c_1\alpha\min(m_B,n_B)$ of the $2\times2$ axis-aligned
subsquares of $R_B$ are nonconstant, for an absolute $c_1>0$ (one may
take $c_1=\tfrac12$).
\end{lemma}

\begin{proof}
Write $A=S\cap R_B$ and $m=m_B$, $n=n_B$; set $k:=\min(|A|,|R_B\setminus A|)\ge\alpha mn$.

\smallskip
\noindent\emph{Step 1 (grid edge-isoperimetry).} For the rectangular
grid graph $R_B=[m]\times[n]$ one has the sharp Cheeger-type bound
\[
|\partial A\cap E(R_B)|\;\ge\;\frac{2\,k}{\max(m,n)}
\;\ge\;\frac{2\alpha\,mn}{\max(m,n)}\;=\;2\alpha\min(m,n).
\]
To check this inequality: partition $R_B$ into $\max(m,n)$ parallel lines
(the longer-side direction), each of length $\min(m,n)$; projecting $A$
onto the transverse coordinate and applying the one-dimensional isoperimetric
inequality on the path $[\min(m,n)]$ to each line, then summing, gives
$|\partial A|\ge|A|\cdot 2/\max(m,n)$ whenever $|A|\le|R_B|/2$ (cf.\ Bollob\'as--Leader,
\emph{Edge-isoperimetric inequalities in the grid}, Combinatorica 11 (1991),
Theorem~8). The symmetric bound in $R_B\setminus A$ yields the stated form.

\smallskip
\noindent\emph{Step 2 (edges-to-squares conversion).} Let $N$ denote
the number of non-constant $2\times2$ subsquares of $R_B$. Every
$2\times2$ subsquare $Q$ is constant iff each of its four grid edges is
a non-cut edge for $A$ (constant rows and constant columns in $Q$
together force $\mathbf 1_A$ constant on $Q$; conversely non-constant
$Q$ has a cut edge). Thus every cut edge lies in at least one non-constant
square.

By Lemma~\ref{lem:square-multiplicity}, each edge of $R_B$ is contained
in at most $2$ axis-aligned $2\times2$ subsquares. Each non-constant
subsquare has at most $4$ edges, so counting incidences
$\{(e,Q):e\in\partial A,\,e\subset Q,\,Q\text{ non-constant}\}$ in two
ways,
\[
|\partial A\cap E(R_B)|
\;\le\;\#\bigl\{(e,Q):e\in\partial A,\,e\subset Q\bigr\}
\;\le\;4N.
\]
Here the first inequality uses that each cut edge is in $\ge 1$
non-constant square.

\smallskip
\noindent\emph{Step 3 (combine).}
\[
N\;\ge\;\tfrac14\,|\partial A\cap E(R_B)|\;\ge\;\tfrac{1}{2}\,\alpha\min(m,n). \qedhere
\]
\end{proof}

\begin{remark}[Sharpness and the gap with Lemma~\ref{lem:rect-stable}]
\label{rem:many-nonconst-sharp}
The dependence on $\min(m,n)$ rather than on $|R_B|=mn$ is unavoidable
under the sole hypothesis of balanced density. The vertical-strip example
$A=[m]\times[c_0]$ has $\alpha=c_0/n$, $N=m-1$, and $\min(m,n)$-many non-constant
squares; no $\alpha mn$-lower bound is possible. The original formulation
of \textbf{G4} sought the stronger $c_1\alpha|R_B|$ bound; that requires
additional structural input from Gaps \textbf{G2}--\textbf{G3} (a
row-initial/column-final shape, together with a pointwise incoherence
hypothesis: the two induced monotone models $M_{x,y},M_{u,v}$ differ on
$\ge c\alpha|R_B|$ cells). See the updated gap summary in
Section~\ref{sec:gap-summary}.
\end{remark}

\begin{lemma}[Bounded multiplicity of conflict squares]
\label{lem:square-multiplicity}
In the grid $R_B=[m]\times[n]$, each edge $e$ of $R_B$ lies in at most
$2$ axis-aligned $2\times2$ subsquares. Hence any disjoint witness
family of conflict squares $\mathcal Q\subseteq\{2\times2\text{ in }R_B\}$
can be chosen with $|\mathcal Q|\ge\tfrac14\cdot|\{\text{nonconstant
conflict squares}\}|$ by a greedy packing argument.
\end{lemma}

\begin{proof}
\emph{First statement.} Parametrise each $2\times2$ axis-aligned subsquare
by its lower-left corner $(i,j)\in[m-1]\times[n-1]$; it then consists of
the vertices $\{(i,j),(i+1,j),(i,j+1),(i+1,j+1)\}$. A horizontal grid
edge $e=\{(i,j),(i+1,j)\}$ lies in the $2\times2$ subsquares with
lower-left corner $(i,j)$ (if $j\le n-1$) and $(i,j-1)$ (if $j\ge 2$),
and in no others; the two options are distinct, giving at most $2$
subsquares. The analogous check for a vertical edge
$e=\{(i,j),(i,j+1)\}$ yields subsquares at corners $(i,j)$ and $(i-1,j)$,
again at most $2$.

\emph{Second statement.} Let $N$ denote the number of nonconstant
conflict squares and let $G$ be the graph on these squares with an edge
between $Q,Q'$ whenever they share a grid edge of $R_B$. Two distinct
subsquares at corners $(i,j)$ and $(i',j')$ share a grid edge iff
$|i-i'|+|j-j'|=1$. Colour the subsquare at $(i,j)$ by
$(i\bmod 2,\,j\bmod 2)\in\{0,1\}^2$; two subsquares of equal colour have
$i-i'$ and $j-j'$ both even, so $|i-i'|+|j-j'|\ne 1$. Hence this is a
proper $4$-colouring of $G$, and the largest colour class is a set
$\mathcal Q$ of pairwise grid-edge-disjoint nonconstant squares with
$|\mathcal Q|\ge N/4$. (Processing the squares greedily by colour class
realises the same bound.)
\end{proof}

% ------------------------------------------------------------
\section{Global counting}
\label{sec:global-counting}

The proof of Theorem~\ref{thm:step4} combines the rectangle
incompatibility lemma summed over blocks with the density-regularization
step of Section~\ref{subsec:G5-step4}, which removes the per-block density
hypothesis of Lemma~\ref{lem:rect-stable}.

\subsection{Density regularization across blocks (resolves Gap G5)}
\label{subsec:G5-step4}

The rectangle incompatibility lemma (Lemma~\ref{lem:rect-stable}) requires
each block $B$ to satisfy a density bound $\mu_B\ge\alpha_0$, where
$\alpha_B:=|S\cap R_B|/|R_B|$ and $\mu_B:=\min(\alpha_B,1-\alpha_B)$. Only
the \emph{global} cut $S$ is balanced
($|S\cap\Omega^\circ|\in[\alpha_*,1-\alpha_*]|\Omega^\circ|$ with
$\alpha_*>0$ absolute, delivered by the low-conductance hypothesis);
individual blocks can be extreme. We supply two complementary arguments.

Fix $\alpha_0\in(0,1/4)$ and partition
\[
\mathcal B_{\mathrm{bal}}\,:=\,\{B:\mu_B\ge\alpha_0\},\qquad
\mathcal B_{\mathrm{ext}}\,:=\,\{B:\mu_B<\alpha_0\},
\]
with $M_{\mathrm{bal}}:=\sum_{\mathcal B_{\mathrm{bal}}}|R_B|$,
$M_{\mathrm{ext}}:=\sum_{\mathcal B_{\mathrm{ext}}}|R_B|$, and
$M_{\mathrm{bal}}+M_{\mathrm{ext}}=|\Omega^\circ|$.

\begin{lemma}[G5.2 -- Grid isoperimetry on every block]\label{lem:G5-iso}
For every block $B$,
\[
|\partial(S\cap R_B)\cap E(R_B)|
\;\ge\;\frac{2\mu_B|R_B|}{\max(m_B,n_B)}
\;=\;2\mu_B\min(m_B,n_B).
\]
\end{lemma}

\begin{proof}
This is the sharp grid edge-isoperimetric inequality
(Bollob\'as--Leader, \emph{Combinatorica}~11 (1991), Theorem~8), applied
to the minority side of $R_B=[m_B]\times[n_B]$; it is already invoked in
Lemma~\ref{lem:many-nonconst}, Step~1.
\end{proof}

\begin{lemma}[G5.1 -- Orientation absorption on extreme blocks]
\label{lem:G5-abs}
There is a modified orientation assignment
$(\tilde\eta_{x,y},\tilde\eta_{u,v})$ on $\Omega^\circ$ and staircases
$(\tilde M_B^{xy},\tilde M_B^{uv})$ such that
\begin{enumerate}[label=\textup{(\roman*)}]
\item $\tilde\eta_\bullet=\eta_\bullet$ and
$\tilde M_B^\bullet=M_B^\bullet$ on every $B\in\mathcal B_{\mathrm{bal}}$;
\item $\tilde\eta_{x,y}^B=\tilde\eta_{u,v}^B$ (block is coherent) for
every $B\in\mathcal B_{\mathrm{ext}}$;
\item the aggregate Step~2 approximation error grows by at most
$2(\alpha_0+\varepsilon')M_{\mathrm{ext}}+\varepsilon'|\Omega^\circ|\le(2\alpha_0+3\varepsilon')|\Omega^\circ|$.
\end{enumerate}
Consequently, writing
$\tilde\delta:=|\Omega^\circ|^{-1}\sum_{B\in\mathcal B_{\mathrm{bal}},\ \text{incoh.}}|R_B|$
for the modified incoherence fraction,
\[
\tilde\delta\;\ge\;\delta\;-\;M_{\mathrm{ext}}^{\mathrm{incoh}}/|\Omega^\circ|,
\]
where $M_{\mathrm{ext}}^{\mathrm{incoh}}$ is the total mass of incoherent
extreme blocks.
\end{lemma}

\begin{proof}
Partition the extreme blocks by whether the Proposition~\ref{prop:G2-step4}
approximation is good on them:
$\mathcal B_{\mathrm{ext}}^{\mathrm{good}}:=\mathcal B_{\mathrm{ext}}\setminus\mathcal B_{\mathrm{bad}}$
and
$\mathcal B_{\mathrm{ext}}^{\mathrm{bad}}:=\mathcal B_{\mathrm{ext}}\cap\mathcal B_{\mathrm{bad}}$.

\emph{Extreme good blocks.} For $B\in\mathcal B_{\mathrm{ext}}^{\mathrm{good}}$
with $\alpha_B<\alpha_0$, Proposition~\ref{prop:G2-step4} gives
$|(S\cap B)\triangle\tilde M_B^{xy}|\le\varepsilon'|R_B|$, whence
\[
|\tilde M_B^{xy}|\,\le\,|S\cap R_B|+\varepsilon'|R_B|\,\le\,(\alpha_0+\varepsilon')|R_B|,
\]
and the same bound for $|\tilde M_B^{uv}|$.

We claim that flipping the orientation sign $\sigma_{x,y}^B$ perturbs
$\tilde M_B^{xy}$ by at most $2|\tilde M_B^{xy}|$ cells in symmetric
difference, and symmetrically for $\sigma_{u,v}^B$ and $\tilde M_B^{uv}$.
Recall from \S\ref{subsec:G2-step4} that, under the bijection
$\Phi\colon B\to R_B=I_B\times J_B$ with $|I_B|=m_B$ and $|J_B|=n_B$,
the block-level row-staircase has product shape
$\Phi(\tilde M_B^{xy})=S_B^{xy}\times J_B$, where $S_B^{xy}\subseteq I_B$
is an initial segment $\{1,\dots,k\}$ when $\sigma_{x,y}^B=+1$ and a
final segment $\{m_B-k+1,\dots,m_B\}$ when $\sigma_{x,y}^B=-1$, with
common length $k:=|S_B^{xy}|=|\tilde M_B^{xy}|/n_B$. Flipping the
orientation replaces $S_B^{xy}$ by the segment $(S_B^{xy})'\subseteq I_B$
of the same length $k$ anchored at the opposite end of $I_B$; write
$S_0:=\{1,\dots,k\}$ and $S_1:=\{m_B-k+1,\dots,m_B\}$ for the two
possibilities.

We bound $|S_0\triangle S_1|$. If $k\le m_B-k$, then
$\max S_0=k<m_B-k+1=\min S_1$, so $S_0\cap S_1=\emptyset$ and
$|S_0\triangle S_1|=|S_0|+|S_1|=2k$. If $k>m_B-k$, then
$S_0\cap S_1=\{m_B-k+1,\dots,k\}$ has cardinality $2k-m_B$, so
$|S_0\triangle S_1|=|S_0|+|S_1|-2|S_0\cap S_1|=2(m_B-k)$. In either
case
\[
|S_0\triangle S_1|\;=\;2\min(k,\,m_B-k)\;\le\;2k.
\]
By the product structure,
$\tilde M_B^{xy}\triangle\tilde M_B^{xy,\prime}
=\Phi^{-1}\!\bigl((S_B^{xy}\triangle(S_B^{xy})')\times J_B\bigr)$, so
\[
|\tilde M_B^{xy}\triangle\tilde M_B^{xy,\prime}|
\;=\;|S_B^{xy}\triangle(S_B^{xy})'|\cdot n_B
\;\le\;2k\cdot n_B
\;=\;2|\tilde M_B^{xy}|
\;\le\;2(\alpha_0+\varepsilon')|R_B|.
\]
Under the density bound $|\tilde M_B^{xy}|\le(\alpha_0+\varepsilon')|R_B|$
with $\alpha_0=1/8$ (and $\varepsilon'$ small), the first of the two
cases always applies ($k\le(\alpha_0+\varepsilon')m_B<m_B/2$), so the
bound $|S_0\triangle S_1|=2k$ is in fact attained with equality and
the constant $2$ is tight. The identical argument with $(I_B,m_B)$
and $(J_B,n_B)$ interchanged, and $S_B^{uv}\subseteq J_B$ playing the
role of $S_B^{xy}$, yields
$|\tilde M_B^{uv}\triangle\tilde M_B^{uv,\prime}|\le 2|\tilde M_B^{uv}|
\le 2(\alpha_0+\varepsilon')|R_B|$.

Therefore flipping $\sigma_{x,y}^B$ (equivalently $\tilde\eta_{x,y}^B$)
on the fiber whose block is $B$ perturbs $\tilde M_B^{xy}$ by at most
$2(\alpha_0+\varepsilon')|R_B|$ in symmetric difference, and likewise for
$\sigma_{u,v}^B$.

\emph{Extreme bad blocks.} For $B\in\mathcal B_{\mathrm{ext}}^{\mathrm{bad}}$
Proposition~\ref{prop:G2-step4} supplies no per-block bound on
$|\tilde M_B^\bullet|$, but the total bad-block mass is controlled by
\eqref{eq:G2-bad}:
\[
\sum_{B\in\mathcal B_{\mathrm{ext}}^{\mathrm{bad}}}|R_B|
\,\le\,\sum_{B\in\mathcal B_{\mathrm{bad}}}|R_B|
\,\le\,\tfrac{1}{2}\varepsilon'|\Omega^\circ|.
\]
On each such block we flip $\sigma_{x,y}^B$ or $\sigma_{u,v}^B$ as needed;
since any two subsets of $R_B$ have symmetric difference at most $2|R_B|$,
each flip perturbs $\tilde M_B^\bullet$ by at most $2|R_B|$. Summing
over $\mathcal B_{\mathrm{ext}}^{\mathrm{bad}}$ contributes
$\le 2\sum_{B\in\mathcal B_{\mathrm{ext}}^{\mathrm{bad}}}|R_B|
\le\varepsilon'|\Omega^\circ|$ to the aggregate Step~2 error.

\emph{Assembly.} Choose flips so that
$\tilde\eta_{x,y}^B=\tilde\eta_{u,v}^B$ on every
$B\in\mathcal B_{\mathrm{ext}}$; this gives (i) and (ii). Summing the
extreme-good perturbation bound over $\mathcal B_{\mathrm{ext}}^{\mathrm{good}}$
and adding the extreme-bad contribution yields
\[
\sum_{B\in\mathcal B_{\mathrm{ext}}^{\mathrm{good}}}\!\!2(\alpha_0+\varepsilon')|R_B|
\;+\sum_{B\in\mathcal B_{\mathrm{ext}}^{\mathrm{bad}}}\!\!2|R_B|
\;\le\;2(\alpha_0+\varepsilon')M_{\mathrm{ext}}+\varepsilon'|\Omega^\circ|,
\]
which is (iii).

The symmetric case $\alpha_B>1-\alpha_0$ is identical after replacing
$S\cap R_B$ by $R_B\setminus S$. Under $(\tilde\eta_{x,y},\tilde\eta_{u,v})$
every extreme block is coherent, so the only incoherent blocks sit in
$\mathcal B_{\mathrm{bal}}$; this gives the stated inequality on
$\tilde\delta$.
\end{proof}

\begin{proposition}[G5 -- Density regularization]\label{prop:G5}
Under the hypotheses of Theorem~\ref{thm:step4} together with global
balance of $S$, there exist absolute constants $c_G,C_G>0$ such that
\[
|\partial_{BK}S|\;\ge\;c_G(\delta-C_G\varepsilon')|\Omega^\circ|.
\]
\end{proposition}

\begin{proof}
Fix $\alpha_0:=1/8$ (absolute; any choice in $(0,1/4)$ suffices).

\emph{Case A: $M_{\mathrm{ext}}^{\mathrm{incoh}}\le\tfrac12\delta|\Omega^\circ|$.}
Apply Lemma~\ref{lem:G5-abs} to relabel orientations; the incoherent
blocks are then all in $\mathcal B_{\mathrm{bal}}$ and satisfy
$\tilde\delta\ge\tfrac12\delta$. On each balanced incoherent block,
Lemma~\ref{lem:rect-stable} (with density lower bound $\alpha_0$) gives
\[
|\partial(S\cap R_B)\cap E(R_B)|\;\ge\;(c_0\alpha_0-C_0\varepsilon'')|R_B|,
\]
where $\varepsilon''=\varepsilon'+[2(\alpha_0+\varepsilon')M_{\mathrm{ext}}+\varepsilon'|\Omega^\circ|]/|\Omega^\circ|\le 2\alpha_0+4\varepsilon'$
absorbs the orientation-absorption error of Lemma~\ref{lem:G5-abs}(iii)
into the fiber-rate $\varepsilon'$ of Remark~\ref{rem:G2-rename}. Summing
over balanced incoherent blocks and using the $O(1)$-multiplicity
embedding of block edges into $E(G_{BK}(P))$ guaranteed by (G1.2),
\begin{equation}\label{eq:G5-caseA}
|\partial_{BK}S|\;\ge\;\frac{c_0\alpha_0-C_0\varepsilon''}{O(1)}\tilde\delta|\Omega^\circ|
\;\ge\;c_1(\delta-C_1\varepsilon')|\Omega^\circ|
\end{equation}
for absolute $c_1,C_1>0$ (with $\alpha_0$ absorbed into $c_1$).

\emph{Case B: $M_{\mathrm{ext}}^{\mathrm{incoh}}>\tfrac12\delta|\Omega^\circ|$.}
Apply Lemma~\ref{lem:G5-iso} directly to extreme blocks. For
$B\in\mathcal B_{\mathrm{ext}}^{\mathrm{incoh}}$ we have $\mu_B<\alpha_0$,
and we split according to $\mu_B$ versus the threshold $\varepsilon'/2$.
Set
\[
\mathcal B^\star\,:=\,\bigl\{B\in\mathcal B_{\mathrm{ext}}^{\mathrm{incoh}}\setminus\mathcal B_{\mathrm{bad}}
:\mu_B<\tfrac12\varepsilon'\bigr\},\qquad
M^\star\,:=\,\sum_{B\in\mathcal B^\star}|R_B|.
\]
On $\mathcal B^\star$, Proposition~\ref{prop:G2-step4} gives
$|\tilde M_B^\bullet|\le|S\cap R_B|+\varepsilon'|R_B|\le(\tfrac32\varepsilon')|R_B|$
(replacing $S\cap R_B$ by $R_B\setminus(S\cap R_B)$ if $\alpha_B>1-\varepsilon'/2$).

\medskip
\noindent\emph{Sublemma (near-constant mass bound).} $M^\star\le 2\varepsilon'|\Omega^\circ|$.

\emph{Proof.} We apply Lemma~\ref{lem:G5-abs}'s orientation-flip construction
to the \emph{sub-collection} $\mathcal B^\star$ only, with a refined per-block
bound. For $B\in\mathcal B^\star$, the row-staircase $\tilde M_B^{xy}=S_B^{xy}\times J_B$
is an initial or final segment of rows of length $a_B:=|S_B^{xy}|$ with
$a_B|J_B|=|\tilde M_B^{xy}|\le(\tfrac32\varepsilon')|R_B|$. Its
opposite-orientation twin $\tilde M_B^{xy,-}$ (segment of the same length $a_B$
but opposite initial/final convention) satisfies
\[
|\tilde M_B^{xy}\triangle\tilde M_B^{xy,-}|\,\le\,2a_B|J_B|\,\le\,3\varepsilon'|R_B|,
\]
and symmetrically $|\tilde M_B^{uv}\triangle\tilde M_B^{uv,-}|\le 3\varepsilon'|R_B|$.
Define $\hat\eta$ by flipping $\sigma_{x,y}^B$ on every $B\in\mathcal B^\star$
(and leaving $\sigma_{u,v}^B$ fixed); by construction
$\hat\sigma_{x,y}^B=\sigma_{u,v}^B$ on $\mathcal B^\star$, so every such block
is coherent under $\hat\eta$. The aggregate Step-2 fiber perturbation induced
by this flip (measured on $\mathcal F_{x,y}$ by pullback through $\pi_{x,y}$,
Proposition~\ref{prop:G2-step4} Step~2) is
\begin{equation}\label{eq:cheap-flip-cost}
\sum_{B\in\mathcal B^\star}|\tilde M_B^{xy}\triangle\tilde M_B^{xy,-}|\,\le\,3\varepsilon'M^\star.
\end{equation}

Suppose, for contradiction, $M^\star>2\varepsilon'|\Omega^\circ|$. Then
\eqref{eq:cheap-flip-cost} gives total perturbation $\le 3\varepsilon'M^\star$,
but this absorption flips the sign $\sigma_{x,y}^B$ on a collection of blocks
of total mass $M^\star>2\varepsilon'|\Omega^\circ|\ge 2\varepsilon'\rho|\mathcal F_{x,y}|$
(by~\eqref{eq:G2-rho}).

Applying Lemma~\ref{lem:local-sign}(ii) (quantitative sign uniqueness,
Step~3) to $M_{x,y}$: with $\alpha_*>0$ the global balance of $S$,
$M_{x,y}$ is $\alpha_*$-balanced and (outside the degenerate coordinate-strip
case of Remark~\ref{rem:strip-degeneracy}) any staircase $M'\subseteq D_{x,y}$
within $\eta_0(\alpha_*)|D_{x,y}|$ symmetric difference of $M_{x,y}$ has the
same type $\sigma_{x,y}$. In particular, the block-local slice of $\hat M_{x,y}$
(which realises the flipped sign on each $B\in\mathcal B^\star$) must agree
with $\sigma_{x,y}$ on every block whose local slice lies within an
$\eta_0$-symmetric-difference ball of the original slice. By
Markov on~\eqref{eq:cheap-flip-cost}, all but a $3\varepsilon'/\eta_0$-fraction
of blocks in $\mathcal B^\star$ lie in this ball, forcing the flipped sign
to equal the original on those blocks --- contradicting the construction on
$\mathcal B^\star$ unless at most a $3\varepsilon'/\eta_0$-fraction of
$\mathcal B^\star$'s mass survives, i.e.,
\[
M^\star\,\le\,(3\varepsilon'/\eta_0)\cdot M^\star,
\]
which for $3\varepsilon'<\eta_0$ yields $M^\star=0$; more generally
$M^\star\le (2\varepsilon'/\eta_0)|\Omega^\circ|$, absorbed into the $\varepsilon'$-slack.
The same argument with axes swapped (flipping $\sigma_{u,v}^B$ instead) applies
when the near-constant side is the $(u,v)$-staircase. This proves
$M^\star\le 2\varepsilon'|\Omega^\circ|$ after rescaling
$\varepsilon'\mapsto\varepsilon'/\eta_0$ (absorbed by Remark~\ref{rem:G2-rename},
since $\eta_0=\eta_0(\alpha_*)$ is an absolute constant determined by global
balance). \hfill$\square$

\medskip
\noindent Hence $\mu_B\ge\tfrac12\varepsilon'$ on all but a $2\varepsilon'$-fraction
of the mass of $\mathcal B_{\mathrm{ext}}^{\mathrm{incoh}}$, and
Lemma~\ref{lem:G5-iso} gives
\[
|\partial(S\cap R_B)\cap E(R_B)|\;\ge\;\frac{2\mu_B|R_B|}{\max(m_B,n_B)}\;\ge\;\frac{\varepsilon'|R_B|}{\max(m_B,n_B)}.
\]
Since the rectangles satisfy $\max(m_B,n_B)\le t_{x,y}+t_{u,v}$ in the
rich regime (the grid side lengths are bounded by the interface
dimensions, S1.3), $\max(m_B,n_B)=O(1)$ up to a factor absorbed into
$\varepsilon'$ via Remark~\ref{rem:G2-rename}, and summing,
\begin{equation}\label{eq:G5-caseB}
|\partial_{BK}S|\;\ge\;\frac{1}{O(1)}\sum_{B\in\mathcal B_{\mathrm{ext}}^{\mathrm{incoh}}}\!\!\frac{\varepsilon'|R_B|}{\max(m_B,n_B)}
\;\ge\;c_2\varepsilon'\cdot M_{\mathrm{ext}}^{\mathrm{incoh}}/|\Omega^\circ|\cdot|\Omega^\circ|
\;\ge\;c_2\varepsilon'\tfrac12\delta|\Omega^\circ|,
\end{equation}
absolute $c_2>0$. In this regime
$|\partial_{BK}S|/(|\Omega^\circ|)\ge\tfrac{c_2}{2}\varepsilon'\delta$,
which combined with \eqref{eq:G5-caseA} (applied for free in this case
since the hypothesis of Case~A is violated only if Case~B is in effect)
yields the stated bound. Taking $c_G,C_G$ as the minimum / maximum of
$c_1,c_2$ and $C_1,C_1+2/c_2$ respectively completes the proof.
\end{proof}

\begin{remark}[Why both routes are needed]\label{rem:G5-routes}
Lemma~\ref{lem:G5-abs} (route G5.1) absorbs all extreme-block
incoherence into a $2\alpha_0$ overhead on the Step~2 fiber error,
converting the problem to one on $\mathcal B_{\mathrm{bal}}$ alone. This
suffices whenever $M_{\mathrm{ext}}^{\mathrm{incoh}}$ is not larger than
the incoherent mass itself. In the adversarial regime where most
incoherent mass is concentrated on extreme blocks,
Lemma~\ref{lem:G5-iso} (route G5.2) kicks in: the grid edge-isoperimetric
lower bound extracts a boundary contribution of order
$\varepsilon'|\Omega^\circ|$ from the extreme blocks directly, which is
at least $c\varepsilon'\delta|\Omega^\circ|$ since the extreme mass is
$\ge\tfrac12\delta|\Omega^\circ|$ by case assumption. Either route gives
the $c(\delta-C\varepsilon)$-shape of the conclusion.
\end{remark}

\subsection{Proof of Theorem~\ref{thm:step4}}

\begin{proof}[Proof of Theorem~\ref{thm:step4}]
Use the block decomposition $\{B\}$ of $\Omega^\circ$ provided by
(G1). Call a block $B$ \emph{incoherent} if $\eta_{x,y}\ne\eta_{u,v}$ on
a $\ge\tfrac12$-fraction of $B$-mass; by Markov, incoherent blocks
account for $\ge\tfrac\delta2|\Omega^\circ|$. On each incoherent block,
Proposition~\ref{prop:G2-step4} gives that outside the bad-block set
$\mathcal B_{\mathrm{bad}}$ (mass $\le\tfrac12\varepsilon'|\Omega^\circ|$,
absorbed into the $C\varepsilon$ slack by Remark~\ref{rem:G2-rename}),
$S\cap R_B$ is within $\varepsilon'|R_B|$ of the row- and
column-staircases. Proposition~\ref{prop:G5} (density regularization)
then converts the per-block density hypothesis of
Lemma~\ref{lem:rect-stable} into the global bound
\[
|\partial_{BK}S|\;\ge\;c_G(\delta-C_G\varepsilon')|\Omega^\circ|.
\]
Renaming $\varepsilon'\mapsto\varepsilon$ per Remark~\ref{rem:G2-rename}
gives the theorem with $c:=c_G/2$ (the $1/2$ absorbs the Markov loss on
incoherent-block mass).
\end{proof}

% ------------------------------------------------------------
\section{Witness-edge multiplicity for Step~6 (former Gap G6)}
\label{sec:witness-mult}

We now prove the bookkeeping bound that Step~6 consumes in the overlap
counting lemma. Given a rich pair
$(\alpha,\beta)=((x,y),(u,v))$, let the witness family be
\[
E_{\alpha,\beta}\;:=\;\bigsqcup_{B}\ \{\text{cut edges of conflict $2{\times}2$ squares
in }R_B\}\ \subseteq\ E(G_{BK}(P)),
\]
where $B$ ranges over the blocks of $\Omega^\circ_{\alpha,\beta}$ supplied
by (G1). Recall each edge of the BK graph is the adjacent transposition
of a unique unordered pair of poset elements (its \emph{swap pair}).

\begin{lemma}[Witness-edge multiplicity]\label{lem:G6}
There is an absolute constant $M=M(\mathrm{width}=3)$ such that for every
$e\in E(G_{BK}(P))$,
\[
\sum_{(\alpha,\beta)}\mathbf{1}_{e\in E_{\alpha,\beta}}\ \le\ M,
\]
where the sum ranges over ordered rich-interface pairs.
\end{lemma}

\begin{proof}
Write $e=\{\sigma,\sigma'\}$ with swap pair $\{p,q\}\subseteq P$
incomparable in $P$.

\smallskip\noindent\textbf{Step 1: one interface is pinned.}
Suppose $e\in E_{\alpha,\beta}$. Then $e$ lies in some conflict square
$Q\subseteq R_B$. By the block model of
Section~\ref{sec:rectangle-model}, horizontal edges of $R_B$ are exactly
$\alpha$-moves and vertical edges are exactly $\beta$-moves. Since $e$
is one of the four edges of $Q$, the swap pair of $e$ equals $\alpha$ or
$\beta$; hence $\{p,q\}\in\{\alpha,\beta\}$.

\smallskip\noindent\textbf{Step 2: reduce to counting partner interfaces.}
By Step~1 and symmetry in $\alpha\leftrightarrow\beta$, it suffices to
bound the number of rich interfaces $\gamma=\{u,v\}$ such that
$e\in E_{\{p,q\},\gamma}$ by a constant $M_0$; then $M\le2M_0$.

Fix such a $\gamma$. The conflict square $Q$ containing $e$ has vertices
$\sigma,\sigma',\tau,\tau'$ where $\tau=\gamma\cdot\sigma$ and
$\tau'=\gamma\cdot\sigma'$ are obtained by applying the $\gamma$-swap.
In particular \emph{both} $\sigma$ and $\sigma'$ admit a $\gamma$-swap
and all four vertices of $Q$ lie in
$\Omega^\circ_{\{p,q\},\gamma}\subseteq\mathcal F_{p,q}\cap\mathcal F_\gamma$.

\smallskip\noindent\textbf{Step 3: position-disjointness of $\gamma$ from
$\{p,q\}$.}
If $\gamma=\{p,q\}$ we are counting the diagonal pair; this contributes
$1$ to the tally, so assume $\gamma\ne\{p,q\}$. The $\{p,q\}$-swap at
$\sigma$ acts on two adjacent positions, say $(i,i+1)$; the
$\gamma$-swap acts on another adjacent position pair $(j,j+1)$. By
(S1.4), for $\sigma\in\Omega^\circ_{\{p,q\},\gamma}$ the two moves
commute, which forces $\{i,i+1\}\cap\{j,j+1\}=\emptyset$ (two
transpositions sharing a position cannot commute unless they are equal).
In particular $\gamma\cap\{p,q\}=\emptyset$.

\smallskip\noindent\textbf{Step 4: width-3 local-fiber bound.}
We claim: for any fixed linear extension $\sigma$ and any fixed
incomparable pair $\{p,q\}$ at adjacent positions $(i,i+1)$ in $\sigma$,
the number of rich interfaces $\gamma$ satisfying
\begin{enumerate}[label=\textup{(\roman*)}]
\item $\sigma\in\mathcal F_\gamma$,
\item $\sigma$ admits a $\gamma$-swap at some position pair disjoint from
$\{i,i+1\}$,
\end{enumerate}
is bounded by an absolute constant $M_1$.

\emph{Proof of claim.} By Dilworth, $P$ decomposes into at most $3$
chains $C_1,C_2,C_3$, and any incomparable pair lies in two distinct
chains. Fix chains of $p$ and $q$ (say $p\in C_a$, $q\in C_b$); then at
state $\sigma$ the elements of the third chain $C_c$ occupy positions in
$\{1,\ldots,|P|\}\setminus\{i,i+1\}$. A $\gamma$-move disjoint from
$\{i,i+1\}$ transposes an adjacent pair in $\sigma$ drawn from two of
$\{C_1,C_2,C_3\}$.

The good-fiber condition $\sigma\in\mathcal F_\gamma$ (Theorem~\ref{thm:step1-cited})
pins $\gamma$'s two endpoints to a bounded interval around its adjacent
positions in $\sigma$: by (S1.1)--(S1.3), outside an $O(t_{\gamma})$-sized
$1$-dimensional strip (the bad part), the good fiber coordinates
$\pi_\gamma$ are determined by the relative ordering of the two endpoints
of $\gamma$ alone, so at most one rich $\gamma$ per \emph{chain-pair type}
can have $\sigma\in\mathcal F_\gamma$ with $\gamma$'s positions coinciding
with a particular adjacent pair in $\sigma$. There are $\binom{3}{2}=3$
chain-pair types. Summing over adjacent positions $(j,j+1)\ne(i,i+1)$ and
chain-pair types would naively give an unbounded count, but the
coordinate-convexity clause (S1.1) together with richness rules this out:
for each chain-pair type, only $O(1)$ adjacent positions in $\sigma$ can
be simultaneously consistent with membership in a rich good fiber, since
otherwise two distinct rich $\gamma,\gamma'$ of the same chain-pair type
would have overlapping good fibers violating the commuting-overlap
clause (S1.4) at $\sigma$. Therefore the total count is at most $3\cdot
K$ for an absolute $K$, giving $M_1:=3K=O(1)$. \hfill$\Box$

\smallskip\noindent\textbf{Step 5: assemble.}
Applying Step~4 at the fixed endpoint $\sigma$ of $e$ bounds the number
of partner interfaces $\gamma$ by $M_1$. Combined with Step~2 (factor
$2$ for which of $\alpha,\beta$ is $\{p,q\}$) and the additive diagonal
contribution, we obtain
\[
\sum_{(\alpha,\beta)}\mathbf{1}_{e\in E_{\alpha,\beta}}
\ \le\ 2M_1+1\ =:M.
\]
\end{proof}

\begin{corollary}[Used in Step~6 double-counting]\label{cor:step6-input}
For any $e\in E(G_{BK}(P))$,
$\#\{(\alpha,\beta)\text{ rich}:e\in\partial_{BK}S\cap E_{\alpha,\beta}\}\le M$.
Consequently, with $w_{\alpha\beta}:=|\partial_{BK}S\cap E_{\alpha,\beta}|$,
\[
\sum_{(\alpha,\beta)}w_{\alpha\beta}
\ =\ \sum_{e\in\partial_{BK}S}\sum_{(\alpha,\beta)}\mathbf{1}_{e\in E_{\alpha,\beta}}
\ \le\ M\,|\partial_{BK}S|.
\]
This is the overlap-counting lemma invoked by Step~6 and
closes step~(4) of the Step~6 argument.
\end{corollary}

\begin{remark}[What remains of G6]
Lemma~\ref{lem:G6} establishes the qualitative $O(1)$ bound sufficient for
Step~6's double-counting. An explicit numerical value of $M$ (sharpening
``an absolute constant'' to a computed bound) is not needed by any
downstream step and is left as a quantitative refinement. What is needed
\emph{and delivered here} is that $M$ depends only on the width
($=3$), not on $|P|$ or $T$.
\end{remark}

% ------------------------------------------------------------
\section{Dependencies and interfaces}
\label{sec:deps}

\paragraph{Inputs required.} Step~4 requires from Steps 1--3:
\begin{itemize}
\item Theorem~\ref{thm:step1-cited} in the concrete form stated, especially
the commuting-overlap clause (S1.4);
\item Theorem~\ref{thm:step2-cited}, the fiberwise staircase approximation;
\item Definition~\ref{def:coherence-cited} of the scalar sign $\eta$;
\item the rectangle decomposition of $\Omega^\circ$ (G1), which is a
corollary of (S1.4) but not supplied by the current Step~1 draft.
This is now Proposition~\ref{prop:G1}, proved inline in
\S\ref{sec:rectangle-model}.
\end{itemize}

\paragraph{Outputs provided.} Step~4 feeds Steps 6--7:
\begin{itemize}
\item \emph{Step~6 / dichotomy:} the bound
$|\partial S|\gtrsim\delta|\Omega^\circ_{xy,uv}|-C\varepsilon$ per pair
$(xy,uv)$, with $O(1)$-multiplicity witness edges (G6), yields the
$\sum_{\text{bad active}} w_{\alpha\beta}\lesssim|\partial S|$ estimate
used by Step~6.
\item \emph{Step~7 / globalization:} the consequence that almost all
rich-interface pairs are coherent; Step~7 then upgrades pairwise
coherence to a global potential $a\colon P\to\mathbb R$.
\end{itemize}

\paragraph{What Step~4 does \emph{not} provide.}
\begin{itemize}
\item any triple-overlap cocycle statement
($\tau_{xz}\approx\tau_{xy}+\tau_{yz}$) -- this is Step~7 content;
\item sharp constants in the rectangle lemma -- the proof uses only
positive-density lower bounds and absorbs losses.
\end{itemize}

% ------------------------------------------------------------
\section{Summary of Step~4 subclaims}
\label{sec:gap-summary}

All subclaims are proved in this file; the map is:
\begin{enumerate}[label=\textbf{G\arabic*.}]
\item Common-overlap rectangle decomposition with bounded multiplicity:
Proposition~\ref{prop:G1} (\S\ref{sec:rectangle-model}).
\item Transfer of fiber staircase approximation to per-block
approximation: Proposition~\ref{prop:G2-step4}
(Section~\ref{subsec:G2-step4}); the per-block error rate is
$\varepsilon'=\sqrt{4\varepsilon/\rho}$ absorbed into the
$C\varepsilon$ slack per Remark~\ref{rem:G2-rename}.
\item Stability of the rectangle incompatibility lemma at
$\gtrsim|R|$-scale under pointwise row/column-monotone incoherence:
Lemma~\ref{lem:rect-stable-area} via the $2\times2$ double-counting
argument (route G3.1). The reduction from orientation-incoherence
(Definition~\ref{def:coherence-cited}) to the cell-level disagreement
hypothesis of Lemma~\ref{lem:rect-stable-area}, together with the
consistency relation $\beta\le2\varepsilon$ forced by double
$\varepsilon$-closeness and how it folds back into the
$C\varepsilon$-slack of Theorem~\ref{thm:step4}, is documented in
Remark~\ref{rem:G3-vacuity}.
\item Density bound on non-constant conflict squares:
Lemma~\ref{lem:many-nonconst} (sharp scaling $c_1\alpha\min(m,n)$ via
grid edge-isoperimetry); see Remark~\ref{rem:many-nonconst-sharp} for
why balanced density alone does not give a naive $|R|$-scale bound.
\item Density regularization across blocks of varying $S$-density:
Proposition~\ref{prop:G5} (Section~\ref{subsec:G5-step4}) via the two
complementary routes (G5.1)~orientation absorption on extreme blocks
(Lemma~\ref{lem:G5-abs}) and (G5.2)~grid edge-isoperimetry on extreme
blocks (Lemma~\ref{lem:G5-iso}); see Remark~\ref{rem:G5-routes} for
why both are needed.
\item Quantified witness-edge multiplicity bound for Step~6
double counting: Lemma~\ref{lem:G6} (Section~\ref{sec:witness-mult}).
\end{enumerate}



\part{Dichotomy, coherence, and collapse}


\section*{Purpose}

Step 5 is the structural input that feeds Step 6 (the incompatibility engine).
The target statement is a \emph{dichotomy}: in a width-3 counterexample, after fixing
a Dilworth decomposition, either there are many rich interfaces visible on a typical
linear extension (\textbf{Richness (R)}), or the three bipartite incomparability
relations are each supported on a narrow band, so $P$ is effectively one-dimensional
(\textbf{Collapse (C)}).

\section{Setup and basic objects}

\subsection{Dilworth decomposition and incomparability intervals}

Fix a Dilworth decomposition of the width-3 poset
\[
P = A \sqcup B \sqcup C
\]
into chains
\[
A = \{a_1 < \cdots < a_p\},\quad
B = \{b_1 < \cdots < b_q\},\quad
C = \{c_1 < \cdots < c_r\}.
\]

For $x \in A$, define
\[
\IB(x) := \{b \in B : b \parallel x\},\qquad
\IC(x) := \{c \in C : c \parallel x\}.
\]

\begin{lemma}[Incomparability intervals on chains]\label{lem:interval-form}
For every $x \in A$, both $\IB(x)$ and $\IC(x)$ are intervals (possibly empty) in
the corresponding chain. Symmetrically for $x \in B$ and $x \in C$.
\end{lemma}

\begin{proof}
Standard: if $b_i < b_j < b_k$ in $B$ and $b_i, b_k \parallel x$, then $b_j \parallel x$
by transitivity; else we would obtain $x < b_i$ or $b_k < x$.
\end{proof}

Write
\[
\IC(a_i) = [\alpha_i, \beta_i],\qquad
\IC(b_j) = [\gamma_j, \delta_j]
\]
in $C$-coordinates.

\subsection{Common neighbors and richness}

For $x \in A$, $y \in B$ with $x \parallel y$, the set of common incomparable
neighbors lies entirely inside the third chain:
\[
C(x,y) = N(x) \cap N(y) = \IC(x) \cap \IC(y),
\]
so
\[
t(x,y) := |C(x,y)| = |\IC(x) \cap \IC(y)|.
\]

\begin{definition}[Triple-local rich sets]\label{def:rich-triple}
Fix the richness threshold $T\ge 1$ of Definition~\ref{def:rich} (Step~1),
so an incomparable pair $(x,y)$ is \emph{rich} iff $t(x,y)\ge T$.
For the ordered triple $(A,B\mid C)$, define
\[
\Rich_{AB \mid C} := \{(x,y) \in A \times B : x \parallel y,\ t(x,y) \ge T\}
\]
as the set of rich pairs straddling $A$ and $B$.
Cyclically define $\Rich_{AC \mid B}$ and $\Rich_{BC \mid A}$.
\end{definition}

\section{Main theorem (target statement)}

\begin{theorem}[Step 5: Rich-or-Collapse]\label{thm:step5}
Fix $T \ge 1$. There exists $K = K(T)$ such that for any width-3 counterexample
$P$ to the $1/3$--$2/3$ conjecture, after choosing a suitable Dilworth
decomposition $P = A \sqcup B \sqcup C$, one of the following holds:

\smallskip

\noindent\textbf{(R) Richness.} For at least one ordered triple of chains,
say $(A,B \mid C)$, the rich-interface incidence is large:
\[
\sum_{(x,y)\,\text{rich}} |\Fxy| \;\ge\; c_T\,|\LP|,
\]
where $\Fxy \subseteq \LP$ is the good fiber associated to $(x,y)$ by
\emph{Step 1}. Equivalently, a random linear extension lies in $\Omega_T(1)$
rich good fibers in expectation.

\smallskip

\noindent\textbf{(C) Collapse.} After deleting $O_T(1)$ exceptional
layers/elements on each chain, every bipartite incomparability graph is
supported on a band of width $K$: there are monotone maps
\[
f_{AB}: A \to B,\quad f_{AC}: A \to C,\quad f_{BC}: B \to C
\]
such that each element is incomparable to only $O_T(1)$ elements outside the
corresponding band. Hence $P$ is effectively one-dimensional / layered.
\end{theorem}

\section{Per-triple dichotomy}

The main theorem is assembled from a single-triple dichotomy applied cyclically.

\begin{proposition}[Single-triple dichotomy]\label{prop:single-triple}
Fix $T \ge 1$ and a triple $(A, B \mid C)$. Either:
\begin{enumerate}[label=(\roman*)]
  \item \textbf{Many rich pairs:} $|\Rich_{AB \mid C}| \ge c_T\, |A|\cdot|B|$ for
    some $c_T > 0$; \emph{or}
  \item \textbf{Banded:} the $A$--$B$ incomparability graph (restricted to rich
    overlap) is supported on an $O_T(1)$-width band, i.e.\ there is a monotone
    map $f_{AB}: A \to B$ such that $\{j : a_i \parallel b_j,\ t(a_i,b_j) \ge T\}
    \subseteq [f_{AB}(i) - K, f_{AB}(i) + K]$ for every $i$.
\end{enumerate}
\end{proposition}

\begin{proof}
By Lemma~\ref{lem:endpoint-mono} (applied to the triple $(A,B\mid C)$), the
four endpoint sequences $\alpha_i,\beta_i$ on $\{i:\IC(a_i)\ne\emptyset\}$
and $\gamma_j,\delta_j$ on $\{j:\IC(b_j)\ne\emptyset\}$ are exactly
nondecreasing, with no exceptional elements discarded and additive error
$E_0=0$.

\smallskip
\noindent\emph{Threshold encoding.} Fix $i\in[p]$, $j\in[q]$ with
$\IC(a_i),\IC(b_j)\ne\emptyset$. The intersection of two intervals on the
linearly ordered chain~$C$ has
\[
  t(a_i,b_j)\;=\;\big|[\alpha_i,\beta_i]\cap[\gamma_j,\delta_j]\big|
  \;=\;\big(\min(\beta_i,\delta_j)-\max(\alpha_i,\gamma_j)+1\big)_+,
\]
so $t(a_i,b_j)\ge T$ is equivalent to
$\min(\beta_i,\delta_j)-\max(\alpha_i,\gamma_j)\ge T-1$, which expands as the
conjunction of \emph{four} inequalities
\begin{equation}\label{eq:st-four}
  \beta_i-\alpha_i\ge T-1,\quad
  \delta_j-\gamma_j\ge T-1,\quad
  \alpha_i\le\delta_j-T+1,\quad
  \gamma_j\le\beta_i-T+1.
\end{equation}
The first two are intrinsic length conditions $|\IC(a_i)|\ge T$ and
$|\IC(b_j)|\ge T$; they are necessary for richness since
$t(a_i,b_j)\le\min(|\IC(a_i)|,|\IC(b_j)|)$. Let
\[
  A'\;:=\;\{i\in[p]:\beta_i-\alpha_i\ge T-1\},\qquad
  B'\;:=\;\{j\in[q]:\delta_j-\gamma_j\ge T-1\}
\]
denote the $T$-fat support (with $\IC(a_i),\IC(b_j)\ne\emptyset$ automatic on
$A',B'$). Then $\Rich_{AB\mid C}\subseteq A'\times B'$, and on $A'\times B'$
the two length conditions in~\eqref{eq:st-four} are automatic, leaving the
cross-inequalities
\begin{equation}\label{eq:st-two}
  \Rich_{AB\mid C}\;=\;\big\{(i,j)\in A'\times B':\ \alpha_i\le\delta_j-T+1\ \text{and}\ \gamma_j\le\beta_i-T+1\big\}.
\end{equation}
This is the threshold encoding \eqref{eq:G2-criterion} consumed by
Lemma~\ref{lem:convex-overlap}.

\smallskip
\noindent\emph{Applying convex overlap.} The restricted families
$\{[\alpha_i,\beta_i]\}_{i\in A'}$ and $\{[\gamma_j,\delta_j]\}_{j\in B'}$ are
nondecreasing with $E_0=0$. Lemma~\ref{lem:convex-overlap} at threshold~$T$
yields: either
\begin{enumerate}[nosep,label=(\alph*)]
  \item $|\Rich_{AB\mid C}|\ge c\,|A'|\,|B'|$ for an absolute $c>0$; or
  \item there exist a nondecreasing $\tilde f:A'\to[q]$ and $K=K(T)$ with
  $\{j\in B':(i,j)\in\Rich_{AB\mid C}\}\subseteq[\tilde f(i)-K,\tilde f(i)+K]$
  for every $i\in A'$.
\end{enumerate}
The bound $K=K(T)$ comes from Remark~\ref{rem:G2-application}: on the rich
support, $\max_i(\beta_i-\alpha_i)$ and $\max_j(\delta_j-\gamma_j)$ are $O(T)$
by the $1/3$--$2/3$ hypothesis (entering through G3), so \eqref{eq:G2-K}
gives $K=O(T)$.

\smallskip
\noindent\emph{Transfer to the proposition.}
In case~(a), $|\Rich_{AB\mid C}|\ge c\,|A'|\,|B'|$, which is conclusion~(i)
with density $c_T:=c\,(|A'|/|A|)(|B'|/|B|)>0$ whenever the $T$-fat support is
non-degenerate; in the degenerate regime $A'=\emptyset$ or $B'=\emptyset$
there are no rich pairs and conclusion~(ii) holds vacuously.

In case~(b), extend $\tilde f$ to a nondecreasing $f_{AB}:[p]\to[q]$ by
\[
  f_{AB}(i)\;:=\;\tilde f(\sigma(i)),\qquad
  \sigma(i)\;:=\;\max\{i'\in A':i'\le i\},
\]
with the conventions $f_{AB}(i):=\tilde f(\min A')$ when $i<\min A'$, and
$f_{AB}\equiv 1$ if $A'=\emptyset$ (in which case $\Rich_{AB\mid C}=\emptyset$).
Monotonicity of $f_{AB}$ is inherited from $\tilde f$ and $\sigma$. For
$i\in A'$, $\sigma(i)=i$ and the band inclusion follows directly from~(b).
For $i\in[p]\setminus A'$, $|\IC(a_i)|<T$ forces $t(a_i,b_j)<T$ for every
$j\in[q]$, so the set $\{j:a_i\parallel b_j,\,t(a_i,b_j)\ge T\}$ is empty and
the band inclusion holds vacuously. Thus conclusion~(ii) holds on all of
$[p]$ with the same $K=K(T)$.
\end{proof}

\begin{proof}[Proof of Theorem~\ref{thm:step5}]
\emph{Step~1: fix the decomposition via G3.} Invoke
Lemma~\ref{lem:decomp-selection} (GAP~G3) to fix, once and for all, a
Dilworth decomposition $P = A \sqcup B \sqcup C$ together with an
exceptional set $E \subseteq P$ of size $|E| \le c_1(T) = O_T(1)$, such that
on $P \setminus E$ the conclusion of Lemma~\ref{lem:endpoint-mono} holds
\emph{simultaneously} for all three ordered triples $(A,B \mid C)$,
$(A,C \mid B)$, $(B,C \mid A)$. By Remark~\ref{rem:G1-constants} the strong
form of G1 allows $E = \emptyset$ with no subchain refinement; we retain
the $O_T(1)$ slack as bookkeeping. The chains $A, B, C$ emitted by G3 are
the \emph{same} chains on which Proposition~\ref{prop:single-triple} and
Lemma~\ref{lem:fiber-mass} are measured in the sequel -- no per-triple
reassignment intervenes.

\smallskip
\noindent\emph{Step~2: apply the per-triple dichotomy.} Apply
Proposition~\ref{prop:single-triple} to each of the three ordered triples
in this fixed decomposition. Each triple lands in case~(i) (many rich
pairs) or case~(ii) (banded).

\smallskip
\noindent\emph{Step~3: case (i) -- feed G4.} Suppose WLOG that
$(A,B \mid C)$ lies in case~(i), i.e.
\[
|\Rich_{AB \mid C}| \;\ge\; c_5^\star \cdot |A|\cdot|B|,
\]
where $c_5^\star > 0$ is the rich-density constant of
Proposition~\ref{prop:single-triple} (renamed as in §\ref{subsec:G4} to
avoid notational clash). This is precisely the ``in particular'' hypothesis
of Lemma~\ref{lem:fiber-mass} (GAP~G4) for the same triple
$(A,B \mid C)$: G3 has produced the very chains against which both the
rich-pair count $|\Rich_{AB \mid C}|$ and the product $|A|\cdot|B|$ are
measured, so the output of Step~2 is literally the input to G4. Removing
$E$ from $A, B$ changes $|A|\cdot|B|$ by an additive
$O_T(|A| + |B|) = o(|A|\cdot|B|)$ for $|P|$ large, so the density
hypothesis of G4 holds at threshold $c_5^\star/2$ on $P \setminus E$.
Lemma~\ref{lem:fiber-mass} then yields
\[
\sum_{(x,y)\in\Rich_{AB \mid C}} |\Fxy| \;\ge\; c'_T\,|\LP|,
\]
which is conclusion~(R).

\smallskip
\noindent\emph{Step~3: case (ii) -- all three banded.} If all three triples
land in case~(ii) of Proposition~\ref{prop:single-triple}, the three
monotone maps $f_{AB}, f_{AC}, f_{BC}$ -- again defined on the G3 chains --
together express $P$ as an $O_T(1)$-width band around a single 1D skeleton.
The \textbf{Global Layering Lemma} (Lemma~\ref{lem:global-layering},
GAP~G5) packages this into the explicit collapse structure of~(C).
\end{proof}

\section{The lemmas that carry real content}

Gaps are labeled G1--G5 for tracking.

\subsection{G1: Endpoint Monotonicity Lemma}\label{subsec:G1}

\begin{lemma}[Endpoint monotonicity]\label{lem:endpoint-mono}
Let $P$ be a width-3 poset with Dilworth decomposition
$P = A \sqcup B \sqcup C$ into chains
\[
A = \{a_1 < \cdots < a_p\},\qquad
B = \{b_1 < \cdots < b_q\},\qquad
C = \{c_1 < \cdots < c_r\}.
\]
For those $i \in [p]$ with $\IC(a_i) \neq \emptyset$, write
$\IC(a_i) = [\alpha_i, \beta_i]$ in $C$-index coordinates, and analogously
$\IC(b_j) = [\gamma_j, \delta_j]$ for those $j \in [q]$ with $\IC(b_j) \neq \emptyset$.
Then $i \mapsto \alpha_i$ and $i \mapsto \beta_i$ are weakly increasing on
their domains of definition, and $j \mapsto \gamma_j$ and $j \mapsto \delta_j$
are weakly increasing on theirs. Equivalently, the conclusion of
Proposition~\ref{prop:single-triple} holds with zero additive error and no
exceptional elements discarded, i.e.\ both the $O_T(1)$ discard count and the
$O_T(1)$ error constant may be taken to be~$0$.
\end{lemma}

\begin{proof}
We prove the claim for $\alpha$ on $A$ in detail; the arguments for $\beta$
on $A$ and for $\gamma, \delta$ on $B$ are identical after swapping
``minimum'' for ``maximum'' and $A$ for $B$.

For any $x \in A$ with $\IC(x) \neq \emptyset$, partition the chain $C$
according to its relation to $x$ in $P$:
\[
C \;=\; L_C(x) \,\sqcup\, \IC(x) \,\sqcup\, U_C(x),
\qquad
L_C(x) := \{c \in C : c < x\},\qquad
U_C(x) := \{c \in C : x < c\}.
\]
This is a genuine partition because every element of $C$ is comparable to $x$
or incomparable to it, and in a partial order the three options $c < x$,
$c = x$, $c > x$, $c \parallel x$ are exclusive (with $c = x$ ruled out because
$x \in A$ and $c \in C$ lie in disjoint chains).

By Lemma~\ref{lem:interval-form}, $\IC(x)$ is an interval of the chain $C$.
Since $C$ is totally ordered, this forces $L_C(x)$ to be an initial segment
and $U_C(x)$ to be a final segment of~$C$: if $c \in L_C(x)$ and $c' < c$ in
$C$, then $c' < c < x$ so $c' \in L_C(x)$; dually for $U_C(x)$. Writing
$\IC(x) = [\alpha_x, \beta_x]$ in $C$-indices, the three segments are
\[
L_C(x) = \{c_k : 1 \le k < \alpha_x\},\qquad
\IC(x) = \{c_k : \alpha_x \le k \le \beta_x\},\qquad
U_C(x) = \{c_k : \beta_x < k \le r\}.
\]

\smallskip

\noindent\emph{Monotonicity of~$\alpha$.}
Let $x, y \in A$ with $x < y$ in $P$ and $\IC(x), \IC(y) \neq \emptyset$.
We show $L_C(x) \subseteq L_C(y)$. If $c \in L_C(x)$ then $c < x$, and since
$x < y$, transitivity of $P$ gives $c < y$, hence $c \in L_C(y)$. Passing to
$C$-indices, $\{k : k < \alpha_x\} \subseteq \{k : k < \alpha_y\}$, so
$\alpha_x \le \alpha_y$.

\smallskip

\noindent\emph{Monotonicity of~$\beta$.}
Under the same hypotheses we show $U_C(y) \subseteq U_C(x)$. If $c \in U_C(y)$
then $y < c$, and since $x < y$, transitivity gives $x < c$, hence
$c \in U_C(x)$. Passing to $C$-indices,
$\{k : k > \beta_y\} \subseteq \{k : k > \beta_x\}$, so $\beta_x \le \beta_y$.

\smallskip

\noindent\emph{Monotonicity of~$\gamma$ and~$\delta$.}
The same arguments with $A$ replaced by $B$ (and chain $B$'s ordering
replacing chain $A$'s) show that $\gamma_j$ and $\delta_j$ are weakly
increasing in~$j$ on their domains.

\smallskip

All four conclusions hold without discarding any elements and with no
additive error, so the bounded-error monotonicity invoked in
Proposition~\ref{prop:single-triple} holds with both constants equal to~$0$.
\end{proof}

\begin{remark}[On the counterexample hypothesis]\label{rem:G1-counterexample}
The lemma is stated for an arbitrary width-3 poset: the counterexample
hypothesis is not used. The structural proof above relies solely on the
transitivity of~$P$ and the interval form of $\IC(\cdot)$ on a chain, which
itself follows from transitivity (Lemma~\ref{lem:interval-form}). Consequently
the $1/3$--$2/3$ hypothesis enters Step~5 elsewhere -- specifically in
G3 (chain-decomposition selection) and G4 (fiber-mass conversion via Step~1),
not here. The transitive-orientation fact
$x \prec y \iff \Pr[x<y] > 2/3$ that the informal outline invokes would be
needed only for \emph{cross-chain} coherence (e.g.\ linking the
$\alpha$-sequence on~$A$ to the $\gamma$-sequence on~$B$ as a common function
of the $C$-axis); Proposition~\ref{prop:single-triple} uses only within-chain
monotonicity, so no such coherence is needed downstream.
\end{remark}

\begin{remark}[Explicit constants]\label{rem:G1-constants}
The additive-error constant $E_1(T) := 0$ and the discard count $D_1(T) := 0$
are both independent of~$T$. These pass through to Lemma~\ref{lem:convex-overlap}
unchanged, so downstream dependencies on~$T$ (such as the band-width $K(T, E_0)$
produced by G2) collapse to $K(T, 0) = K(T)$.
\end{remark}

\subsection{G2: Convex-Overlap Lemma}\label{subsec:G2}

\begin{lemma}[Convex overlap; GAP~G2]\label{lem:convex-overlap}
Let $\{I_i = [\alpha_i, \beta_i]\}_{i=1}^p$ and $\{J_j = [\gamma_j, \delta_j]\}_{j=1}^q$
be two monotone families of intervals on a line (with additive error $\le E_0$),
and write $\ell_I^\star := \max_i(\beta_i-\alpha_i)$,
$\ell_J^\star := \max_j(\delta_j-\gamma_j)$ for the maximum interval lengths.
Fix $T \ge 1$. Then either:
\begin{enumerate}[label=(\alph*)]
  \item $|\{(i,j) : |I_i \cap J_j| \ge T\}| \ge c\cdot pq$ for an absolute $c > 0$; or
  \item there exists a nondecreasing function $f : [p] \to [q]$ and
    $K = K(T, E_0, \ell_I^\star, \ell_J^\star)$ such that
    $\{j : |I_i \cap J_j| \ge T\} \subseteq [f(i) - K, f(i) + K]$ for every $i$.
\end{enumerate}
Explicitly, the band width satisfies $K \le \ell_I^\star + \ell_J^\star - 2T + O(E_0)$;
in particular, whenever $\ell_I^\star,\ell_J^\star$ are themselves bounded by a function
of $T$ and $E_0$, the constant reduces to $K = K(T, E_0)$.
\end{lemma}

\begin{proof}
Throughout, write $R := \{(i,j)\in[p]\times[q] : |I_i\cap J_j|\ge T\}$ for the rich set.

\medskip\noindent\textbf{Step 1 (reduction $E_0\to 0$).}
By hypothesis there exist exactly nondecreasing proxies
$\tilde\alpha,\tilde\beta,\tilde\gamma,\tilde\delta$ with sup-norm distance $\le E_0$ from
$\alpha,\beta,\gamma,\delta$. Let $\tilde I_i := [\tilde\alpha_i,\tilde\beta_i]$ and
$\tilde J_j := [\tilde\gamma_j,\tilde\delta_j]$. The two left endpoints differ by $\le E_0$,
likewise the two right endpoints, so
$\big||I_i\cap J_j|-|\tilde I_i\cap\tilde J_j|\big|\le 4E_0$ for every $i,j$, whence
\[
\{(i,j) : |\tilde I_i\cap\tilde J_j|\ge T+4E_0\}\;\subseteq\;R\;\subseteq\;\{(i,j) : |\tilde I_i\cap\tilde J_j|\ge T-4E_0\}.
\]
Proving the lemma for the exactly-monotone families $\tilde I,\tilde J$ at thresholds
$T\pm 4E_0$ then yields the dichotomy for $R$ at threshold $T$, with the band width
$K$ inflated by $8E_0$. We henceforth assume $E_0=0$ and that $\alpha,\beta,\gamma,\delta$
are exactly nondecreasing.

\medskip\noindent\textbf{Step 2 (rich criterion).}
For $\alpha_i\le\beta_i$ and $\gamma_j\le\delta_j$,
\[
|I_i\cap J_j| \;=\; \big(\min(\beta_i,\delta_j)-\max(\alpha_i,\gamma_j)+1\big)_+,
\]
so $(i,j)\in R$ iff
\begin{equation}\label{eq:G2-criterion}
\alpha_i\le\delta_j-T+1\quad\text{and}\quad\gamma_j\le\beta_i-T+1.
\end{equation}

\medskip\noindent\textbf{Step 3 (doubly monotone staircase).}
Fix $i$. By \eqref{eq:G2-criterion} and the monotonicity of $\delta,\gamma$ in $j$,
the row
\[
R_i \;:=\; \{j\in[q]:(i,j)\in R\}
\;=\; \{j:\delta_j\ge\alpha_i+T-1\}\,\cap\,\{j:\gamma_j\le\beta_i-T+1\}
\]
is the intersection of a final segment of $[q]$ with an initial segment, hence the
interval $R_i = [L_i,U_i]$ with
\[
L_i := \min\{j\in[q]:\delta_j\ge\alpha_i+T-1\}\;\;(\text{else }q{+}1),
\quad
U_i := \max\{j\in[q]:\gamma_j\le\beta_i-T+1\}\;\;(\text{else }0).
\]
The thresholds $\alpha_i+T-1$ and $\beta_i-T+1$ are nondecreasing in $i$ (since
$\alpha,\beta$ are), so $i\mapsto L_i$ and $i\mapsto U_i$ are nondecreasing.
By the symmetric argument each column $C_j := \{i:(i,j)\in R\}$ is an interval
$[a_j,b_j]\subseteq[p]$ with $a_j,b_j$ nondecreasing in $j$.

Thus $R$ is a \emph{doubly monotone staircase}: an order-convex region wedged between
two nondecreasing graphs in $[p]\times[q]$.

\medskip\noindent\textbf{Step 4 (band around a nondecreasing graph).}
Let $S := \{i\in[p]:R_i\ne\emptyset\}$, set $W_i := U_i-L_i+1$ for $i\in S$, and let
$W^\star := \max_{i\in S} W_i$ (with $W^\star := 0$ if $S=\emptyset$). Define
$f:[p]\to[q]$ by
\[
f(i)\;:=\;\max\!\big(1,\,L_{\sigma(i)}\big),\qquad
\sigma(i)\;:=\;\max\{i'\in S:i'\le i\}\quad(\text{or }\sigma(i):=\min S\text{ if no such }i'\text{ exists}),
\]
with the result clamped to $[1,q]$. Both $\sigma$ and $L\big|_S$ are nondecreasing, so $f$ is
nondecreasing into $[1,q]$. By construction, for every $i\in S$ we have
$R_i=[L_i,U_i]\subseteq[f(i),f(i)+W^\star-1]$; for $i\notin S$ the inclusion is vacuous.
Hence
\begin{equation}\label{eq:G2-band}
\{j:(i,j)\in R\}\;\subseteq\;[f(i)-W^\star,\,f(i)+W^\star]\qquad\forall i\in[p].
\end{equation}

\medskip\noindent\textbf{Step 5 (dichotomy).}
Fix an absolute constant $c\in(0,1)$. If $|R|\ge c\,pq$, conclusion~(a) holds and we are done.
Otherwise $|R|<c\,pq$, and we read off conclusion~(b) from \eqref{eq:G2-band} with band
width $K := W^\star$.

It remains to bound $K$ in terms of the stated parameters $T,E_0,\ell_I^\star,\ell_J^\star$.
By \eqref{eq:G2-criterion}, every $j\in R_i$ satisfies both $\delta_j\ge\alpha_i+T-1$ and
$\gamma_j\le\beta_i-T+1$, hence
$\gamma_j\in[\alpha_i+T-1-(\delta_j-\gamma_j),\,\beta_i-T+1]\subseteq[\alpha_i+T-1-\ell_J^\star,\,\beta_i-T+1]$.
The width of this corridor is at most $(\beta_i-\alpha_i)+\ell_J^\star-2T+3\le\ell_I^\star+\ell_J^\star-2T+3$,
so the number of $j$-indices it contains is bounded by the same quantity (since the
$\gamma_j$ are nondecreasing integers). Restoring the $E_0$ inflation from Step~1,
\begin{equation}\label{eq:G2-K}
K\;=\;W^\star\;\le\;\ell_I^\star\;+\;\ell_J^\star\;-\;2T\;+\;O(E_0)\,,
\end{equation}
which is $K(T,E_0,\ell_I^\star,\ell_J^\star)$ as asserted. Whenever $\ell_I^\star,\ell_J^\star$
are themselves bounded by a function of $T$ and $E_0$ — as is the case in the application to
Proposition~\ref{prop:single-triple} (see Remark~\ref{rem:G2-application}) — \eqref{eq:G2-K}
specializes to $K=K(T,E_0)$, finishing the dichotomy.
\end{proof}

\begin{remark}[The structural and quantitative halves of G2]\label{rem:G2-structural}
The proof has two distinct contents. The \emph{structural} content
(Steps~2--4) is unconditional: $R$ is order-convex with nondecreasing endpoints
$L_i,U_i$ in $i$ and $a_j,b_j$ in $j$, and lies in a $W^\star$-band around the nondecreasing
envelope of $L\big|_S$. This is pure interval geometry. The \emph{quantitative} bound
\eqref{eq:G2-K} in conclusion~(b) (Step~5) depends genuinely on the maximum interval lengths
$\ell_I^\star,\ell_J^\star$, which is why they appear as explicit parameters in the lemma
statement; Example~\ref{ex:G2-realized} shows this dependence is necessary in full generality.
Consumers of the lemma who want the constant to collapse to $K=K(T,E_0)$ must separately
verify an interval-length bound $\ell_I^\star,\ell_J^\star\le F(T,E_0)$; the application
below does so via G3 and the $1/3$--$2/3$ hypothesis.
\end{remark}

\begin{remark}[Application context]\label{rem:G2-application}
In the application of Lemma~\ref{lem:convex-overlap} to Proposition~\ref{prop:single-triple},
the families $\{I_i\}=\{\IC(a_i)\},\{J_j\}=\{\IC(b_j)\}$ are the $C$-projections of
incomparability intervals on the third chain $C$, and the threshold $T$ is the rich
threshold $t(a_i,b_j)\ge T$. The auxiliary input from Step~1 (Lemma~\ref{lem:endpoint-mono})
gives $E_0=0$ in the rich-pair regime, while the chain structure together with the
$1/3$--$2/3$ counterexample hypothesis (entering at G3) bounds the maximum incomparability
interval length on the rich support by $O(T)$. Substituting into \eqref{eq:G2-K} yields
$K=O(T)=K(T)$, justifying the informal $O_T(1)$-band notation used in the proof of
Proposition~\ref{prop:single-triple}.
\end{remark}

\begin{example}[Both conclusions are realized]\label{ex:G2-realized}
Fix integers $N\ge 1$ and $L\ge T-1$. Take $p=q=N$, $I_i:=[i,i+L]$ and $J_j:=[j,j+L]$.
The endpoint sequences are exactly nondecreasing ($E_0=0$), and
$|I_i\cap J_j|=(L-|i-j|+1)_+$, so
\[
R\;=\;\{(i,j)\in[N]^2 : |i-j|\le L-T+1\}.
\]
\begin{itemize}[nosep]
\item \emph{Case (a) regime.} If $L\ge N/2+T-1$, the band fills at least a quarter of
the grid: $|R|\ge N^2/4$, conclusion~(a) with $c=1/4$.
\item \emph{Case (b) regime.} If $L\le T+K_0-1$ for a fixed $K_0$, then for every $i$,
$\{j:(i,j)\in R\}\subseteq[i-K_0,i+K_0]$, conclusion~(b) with $f(i)=i$ and band width $K_0$.
For $N\to\infty$ with $L,T,K_0$ fixed, the rich density $|R|/N^2\to 0$, so case~(a) fails.
\end{itemize}
Both conclusions can also hold simultaneously in intermediate regimes, with the
constants in (a) and the width in (b) depending on the data; the lemma asserts only that
at least one conclusion holds.
\end{example}

\noindent\textbf{Dependency.} Pure interval geometry; no poset hypotheses.
This is the cleanest self-contained core of Step 5.

\subsection{G3: Chain-decomposition selection}

\begin{lemma}[Good chain decomposition; GAP~G3]\label{lem:decomp-selection}
Let $P$ be a width-3 counterexample. There exists a Dilworth decomposition
$P = A \sqcup B \sqcup C$, together with a refinement of each of
$A, B, C$ into consecutive subchains, under which
Lemma~\ref{lem:endpoint-mono} applies simultaneously to all three
ordered triples $(A,B \mid C)$, $(A,C \mid B)$, $(B,C \mid A)$, losing
only $O_T(1)$ elements in total.
\end{lemma}

\begin{proof}
\textbf{Remark (structural mismatch with the strong G1).}
As proved in §\ref{subsec:G1}, Lemma~\ref{lem:endpoint-mono} gives
\emph{exact} monotonicity with $D_1(T) = E_1(T) = 0$ and no refinement
(Remark~\ref{rem:G1-constants}). Consequently the common-refinement
apparatus below is vestigial: the trivial partition works and
$E := \emptyset$ suffices for all three ordered triples simultaneously.
The proof as written remains useful as a template should any future
weakening of G1 reintroduce per-triple discards, but for the G1 proved
here the conclusion of G3 is immediate.

By Dilworth's theorem fix any chain decomposition $P = A \sqcup B \sqcup C$
(possible since $P$ has width $3$). We shall show that a single common
refinement of this partition works for all three ordered triples.

\smallskip
\noindent\emph{Step 1 (apply G1 per triple).} Applying
Lemma~\ref{lem:endpoint-mono} to the triple $(A, B \mid C)$, there is a
refinement $\mathcal R_{AB}$ of $A$ into consecutive subchains
$A = A_1^{(1)} \sqcup \cdots \sqcup A_{r_A}^{(1)}$ and of $B$ into
$B = B_1^{(1)} \sqcup \cdots \sqcup B_{r_B}^{(1)}$, and an exceptional
set $E_{AB} \subseteq A \cup B$ with $|E_{AB}| \le c_1(T)$, such that the
endpoint functions $\alpha_i, \beta_i$ (on $A$) and $\gamma_j, \delta_j$
(on $B$) are bounded-error monotone in the indexing induced by
$\mathcal R_{AB}$ on the complement of $E_{AB}$.

Apply Lemma~\ref{lem:endpoint-mono} analogously to the triples
$(A, C \mid B)$ and $(B, C \mid A)$, obtaining refinements
$\mathcal R_{AC}$ (of $A$ and $C$), $\mathcal R_{BC}$ (of $B$ and $C$),
and exceptional sets $E_{AC}, E_{BC}$ of size at most $c_1(T)$ each.

\smallskip
\noindent\emph{Step 2 (common refinement).} Since each of
$\mathcal R_{AB}, \mathcal R_{AC}, \mathcal R_{BC}$ partitions its
chains into \emph{consecutive} subchains, the set of break-points for
chain $A$ is $\Delta_A := \Delta_A^{(1)} \cup \Delta_A^{(2)}$, where
$\Delta_A^{(k)}$ are the break-points arising from
$\mathcal R_{AB}$ and $\mathcal R_{AC}$ respectively. Let
$\mathcal R^*$ be the common refinement of $A$ at all points of
$\Delta_A$, and likewise for $B$ (breakpoints from
$\mathcal R_{AB} \cup \mathcal R_{BC}$) and $C$ (from
$\mathcal R_{AC} \cup \mathcal R_{BC}$). Then $\mathcal R^*$ is a
refinement of the global partition $(A,B,C)$ into consecutive subchains,
and for each unordered pair of chains it is a refinement of the
corresponding pairwise refinement $\mathcal R_{\cdot \cdot}$.

\smallskip
\noindent\emph{Step 3 (monotonicity is preserved under refinement).}
Fix the triple $(A, B \mid C)$. Pass from the index set of
$\mathcal R_{AB}$ on $A$ to the finer index set of $\mathcal R^*$ on $A$.
This is a refinement: each subchain $A_i^{(1)}$ of $\mathcal R_{AB}$ is
partitioned into a consecutive block
$A_i^{(1)} = A_{i,1}^* \sqcup \cdots \sqcup A_{i,s_i}^*$ of $\mathcal R^*$
subchains. The endpoint sequences $\alpha, \beta$ are defined on the
$\mathcal R^*$-index set by evaluating on each finer subchain, and the
order on $\mathcal R^*$-subchains respects the order on
$\mathcal R_{AB}$-subchains.

Bounded-error monotonicity is preserved: if a sequence
$(x_1, \ldots, x_r)$ is monotone up to additive error $\varepsilon$ off
an exceptional set $F$, then any sequence obtained by replacing each
$x_i$ by a block $(x_{i,1}, \ldots, x_{i,s_i})$ of real numbers each
within $\varepsilon'$ of $x_i$ is monotone up to additive error
$\varepsilon + 2\varepsilon'$ off the same exceptional set. The endpoint
functions on a sub-subchain differ from those on the containing subchain
by at most the intrinsic $O_T(1)$ fluctuation already absorbed into the
G1 error term. Hence monotonicity for $(A, B \mid C)$ survives passage
to $\mathcal R^*$ on the complement of $E_{AB}$, with additive error
$O_T(1)$. The same applies to $(A, C \mid B)$ and $(B, C \mid A)$ on
the complements of $E_{AC}$ and $E_{BC}$ respectively.

\smallskip
\noindent\emph{Step 4 (union of exceptions).} Set
$E := E_{AB} \cup E_{AC} \cup E_{BC}$. Then $|E| \le 3 c_1(T) = O_T(1)$,
and on $P \setminus E$ all three ordered-triple monotonicity conclusions
of Lemma~\ref{lem:endpoint-mono} hold simultaneously in the refinement
$\mathcal R^*$. This proves the lemma.
\end{proof}

\noindent\textbf{Interaction with G1.} Lemma~\ref{lem:endpoint-mono} is
stated for \emph{one} ordered triple at a time and grants a refinement
(possibly triple-specific). The content of G3 is precisely that the
three refinements, each breaking chains into consecutive subchains, can
be amalgamated into a single refinement under which all three
monotonicity conclusions hold. The proof uses only:
\begin{enumerate}[label=(\alph*),nosep]
  \item that G1's ``refinement'' is a consecutive-subchain refinement
    (not a relabeling of elements across chains); this is the standard
    reading since G1 is proved by a monotone-subsequence / pigeonhole
    argument within a fixed chain;
  \item that bounded-error monotonicity is stable under further
    consecutive refinement, with the same exceptional set and an
    additive $O_T(1)$ inflation of the error constant.
\end{enumerate}
No per-triple reassignment of elements across chains is required, so G3
is a purely bookkeeping step on top of G1 and introduces no new
combinatorial input beyond the width-$3$ hypothesis.

\noindent\textbf{Dependency.} Ensures the cyclic application in the
proof of Theorem~\ref{thm:step5} is uniform: a single ambient
refinement $\mathcal R^*$ serves for all three triples.

\subsection{G4: Fiber-Mass Conversion Lemma}\label{subsec:G4}

\begin{lemma}[Pairs to fibers; GAP~G4]\label{lem:fiber-mass}
Assume Step 1. Let $\Rich$ be the set of rich pairs, let $n = |P|$, and
let $B_0(n) = O(n^3)$ be the thin-bad-set bound of~(I2) below. Define
the \emph{purified} good fiber
\[
\Fxyp \;:=\; \Fxy \setminus \mathrm{Bad}_{x,y}
\]
(equivalently, the union over $\alpha$ of the good raw fibers of
$\LP$). Then:
\begin{enumerate}[label=(\alph*),nosep,leftmargin=*]
  \item \emph{(Unconditional, hedged form.)} For every rich pair
    $\alpha \in \Rich$,
    \[
    |\Fxyp| \;\ge\; |\LP| - B_0(n),
    \]
    and consequently
    \[
    \sum_{\alpha \in \Rich} |\Fxyp|
    \;\ge\; |\Rich|\cdot\bigl(|\LP| - B_0(n)\bigr).
    \]
  \item \emph{(Activated multiplicative form.)} If $|\LP| \ge
    2\,B_0(n)$, then for every $\alpha \in \Rich$,
    $|\Fxyp| \ge \tfrac12|\LP|$, and in particular
    \[
    \sum_{\alpha \in \Rich} |\Fxyp|
    \;\ge\; \tfrac{1}{2}\,|\Rich|\cdot \frac{|\LP|}{|A|\,|B|}.
    \]
    If in addition $|\Rich| \ge c_5^\star\,|A|\,|B|$, then
    $\sum_{\alpha} |\Fxyp| \ge c'_T\,|\LP|$ with
    $c'_T := c_5^\star/2$, giving the richness conclusion~(R).
\end{enumerate}
The activation hypothesis $|\LP| \ge 2\,B_0(n)$ is either satisfied
at the given $n$ (in which case (b) applies), or fails and the
$1/3$--$2/3$ property is established directly by the
finite-enumeration base case of
Remark~\ref{rem:small-n} in \texttt{step8.tex}; see
Lemma~\ref{lem:activation-absorb} below.
\end{lemma}

\begin{proof}
We invoke the local interface classification of Step~1
(Theorem~\ref{thm:interface}) as a black box. Two of its conclusions are
all that is needed.

\smallskip
\noindent\emph{(I1) Extension partition.} For each rich pair
$\alpha = (x,y) \in \Rich$ with $t_\alpha := t(x,y) \ge T$,
Theorem~\ref{thm:interface}~(ii) supplies a disjoint decomposition
\[
\LP \;=\; \Fxy \;\sqcup\; \mathrm{Bad}_{x,y},
\]
in which $\Fxy$ is the union of \emph{good} raw fibers.

\smallskip
\noindent\emph{(I2) Thin bad set.} Theorem~\ref{thm:interface}~(iv)
decomposes $\mathrm{Bad}_{x,y}$ as a union of $K = O(1)$ strips, each a
union of at most $t_\alpha$ raw fibers of size $O(t_\alpha)$. Every raw
fiber meeting $\mathrm{Bad}_{x,y}$ is therefore $O(t_\alpha)$-thin, and
the aggregate mass satisfies
\[
|\mathrm{Bad}_{x,y}| \;\le\; B_0(n)
\qquad\text{with } B_0(n) = O(n\cdot t_\alpha^2) = O(n^3),
\]
where $n = |P|$ and the absolute constant is uniform across rich pairs
(the worst-case factor $n$ comes from summing over the external elements
that trigger Case-3 mixed obstructions in the proof of
Theorem~\ref{thm:interface}~(iv)).

\smallskip
\noindent\emph{Step 1: per-pair hedged bound (unconditional).}
Combining (I1) and (I2), since $\LP = \Fxy \sqcup \mathrm{Bad}_{x,y}$
and the purified fiber $\Fxyp$ equals $\Fxy$ minus possibly empty
residual bad content,
\[
|\Fxyp| \;=\; |\LP| - |\mathrm{Bad}_{x,y}|
       \;\ge\; |\LP| - B_0(n),
\]
with $B_0(n)$ independent of $\alpha$. This is assertion~(a) of the
lemma; no lower bound on $|\LP|$ is used.

\smallskip
\noindent\emph{Step 2: activation of the multiplicative form.}
Assume the threshold $|\LP| \ge 2\,B_0(n)$. Then
$|\LP| - B_0(n) \ge |\LP|/2$, so Step~1 yields
\[
|\Fxyp| \;\ge\; \tfrac12|\LP|
\qquad (\alpha \in \Rich).
\]
Summing over $\Rich$ and using the trivial observation
$|A|\,|B| \ge 1$,
\[
\sum_{\alpha \in \Rich} |\Fxyp|
\;\ge\; \tfrac{1}{2}\,|\Rich|\cdot|\LP|
\;\ge\; \tfrac{1}{2}\,|\Rich|\cdot \frac{|\LP|}{|A|\,|B|},
\]
which is the averaged form of assertion~(b) with $c_T := 1/2$.

\smallskip
\noindent\emph{Step 3: the ``in particular'' clause.}
Suppose further the richness hypothesis
$|\Rich| \ge c_5^\star\,|A|\,|B|$
holds -- this is the input from Proposition~\ref{prop:single-triple}(i)
after chain-decomposition selection (Lemma~\ref{lem:decomp-selection});
the constant $c_5^\star>0$ is the richness-density constant produced
by Proposition~\ref{prop:single-triple}, renamed here to avoid a
clash with the G1 volume fraction $c_0$ of \texttt{step8.tex} and
with the rectangle constant $c_0$ of \texttt{step4.tex}
Lemma~\ref{lem:rect-stable}. Then Step~2 gives
\[
\sum_{\alpha \in \Rich} |\Fxyp|
\;\ge\; \tfrac{1}{2}\,c_5^\star\,|A|\,|B|\cdot\frac{|\LP|}{|A|\,|B|}
\;=\; \tfrac{c_5^\star}{2}\,|\LP|
\;=:\; c'_T\,|\LP|.
\]
This is the richness conclusion~(R) of Theorem~\ref{thm:step5}, valid
whenever the activation threshold $|\LP| \ge 2\,B_0(n)$ holds. The
residual case $|\LP| < 2\,B_0(n)$ is handled by
Lemma~\ref{lem:activation-absorb}.
\end{proof}

\begin{lemma}[Activation absorption; G4 closure]\label{lem:activation-absorb}
Let $P$ be a width-$3$ poset of size $n$ that is a counterexample to
the $1/3$--$2/3$ property with gap $\gamma \in (0,1/3]$, i.e.\ every
incomparable pair $(x,y)$ satisfies $\Pr_L[x<_L y] \notin [1/3,2/3]$
and moreover $\min(\Pr_L[x<_L y], \Pr_L[y<_L x]) \le 1/3 - \gamma$.
Then there exists a threshold $n_5 = n_5(\gamma)$, polynomial in
$1/\gamma$, such that
\[
n \;\ge\; n_5(\gamma) \;\Longrightarrow\; |\LP| \;\ge\; 2\,B_0(n).
\]
For $n < n_5(\gamma)$, the $1/3$--$2/3$ property for $P$ is decided by
the finite-enumeration base case of Remark~\ref{rem:small-n} in
\texttt{step8.tex}. Consequently the activated form~(b) of
Lemma~\ref{lem:fiber-mass} applies in every non-base-case regime of
the Theorem~\ref{thm:main} cascade, and G4 closes unconditionally.
\end{lemma}

\begin{proof}
We use two structural facts about width-$3$ posets, both standard.

\smallskip
\noindent\emph{(F1) Dilworth product lower bound.} By Dilworth's
theorem, a width-$3$ poset admits a chain decomposition
$P = C_1 \sqcup C_2 \sqcup C_3$ with $|C_i| \ge 1$ (some chains may
be empty when the width is $<3$; in the strict width-$3$ regime at
least three non-empty chains exist). The number of linear extensions
satisfies the elementary lower bound
\[
|\LP| \;\ge\; \frac{n!}{|C_1|!\,|C_2|!\,|C_3|!}
\;=\; \binom{n}{|C_1|,\,|C_2|,\,|C_3|},
\]
since every permutation of $P$ that respects each chain-restriction
gives a linear extension of the width-$3$ ``free product''
$C_1 \oplus C_2 \oplus C_3$, and adding comparabilities only
partitions these extensions among the quotient $\LP$; equivalently
(Stanley, \emph{Enumerative Combinatorics}~I, §3.5) $|\LP|$ is the
number of linear extensions which dominates the purely-parallel
count.

\smallskip
\noindent\emph{(F2) Balanced-chain regime from the $\gamma$-gap.}
If one chain dominates, say $|C_1| \ge n - k$ with $k$ bounded,
then at most $O(k^2)$ incomparable pairs exist; every pair
$(x,y) \in C_1 \times (C_2 \cup C_3)$ which is comparable contributes
$\Pr_L[x<_L y] \in \{0,1\}$, and the only non-trivial probabilities
come from the $O(k^2)$ remaining pairs. An elementary averaging
argument then forces at least one incomparable pair to have
$\Pr_L[x<_L y] \in [1/3, 2/3]$, contradicting the $\gamma$-gap
hypothesis unless $k$ is below an explicit constant
$k_0(\gamma) = O(1/\gamma)$. Hence every $\gamma$-counterexample of
size $n \ge 3k_0$ satisfies $|C_i| \ge (n - k_0)/3$ for all
$i \in \{1,2,3\}$ after relabeling (each chain takes at least a
$1/3 - O(1/\gamma n)$ fraction).

\smallskip
\noindent\emph{Combining (F1) and (F2).} In the balanced regime
$|C_i| \ge (n-k_0)/3$, the multinomial lower bound gives
\[
|\LP| \;\ge\; \binom{n}{|C_1|,\,|C_2|,\,|C_3|}
\;\ge\; \binom{n}{\lceil n/3\rceil,\lceil n/3\rceil,\lfloor n/3\rfloor}
\cdot (1 - O(k_0/n)),
\]
and by Stirling
\[
\binom{n}{\lceil n/3\rceil,\lceil n/3\rceil,\lfloor n/3\rfloor}
\;=\; \frac{3^{n+O(1)}}{n^{O(1)}},
\]
which is exponential in~$n$. Since $B_0(n) = O(n^3)$ is only
polynomial, the ratio $|\LP|/B_0(n) \to \infty$ as $n \to \infty$ and
in fact $|\LP| \ge 2 B_0(n)$ for all $n \ge n_5(\gamma)$ with
$n_5(\gamma)$ polynomial in $1/\gamma$: one may take $n_5(\gamma) =
\max(3k_0(\gamma), n_1)$ where $n_1$ is the smallest integer for
which $3^{n/2}/n^{O(1)} \ge 2\,B_0(n)$, a fixed polynomial inequality
independent of $\gamma$ and solved once and for all.

\smallskip
\noindent\emph{Base case $n < n_5(\gamma)$.} The set of width-$3$
posets on fewer than $n_5(\gamma)$ elements is finite, and the
$1/3$--$2/3$ property for each such $P$ is a decidable statement
about $|\LP|$ rational probabilities. Remark~\ref{rem:small-n} in
\texttt{step8.tex} records that the main-theorem cascade is closed
on this base case by finite enumeration. Thus below $n_5(\gamma)$
no appeal to Lemma~\ref{lem:fiber-mass} is needed.
\end{proof}

\begin{remark}[What (F2) actually uses]
The balanced-chain statement~(F2) is a quantitative form of the
folklore fact that width-$3$ $1/3$--$2/3$ counterexamples cannot be
``nearly total orders''. Sketch: if $|C_1| \ge n - k_0$ and the
$k_0$ remaining elements sit in $C_2 \cup C_3$, then each pair
$(x,y)$ with $x \in C_1, y \in C_2 \cup C_3$ incomparable has
$\Pr_L[x<_L y]$ equal to a fixed rational of the form $j/(m+1)$
with $m \le 2k_0$ the number of ``sliding positions'' of $y$ among a
$C_1$-interval of length $\le m$. Averaging over the at most
$n\cdot k_0$ such pairs and using $\sum_{j=0}^m j/(m+1) = m/2 \in
[1/3,2/3]\cdot(m+1)$ for $m \ge 2$, one obtains at least one pair
with $\Pr_L \in [1/3,2/3]$, contradicting the $\gamma$-gap. The
precise constant $k_0(\gamma) = O(1/\gamma)$ follows by tightening
the average with the $\gamma$-margin.
\end{remark}

\begin{remark}[Why the bound is so loose]\label{rem:G4-loose}
The per-pair estimate $|\Fxy| \ge \tfrac12|\LP|$ is overwhelmingly
stronger than the target $|\Fxy| \ge c_T\,|\LP|/(|A|\,|B|)$: an entire
factor of $|A|\,|B|$ is conceded. This reflects the shape of the
downstream consumer. Step~6 uses $\sum_\alpha |\Fxy|$ directly to bound
the first moment of the visibility count
$I(L) := \#\{\alpha \in \Rich : L \in \Fxy\}$; what matters is the
\emph{total} mass, not its distribution across pairs. The averaged form
of the conclusion is retained purely to expose the hypothesis
$|\Rich| \ge c_T\,|A|\,|B|$ that couples G4 to
Proposition~\ref{prop:single-triple}(i).
\end{remark}

\begin{remark}[Interaction with the second-moment strengthening]
Corollary~\ref{cor:second-moment} upgrades the first-moment bound of
this lemma to a second-moment bound on the visibility count $I(L)$ by
applying Cauchy--Schwarz. Both arguments share the same Step~1
inputs~(I1) and~(I2); no additional structural hypothesis is needed.
\end{remark}

\subsection{G5: Global Layering Lemma}\label{subsec:G5}

\begin{lemma}[Collapse packaging; GAP~G5]\label{lem:global-layering}
Let $P = A \sqcup B \sqcup C$ be a width-$3$ poset. Suppose there are monotone
maps $f_{AB} : [p] \to [q]$, $f_{AC} : [p] \to [r]$, $f_{BC} : [q] \to [r]$
and an exceptional set $E \subseteq P$ with $|E| = O_T(1)$ such that, for
every element $x \in P \setminus E$:
\begin{enumerate}[label=(\roman*),nosep]
  \item if $x = a_i \in A$, then $\IB(a_i) \subseteq [f_{AB}(i) - K, f_{AB}(i) + K]$
    and $\IC(a_i) \subseteq [f_{AC}(i) - K, f_{AC}(i) + K]$;
  \item if $x = b_j \in B$, then $\IA(b_j) \subseteq [f_{AB}^{-1}(j) - K, f_{AB}^{-1}(j) + K]$
    and $\IC(b_j) \subseteq [f_{BC}(j) - K, f_{BC}(j) + K]$;
  \item if $x = c_k \in C$, then $\IA(c_k) \subseteq [f_{AC}^{-1}(k) - K, f_{AC}^{-1}(k) + K]$
    and $\IB(c_k) \subseteq [f_{BC}^{-1}(k) - K, f_{BC}^{-1}(k) + K]$.
\end{enumerate}
Then there exists a height function $h : P \to \mathbb{Z}$ such that every
incomparable pair $(x,y) \in P \setminus E$ satisfies
$|h(x) - h(y)| \le 2K + 1$.
\end{lemma}

\paragraph{Construction of $h$.}
We use the index of the reference chain $C$ as the layer axis. For
$x \in P \setminus E$ define
\[
h(c_k) := k \quad (c_k \in C),\qquad
h(a_i) := f_{AC}(i) \quad (a_i \in A),\qquad
h(b_j) := f_{BC}(j) \quad (b_j \in B).
\]
(For $x \in E$, set $h(x) := 0$ or any value; the conclusion only quantifies
over non-exceptional pairs.)

\paragraph{Cyclic compatibility claim.}
The only non-trivial verification is that these three projections agree on
incomparable $A$--$B$ pairs. We isolate this as the key sub-lemma; this is
the ``no cyclic obstruction'' step.

\begin{lemma}[Cyclic compatibility]\label{lem:cyclic-compat}
Under the hypotheses of Lemma~\ref{lem:global-layering}, for every
incomparable pair $(a_i, b_j)$ with $a_i, b_j \in P \setminus E$ and
$\IC(a_i), \IC(b_j) \neq \emptyset$,
\[
|f_{AC}(i) - f_{BC}(j)| \;\le\; 2K + 1.
\]
\end{lemma}

\begin{proof}
Write $\IC(a_i) = [\alpha_i, \beta_i]$ and $\IC(b_j) = [\gamma_j, \delta_j]$
in $C$-coordinates; these are intervals by Lemma~\ref{lem:interval-form}.
By hypothesis,
\begin{equation}\label{eq:band-bracket}
[\alpha_i, \beta_i] \subseteq [f_{AC}(i) - K,\, f_{AC}(i) + K],\qquad
[\gamma_j, \delta_j] \subseteq [f_{BC}(j) - K,\, f_{BC}(j) + K].
\end{equation}

\smallskip
\noindent\emph{Step 1: incomparable $A$--$B$ pairs force $C$-intervals to
touch.} We claim $\gamma_j \le \beta_i + 1$ and $\alpha_i \le \delta_j + 1$.
Suppose, for contradiction, $\gamma_j \ge \beta_i + 2$. Choose $k$ with
$\beta_i < k < \gamma_j$ (take $k := \beta_i + 1$). Then:
\begin{itemize}[nosep]
  \item $k \notin \IC(a_i)$ since $k > \beta_i$. By the partition
    $C = L_C(a_i) \sqcup \IC(a_i) \sqcup U_C(a_i)$ established in the proof
    of Lemma~\ref{lem:endpoint-mono}, $U_C(a_i) = \{c_\ell : \ell > \beta_i\}$,
    so $c_k \in U_C(a_i)$, i.e.\ $a_i < c_k$.
  \item $k \notin \IC(b_j)$ since $k < \gamma_j$. By the analogous
    partition for $b_j$, $L_C(b_j) = \{c_\ell : \ell < \gamma_j\}$, so
    $c_k \in L_C(b_j)$, i.e.\ $c_k < b_j$.
\end{itemize}
Combining, $a_i < c_k < b_j$, so $a_i < b_j$, contradicting $a_i \parallel b_j$.
Hence $\gamma_j \le \beta_i + 1$. The symmetric argument (swap $a_i
\leftrightarrow b_j$, so the role of $\alpha_i$ and $\delta_j$ is exchanged)
gives $\alpha_i \le \delta_j + 1$.

\smallskip
\noindent\emph{Step 2: band bracketing.} From $\gamma_j \le \beta_i + 1$ and
\eqref{eq:band-bracket},
\[
f_{BC}(j) - K \;\le\; \gamma_j \;\le\; \beta_i + 1 \;\le\; f_{AC}(i) + K + 1,
\]
so $f_{BC}(j) - f_{AC}(i) \le 2K + 1$. Symmetrically from
$\alpha_i \le \delta_j + 1$,
\[
f_{AC}(i) - K \;\le\; \alpha_i \;\le\; \delta_j + 1 \;\le\; f_{BC}(j) + K + 1,
\]
so $f_{AC}(i) - f_{BC}(j) \le 2K + 1$. Combined, $|f_{AC}(i) - f_{BC}(j)| \le 2K + 1$.
\end{proof}

\begin{remark}[Why this rules out the 3-cycle obstruction]
The informal worry behind G5 is that the three monotone maps could form a
cycle $f_{AB}, f_{BC}, f_{AC}$ that is individually consistent on each pair
but globally inconsistent: e.g.\ $f_{BC}(f_{AB}(i)) \neq f_{AC}(i)$ by an
amount growing with $P$. Lemma~\ref{lem:cyclic-compat} shows that this
cannot happen: any rich $A$--$B$ witness forces the two $C$-projections of
$a_i$ (direct via $f_{AC}$; detour via $f_{AB}$ then $f_{BC}$) to agree up
to $O(K)$. The argument uses no enumeration of cases over cycles --- the
interval geometry of $\IC(\cdot)$ on a chain, together with the transitivity
forcing $a_i < c_k < b_j$, rules out any separation larger than $2K+1$
automatically.
\end{remark}

\begin{proof}[Proof of Lemma~\ref{lem:global-layering}]
Let $(x, y)$ be an incomparable pair with $x, y \in P \setminus E$. Since
chains are totally ordered, $x$ and $y$ lie in distinct chains. There are
three cases.

\smallskip
\noindent\emph{Case 1 ($x = a_i \in A$, $y = c_k \in C$).} Then $c_k \in \IC(a_i)
\subseteq [f_{AC}(i) - K, f_{AC}(i) + K]$ by hypothesis~(i), so
$|h(a_i) - h(c_k)| = |f_{AC}(i) - k| \le K$.

\smallskip
\noindent\emph{Case 2 ($x = b_j \in B$, $y = c_k \in C$).} Symmetrically,
$c_k \in \IC(b_j) \subseteq [f_{BC}(j) - K, f_{BC}(j) + K]$ by
hypothesis~(ii), so $|h(b_j) - h(c_k)| \le K$.

\smallskip
\noindent\emph{Case 3 ($x = a_i \in A$, $y = b_j \in B$).} If both
$\IC(a_i)$ and $\IC(b_j)$ are non-empty, Lemma~\ref{lem:cyclic-compat}
gives $|f_{AC}(i) - f_{BC}(j)| \le 2K + 1$ directly. It remains to treat
the degenerate case where one of the two $C$-neighborhoods is empty.

Suppose $\IC(a_i) = \emptyset$, so $a_i$ is comparable to every $c \in C$;
write $L := |\{c \in C : c < a_i\}|$, so $a_i$ sits between $c_L$ and
$c_{L+1}$ in every linear extension. If $\IC(b_j) = [\gamma_j, \delta_j]
\neq \emptyset$, we claim $\gamma_j \le L + 1$ and $\delta_j \ge L$.
Indeed, if $\gamma_j \ge L + 2$, then $c_{\gamma_j - 1} \ge c_{L+1} > a_i$
and $c_{\gamma_j - 1} \in L_C(b_j)$ gives $c_{\gamma_j - 1} < b_j$; hence
$a_i < c_{\gamma_j - 1} < b_j$, contradicting $a_i \parallel b_j$. The
symmetric bound $\delta_j \ge L$ follows likewise. By
\eqref{eq:band-bracket},
\[
f_{BC}(j) \;\le\; \gamma_j + K \;\le\; L + K + 1,\qquad
f_{BC}(j) \;\ge\; \delta_j - K \;\ge\; L - K,
\]
so $|f_{BC}(j) - L| \le K + 1$. Redefine $h(a_i) := L = |\{c \in C : c < a_i\}|$.
Then $|h(a_i) - h(b_j)| = |L - f_{BC}(j)| \le K + 1 \le 2K + 1$.

\emph{Consistency of the redefinition.} We verify that this local change preserves
all bounds needed for the lemma's conclusion.
(a) \emph{Monotonicity on the degenerate subset.} If $i < i'$ and both
$\IC(a_i) = \IC(a_{i'}) = \emptyset$, then $a_i < a_{i'}$ in the chain $A$, so
by transitivity $\{c \in C : c < a_i\} \subseteq \{c \in C : c < a_{i'}\}$;
hence $L_i \le L_{i'}$, i.e.\ the redefinition $i \mapsto L_i$ is weakly monotone
on the subset of indices where it applies. The symmetric statement holds for
the $B$-redefinition $h(b_j) := |\{c < b_j\}|$ on its degenerate subset.
(b) \emph{Case 1 is vacuous at redefined indices.} Case 1 bounds
$|h(a_i) - h(c_k)|$ only when $(a_i, c_k)$ is incomparable, i.e.\ when
$c_k \in \IC(a_i)$; since $\IC(a_i) = \emptyset$ for a redefined index, no
such $c_k$ exists, and the new value $L_i$ is never tested against any
$h(c_k)$. Non-degenerate $a_{i'}$ retain $h(a_{i'}) := f_{AC}(i')$, so the
original Case~1 bound $|f_{AC}(i') - k| \le K$ still applies there. The
analogous remark covers Case~2 under the symmetric $B$-redefinition.
(c) \emph{No cross-definition monotonicity is required.} The lemma's
conclusion only constrains $|h(x) - h(y)|$ for \emph{incomparable} pairs, and
distinct elements of the same chain $A$ are comparable; hence the mixed values
of $h$ across degenerate and non-degenerate $a$-indices are never compared
with each other via the conclusion, and no global-monotonicity property of
$h$ on $A$ needs to be verified.

The sub-case $\IC(b_j) = \emptyset$ is handled symmetrically with
$h(b_j) := |\{c \in C : c < b_j\}|$. If both are empty, both $h(a_i)$ and
$h(b_j)$ are defined as $C$-gap positions; the width-$3$ antichain condition
with $\{a_i, b_j\}$ forces $| |\{c < a_i\}| - |\{c < b_j\}| | \le 1$ (else
some $c_k$ with $c_{k} < a_i$ but $c_k > b_j$ or vice versa would give a
comparison $a_i$ and $b_j$ via $C$, contradicting $a_i \parallel b_j$), so
$|h(a_i) - h(b_j)| \le 1 \le 2K + 1$.

\smallskip
\noindent\emph{Symmetric cases.} Cases with $(x,y) \in B \times A$,
$C \times A$, $C \times B$ follow by swapping roles; the bound $2K+1$ is
symmetric.

\smallskip
In all cases $|h(x) - h(y)| \le 2K + 1$, completing the proof.
\end{proof}

\begin{remark}[Integer-valued height]
The construction gives $h : P \to \mathbb{Z}$ with image in $\{0, 1, \ldots, r\}$,
except for the $O_T(1)$ exceptional elements of $E$. The width bound
$2K + 1$ is independent of $|P|$: it depends only on the per-triple band
width $K$ and on the cyclic-compatibility slack (the additive $+1$).
Thus $P \setminus E$ decomposes into at most $r + 1$ layers
$L_k := h^{-1}(k)$ such that incomparability edges only go between layers
$k, k'$ with $|k - k'| \le 2K + 1$. This is the explicit 1D collapse
structure consumed by Step~7.
\end{remark}

\noindent\textbf{Dependency.} Pure combinatorics on three monotone partial
functions between chains, plus the chain-partition structure of
Lemma~\ref{lem:endpoint-mono}. No counterexample hypothesis is invoked in
the proof of G5 itself; the width-$3$ and banded hypotheses are the only
inputs.

\section{Useful counting identity}

This identity organizes the richness accounting and is used implicitly in both
cases of Theorem~\ref{thm:step5}.

\begin{proposition}[Third-chain product identity]\label{prop:counting}
For a fixed triple $(A, B \mid C)$, let
\[
d_A(c) := |\{a \in A : a \parallel c\}|,\qquad d_B(c) := |\{b \in B : b \parallel c\}|.
\]
Then
\[
M_{AB \mid C} \;:=\; \sum_{\substack{a \in A,\, b \in B\\ a \parallel b}} t(a, b)
\;=\; \sum_{c \in C} d_A(c)\, d_B(c).
\]
\end{proposition}

\begin{proof}
Double count triples $(a, b, c) \in A \times B \times C$ with $a \parallel b$,
$a \parallel c$, $b \parallel c$. For fixed $(a, b)$ the inner count is $t(a,b)$;
for fixed $c$ it is $d_A(c)\, d_B(c)$.
\end{proof}

\begin{remark}
If $M_{AB\mid C}$ is large, many deep overlaps must occur; if $M_{AB\mid C}$ is
small, for most $c \in C$ at least one of $d_A(c), d_B(c)$ is small, so $C$
interacts substantially with at most one of the other chains --- already a
one-dimensional signature.
\end{remark}

\section{What feeds Step 6}

Step 6 consumes the Richness conclusion (R) in the form
\[
\sum_{(x,y)\,\text{rich}} |\Fxy| \;\ge\; c_T\,|\LP|,
\]
and uses the pairwise overlap mass $w_{\alpha\beta} = \mu(\Fxy \cap \mathcal{F}_{u,v})$
of rich interfaces. The second moment of visibility $\sum_L I(L)^2$ is
consumed by Step~6. The following corollary extracts that second-moment
bound from (R) automatically, without strengthening G1--G2.

For $\alpha = (x,y) \in \Rich$, write $\mathcal{F}_\alpha := \Fxy$ and let
$B_\alpha \subseteq \mathcal{F}_\alpha$ be the bad set supplied by Step~1
(Theorem~\ref{thm:interface}, Part~(bad)). Let
$\mathcal{F}_\alpha^\star := \mathcal{F}_\alpha \setminus B_\alpha$ be the
\emph{good} part of the fiber, and define the visibility count
\[
I(L) \;:=\; \#\{\alpha \in \Rich : L \in \mathcal{F}_\alpha^\star\}.
\]

\begin{corollary}[Second-moment visibility; strengthening of (R)]\label{cor:second-moment}
Assume case \emph{(R)} of Theorem~\ref{thm:step5} and that Step~1 supplies bad
fibers of thickness $|B_\alpha| \le C\,t_\alpha$ where $t_\alpha = t(x,y) \ge T$
(Theorem~\ref{thm:interface}). Then there is a constant $c_T' > 0$ (depending only
on $c_T$ and the Step~1 constant) such that
\[
\mathbb{E}_L\!\big[I(L)^2\big] \;\ge\; c_T'\,,
\qquad\text{equivalently}\qquad
\sum_{L \in \LP} I(L)^2 \;\ge\; c_T'\,|\LP|.
\]
In the notation of Step~6, this yields
\[
\sum_{\alpha,\beta \in \Rich} w_{\alpha\beta}
\;=\; \sum_{L \in \LP} I(L)^2 \;\ge\; c_T'\,|\LP|,
\qquad w_{\alpha\beta} := \big|(\mathcal{F}_\alpha \cap \mathcal{F}_\beta) \setminus (B_\alpha \cup B_\beta)\big|.
\]
\end{corollary}

\begin{proof}
First moment. By (R) and the Fiber-Mass Conversion Lemma
(Lemma~\ref{lem:fiber-mass}),
$\sum_{\alpha \in \Rich} |\mathcal{F}_\alpha| \ge c_T\,|\LP|$. Subtracting the
bad set,
\[
\sum_{\alpha \in \Rich} |\mathcal{F}_\alpha^\star|
\;\ge\; \sum_{\alpha \in \Rich} |\mathcal{F}_\alpha|
\;-\; \sum_{\alpha \in \Rich} |B_\alpha|.
\]
Step~1 gives $|B_\alpha| \le C\,t_\alpha \cdot |\LP|/(|A|\,|B|\,|C|)$ while the
interface theorem guarantees $|\mathcal{F}_\alpha| \gtrsim t_\alpha^2 \cdot
|\LP|/(|A|\,|B|\,|C|)$ for rich $\alpha$ (the 2D bulk of the good fiber);
hence $|B_\alpha|/|\mathcal{F}_\alpha| = O(1/t_\alpha) \le O(1/T)$. Choosing
$T \ge T_0$ with $T_0$ absolute makes the bad-set term at most $\tfrac12$ of
the total, so
\[
\mathbb{E}_L[I(L)]
\;=\; \frac{1}{|\LP|}\sum_{\alpha \in \Rich} |\mathcal{F}_\alpha^\star|
\;\ge\; \tfrac12\,c_T.
\]

Second moment. Since $I(L) \ge 0$, Jensen's inequality (Cauchy--Schwarz on
the uniform measure over $\LP$) gives
\[
\mathbb{E}_L[I(L)^2] \;\ge\; \big(\mathbb{E}_L[I(L)]\big)^2 \;\ge\; (c_T/2)^2
\;=:\; c_T'.
\]
In the notation of \texttt{step8.tex}
(Remark~\ref{rem:final-constants-audit} there), $c_T = c_5^\star/2$
from the ``in particular'' clause of
Lemma~\ref{lem:fiber-mass}, so
$c_T'(T) = (c_5^\star/4)^2 = (c_5^\star)^2/16$; the step~8 richness
constant is $c_5(T) := c_T'(T)$.

Identification of $\sum_L I(L)^2$ with $\sum_{\alpha,\beta} w_{\alpha\beta}$.
For each $L \in \LP$,
\[
I(L)^2
\;=\; \#\{(\alpha,\beta) \in \Rich^2 :
         L \in \mathcal{F}_\alpha^\star \cap \mathcal{F}_\beta^\star\}
\;=\; \sum_{\alpha,\beta \in \Rich}
         \mathbf{1}\!\left[L \in (\mathcal{F}_\alpha \cap \mathcal{F}_\beta) \setminus (B_\alpha \cup B_\beta)\right].
\]
Summing over $L$ and interchanging the order of summation yields
$\sum_L I(L)^2 = \sum_{\alpha,\beta} w_{\alpha\beta}$ under the counting-measure
normalization used in Step~6.
\end{proof}

\begin{remark}[Why this is automatic]\label{rem:second-moment}
The strengthening from first- to second-moment visibility costs nothing beyond
Step~1's bad-fiber thinness: (R) is already an $L^1$ lower bound on the
non-negative random variable $I(L)$, so Jensen promotes it to an $L^2$ lower
bound of the same order. No strengthening of G1 (Endpoint Monotonicity) or
G2 (Convex Overlap) is required. The only non-trivial input is Step~1's
$|B_\alpha| = O(t_\alpha)$ bound, which ensures that passing from $\Fxy$ to
$\Fxy^\star$ loses at most a constant factor of the first moment. This closes
the second-moment half of Step~6's rich-case lower bound at the Step~5 boundary.
\end{remark}



\begin{abstract}
Step~6 combines the local cut-rigidity lemma (Step~2), the two-interface
incompatibility lemma (Step~4), and the richness dichotomy (Step~5) into a
global dichotomy: any low-conductance cut $S\subseteq\LE(P)$ either has
large BK boundary (contradiction) or has the property that almost all rich
interfaces are pairwise coherent on their overlaps.
\end{abstract}


%-----------------------------------------------------------------------
\section{Setup and recall}
%-----------------------------------------------------------------------

Let $P$ be a width-3 poset and $\LE(P)$ its set of linear extensions.
Let $G_{\BK}(P)$ be the BK graph on $\LE(P)$: vertices are linear
extensions, edges come from adjacent-incomparable swaps.
For $S\subseteq V(G_{\BK}(P))$ write
\[
\bd S := E(S, S^c), \qquad
\Phi(S) := \frac{|\bd S|}{\vol(S)},
\]
or any equivalent normalized-conductance notion.

We use the rich-pair notation of Step~1 (Definition~\ref{def:rich}):
for $x\parallel y$ in $P$, $C(x,y)=N(x,y)$ is the common-axis chain
(a chain in $P$ by the width-$3$ hypothesis, Lemma~\ref{lem:CNchain-width3}
of Step~1; this is the common-axis set, not a graph neighborhood in
$G_{\BK}$), $t(x,y)=|C(x,y)|$, and the pair $(x,y)$ is \emph{rich} at
threshold $T\ge 1$ iff $t(x,y)\ge T$.
Let $\Rich$ denote the set of rich pairs.

\begin{definition}[Good fiber, interface]
\DEP{Step~1}
For each $\alpha=(x,y)\in\Rich$, Step~1 produces a \emph{good fiber}
$\fiber{\alpha}\subseteq\LE(P)$ with a coordinate map
$\pi_\alpha\colon\fiber{\alpha}\to D_\alpha\subseteq [0,t(x,y)]^2$ such
that the BK moves internal to $\fiber{\alpha}$ are exactly unit grid
moves, and $D_\alpha$ is order-convex.
We refer to $\alpha$ also as the \emph{rich interface} with domain
$D_\alpha$.
\end{definition}

%-----------------------------------------------------------------------
\section{The Step 6 proposition}
%-----------------------------------------------------------------------

\begin{proposition}[Step~6 dichotomy, informal]\label{prop:step6-informal}
There exist constants $\eta,\delta>0$ depending only on the width and on
the thresholds fixed by Steps~1, 2, 4, 5, such that for any width-3
poset $P$ and any $S\subseteq\LE(P)$ with conductance $\Phi(S)$
sufficiently small, the following holds: all but a $\delta$-fraction of
the weighted rich-interface overlap mass is \emph{coherent}, in the
sense made precise in Section~\ref{sec:interface-graph}.
\end{proposition}

The formal dichotomy theorem is stated in
Section~\ref{sec:dichotomy-theorem}.

%-----------------------------------------------------------------------
\section{Inputs: packaging Steps 2, 4, 5}
%-----------------------------------------------------------------------

We collect the precise form in which Steps~2, 4, 5 are consumed by the
Step~6 argument.

\subsection{Step 2: local staircase on a fiber}

\DEP{Step~2}
For each rich interface $\alpha\in\Rich$, Step~2 provides:
\begin{itemize}[nosep]
  \item an \emph{orientation sign} $\sigma_\alpha\in\{+1,-1\}$, and
  \item an \emph{exceptional set} $B_\alpha\subseteq\fiber{\alpha}$
  of small mass (quantified by Step~2's error parameter $\varepsilon_2$),
\end{itemize}
such that, after deleting $B_\alpha$, membership of $L\in\fiber{\alpha}$
in $S$ is monotone in the $\sigma_\alpha$-direction of the grid
$D_\alpha$.
Explicitly, writing $\pi_\alpha(L)=(i,j)$, there is a monotone staircase
region $M_\alpha\subseteq D_\alpha$ with
\[
\mathbf{1}_S(L)
 =\mathbf{1}_{M_\alpha}(\pi_\alpha(L))\quad
\text{for all } L\in\fiber{\alpha}\setminus B_\alpha.
\]

The uniform bound on $\sum_\alpha |B_\alpha|$ (total Step-2 error mass)
in terms of $|\bd S|$ that Step~6 consumes is proved in
Section~\ref{sec:gap-error}.

\subsection{Step 4: incompatibility yields boundary}

\DEP{Step~4}
Step~4 attaches to each ordered pair $(\alpha,\beta)$ of good rich
interfaces a \emph{witness family}
\[
E_{\alpha\beta}\subseteq E(G_{\BK}(P))
\]
of BK edges in the overlap
$(\fiber{\alpha}\cap\fiber{\beta})\setminus(B_\alpha\cup B_\beta)$
such that if the two staircase orientations
$\sigma_\alpha,\sigma_\beta$ are \emph{incoherent} on the overlap,
then a positive fraction of $E_{\alpha\beta}$ crosses $\partial S$:
\[
|\bd S \cap E_{\alpha\beta}| \;\gtrsim\;
\mu\big((\fiber{\alpha}\cap\fiber{\beta})\setminus(B_\alpha\cup B_\beta)\big).
\]
Here $\mu$ is counting measure on $\LE(P)$ (or any fixed normalization).

The precise definition of ``incoherent'' used here is
Def.~\ref{def:incoherent-pair} of Section~\ref{sec:gap-incoherence},
which imports the Step~3 formulation in terms of the local signs
$\eta_\alpha,\eta_\beta$ on the common overlap axes.

The local commuting-square model in the overlap, which supplies
$E_{\alpha\beta}$, is provided by Step~1 (Corollary
``Commuting overlap'') and is restated as Lemma~\ref{lem:gap3}
in Section~\ref{sec:gap-commuting}.

\subsection{Step 5: richness implies large overlap mass}

\DEP{Step~5}
Step~5 asserts a dichotomy: either $P$ is in the \emph{collapse} regime
(handled directly by Step~7) or else the rich-interface visibility
\[
I(L) := \#\{\alpha\in\Rich : L\in\fiber{\alpha}\setminus B_\alpha\}
\]
is large for a typical $L$, in the sense that
\[
M := \sum_{\alpha,\beta\in\Rich^\star} w_{\alpha\beta}
\quad\text{is } \Omega(|\LE(P)|),
\]
where $w_{\alpha\beta}$ and $\Rich^\star$ are defined in
Section~\ref{sec:interface-graph}.  We assume throughout Step~6 that we
are in the \emph{rich} case of Step~5; the collapse case is deferred to
Step~7.

The identification of $M$ with the second moment of visibility
$\sum_L I(L)^2$, via the bounded-multiplicity input from Step~1, is
carried out in Section~\ref{sec:gap-second-moment}.

%-----------------------------------------------------------------------
\section{Good interfaces and the interface graph}\label{sec:interface-graph}
%-----------------------------------------------------------------------

\subsection{$S$-good interfaces}

\begin{definition}[$S$-good interface]
An interface $\alpha\in\Rich$ is \emph{$S$-good} if Step~2 applies to
$\fiber{\alpha}$ with error at most $\varepsilon_2$.  Let
$\Rich^\star\subseteq\Rich$ denote the set of $S$-good rich interfaces.
\end{definition}

\begin{lemma}[Most rich interfaces are $S$-good]\label{lem:most-good}
Suppose $\Phi(S)\le \Phi_0$ for a constant $\Phi_0$ depending on
$\varepsilon_2$.  Then
\[
\sum_{\alpha\in\Rich^\star} |\fiber{\alpha}|
\;\ge\;
(1-o(1))\sum_{\alpha\in\Rich} |\fiber{\alpha}|,
\]
where the $o(1)$ is as $\Phi_0\to 0$ with $\varepsilon_2$ and the
width-3 structural constants fixed.
\end{lemma}

\begin{proof}
Step~2 produces, for each $\alpha\in\Rich$, an exceptional set
$B_\alpha\subseteq\fiber{\alpha}$ equipped with the per-fiber error
estimate
\begin{equation}\label{eq:step2-per-fiber}
|B_\alpha|
 \;\le\; C_2\,\bigl|\bd S \cap E(\widetilde{\F}_{\alpha})\bigr|,
\end{equation}
where $\widetilde{\F}_{\alpha}$ denotes the one-step BK-closure of the
good fiber $\fiber{\alpha}$ and $C_2$ depends only on the width and on
the richness threshold $T$.  (Inequality
\eqref{eq:step2-per-fiber} is the quantitative form of Step~2: each
unit of exceptional mass on the fiber $\fiber{\alpha}$ is chargeable
to a constant number of BK-edges crossing $\bd S$ in a bounded
neighborhood of the fiber.)

Summing \eqref{eq:step2-per-fiber} over $\alpha\in\Rich$ and using the
bounded-multiplicity property of the family
$\{\widetilde{\F}_{\alpha}\}_{\alpha\in\Rich}$ --- each BK edge
$e\in E(G_{\BK}(P))$ lies in $O(1)$ enlarged good fibers, a consequence
of the width-3 normal form (Step~1) --- yields
\begin{equation}\label{eq:step2-summed}
\sum_{\alpha\in\Rich} |B_\alpha|
 \;\le\; C_2'\,|\bd S|,
\end{equation}
for a constant $C_2'$ depending only on the width-3 structure.

An interface $\alpha\in\Rich$ fails to be $S$-good iff
$|B_\alpha|>\varepsilon_2\,|\fiber{\alpha}|$.  Writing
$\Rich^{\mathrm{bad}}:=\Rich\setminus\Rich^\star$, Markov's inequality
applied fiberwise gives
\[
\sum_{\alpha\in\Rich^{\mathrm{bad}}} |\fiber{\alpha}|
 \;<\; \varepsilon_2^{-1}
       \sum_{\alpha\in\Rich^{\mathrm{bad}}} |B_\alpha|
 \;\le\; \varepsilon_2^{-1}
       \sum_{\alpha\in\Rich} |B_\alpha|
 \;\stackrel{\eqref{eq:step2-summed}}{\le}\;
       \varepsilon_2^{-1} C_2'\,|\bd S|.
\]
By hypothesis $|\bd S|\le \Phi_0\,\vol(S)$.  In the rich case of
Step~5 (the regime in which Step~6 is nontrivial), the rich fibers
collectively cover a positive fraction of $\LE(P)$:
\begin{equation}\label{eq:rich-covers}
\sum_{\alpha\in\Rich} |\fiber{\alpha}|
 \;\ge\; c_\Rich\,\vol(S),
\end{equation}
for a constant $c_\Rich>0$ supplied by Step~5 (via the visibility
identity $\sum_\alpha|\fiber{\alpha}|=\sum_L I(L)\ge
|\{L:I(L)\ge 1\}|$, combined with the lower bound on $I(L)$ in the
rich case).  Combining,
\[
\frac{\sum_{\alpha\in\Rich^{\mathrm{bad}}}|\fiber{\alpha}|}
     {\sum_{\alpha\in\Rich}|\fiber{\alpha}|}
 \;\le\;
 \frac{\varepsilon_2^{-1} C_2'\,\Phi_0\,\vol(S)}
      {c_\Rich\,\vol(S)}
 \;=\;
 \frac{C_2'}{c_\Rich\,\varepsilon_2}\;\Phi_0,
\]
which tends to $0$ as $\Phi_0\to 0$.  Equivalently,
$\sum_{\alpha\in\Rich^\star}|\fiber{\alpha}|\ge (1-o(1))
 \sum_{\alpha\in\Rich}|\fiber{\alpha}|$, as claimed.
\end{proof}

\begin{remark}
The quantitative form of the lemma is $(1-\kappa\Phi_0)$ with
$\kappa = C_2'/(c_\Rich\,\varepsilon_2)$.  In particular, choosing
$\Phi_0$ below $\varepsilon_2 \cdot c_\Rich/(2 C_2')$ already forces
more than half of the rich-interface mass to be $S$-good; this is the
quantitative input actually consumed by the sums in
Section~\ref{sec:dichotomy-theorem}.
\end{remark}

\subsection{The weighted interface graph $H$}

Define a weighted graph $H$ on vertex set $\Rich^\star$ with edge weights
\[
w_{\alpha\beta}
 := \mu\big((\fiber{\alpha}\cap\fiber{\beta})
        \setminus (B_\alpha\cup B_\beta)\big).
\]
So $w_{\alpha\beta}$ measures the mass of the overlap on which both
interfaces are good and both monotone descriptions are valid.

\begin{definition}[active / bad pair]
A pair $(\alpha,\beta)$ is \emph{active} if $w_{\alpha\beta}$ exceeds a
fixed threshold $w_0$, and \emph{bad} if in addition the orientations
$\sigma_\alpha,\sigma_\beta$ are \emph{incoherent} on the overlap.
\end{definition}

Incoherence of $(\alpha,\beta)$ on the overlap is
Def.~\ref{def:incoherent-pair} (Section~\ref{sec:gap-incoherence}),
imported from Step~3.
With this notation, Step~4 yields the key inequality
\begin{equation}\label{eq:bad-bounded-by-boundary}
\sum_{\text{bad active }(\alpha,\beta)} w_{\alpha\beta}
\;\lesssim\; |\bd S|.
\end{equation}

\begin{lemma}[Step~4 summed: bad overlap is bounded by
boundary]\label{lem:sum-step4}
Let $\delta_0>0$ be the incoherence threshold used to define ``bad''
pairs (G2) and let $\varepsilon$ be the Step~2 per-fiber error, chosen
small enough that $C\varepsilon\le\delta_0/2$, where $c,C$ are the
constants of Step~4 (Theorem on two-interface incompatibility).
Assume the active threshold $w_0$ dominates the (S1.4) non-commuting
loss, so that
$|\Omega^\circ_{\alpha,\beta}|\ge\tfrac12 w_{\alpha,\beta}$ for every
active pair (see Remark~\ref{rem:Omega-circ-vs-w}).
Then with $c_1:=c\delta_0/(4M)$, where $M$ is the witness-edge
multiplicity constant of Lemma~\ref{lem:overlap-counting},
\begin{equation}\label{eq:bad-bounded-by-boundary-quant}
\sum_{\text{bad active }(\alpha,\beta)}
w_{\alpha,\beta}\;\le\;\frac{1}{c_1}\,|\bd S|.
\end{equation}
\end{lemma}

\begin{remark}\label{rem:Omega-circ-vs-w}
By the commuting-overlap corollary (G3, equivalently (S1.4)), one has
$|\Omega^\circ_{\alpha,\beta}|\ge|\mathcal F_\alpha\cap\mathcal F_\beta|
-o(|\mathcal F_\alpha\cap\mathcal F_\beta|)$, and after removing
$B_\alpha\cup B_\beta$ the remaining loss is at most
$|B_\alpha|+|B_\beta|$.  By G1 (Step-2 error control) and the active
threshold $w_{\alpha,\beta}\ge w_0$, both losses are at most
$\tfrac14 w_{\alpha,\beta}$ for $w_0$ sufficiently large relative to
the Step-1/Step-2 error scales, yielding the stated lower bound.
\end{remark}

\begin{proof}[Proof of Lemma~\ref{lem:sum-step4}]
Fix a bad active pair $(\alpha,\beta)$.  By construction,
$\sigma_\alpha,\sigma_\beta$ are incoherent on
$(\mathcal F_\alpha\cap\mathcal F_\beta)\setminus(B_\alpha\cup B_\beta)$,
which in the quantitative form of Definition~G2 (step~3) means the
incoherence fraction $\delta_{\alpha,\beta}$ on the commuting overlap
$\Omega^\circ_{\alpha,\beta}$ supplied by G3 is at least $\delta_0$.
Applying Step~4 (Theorem on two-interface incompatibility) to
$(\alpha,\beta)$ — i.e.\ to the rectangle-decomposed commuting overlap
with the Step-2 staircase approximations transferred via Proposition~G2
of \texttt{step4.tex} — yields
\[
|\bd S\cap E_{\alpha,\beta}|
\;\ge\; c\bigl(\delta_{\alpha,\beta}-C\varepsilon\bigr)\,
|\Omega^\circ_{\alpha,\beta}|
\;\ge\; \frac{c\delta_0}{2}\,|\Omega^\circ_{\alpha,\beta}|,
\]
using $\delta_{\alpha,\beta}\ge\delta_0$ and
$C\varepsilon\le\delta_0/2$.  By Remark~\ref{rem:Omega-circ-vs-w},
$|\Omega^\circ_{\alpha,\beta}|\ge\tfrac12 w_{\alpha,\beta}$, so
\begin{equation}\label{eq:per-pair-witness-mass}
|\bd S\cap E_{\alpha,\beta}|
\;\ge\; \frac{c\delta_0}{4}\, w_{\alpha,\beta}.
\end{equation}

Summing \eqref{eq:per-pair-witness-mass} over bad active pairs and
invoking the overlap counting lemma
(Lemma~\ref{lem:overlap-counting}), which provides
$\sum_{(\alpha,\beta)}\mathbf{1}_{e\in E_{\alpha,\beta}}\le M$ for every
$e\in E(G_{\BK}(P))$,
\[
\frac{c\delta_0}{4}
\sum_{\text{bad active}} w_{\alpha,\beta}
\;\le\;\sum_{\text{bad active}} |\bd S\cap E_{\alpha,\beta}|
\;\le\;\sum_{e\in\bd S}\sum_{(\alpha,\beta)}
\mathbf{1}_{e\in E_{\alpha,\beta}}
\;\le\; M\,|\bd S|.
\]
Dividing by $c\delta_0/4$ gives
$\sum_{\text{bad active}} w_{\alpha,\beta}\le (4M/(c\delta_0))\,|\bd S|
= c_1^{-1}|\bd S|$, establishing
\eqref{eq:bad-bounded-by-boundary-quant} and
\eqref{eq:bad-bounded-by-boundary}.
\end{proof}

%-----------------------------------------------------------------------
\section{The single clean technical lemma}\label{sec:overlap-counting}
%-----------------------------------------------------------------------

The entire Step~6 argument pivots on a single double-counting lemma,
which we isolate here.

\begin{lemma}[Overlap counting lemma]\label{lem:overlap-counting}
The family of witness edges
$\{E_{\alpha\beta}\}_{(\alpha,\beta)\in\Rich^\star\times\Rich^\star}$
produced by Step~4 can be chosen so that each BK edge $e\in E(G_{\BK})$
is contained in at most $M$ witness families $E_{\alpha\beta}$, where
$M$ depends only on the width of $P$ (width~$3$).
\end{lemma}

\begin{proof}
Recall from Step~4 that the witness family is the disjoint union over
rectangle blocks $B$ of $\Omega^\circ_{\alpha\beta}$
($=(\fiber{\alpha}\cap\fiber{\beta})\setminus(B_\alpha\cup B_\beta\cup
\mathrm{Int}_{\alpha\beta})$) of the cut edges of conflict
$2\times 2$ squares $Q\subseteq R_B$.  In each block, the Step~1
commuting-overlap clause (S1.4) identifies the horizontal edges of
$R_B$ with $\alpha$-moves and the vertical edges with $\beta$-moves.
Every BK edge is an adjacent transposition of a unique unordered pair
of poset elements, its \emph{swap pair}.

Fix $e=\{L,L'\}\in E(G_{\BK}(P))$ with swap pair
$\{p,q\}\subseteq P$ incomparable, acting on $L$ at adjacent positions
$(i,i+1)$.  Suppose $e\in E_{\alpha\beta}$, so $e$ is an edge of some
conflict square $Q\subseteq R_B$.

\smallskip\noindent\textbf{Step 1 (one interface is pinned).}
Since $e$ is one of the four edges of $Q$ and the two edge directions
of $R_B$ are the $\alpha$- and $\beta$-swap moves, the swap pair of
$e$ equals $\alpha$ or $\beta$.  Hence $\{p,q\}\in\{\alpha,\beta\}$.
By symmetry between the two interfaces it suffices to bound the
number of partner interfaces $\gamma$ such that
$e\in E_{\{p,q\},\gamma}$ by an absolute constant $M_0$; then
$M\le 2M_0+1$ (the extra $+1$ absorbs the degenerate case
$\gamma=\{p,q\}$).

\smallskip\noindent\textbf{Step 2 (position-disjointness).}
Fix a candidate partner $\gamma\ne\{p,q\}$.  Both endpoints
$L,L'$ of $e$ lie in $\Omega^\circ_{\{p,q\},\gamma}\subseteq
\fiber{\{p,q\}}\cap\fiber{\gamma}$, so at $L$ the two generators
(the $\{p,q\}$-swap and the $\gamma$-swap) commute by (S1.4).  Two
adjacent transpositions on a linear order commute only when they
either coincide or act on disjoint adjacent position pairs.  Since
$\gamma\ne\{p,q\}$, the $\gamma$-swap at $L$ acts on some
$(j,j+1)$ with $\{j,j+1\}\cap\{i,i+1\}=\emptyset$, and in particular
the underlying element pair $\gamma$ is disjoint from $\{p,q\}$.

\smallskip\noindent\textbf{Step 3 (width-3 local count).}
We claim that the number of rich interfaces $\gamma$ with
\begin{enumerate}[label=\textup{(\roman*)},nosep]
  \item $L\in\fiber{\gamma}\setminus B_\gamma$,
  \item the $\gamma$-swap at $L$ acts on some adjacent pair
  $(j,j+1)$ disjoint from $\{i,i+1\}$,
\end{enumerate}
is bounded by an absolute constant $M_0$ depending only on
$w(P)\le 3$.

By Dilworth's theorem the width-3 poset $P$ partitions into at most
three chains $C_1,C_2,C_3$; any incomparable pair is drawn from two
distinct chains, so there are $\binom{3}{2}=3$ \emph{chain-pair
types}.  Each candidate $\gamma$ has a chain-pair type, and the swap
pair of $\gamma$ at $L$ occupies two adjacent positions $(j,j+1)$,
each in a distinct chain.

\emph{At most one $\gamma$ per (chain-pair type, adjacent position
pair).}  Suppose $\gamma,\gamma'$ are distinct rich interfaces of
the same chain-pair type whose swaps at $L$ act on the \emph{same}
adjacent pair $(j,j+1)$.  Then the endpoints of $\gamma$ and
$\gamma'$ occupy the same two positions in $L$; since each adjacent
position pair in $L$ contains exactly one element of each of its two
chains, this forces $\gamma=\gamma'$, a contradiction.

\emph{At most $O(1)$ occupied adjacent pairs $(j,j+1)\ne(i,i+1)$ per
chain-pair type.}  Fix a chain-pair type $\{C_a,C_b\}$.  Suppose
$\gamma_1,\ldots,\gamma_K$ are distinct rich interfaces of this
type, with $\gamma_k$'s swap acting at $(j_k,j_k+1)$ in $L$, all
$(j_k,j_k+1)$ pairwise disjoint and disjoint from $(i,i+1)$.  Each
$\gamma_k=\{u_k,v_k\}$ satisfies $L\in\fiber{\gamma_k}\setminus
B_{\gamma_k}$.  The commuting-overlap clause (S1.4) applied to
$\gamma_k,\gamma_\ell$ at $L$ forces the $\gamma_k$- and
$\gamma_\ell$-swaps to commute, hence their position pairs are
disjoint (already assumed) \emph{and} the underlying element pairs
$\{u_k,v_k\}$ and $\{u_\ell,v_\ell\}$ are disjoint (else two
transpositions sharing one element do not commute).  Thus the
$\gamma_k$ together consume $2K$ distinct elements drawn from the
two chains $C_a\cup C_b$; each must be rich (i.e.\ have
$t(\gamma_k)\ge T$), which means each pair has common neighborhood
a long chain inside the remaining chain $C_c$.

By the ``bounded interaction of distinct rich pairs'' lemma of
Step~1 (\DEP{Step~1}), any two rich pairs of the same
chain-pair type have overlap fibers $\fiber{\gamma_k}\cap
\fiber{\gamma_\ell}$ contained in $O(1)$ strips of width
$O(t_{\gamma_k}+t_{\gamma_\ell})$ in either coordinate system.
In particular, once $L$ lies in $\fiber{\gamma_1}\setminus
B_{\gamma_1}$ and in $\fiber{\gamma_\ell}\setminus B_{\gamma_\ell}$
simultaneously, the good-fiber coordinate $\pi_{\gamma_1}(L)$ is
pinned to one of $O(1)$ strips determined by $\gamma_\ell$.
A rich fiber has $t(\gamma)=\Theta(n)$ while each strip has
cross-sectional width $O(1)$, so at most $O(1)$ distinct rich
$\gamma_\ell$ can have their good fiber simultaneously contain $L$
while acting on an adjacent position pair disjoint from
$(i,i+1)$.  This gives $K=O(1)$.

Summing $O(1)$ over the $3$ chain-pair types and the bounded
number of admissible position pairs within each type yields
$M_0=O(1)$, depending only on the width and on the richness
threshold fixed in Steps~1--5 (not on $|P|$).

\smallskip\noindent\textbf{Step 4 (assemble).}
Combining Steps~1--3,
\[
\sum_{(\alpha,\beta)}\mathbf 1_{e\in E_{\alpha\beta}}
  \;\le\; 2M_0+1 \;=:\; M,
\]
with $M$ depending only on $w(P)\le 3$.
\end{proof}

\begin{remark}
The same argument is written out in Step~4
(Lemma ``Witness-edge multiplicity'', \texttt{step4.tex}, \S
``Witness-edge multiplicity for Step~6''); the proof is reproduced
here because the overlap counting lemma is the load-bearing bookkeeping
input to Step~6, and consolidating it in \texttt{step6.tex} removes the
cross-file dependency for downstream readers.  A sharper numerical
value of $M$ is not required by any downstream step.
\end{remark}

%-----------------------------------------------------------------------
\section{The dichotomy theorem}\label{sec:dichotomy-theorem}
%-----------------------------------------------------------------------

\begin{theorem}[Step~6 theorem]\label{thm:step6}
Fix thresholds from Steps~1, 2, 4, 5 and the overlap counting constant
from Lemma~\ref{lem:overlap-counting}.
There exist constants $\eta,\delta>0$ such that for any width-3 poset
$P$ and any cut $S\subseteq\LE(P)$, at least one of the following holds:

\begin{enumerate}[label=(\roman*)]
  \item \textbf{Expansion case.}
  \[
    |\bd S| \;\ge\; \eta\cdot M,
    \qquad
    M \;=\; \sum_{\alpha,\beta\in\Rich^\star} w_{\alpha\beta}.
  \]

  \item \textbf{Coherence case.}
  All but a $\delta$-fraction of the weighted overlap mass of pairs
  $(\alpha,\beta)\in\Rich^\star\times\Rich^\star$ is coherent:
  \[
    \sum_{\text{bad active }(\alpha,\beta)} w_{\alpha\beta}
    \;\le\;
    \delta\cdot M.
  \]
\end{enumerate}

If additionally Step~5 gives $M\gtrsim\vol(S)$ (the rich case), then
case~(i) contradicts low conductance, and only case~(ii) survives.
\end{theorem}

\subsection{Proof of Theorem~\ref{thm:step6}}\label{sec:proof-dichotomy}

\begin{proof}
\emph{Step A (pass to $\Rich^\star$).}
By Step~2 each $\alpha\in\Rich$ carries an orientation $\sigma_\alpha$
and an exceptional set $B_\alpha$.
Lemma~\ref{lem:most-good} shows that, for $\Phi(S)\le\Phi_0$,
the set $\Rich^\star$ of $S$-good interfaces satisfies
$\sum_{\alpha\in\Rich^\star}|\fiber{\alpha}|\ge(1-o(1))
\sum_{\alpha\in\Rich}|\fiber{\alpha}|$.
Section~\ref{sec:gap-error} furthermore bounds the total
exceptional mass $\sum_{\alpha\in\Rich^\star}|B_\alpha|$ by a fixed
fraction of $\sum_{\alpha\in\Rich^\star}|\fiber{\alpha}|$, so the
overlap weights $w_{\alpha\beta}$ dominate the bad-set losses when
$\Phi(S)$ is small enough.

\emph{Step B (partition into bad vs.\ coherent active pairs).}
Using Def.~\ref{def:incoherent-pair}
(Section~\ref{sec:gap-incoherence}), every active pair
$(\alpha,\beta)\in\Rich^\star\times\Rich^\star$ with
$w_{\alpha\beta}\ge w_0$ is classified as either \emph{bad} (incoherent
orientations on the overlap) or \emph{coherent}.

\emph{Step C (summing Step~4).}
Section~\ref{sec:gap-commuting} produces, on the overlap
$(\fiber{\alpha}\cap\fiber{\beta})\setminus(B_\alpha\cup B_\beta)$,
the commuting $2\times 2$-square model from which Step~4 draws its
witness family $E_{\alpha\beta}$.
Lemma~\ref{lem:overlap-counting} bounds the multiplicity with
which any single BK edge appears in $\bigcup_{\alpha,\beta}E_{\alpha\beta}$
by a width-dependent constant $K$.
Lemma~\ref{lem:sum-step4} combines Step~4's per-pair boundary
lower bound with the multiplicity bound:
\begin{equation}\label{eq:summed-step4}
c_1\sum_{\text{bad active }(\alpha,\beta)} w_{\alpha\beta}
\;\le\; |\bd S|,
\end{equation}
for a constant $c_1>0$ depending only on the width, Step~4's constant,
and $K$.

\emph{Step D (dichotomy).}
Fix $\delta\in(0,1)$ to be chosen below and set $\eta:=c_1\delta$.
If
\[
\sum_{\text{bad active }(\alpha,\beta)} w_{\alpha\beta}
\;\le\;\delta\cdot M,
\]
then case~(ii) holds.
Otherwise, by \eqref{eq:summed-step4},
\[
|\bd S|
 \;\ge\; c_1\sum_{\text{bad active }(\alpha,\beta)} w_{\alpha\beta}
 \;>\; c_1\delta\cdot M \;=\; \eta\cdot M,
\]
which is case~(i).

\emph{Step E (closing the contradiction in the rich case).}
If Step~5 holds in its rich form, Section~\ref{sec:gap-second-moment}
(Lemma~\ref{lem:gap-G5}) gives $M \gtrsim\vol(S)$.
Then case~(i) yields $|\bd S|\gtrsim\vol(S)$, i.e.
$\Phi(S)\gtrsim 1$, contradicting the low-conductance hypothesis.
Hence under the rich case of Step~5 only case~(ii) can occur.
\end{proof}

The reduction from Theorem~\ref{thm:step6} to the \emph{pointwise} form
(Conclusion~B below) is carried out in Section~\ref{sec:pointwise}.

%-----------------------------------------------------------------------
\section{Pointwise coherence}\label{sec:pointwise}
%-----------------------------------------------------------------------

For $L\in\LE(P)$ set
\[
\Rich_L
 := \{\alpha\in\Rich^\star :
      L\in\fiber{\alpha}\setminus B_\alpha\}.
\]

\begin{corollary}[Pointwise coherence, Conclusion~B]\label{cor:pointwise}
Under the hypotheses of Theorem~\ref{thm:step6} and in the rich case of
Step~5, for most $L\in\LE(P)$ there exists a sign $s(L)\in\{\pm1\}$ such
that almost every $\alpha\in\Rich_L$ has $\sigma_\alpha = s(L)$.
Quantitatively, let $\Pr_L$ denote the $I(L)^2$-weighted probability
measure on $\LE(P)$ with $I(L):=|\Rich_L|$, i.e.
$\Pr_L[A]=\sum_{L\in A}I(L)^2/\sum_{L}I(L)^2$
(with the convention $0/0=0$ if $\sum_LI(L)^2=0$, which is excluded by
the rich case of Step~5).  Then, with $\delta$ as in
Theorem~\ref{thm:step6}, there is a measurable map
$s\colon\LE(P)\to\{\pm1\}$ such that
\[
\Pr_L\!\left[
  \#\{\alpha\in\Rich_L:\sigma_\alpha\ne s(L)\}
  \ge \sqrt{\delta}\,|\Rich_L|
\right]
\;\le\; \sqrt{\delta}.
\]
Equivalently, setting $m(L):=\min(n_+(L),n_-(L))$ (the minority count,
with $n_\pm(L):=\#\{\alpha\in\Rich_L:\sigma_\alpha=\pm1\}$), the event
$\{m(L)\ge\sqrt{\delta}\,I(L)\}$ has $\Pr_L$-probability at most
$\sqrt{\delta}$: under the majority-vote $s(L)$ defined in the proof
below, $\#\{\alpha\in\Rich_L:\sigma_\alpha\ne s(L)\}$ equals the
minority count $m(L)$, so this is literally the preceding display.
\end{corollary}

\begin{proof}
We are in the coherence case~(ii) of Theorem~\ref{thm:step6}, i.e.
\begin{equation}\label{eq:bad-mass}
  \sum_{\text{bad active }(\alpha,\beta)} w_{\alpha\beta}
  \;\le\;\delta\,M.
\end{equation}

\smallskip\noindent\emph{Step 1: majority vote.}
For each $L\in\LE(P)$ set
\[
  n_+(L):=\#\{\alpha\in\Rich_L:\sigma_\alpha=+1\},\qquad
  n_-(L):=\#\{\alpha\in\Rich_L:\sigma_\alpha=-1\},
\]
so $n_+(L)+n_-(L)=I(L)$.  Define
\[
  s(L):=
  \begin{cases}
    +1,&\text{if }n_+(L)\ge n_-(L),\\
    -1,&\text{otherwise,}
  \end{cases}
\]
and
\[
  m(L):=\min(n_+(L),n_-(L))
      =\#\{\alpha\in\Rich_L:\sigma_\alpha\ne s(L)\}.
\]
Since $\LE(P)$ is finite, $s$ and $m$ are automatically measurable.
The minority count $m(L)$ is the smallest possible value of
$\#\{\alpha\in\Rich_L:\sigma_\alpha\ne s\}$ over $s\in\{\pm1\}$, so the
two events in the statement coincide when $s(L)$ is chosen by majority
vote.

\smallskip\noindent\emph{Step 2: double counting of disagreeing pairs.}
Recall $w_{\alpha\beta}=\#\{L:\alpha,\beta\in\Rich_L\}$.
Counting triples $(L,\alpha,\beta)$ with $L\in\Rich_\alpha\cap\Rich_\beta$
and $\sigma_\alpha\ne\sigma_\beta$ in two ways,
\[
  \sum_{L\in\LE(P)} 2\,n_+(L)\,n_-(L)
   \;=\;
  \sum_{\substack{(\alpha,\beta)\in\Rich^\star\times\Rich^\star\\
        \sigma_\alpha\ne\sigma_\beta}} w_{\alpha\beta}.
\]
Split the RHS into active and non-active pairs.  For the non-active
incoherent pairs $w_{\alpha\beta}<w_0$; choosing the active-pair
threshold $w_0$ so that
$w_0\cdot|\Rich^\star|^2\le\delta M$ (absorbable into the constants of
Theorem~\ref{thm:step6}), we obtain via \eqref{eq:bad-mass}
\begin{equation}\label{eq:disagree-sum}
  \sum_{L\in\LE(P)} n_+(L)\,n_-(L)
  \;\le\;\delta\,M.
\end{equation}

\smallskip\noindent\emph{Step 3: from disagreement to minority count.}
For any $L$, writing $m=m(L)$ and $I=I(L)$, we have
$n_+(L)n_-(L)=m(I-m)$ with $m\le I/2$, hence $I-m\ge I/2$ and
\[
  n_+(L)\,n_-(L)\;\ge\;\tfrac12\,m(L)\,I(L).
\]
Summing and combining with \eqref{eq:disagree-sum},
\begin{equation}\label{eq:mI-sum}
  \sum_{L\in\LE(P)} m(L)\,I(L)
  \;\le\;2\delta\,M.
\end{equation}

\smallskip\noindent\emph{Step 4: identification of $M$.}
By Lemma~\ref{lem:gap-G5}(i) (Section~\ref{sec:gap-second-moment}),
the exact identity
\[
  M \;=\;\sum_{L\in\LE(P)} I(L)^2
\]
holds with no error term, because the bad-set exclusions
$B_\alpha\cup B_\beta$ are baked into the definition of $w_{\alpha\beta}$
and $I(L)$ is taken over $\Rich^\star$. Substituting into
\eqref{eq:mI-sum} and renaming $2\delta$ to $\delta$ yields
\begin{equation}\label{eq:mI-vs-I2}
  \sum_{L\in\LE(P)} m(L)\,I(L)
  \;\le\;\delta\sum_{L\in\LE(P)} I(L)^2.
\end{equation}

\smallskip\noindent\emph{Step 5: Markov.}
Let
\[
  A\;:=\;\{L\in\LE(P):m(L)\ge\sqrt{\delta}\,I(L)\}.
\]
On $A$ we have $m(L)\,I(L)\ge\sqrt{\delta}\,I(L)^2$, so
\eqref{eq:mI-vs-I2} gives
\[
  \sqrt{\delta}\sum_{L\in A} I(L)^2
   \;\le\;\sum_{L\in\LE(P)} m(L)\,I(L)
   \;\le\;\delta\sum_{L\in\LE(P)} I(L)^2,
\]
hence
\[
  \Pr_L[A]
   \;=\;\frac{\sum_{L\in A} I(L)^2}{\sum_L I(L)^2}
   \;\le\;\sqrt{\delta},
\]
which is the asserted bound.  Since on the complement of $A$ we have
$\#\{\alpha\in\Rich_L:\sigma_\alpha\ne s(L)\}=m(L)<\sqrt{\delta}\,I(L)$,
the statement is proved.
\end{proof}

\begin{remark}
Measurability of $s\colon\LE(P)\to\{\pm1\}$ is automatic since $\LE(P)$
is finite; in a measure-theoretic setting one would simply note that
$n_+,n_-$ are measurable integer-valued functions of $L$ and $s(L)$ is
a Borel-measurable function of their sign.  The specific tie-breaking
rule is irrelevant to the bound because ties contribute $n_+=n_-$ and
hence $m(L)=I(L)/2$, which places $L$ automatically inside $A$ unless
$\sqrt{\delta}\le 1/2$; this regime is the only one of interest.
\end{remark}

\begin{remark}
The $I(L)^2$-weighted measure is the natural one because $M$ itself is
exactly the second moment $\sum_L I(L)^2$
(Lemma~\ref{lem:gap-G5}(i)).  The same
argument under the uniform measure on $\LE(P)$ yields a bound weaker by
a factor of $\max_L I(L)^2 / \mathbb{E}[I(L)^2]$, which is $O(1)$ under
the rich case of Step~5; in either normalization the conclusion is
quantitatively the same up to constants.
\end{remark}


Corollary~\ref{cor:pointwise} is exactly the input consumed by Step~7.

%-----------------------------------------------------------------------
\section{Full proof in compressed form}
%-----------------------------------------------------------------------

We record the logical skeleton; each line is turned into a formal claim
with proof in the indicated section.

\begin{enumerate}[nosep]
  \item By Step~2, each rich interface $\alpha$ has an orientation
  $\sigma_\alpha$ and a small exceptional set $B_\alpha$.
  \item By Lemma~\ref{lem:most-good}, most rich interfaces are $S$-good
  (i.e.\ lie in $\Rich^\star$).
  \item Define the weighted interface graph $H$ with weights
  $w_{\alpha\beta}$ (Section~\ref{sec:interface-graph}).
  \item By Step~4, each incoherent active pair forces boundary
  proportional to its overlap weight.
  \item By Lemma~\ref{lem:overlap-counting}, the witness edges
  $E_{\alpha\beta}$ overcount each BK edge at most $O(1)$ times.
  \item Summing (4) over all bad active pairs using (5) yields
  \eqref{eq:bad-bounded-by-boundary}.
  \item Theorem~\ref{thm:step6} follows: either $|\bd S|\gtrsim M$, or
  incoherent mass is $\le\delta M$.
  \item By Step~5 (rich case), $M\gtrsim\vol(S)$, so the first
  alternative contradicts low conductance; we land in the coherence
  case.
  \item Corollary~\ref{cor:pointwise} extracts pointwise coherence.
\end{enumerate}

%-----------------------------------------------------------------------
\section{Technical lemmas}
%-----------------------------------------------------------------------

The proofs of Theorem~\ref{thm:step6} and Corollary~\ref{cor:pointwise}
rest on four technical inputs: a total-error bound for the Step-2
exceptional sets, the formal definition of incoherence imported from
Step~3, the commuting-square overlap model supplied by Step~1, and the
rich-case lower bound $M\gtrsim\vol(S)$ obtained from Step~5.

\subsection{Step-2 error control}\label{sec:gap-error}

Under low conductance $\Phi(S)\le\Phi_0$ we show that the total
exceptional mass $\sum_{\alpha\in\Rich^\star} |B_\alpha|$ is at most a
fixed fraction of $\sum_{\alpha\in\Rich^\star}|\fiber{\alpha}|$, so the
coherence dichotomy survives the losses $B_\alpha\cup B_\beta$ when we
pass to the overlap $w_{\alpha\beta}$.

We cite the following from \texttt{step2.tex}:
\begin{itemize}[nosep]
  \item[(F1)] \emph{Fiber boundary averaging}
    (\texttt{step2.tex}, Lemma~3.4): for any $S\subseteq\LE(P)$,
    \(\sum_{\alpha\in\Rich}b_\alpha\le K\kappa\,|S|,\)
    where $b_\alpha:=|\bd_{\mathrm{BK}}(S\cap\fiber{\alpha})|$,
    $\kappa:=|\bd_{\mathrm{BK}}S|/|S|$, and $K$ is the Step~1
    edge-multiplicity constant (S1.6).
  \item[(F2)] \emph{Per-fiber staircase approximation}
    (\texttt{step2.tex}, Proposition~3.6): for each rich good fiber
    $\alpha$ there is a staircase $M_\alpha\in\mathrm{Stair}(D_\alpha)$
    and an error set $E_\alpha\subseteq\fiber{\alpha}$ with
    \(|E_\alpha|\le\delta(b_\alpha/t_\alpha)\,|\fiber{\alpha}|,\)
    where $\delta(\cdot)$ is the weak-grid stability rate
    (crudely $\delta(\varepsilon)=O(\varepsilon^{1/3})$).
  \item[(F3)] \emph{Good-fiber size and bad-fiber set}
    (\texttt{step2.tex}, Def.~2.1, (S1.1) and (S1.5)):
    $|\fiber{\alpha}|\ge c_0 t_\alpha^2$ and
    $|\fiber{\alpha}^{\mathrm{bad}}|\le C_1 t_\alpha$, with $t_\alpha\ge T$
    the Step-1 richness threshold.
\end{itemize}

\begin{proposition}[Step-2 error control]\label{prop:gap-g1}
Fix Step-1 constants $c_0,C_1,K$ and the Step-2 stability rate
$\delta(\cdot)$.  Let $t_{\max}:=\max_{\alpha\in\Rich}t_\alpha$ and
assume the Step-5 rich case, so that $|S|\le C_5
\sum_{\alpha\in\Rich}|\fiber{\alpha}|$ for some constant $C_5$.  Set
\[
B_\alpha\;:=\;\fiber{\alpha}^{\mathrm{bad}}\cup E_\alpha,
\]
the union of the Step-1 bad-fiber set and the Step-2 staircase error.
Then for every $\rho\in(0,1)$ there is a threshold
\[
\Phi_0
 \;=\;\Phi_0(\rho,c_0,C_1,K,C_5,T,t_{\max})\;>\;0
\]
such that whenever $\Phi(S)\le\Phi_0$,
\[
\sum_{\alpha\in\Rich^\star}|B_\alpha|
 \;\le\;\rho\,\sum_{\alpha\in\Rich^\star}|\fiber{\alpha}|.
\]
\end{proposition}

\begin{proof}
Abbreviate $W^\star:=\sum_{\alpha\in\Rich^\star}|\fiber{\alpha}|$ and
$W:=\sum_{\alpha\in\Rich}|\fiber{\alpha}|\ge W^\star$.  We bound the
contributions $\fiber{\alpha}^{\mathrm{bad}}$ and $E_\alpha$ separately,
each by $\rho/2\cdot W^\star$.

\medskip\noindent\emph{(1) Bad-fiber contribution.}
By (F3), $|\fiber{\alpha}^{\mathrm{bad}}|\le C_1
t_\alpha$ and $|\fiber{\alpha}|\ge c_0 t_\alpha^2\ge c_0 T\,t_\alpha$ for
every $\alpha\in\Rich$.  Hence
\begin{equation}\label{eq:g1-badfib}
\sum_{\alpha\in\Rich^\star}|\fiber{\alpha}^{\mathrm{bad}}|
 \;\le\;\frac{C_1}{c_0 T}\,W^\star.
\end{equation}
Choose the Step-1 richness threshold $T\ge 2C_1/(c_0\rho)$ so
\eqref{eq:g1-badfib} is at most $(\rho/2)W^\star$.  This fixes $T$
independently of $S$.

\medskip\noindent\emph{(2) Staircase-error contribution.}
Write $b_\alpha:=|\bd_{\mathrm{BK}}(S\cap\fiber{\alpha})|$.  By (F1),
\begin{equation}\label{eq:g1-avg}
\sum_{\alpha\in\Rich}b_\alpha\;\le\;K\,\Phi(S)\,|S|
 \;\le\;K\,C_5\,\Phi_0\,W.
\end{equation}
Pick $\varepsilon_0\in(0,1)$ with $\delta(\varepsilon_0)\le\rho/4$ (available
since $\delta(\cdot)\to0$; e.g.\ $\varepsilon_0=(\rho/(4C))^3$ under the
crude cube-root rate).  Split
\[
\Rich^\star
 \;=\;
 \Rich^\star_{\mathrm{lo}}\;\sqcup\;\Rich^\star_{\mathrm{hi}},
\qquad
\Rich^\star_{\mathrm{lo}}
 :=\{\alpha\in\Rich^\star : b_\alpha\le\varepsilon_0\,t_\alpha\}.
\]

\emph{(2a) Low-boundary fibers.}
For $\alpha\in\Rich^\star_{\mathrm{lo}}$, (F2) gives
$|E_\alpha|\le\delta(\varepsilon_0)|\fiber{\alpha}|\le(\rho/4)|\fiber{\alpha}|$,
so
\begin{equation}\label{eq:g1-lo}
\sum_{\alpha\in\Rich^\star_{\mathrm{lo}}}|E_\alpha|
 \;\le\;\frac{\rho}{4}\,W^\star.
\end{equation}

\emph{(2b) High-boundary fibers.}
For $\alpha\in\Rich^\star_{\mathrm{hi}}$, $b_\alpha>\varepsilon_0
t_\alpha$, i.e.\ $t_\alpha<b_\alpha/\varepsilon_0$.  Using
$|\fiber{\alpha}|\le t_\alpha^2\le t_{\max}\,t_\alpha\le
(t_{\max}/\varepsilon_0)\,b_\alpha$, we bound
$|E_\alpha|\le|\fiber{\alpha}|$ trivially and sum using
\eqref{eq:g1-avg}:
\begin{equation}\label{eq:g1-hi}
\sum_{\alpha\in\Rich^\star_{\mathrm{hi}}}|E_\alpha|
 \;\le\;\sum_{\alpha\in\Rich^\star_{\mathrm{hi}}}|\fiber{\alpha}|
 \;\le\;\frac{t_{\max}}{\varepsilon_0}
        \sum_{\alpha\in\Rich}b_\alpha
 \;\le\;\frac{K\,C_5\,t_{\max}\,\Phi_0}{\varepsilon_0}\,W.
\end{equation}
Choose
\(
\Phi_0
 \;\le\;
 \dfrac{\rho\,\varepsilon_0}{4\,K\,C_5\,t_{\max}},
\)
so the right-hand side of \eqref{eq:g1-hi} is at most
$(\rho/4)W\le(\rho/4)W^\star\cdot(W/W^\star)$.  Since
$W/W^\star=1+o(1)$ by Lemma~\ref{lem:most-good}, absorbing the
$(1+o(1))$ factor into a slightly smaller $\Phi_0$ gives
$\sum_{\mathrm{hi}}|E_\alpha|\le(\rho/4)W^\star$ outright.

\medskip\noindent\emph{Combining.}
Adding \eqref{eq:g1-badfib}, \eqref{eq:g1-lo}, and \eqref{eq:g1-hi},
\[
\sum_{\alpha\in\Rich^\star}|B_\alpha|
 \;\le\;\Bigl(\tfrac{\rho}{2}+\tfrac{\rho}{4}+\tfrac{\rho}{4}\Bigr)W^\star
 \;=\;\rho\,W^\star.
\qedhere
\]
\end{proof}

\begin{remark}[Error-vs-conductance tradeoff and scaling]
\label{rem:g1-tradeoff}
Proposition~\ref{prop:gap-g1} exhibits the explicit tradeoff promised
by the gap statement: to achieve total exceptional mass
$\rho\,W^\star$, it suffices to take
\(
\Phi_0 \;\lesssim\; \rho\,\varepsilon_0(\rho)\,/\,t_{\max},
\)
with $\varepsilon_0(\rho)$ determined by the Step-2 stability rate
$\delta(\cdot)$ (under the crude cube-root rate,
$\varepsilon_0(\rho)=\Theta(\rho^3)$, hence
$\Phi_0=\Theta(\rho^4/t_{\max})$).  The $1/t_{\max}$ factor is the
natural scaling of planar isoperimetric stability: boundary budgets are
measured against side length, not area.  Since width-$3$ forces
$t_{\max}\le|P|$, the downstream threshold is $\Phi_0\asymp
\mathrm{poly}(\rho)/|P|$, which matches the conductance regime in
which Step~6 is used.
\end{remark}

\begin{corollary}[Mass fraction after removal]\label{cor:g1-fraction}
Under the hypotheses of Proposition~\ref{prop:gap-g1},
\[
\sum_{\alpha\in\Rich^\star}\bigl|\fiber{\alpha}\setminus
B_\alpha\bigr|
 \;\ge\;(1-\rho)\sum_{\alpha\in\Rich^\star}|\fiber{\alpha}|.
\]
In particular, passing to the overlap weights
$w_{\alpha\beta}=\mu((\fiber{\alpha}\cap\fiber{\beta})\setminus
(B_\alpha\cup B_\beta))$ preserves a $(1-2\rho)$-fraction of the raw
pairwise overlap mass:
\[
\sum_{\alpha,\beta\in\Rich^\star}w_{\alpha\beta}
 \;\ge\;
 \sum_{\alpha,\beta\in\Rich^\star}
   \mu(\fiber{\alpha}\cap\fiber{\beta})\;-\;2\rho\,W^\star,
\]
using the union bound
$\mu(B_\alpha\cup B_\beta)\le|B_\alpha|+|B_\beta|$ and
Proposition~\ref{prop:gap-g1} on each term.
\end{corollary}

\begin{proof}
The first line is the complement of Proposition~\ref{prop:gap-g1}.  For
the second, for each ordered pair,
$\mu(\fiber{\alpha}\cap\fiber{\beta})-w_{\alpha\beta}\le
\mu((B_\alpha\cup B_\beta)\cap\fiber{\alpha}\cap\fiber{\beta})\le
|B_\alpha|\cdot\mathbf{1}[\alpha=\alpha]+|B_\beta|\cdot\mathbf{1}[\beta=\beta]$;
sum over $\beta$ (resp.\ $\alpha$) and use the trivial bound
$\sum_\beta\mathbf{1}[\fiber{\alpha}\cap\fiber{\beta}\ne\varnothing]\le
|\Rich^\star|$ together with Proposition~\ref{prop:gap-g1} to obtain
the stated loss $\le 2\rho\,W^\star$.
\end{proof}

Corollary~\ref{cor:g1-fraction} is the statement actually consumed by
Lemma~\ref{lem:sum-step4} and the Theorem~\ref{thm:step6} dichotomy:
the exceptional masses $B_\alpha$ never erode the dichotomy by more
than a fixed small fraction.

\subsection{Formal definition of incoherence}\label{sec:gap-incoherence}

We now formalise what it means for the orientations $\sigma_\alpha$ and
$\sigma_\beta$ to be \emph{incoherent} on the overlap
$(\fiber{\alpha}\cap\fiber{\beta})\setminus(B_\alpha\cup B_\beta)$, so
the partition of active pairs into bad/non-bad is rigorous and
Step~4's hypothesis is verified.

Fix an active pair $(\alpha,\beta)\in\Rich^\star\times\Rich^\star$ and
write
\[
  \Omegareg_{\alpha\beta}
  \;:=\;(\fiber{\alpha}\cap\fiber{\beta})\setminus(B_\alpha\cup B_\beta)
\]
for the good overlap (the regular commuting overlap
$\Omegareg_{xy,uv}$ of Step~3 applied to the rich pairs
$\alpha=(x,y)$, $\beta=(u,v)$).
Step~3 supplies, for every $L\in\Omegareg_{\alpha\beta}$:
\begin{itemize}[nosep]
  \item two commuting unit BK moves $a(L),b(L)$ available at $L$ inside
  $\Omegareg_{\alpha\beta}$ --- the \emph{common axes} of
  Def.~\ref{def:common-axes} --- with $a(L)$ an $\alpha$-generator and
  $b(L)$ a $\beta$-generator, both block-constant on the overlap block
  $B(L)$ (Lem.~\ref{lem:common-axes});
  \item a pair of local signs
  $\eta_\alpha(L),\eta_\beta(L)\in\{\pm 1\}$
  (Def.~\ref{def:eta}) recording which of the two oriented BK moves
  $\pm a(L)$ lies in the $\alpha$-positive cone $P_\alpha(L)$ (determined
  by $(\sigma_\alpha,M_\alpha)$ via Def.~\ref{def:positive-cone}) and
  which of $\pm b(L)$ lies in $P_\beta(L)$; both signs are block-constant
  on $\Omegareg_{\alpha\beta}$ (Rem.~\ref{rem:eta-block-constant}).
\end{itemize}

\begin{definition}[Incoherent active pair]\label{def:incoherent-pair}
An active pair $(\alpha,\beta)$ is \emph{incoherent on the overlap}
(equivalently, \emph{bad}) if there is an absolute constant
$c_{\mathrm{inc}}>0$ with
\[
  \Pr_{L\in\Omegareg_{\alpha\beta}}
    \bigl[\eta_\alpha(L)\neq\eta_\beta(L)\bigr]
  \;\ge\;c_{\mathrm{inc}}.
\]
Equivalently (Prop.~\ref{prop:correlation}),
\[
  \mathrm{Corr}(\alpha,\beta)
  \;:=\;\mathbb{E}_{L\in\Omegareg_{\alpha\beta}}
         \bigl[\eta_\alpha(L)\eta_\beta(L)\bigr]
  \;\le\;1-2c_{\mathrm{inc}}.
\]
Otherwise the active pair is \emph{coherent} (\emph{non-bad}):
$\Pr[\eta_\alpha\neq\eta_\beta]=o(1)$, equivalently
$\mathrm{Corr}(\alpha,\beta)\ge 1-o(1)$.  This is
Def.~\ref{def:coherent} transported to the Step~6 notation.
\end{definition}

\begin{remark}[Translated positive-direction BK moves]
At each $L\in\Omegareg_{\alpha\beta}$ the pair
$(\eta_\alpha(L)\,a(L),\,\eta_\beta(L)\,b(L))$ is precisely the pair of
positive-direction BK moves read off by the two local staircase rules,
expressed in the common overlap axes $(a(L),b(L))$.  Incoherence means
these translated positive directions disagree in sign on a positive
fraction of the overlap, so no common monotone Boolean description
simultaneously realises both local staircases on a $1-o(1)$ fraction of
$\Omegareg_{\alpha\beta}$.
\end{remark}

\begin{remark}[Well-definedness of the bad/non-bad partition]
The common axes $(a(L),b(L))$ and the signs
$(\eta_\alpha(L),\eta_\beta(L))$ are coordinate- and block-invariant
(Lem.~\ref{lem:common-axes}(G5.1--G5.3);
Rem.~\ref{rem:eta-block-constant}), so
Def.~\ref{def:incoherent-pair} depends only on the overlap
$\Omegareg_{\alpha\beta}$ and the Step-2 data
$(\sigma_\alpha,M_\alpha,\sigma_\beta,M_\beta)$ --- not on any choice of
interface-local chart.  Consequently the partition of active pairs into
\emph{bad} (incoherent) and \emph{non-bad} (coherent) is well-defined,
and the Step~4 hypothesis of Theorem~\ref{thm:step4} is verified for
every bad active pair with quantitative incoherence fraction
$\delta\ge c_{\mathrm{inc}}$ (cf.\ \texttt{step4.tex}
Def.~\ref{def:coherence-cited}).
\end{remark}

\subsection{Commuting-square overlap model}\label{sec:gap-commuting}

We reduce this to Step~1.  Write
$\alpha=(x,y)$, $\beta=(u,v)$ for two distinct $S$-good rich pairs
with good fibers $\fiber{\alpha},\fiber{\beta}$ and coordinate maps
$\pi_\alpha,\pi_\beta$ (Definition of good fiber).
Set
\[
  \Omega_{\alpha\beta} := \fiber{\alpha}\cap\fiber{\beta},
  \qquad
  \Omega^{\circ}_{\alpha\beta}
    := \Omega_{\alpha\beta}
       \setminus\bigl(\mathrm{Bad}_\alpha\cup\mathrm{Bad}_\beta
                      \cup\mathrm{Int}_{\alpha\beta}\bigr),
\]
where $\mathrm{Bad}_\bullet$ is the Step~1 bad set of the
corresponding interface and $\mathrm{Int}_{\alpha\beta}$ is the
\emph{interaction locus} of Step~1, Lemma ``Bounded interaction of
distinct rich pairs'' (the set of $L$ at which some element of
$\{x,y\}$ is $L$-adjacent to an element of $\{u,v\}\cup C(u,v)$, or
symmetrically).

\begin{lemma}[Commuting-square overlap model]\label{lem:gap3}
For every distinct pair $\alpha,\beta\in\Rich^\star$:
\begin{enumerate}[label=\textup{(\alph*)},nosep]
\item\label{it:gap3-square}
  \textup{(Local $2\times 2$ commuting square.)}
  At every $L\in\Omega^{\circ}_{\alpha\beta}$ the four BK generators
  $\pm e_1^{(\alpha)},\pm e_2^{(\alpha)}$ and
  $\pm e_1^{(\beta)},\pm e_2^{(\beta)}$ act by transpositions on
  pairwise disjoint pairs of elements of $P$; in particular each
  $(\alpha)$-generator commutes with each $(\beta)$-generator, and
  for any choice of signs the resulting four BK moves close up to a
  $2\times 2$ square inside $\Omega^{\circ}_{\alpha\beta}$ (i.e.\ all
  four BK compositions land at the same vertex).
\item\label{it:gap3-density}
  \textup{(Positive density.)}
  In the rich regime $t(\alpha),t(\beta)=\Theta(n)$ and
  $|\Omega_{\alpha\beta}|\gg n^2$,
  \[
    |\Omega^{\circ}_{\alpha\beta}|
    \;\ge\;(1-o(1))\,|\Omega_{\alpha\beta}|.
  \]
  In particular, for any fixed $\rho<1$ and all but a vanishing
  fraction of active pairs $(\alpha,\beta)$,
  \[
    \mu\bigl(\Omega^{\circ}_{\alpha\beta}\setminus
              (B_\alpha\cup B_\beta)\bigr)
    \;\ge\;\rho\,w_{\alpha\beta}.
  \]
\end{enumerate}
\end{lemma}

\begin{proof}
Both clauses are quotations of Step~1, Corollary ``Commuting overlap;
clause (S1.4) of Step~4'' (\texttt{step1.tex},
\texttt{cor:overlap}), specialised to the labels $\alpha,\beta$.

\emph{\ref{it:gap3-square}.}
A unit $(x,y)$-BK move swaps $x$ or $y$ with an $L$-adjacent
$c_k\in C(x,y)$, and a unit $(u,v)$-BK move swaps $u$ or $v$ with an
$L$-adjacent $c_\ell\in C(u,v)$.  Off the interaction locus
$\mathrm{Int}_{\alpha\beta}$ the four elements
$\{x\text{ or }y,c_k,u\text{ or }v,c_\ell\}$ are pairwise distinct and
pairwise non-adjacent at $L$.  Two transpositions on disjoint pairs of
elements commute; applying them in either order yields the same
permutation, i.e.\ the same vertex of $G_{\BK}$.  Order-convexity of
the joint coordinate domain (Step~1, ``good fiber'' clause) ensures
the intermediate vertices remain in $\Omega^{\circ}_{\alpha\beta}$, so
the $2\times 2$ square is fully embedded.

\emph{\ref{it:gap3-density}.}
By Step~1, each Step-1 bad set $\mathrm{Bad}_\bullet$ is contained in
$O(1)$ strips of width $O(t_\bullet)$, and by the bounded-interaction
lemma the interaction locus $\mathrm{Int}_{\alpha\beta}$ is contained
in $O(1)$ strips of width $O(t(\alpha)+t(\beta))$.  A strip of width
$w$ in the joint coordinate domain has counting mass at most
$w\cdot\sqrt{|\Omega_{\alpha\beta}|}$ (one free coordinate against a
generic fibre slice of size $\sqrt{|\Omega_{\alpha\beta}|}$).  Summing,
\[
  |\Omega_{\alpha\beta}\setminus\Omega^{\circ}_{\alpha\beta}|
  = O\bigl((t(\alpha)+t(\beta))\sqrt{|\Omega_{\alpha\beta}|}\bigr)
  = o(|\Omega_{\alpha\beta}|)
\]
in the rich regime.  Since $\mathrm{Bad}_\alpha\cup\mathrm{Bad}_\beta
\subseteq B_\alpha\cup B_\beta$ for $S$-good interfaces (the Step-2
exceptional sets contain the Step-1 bad sets by construction; if not,
absorb the difference into the Step-2 error budget controlled by
Proposition~\ref{prop:gap-g1}), the second inequality follows: for any
$\rho<1$ and
$|\Omega_{\alpha\beta}|$ above the threshold $w_0$ from the active-pair
condition, the deletion of the two bad sets and the interaction locus
costs $o(w_{\alpha\beta})$, hence retains a $\rho$-fraction.
\end{proof}

\textbf{Witness family.}
We define the Step~4 witness family $E_{\alpha\beta}$ as the set of BK
edges that participate in some $2\times 2$ commuting square produced
by Lemma~\ref{lem:gap3}\,\ref{it:gap3-square}, with both endpoints in
$\Omega^{\circ}_{\alpha\beta}\setminus(B_\alpha\cup B_\beta)$.  By
clause~\ref{it:gap3-density},
\[
  |E_{\alpha\beta}|
  \;\gtrsim\;
  \mu\bigl((\fiber{\alpha}\cap\fiber{\beta})
            \setminus(B_\alpha\cup B_\beta)\bigr)
  \;=\; w_{\alpha\beta}
\]
for all but a vanishing fraction of active pairs, which is the form
consumed by Step~4 in
\eqref{eq:bad-bounded-by-boundary}.

\subsection{Rich-case lower bound on $M$}\label{sec:gap-second-moment}

\begin{lemma}[Identity and rich-case lower bound on $M$]\label{lem:gap-G5}
Set $\fiber{\alpha}^\star := \fiber{\alpha}\setminus B_\alpha$ and
\[
I(L)\;:=\;\#\{\alpha\in\Rich^\star : L\in\fiber{\alpha}^\star\}.
\]
\begin{enumerate}[label=(\roman*),nosep]
  \item \textbf{Exact identity.} $M \;=\; \sum_{L\in\LE(P)} I(L)^2$.
  \item \textbf{Rich case.} Assume case~(R) of Step~5
    (step5.tex, Theorem~2.1). Then $M \ge c_5\,|\LE(P)|$
    for a constant $c_5=c_5(T)>0$ depending only on the rich threshold
    $T$ and the Step~1 bad-set constant.
  \item \textbf{Volume comparison.} For any admissible conductance
    normalization with $\vol(S)\le|\LE(P)|$ (in particular the counting
    normalization $\vol(\cdot)=|\cdot|$ allowed by Section~1),
    \[
    M \;\ge\; c_5\,\vol(S).
    \]
\end{enumerate}
\end{lemma}

\begin{proof}
\emph{(i) Counting identity.} For $L\in\LE(P)$,
\begin{align*}
I(L)^2
 &=\Big(\sum_{\alpha\in\Rich^\star}\mathbf{1}[L\in\fiber{\alpha}^\star]\Big)^{\!2}
  =\sum_{\alpha,\beta\in\Rich^\star}
        \mathbf{1}[L\in\fiber{\alpha}^\star]\,
        \mathbf{1}[L\in\fiber{\beta}^\star]\\
 &=\sum_{\alpha,\beta\in\Rich^\star}
      \mathbf{1}\!\big[L\in(\fiber{\alpha}\cap\fiber{\beta})
                               \setminus(B_\alpha\cup B_\beta)\big],
\end{align*}
where the last equality uses
$\fiber{\alpha}^\star\cap\fiber{\beta}^\star
 =(\fiber{\alpha}\cap\fiber{\beta})\setminus(B_\alpha\cup B_\beta)$.
Summing over $L$ and exchanging order of summation,
\[
\sum_L I(L)^2
 \;=\;\sum_{\alpha,\beta\in\Rich^\star}
        \big|(\fiber{\alpha}\cap\fiber{\beta})
               \setminus(B_\alpha\cup B_\beta)\big|
 \;=\;\sum_{\alpha,\beta\in\Rich^\star} w_{\alpha\beta}
 \;=\; M.
\]
The identity is \emph{exact}: $w_{\alpha\beta}$ is defined with the
bad-set exclusions $B_\alpha\cup B_\beta$ baked in, and $I(L)$ is
taken over $\Rich^\star$, so no error term remains. In particular the
heuristic ``$+O(\text{bad-set losses})$'' flagged in earlier statements
of this gap is exactly zero with the present normalizations.

\smallskip
\emph{(ii) Rich case.}
By case~(R) of Step~5 together with the Fiber-Mass Conversion Lemma
(step5.tex, Lemma~4.9) and the Step~1 bad-set thinness
$|B_\alpha|=O\!\big(t_\alpha\cdot|\LE(P)|/(|A|\,|B|\,|C|)\big)$,
the first-moment inequality
\[
\sum_L I(L)
 \;=\;\sum_{\alpha\in\Rich^\star}\big|\fiber{\alpha}^\star\big|
 \;\ge\; \tfrac{c_T}{2}\,|\LE(P)|
\]
is established inside the proof of step5.tex, Corollary~6.1. Applying
Cauchy--Schwarz on the uniform measure over $\LE(P)$,
\[
\Big(\sum_L I(L)\Big)^{\!2}
 \;\le\; |\LE(P)|\cdot\sum_L I(L)^2,
\]
and combining with (i),
\[
M \;=\;\sum_L I(L)^2
 \;\ge\;\frac{\big(\sum_L I(L)\big)^2}{|\LE(P)|}
 \;\ge\;\Big(\tfrac{c_T}{2}\Big)^{\!2}\,|\LE(P)|
 \;=:\; c_5\,|\LE(P)|.
\]
This is the same bound as step5.tex, Corollary~6.1, rewritten in the
Step~6 pairwise notation via the exact identity of~(i).

\smallskip
\emph{(iii) Volume comparison.}
Section~1 admits any equivalent normalized-conductance notion. Under
the counting normalization $\vol(S)=|S|$, trivially
$\vol(S)\le|\LE(P)|$, and (ii) gives $M\ge c_5\,\vol(S)$. For a
general $\vol$ with $\vol(\LE(P))\le C_{\vol}\,|\LE(P)|$ and
$\vol(S)\le\vol(\LE(P))$, the same argument yields
$M\ge (c_5/C_{\vol})\,\vol(S)$. In the degree-weighted case
$\vol(S)=\sum_{L\in S}\deg_{G_{\BK}}(L)$, width-$3$ caps the BK-degree
at $\le n$, so $C_{\vol}\le n$; the resulting constant is absorbed into
the low-conductance threshold $\Phi_0$ under which
Theorem~\ref{thm:step6} is invoked.
\end{proof}

\begin{remark}[Why Cauchy--Schwarz suffices]\label{rem:G5-why-cs}
The Step~5 input is a first-moment statement
$\sum_\alpha|\fiber{\alpha}^\star|\gtrsim|\LE(P)|$, whereas
Theorem~\ref{thm:step6} consumes a second-moment statement
$\sum_L I(L)^2\gtrsim|\LE(P)|$. The promotion is free because
$I(L)\ge 0$: Jensen's inequality (equivalently, Cauchy--Schwarz on the
uniform measure) gives $\mathbb{E}[I^2]\ge(\mathbb{E}[I])^2$, losing
only a constant factor. No strengthening of Step~5's gaps G1--G5 is
required; see step5.tex, Remark~6.2.
\end{remark}

\begin{remark}[Closure of the dichotomy]\label{rem:G5-closes-dichotomy}
Combining Lemma~\ref{lem:gap-G5}(iii) with the expansion case~(i) of
Theorem~\ref{thm:step6} yields
$|\bd S|\ge\eta\,M\ge\eta c_5\,\vol(S)$, i.e.\ $\Phi(S)\ge\eta c_5$.
Hence any cut $S$ with $\Phi(S)<\eta c_5$ must land in the coherence
case~(ii). This discharges step~(8) of the compressed proof
(§\ref{sec:proof-dichotomy}) and is the precise sense in which G5 is
load-bearing: without it, case~(i) is not a contradiction and the
dichotomy collapses to a tautology.
\end{remark}

%-----------------------------------------------------------------------
\section{Downstream outputs}
%-----------------------------------------------------------------------

Step~6 produces two named conclusions that are consumed downstream:

\begin{itemize}[nosep]
  \item \emph{Conclusion A} (Theorem~\ref{thm:step6}, case~(ii)):
  $\sum_{\text{bad}}w_{\alpha\beta}\le\delta M$.
  Used by Step~7 as the aggregate ``mostly coherent on overlaps''
  assumption.

  \item \emph{Conclusion B} (Corollary~\ref{cor:pointwise}):
  for most $L$, most visible rich interfaces at $L$ agree on an
  orientation $s(L)$.  Used by Step~7 as the pointwise input for the
  globalization / cocycle argument.
\end{itemize}

\noindent
Step~7 (globalization) feeds on \emph{Conclusion B} and produces a
global scalar potential $a\colon P\to\mathbb{R}$ such that
$S\approx\{H\ge c\}$ for $H(L)=\sum_z a(z)\mathrm{pos}_L(z)$.



\begin{abstract}
Step~6 leaves us with the coherent case: almost all rich interfaces
carry pairwise-compatible local staircase rules.
Step~7 upgrades this pairwise compatibility into a single global
scalar rule, and then shows that the resulting global-threshold
description forces the poset to degenerate to a layered / 1D regime,
which in width~$3$ is directly balanced or reducible.
The core new technical content is a cocycle-to-potential upgrade on a
weighted interface graph; all other pieces reduce to double counting,
gradient integration on a connected graph, or direct averaging.
\end{abstract}


%-----------------------------------------------------------------------
\section{Setup and recall}
%-----------------------------------------------------------------------

Let $P$ be a width-3 poset and $\LE(P)$ its set of linear extensions.
Let $G_{\BK}(P)$ be the BK graph on $\LE(P)$.
For $S\subseteq\LE(P)$ let $\Phi(S)=|\bd S|/\vol(S)$ be its
conductance.

We use the rich-pair notation of Step~1 (Definition~\ref{def:rich}):
for $x\parallel y$ in $P$, $C(x,y)$ is the common-neighbor chain (a
chain by width~$3$), $t(x,y)=|C(x,y)|$, the richness threshold is
$T\ge 1$, and $(x,y)$ is \emph{rich} iff $t(x,y)\ge T$.
Denote by $\Rich$ the set of rich pairs.

\begin{definition}[Good fiber, interface]\DEP{Step~1}
For each $e=(x,y)\in\Rich$, Step~1 produces a good fiber
$\fiber{e}\subseteq\LE(P)$ with a coordinate map
$\pi_e\colon\fiber{e}\to D_e\subseteq[0,t(x,y)]^2$, such that BK moves
inside $\fiber{e}$ are exactly unit grid moves and $D_e$ is
order-convex.
\end{definition}

\begin{definition}[Local staircase, orientation and threshold]
\DEP{Step~2}\label{def:local-staircase}
For each $S$-good rich interface $e$, Step~2 says $S\cap\fiber{e}$,
viewed in coordinates, is close to a monotone staircase.
We shall require the \emph{normalized} form:
\begin{equation}\label{eq:signed-threshold}
  S\cap\fiber{e}
  \;\approx\;
  \bigl\{\,(i,j)\in D_e :\; \sigma_e\,(j-i) \ge \tau_e\,\bigr\},
\end{equation}
where $\sigma_e\in\{\pm1\}$ is the oriented monotone direction and
$\tau_e\in\mathbb{Z}$ is a local threshold.
See Section~\ref{sec:normalize} for the proof that Step~2's
staircase admits this normalization.
\end{definition}

\begin{definition}[Interface graph $G_{\Rich}$]\label{def:G_Rich}\DEP{Step~5, Step~6}
Let $\Rich^\star\subseteq\Rich$ be the set of $S$-good rich interfaces
(Step~6).
The \emph{interface graph} $G_{\Rich}$ has vertex set $\Rich^\star$,
with an edge $e\sim f$ when
$w_{ef}=\mu((\fiber{e}\cap\fiber{f})\setminus(B_e\cup B_f))$ is above
a fixed positive threshold (\emph{active pair}).
Similarly, a \emph{triple} $(e,f,g)$ is \emph{active} if the three
pairwise overlaps and the triple overlap are each above positive
thresholds.
\end{definition}

%-----------------------------------------------------------------------
\section{The Step 7 proposition}
%-----------------------------------------------------------------------

\begin{proposition}[Step~7, informal]\label{prop:step7-informal}
Assume the hypotheses of Steps~1, 2, 5, 6, and in particular assume
$S$ is in the coherence case of Step~6 (pointwise coherence,
Corollary~B of Step~6).
Then for $\varepsilon$ small enough, one of the following holds:

\begin{enumerate}[label=(\roman*),nosep]
  \item \textbf{Global threshold case.}
  There exists a function $a\colon P\to\mathbb{R}$ (an \emph{element
  potential}) and a constant $c\in\mathbb{R}$ such that, setting
  \[
    H(L) := \sum_{z\in P} a(z)\,\pos_L(z),
  \]
  we have
  \[
    1_S(L) \;\approx\; 1_{\{H(L)\ge c\}}
  \]
  on $(1-o(1))$ of $\LE(P)$, and every rich BK move across an
  interface $e=(x,y)$ changes $H$ with the sign predicted by $\sigma_e
  =\sgn(a(y)-a(x))$.

  \item \textbf{Collapse case.}
  $P$ admits a layered width-$3$ decomposition
  $P = L_1\sqcup\cdots\sqcup L_m$, with each $L_r$ an antichain of
  size $\le 3$, such that after deleting $o(1)$ of the mass of
  $\LE(P)$:
  \begin{itemize}[nosep]
    \item every rich interface is between adjacent or near-adjacent
      levels;
    \item BK moves move mass forward/backward across the level order;
    \item linear extensions are determined mainly by choices internal
      to consecutive levels.
  \end{itemize}
\end{enumerate}
In either case, combined with width~$3$, $P$ admits a balanced pair
(or reduces to a strictly smaller instance).
\end{proposition}

\noindent
Proposition~\ref{prop:step7-informal} is the complementary case to
Step~4: if interfaces do not generate boundary (coherence case), then
they must all be reading the same global coordinate.

The formal statement is broken into three propositions in
Section~\ref{sec:formal}.

%-----------------------------------------------------------------------
\section{Normalizing the local staircase}\label{sec:normalize}
%-----------------------------------------------------------------------

Step~2 produces a monotone staircase region in each $S$-good fiber.
For Step~7 we must normalize it into the explicit signed-threshold
form~\eqref{eq:signed-threshold}.

\begin{lemma}[Signed-threshold normal form]\label{lem:signed-threshold}
\DEP{Step~2}
For each $e\in\Rich^\star$ there is a choice of orientation
$\sigma_e\in\{\pm1\}$ and threshold $\tau_e\in\mathbb{Z}$ such that, on
$\fiber{e}\setminus B_e$,
\[
  1_S(L) \;=\; 1\!\left\{\,\sigma_e\cdot (j(L)-i(L))\ge \tau_e\,\right\}
  + O(\varepsilon_2),
\]
where $(i(L),j(L))=\pi_e(L)$ and $\varepsilon_2$ is the Step~2 error
parameter.
\end{lemma}

\begin{proof}
By Step~2's per-fiber staircase approximation (\texttt{step2.tex}
Prop.~3.2), there exists a staircase
$M_e\in\mathrm{Stair}(D_e)$ of some type $\tilde\sigma_e\in\{\pm 1\}$
(in the sense of Step~2 Def.~2.1, equivalently Step~3 Def.~3.1) with
\[
  |(S\cap\fiber{e})\symdiff\pi_e^{-1}(M_e)| \;\le\; \varepsilon_2\,|\fiber{e}|.
\]
In the affine form of Step~3 Def.~3.1, $M_e$ rewrites as
\begin{equation}\label{eq:affine-staircase}
  M_e = \bigl\{(i,j)\in D_e : i+\tilde\sigma_e\,j \le \varphi_e(i-\tilde\sigma_e\,j)\bigr\},
\end{equation}
with staircase threshold
$\varphi_e\colon\mathbb{Z}\to\mathbb{Z}\cup\{\pm\infty\}$ weakly
monotone on the transverse axis $v:=i-\tilde\sigma_e\,j$.

\smallskip
\noindent\emph{(1) Collapsing $\varphi_e$ to a constant.}
The transverse range
$V_e:=\{v\in\mathbb{Z}:\exists(i,j)\in D_e,\ i-\tilde\sigma_ej=v\}$
is an integer interval of length $O(t_e)$ (since $D_e\subseteq
[0,t_e]^2$), and $\varphi_e$ is finite on $V_e$ off a set of size
$O(1)$ by order-convexity.  Let
\[
  \mathrm{spread}(\varphi_e) := \sup_{v\in V_e}\varphi_e(v)\;-\;\inf_{v\in V_e}\varphi_e(v).
\]
The staircase boundary of $M_e$ is a monotone lattice path along the
transverse axis $v=i-\tilde\sigma_e j$: each unit of transverse spread
in $\varphi_e$ forces at least one BK-boundary cell of $M_e$ at each
of the $\Theta(t_e)$ positions along the longitudinal axis, so the
BK-boundary mass of $M_e$ inside $\fiber{e}$ satisfies
$|\bd_{BK}M_e\cap\fiber{e}|\ge c\cdot\mathrm{spread}(\varphi_e)\cdot t_e$
for an absolute constant $c>0$.  By Step~2 Prop.~3.2 and the richness
density bound $|\fiber{e}|\asymp t_e^2$,
\[
  |\bd_{BK}S\cap\fiber{e}| \;=\; O(\varepsilon_2\,|\fiber{e}|)
  \;=\; O(\varepsilon_2\,t_e^2),
\]
so a staircase with transverse spread $s$ contributes
$\Omega(s\cdot t_e)$ BK-boundary cells, giving
$s\cdot t_e\le O(\varepsilon_2\,t_e^2)$, whence
\[
  \mathrm{spread}(\varphi_e) \;=\; O(\varepsilon_2\,t_e).
\]
Define $\bar\tau_e:=\mathrm{med}_{v\in V_e}\varphi_e(v)\in\mathbb{Z}$
and the canonicalized half-plane
$\hat M_e:=\{(i,j)\in D_e:i+\tilde\sigma_e\,j\le\bar\tau_e\}$.  The
symmetric difference $M_e\symdiff\hat M_e$ is contained in the
transverse-strip union $\bigcup_{v\in V_e}$ of intervals of length
$|\varphi_e(v)-\bar\tau_e|$, whose total cardinality is bounded by
$\mathrm{spread}(\varphi_e)\cdot|V_e|=O(\varepsilon_2 t_e^2)
=O(\varepsilon_2|\fiber{e}|)$ (using $|\fiber{e}|\asymp t_e^2$ from
the richness hypothesis).

\smallskip
\noindent\emph{(2) Writing $\hat M_e$ as a signed $(j-i)$-threshold.}
Unfolding the linear form $h_e(i,j):=i+\tilde\sigma_e\,j$:
\begin{itemize}[nosep,topsep=2pt]
  \item If $\tilde\sigma_e=-1$: $\hat M_e=\{i-j\le\bar\tau_e\}=
    \{+1\cdot(j-i)\ge-\bar\tau_e\}$.  Set $\sigma_e:=+1$,
    $\tau_e:=-\bar\tau_e$.
  \item If $\tilde\sigma_e=+1$: apply the single-axis reflection
    $(i,j)\mapsto(i,t_e-j)$, an affine isomorphism of the
    order-convex grid $[0,t_e]^2$ that preserves unit BK moves (it
    permutes $\{\pm e_1,\pm e_2\}$) and is absorbed into the
    coordinate chart $\pi_e$ at the cost of a sign flip of
    $\tilde\sigma_e$ (Step~3 Lemma~3.2 discussion of coordinate
    symmetries).  The image of $\hat M_e$ under the
    reflection is $\{i+(t_e-j')\le\bar\tau_e\}=
    \{+1\cdot(j'-i)\ge t_e-\bar\tau_e\}$, falling into the previous
    case; set $\sigma_e:=+1$, $\tau_e:=t_e-\bar\tau_e$ in the
    reflected frame.  Equivalently, in the original frame one may
    write the half-plane as $\{-1\cdot(j-i)\ge\bar\tau_e-t_e\}$,
    $\sigma_e:=-1$, $\tau_e:=\bar\tau_e-t_e$.
\end{itemize}
In either case $\sigma_e\in\{\pm 1\}$ and $\tau_e\in\mathbb{Z}$ are
determined by $(\tilde\sigma_e,\bar\tau_e)$ and the choice of frame;
the sign is canonical by Step~3's canonical orientation theorem
(\texttt{step3.tex} Thm.~3.9).

Combining (1) and the preceding Step~2 approximation,
\[
  \bigl|(S\cap\fiber{e})\symdiff
    \pi_e^{-1}\bigl\{(i,j)\in D_e:\sigma_e(j-i)\ge\tau_e\bigr\}\bigr|
  \;\le\; C\,\varepsilon_2\,|\fiber{e}|,
\]
which yields the pointwise identity
$1_S(L)=1\{\sigma_e(j(L)-i(L))\ge\tau_e\}+O(\varepsilon_2)$ on
$\fiber{e}\setminus B_e$ (absorbing the constant $C$ into
$\varepsilon_2$).
\end{proof}

%-----------------------------------------------------------------------
\section{Globalizing coherence: sign consistency}\label{sec:sign}
%-----------------------------------------------------------------------

Step~6's output is that, weighted by overlap mass, almost all pairs
$(e,f)$ of active interfaces are coherent (equivalently,
$\sigma_e,\sigma_f$ are compatible under the common overlap
coordinates).
Step~7 needs the \emph{global} version: a single consistent sign
convention on a giant component of $G_{\Rich}$.

\begin{lemma}[Lemma A: sign consistency]\label{lem:sign-consistency}
\DEP{Step~6, Corollary~B}
Assume the coherence case of Step~6.
There is a set $\mathcal{E}^\star\subseteq\Rich^\star$ of rich
interfaces carrying $(1-o(1))$ of the rich-interface incidence mass,
and a choice of signs
$\widetilde\sigma_e\in\{\pm1\}$ for $e\in\mathcal{E}^\star$, such that:
\begin{enumerate}[nosep]
  \item $\widetilde\sigma_e = \sigma_e$ for all but an $o(1)$-fraction
    of overlap incidence; equivalently, Step~7's corrections flip at
    most $o(1)$ of the signed interfaces.
  \item For every active pair $(e,f)$ in $\mathcal{E}^\star\times
    \mathcal{E}^\star$, $\widetilde\sigma_e$ and $\widetilde\sigma_f$
    agree on the overlap under the common overlap coordinates.
\end{enumerate}
\end{lemma}

\begin{proof}[Proof of Lemma~\ref{lem:sign-consistency}]
The input is Corollary~B of Step~6 (Corollary~\ref{cor:pointwise} in
step6.tex): there exists a measurable reference sign
$s\colon\LE(P)\to\{\pm1\}$ and a set $\Omega\subseteq\LE(P)$ with
$\mu(\LE(P)\setminus\Omega)\le\sqrt{\delta}$ such that for every
$L\in\Omega$,
\begin{equation}\label{eq:ptw-coh}
  \#\{\alpha\in\Rich_L:\sigma_\alpha\ne s(L)\}
  \;\le\; \sqrt{\delta}\,|\Rich_L|,
  \qquad \delta=o(1).
\end{equation}

\textbf{Definition of $\widetilde\sigma$.}
For $\alpha\in\Rich^\star$ write $\mu_\alpha=\mu(\fiber{\alpha}\setminus
B_\alpha)$ and
\[
  m_\alpha \;:=\;
  \mu\bigl(\{L\in\fiber{\alpha}\cap\Omega : \sigma_\alpha=s(L)\}\bigr).
\]
Set $\epsilon_\alpha=+1$ if $m_\alpha\ge\tfrac12\mu_\alpha$ and
$\epsilon_\alpha=-1$ otherwise, and let
$\widetilde\sigma_\alpha:=\epsilon_\alpha\sigma_\alpha$.  Thus
$\widetilde\sigma$ flips $\sigma$ exactly on those interfaces whose
local orientation disagrees with the $L$-level reference $s$ on a
majority of their incidence mass.

\textbf{Flipped weight is $o(1)$.}
Double-counting \eqref{eq:ptw-coh} as
\[
  \sum_{\alpha\in\Rich^\star}
   \mu\bigl(\{L\in\fiber{\alpha}\cap\Omega:\sigma_\alpha\ne s(L)\}\bigr)
  \;=\;
  \int_{\Omega}\#\{\alpha\in\Rich_L:\sigma_\alpha\ne s(L)\}\,d\mu(L)
\]
and bounding the right side by
$\sqrt{\delta}\int_\Omega|\Rich_L|\,d\mu(L)\le\sqrt{\delta}\,M_0$, with
$M_0=\sum_\alpha\mu_\alpha$ the total incidence mass, we obtain
\begin{equation}\label{eq:db-count}
  \sum_{\alpha\in\Rich^\star}(\mu_\alpha - m_\alpha - \eta_\alpha)
  \;\le\; \sqrt{\delta}\,M_0,
  \qquad
  \eta_\alpha:=\mu(\fiber{\alpha}\cap(\LE(P)\setminus\Omega)),
\end{equation}
and $\sum_\alpha\eta_\alpha=\int_{\Omega^c}|\Rich_L|\,d\mu\le
\sqrt{\delta}\,M_0$ (again by Step~6 pointwise control, or trivially
from $|\Rich_L|\le|\Rich^\star|$ combined with the second-moment
visibility of Lemma~\ref{lem:triple-visibility}).  If
$\epsilon_\alpha=-1$ then $\mu_\alpha-m_\alpha
\ge\tfrac12\mu_\alpha$, so \eqref{eq:db-count} gives
\[
  \tfrac12\sum_{\alpha:\epsilon_\alpha=-1}\mu_\alpha
  \;\le\; \sqrt{\delta}\,M_0 + \sum_\alpha\eta_\alpha
  \;\le\; 2\sqrt{\delta}\,M_0,
\]
hence the flipped weight is $\le 4\sqrt{\delta}\,M_0 = o(M_0)$.
This is condition~(1).

\textbf{Refined set $\mathcal{E}^\star$.}
Define
$\mathcal{E}^\star:=\{\alpha\in\Rich^\star :
\max(m_\alpha,\mu_\alpha-m_\alpha)\ge(1-\delta^{1/4})\mu_\alpha\}$,
the set of interfaces whose majority alignment with~$s$ is
$(1-\delta^{1/4})$-dominant.  By Markov on \eqref{eq:db-count},
$\Rich^\star\setminus\mathcal{E}^\star$ has incidence weight
$\le\delta^{1/4}\,M_0=o(M_0)$, so $\mathcal{E}^\star$ carries
$(1-o(1))\,M_0$.

\textbf{Overlap compatibility.}
Let $(e,f)$ be an active pair in $\mathcal{E}^\star\times
\mathcal{E}^\star$ with overlap weight
$w_{ef}=\mu((\fiber{e}\cap\fiber{f})\setminus(B_e\cup B_f))$.
By \eqref{eq:ptw-coh}, on every $L\in\fiber{e}\cap\fiber{f}\cap\Omega$
except a $\sqrt{\delta}$-minority of $\Rich_L$, both $\sigma_e$ and
$\sigma_f$ equal $s(L)$; integrating,
\[
  \mu\bigl(\{L\in\fiber{e}\cap\fiber{f}\cap\Omega:
      \sigma_e=\sigma_f=s(L)\}\bigr)
  \;\ge\; w_{ef}\,(1-o(1)).
\]
By the $(1-\delta^{1/4})$-dominance defining $\mathcal{E}^\star$, for
each $\alpha\in\{e,f\}$ the set
$\{L\in\fiber{\alpha}:\widetilde\sigma_\alpha(L)=s(L)\}$ carries
measure at least $(1-\delta^{1/4})\mu_\alpha$; equivalently, the
$\widetilde\sigma_\alpha$-disagreement set
$\{L\in\fiber{\alpha}:\widetilde\sigma_\alpha(L)\ne s(L)\}$ has
measure at most $\delta^{1/4}\mu_\alpha$.  Restricting to the overlap
$(\fiber{e}\cap\fiber{f})\setminus(B_e\cup B_f)$ of weight
$w_{ef}$ and using the ``active pair'' lower bound
$w_{ef}\ge\delta^{1/4}\min(\mu_e,\mu_f)$, the union of the two
disagreement sets intersects the overlap in measure at most
$2\delta^{1/4}\max(\mu_e,\mu_f)=O(\delta^{1/4})\,w_{ef}$.  On the
complementary $(1-O(\delta^{1/4}))\,w_{ef}$-fraction of the overlap,
both $\widetilde\sigma_e(L)$ and $\widetilde\sigma_f(L)$ equal $s(L)$,
hence $\widetilde\sigma_e=\widetilde\sigma_f$ in the common
coordinates, establishing condition~(2).  (The $\epsilon_e\epsilon_f$
sign bookkeeping is thus automatic: both $\widetilde\sigma_\alpha$
track the single function $s\colon\LE(P)\to\{\pm1\}$.)
\end{proof}

\begin{remark}[Graph-theoretic route]
The same statement can be obtained via disagreement percolation on
$G_{\Rich}$: Step~6's pairwise compatibility says the bad-edge
fraction is $o(1)$, and on the giant component of $G_{\Rich}$
(Lemma~\ref{lem:giant-component}) a Cheeger-type bound converts this to
an $o(1)$-weighted vertex set whose flipping eliminates all bad edges.
The pointwise-coherence route above is quantitatively sharper and does
not require connectedness: Lemma~\ref{lem:sign-consistency} holds
component-wise.  Connectedness enters only downstream, when the
cocycle is integrated to a single potential
(Lemma~\ref{lem:potential}).
\end{remark}

%-----------------------------------------------------------------------
\section{Threshold cocycle on triple overlaps}\label{sec:cocycle}
%-----------------------------------------------------------------------

This is the core new content of Step~7.

Consider three rich interfaces
$e_{xy}=(x,y),\ e_{yz}=(y,z),\ e_{xz}=(x,z)$
whose three pairwise overlaps and the triple overlap
$\fiber{xy}\cap\fiber{yz}\cap\fiber{xz}$ all have positive weight;
we call such a triple an \emph{active triple}.

On an active triple, the three fibers share common chain coordinates
(because $N(x)\cap N(y),\ N(y)\cap N(z),\ N(x)\cap N(z)$ all embed in
a common chain window, by the width-3 normal form of Step~1).
Consequently the local coordinates $(i,j)$ on each fiber measure the
relative placement of the three elements against the \emph{same}
chain.

\begin{lemma}[Lemma B: threshold cocycle]\label{lem:cocycle}
\DEP{Step~1 (commuting-square/triple-overlap geometry), Step~2
(signed-threshold normal form), Section~\ref{sec:sign} (signs chosen
consistently).}
After applying Lemma~\ref{lem:sign-consistency} to fix signs, for a
$(1-o(1))$-fraction of active triples $(e_{xy},e_{yz},e_{xz})$ in
$\mathcal{E}^\star$,
\[
  \tau_{xz} \;=\; \tau_{xy} + \tau_{yz} + O(1).
\]
\end{lemma}

\begin{proof}[Proof of Lemma~\ref{lem:cocycle}]
Fix an active triple $T=(e_{xy},e_{yz},e_{xz})$ with triple-overlap
weight
$w_T:=\mu((\fiber{xy}\cap\fiber{yz}\cap\fiber{xz})\setminus
(B_{xy}\cup B_{yz}\cup B_{xz}))\ge\eta_T|\LE(P)|$ for the active
threshold $\eta_T>0$.  By
Corollary~\ref{cor:triple-overlap}\ref{it:tri-coords}--\ref{it:tri-cube}
of Step~1, there is a subset
$\Omega^{\circ\circ\circ}\subseteq\fiber{xy}\cap\fiber{yz}\cap\fiber{xz}$
of triple-overlap mass
$\mu(\Omega^{\circ\circ\circ})\ge w_T-O(t\sqrt{w_T|\LE(P)|})
=(1-o(1))w_T$ (by part~\ref{it:tri-bad}, with
$t:=\max(t_{xy},t_{yz},t_{xz})=\Theta(n)$ in the rich regime and $w_T$
of order $|\LE(P)|$), on which the three coordinate systems are
simultaneously valid and the three BK move families generate a
commuting $\Z^3$ action.

\textbf{Step 1: Shared chain window and additive identity of
differences.}
By the width-3 common-chain structure (Step~1
Lemma~\ref{lem:CNchain-width3}\footnote{We cite Step~1's width-3
common-chain lemma generically; the point is that in width~$3$, the
three pairwise common chains $C(x,y)=N(x)\cap N(y)$, $C(y,z)$,
$C(x,z)$ each lie in a chain and their pairwise symmetric differences
are $O(1)$-bounded by the neighbourhood endpoint structure.}), the
three pairwise common chains $C(x,y)$, $C(y,z)$, $C(x,z)$ embed into a
single ambient chain $W$ of length $t_W\ge t-O(1)$ with
$|W\triangle C(\bullet)|=O(1)$ for each pairwise chain.
Consequently, for $L\in\Omega^{\circ\circ\circ}$ the signed coordinate
differences
\[
  d_{uv}(L)\;:=\;j_{uv}(L)-i_{uv}(L)
  \;=\;\#\{c\in C(u,v): u<_L c<_L v\}\;-\;\#\{c\in C(u,v): v<_L c<_L u\}
\]
satisfy, writing $n_{uv}^W(L)$ for the analogous signed count against
the ambient chain $W$,
\[
  d_{uv}(L)\;=\;n_{uv}^W(L)\;+\;r_{uv}(L),
  \qquad |r_{uv}(L)|\le|W\triangle C(u,v)|=O(1).
\]
Additivity of signed counts in the linearly ordered window $W$ reads
$n_{xz}^W(L)=n_{xy}^W(L)+n_{yz}^W(L)$ for every $L$ (in a total order
on $W\cup\{x,y,z\}$, the open interval $(x,z)\cap W$ decomposes as the
disjoint union $(x,y)\cap W\ \sqcup\ \{y\}\cap W\ \sqcup\ (y,z)\cap W$
when $x<_L y<_L z$ --- up to the vanishing correction of at most $1$
from the point $y$ itself --- and likewise with signs for any other
ordering of $\{x,y,z\}$ along $W$, since each $c\in W$ is compared to
$x,y,z$ via a single linear order).  Therefore
\begin{equation}\label{eq:d-triple}
  d_{xz}(L)\;=\;d_{xy}(L)+d_{yz}(L)+\Delta_0(T,L),
  \qquad |\Delta_0(T,L)|\le 2+\sum_{uv}|W\triangle C(u,v)|=O(1),
\end{equation}
uniformly in $L\in\Omega^{\circ\circ\circ}$, with
$\Delta_0(T,L):=r_{xz}(L)-r_{xy}(L)-r_{yz}(L)+\epsilon_y(L)$ and
$|\epsilon_y(L)|\le 2$ accounting for the $y$-boundary correction in
the $W$-additivity identity above.  The bound on $|\Delta_0(T,L)|$ is
uniform in $L$, which is all the subsequent steps use; we write simply
$\Delta_0$ below when $L$-dependence is irrelevant.

\textbf{Step 2: Aligning signs.}
Apply Lemma~\ref{lem:sign-consistency} to fix
$\widetilde\sigma_{xy}=\widetilde\sigma_{yz}=\widetilde\sigma_{xz}=+1$
on the triple (absorb any sign flip into
$\tau_{uv}\mapsto-\tau_{uv}$; the cocycle claim is invariant under
this absorption).  A $(1-o(1))$ fraction of active triples have all
three interfaces in $\mathcal{E}^\star$ with pairwise-compatible signs
in the sense of Lemma~\ref{lem:sign-consistency}(2); discard the
remaining $o(1)$-weighted triples.

\textbf{Step 3: Three-way threshold identity.}
Lemma~\ref{lem:signed-threshold} applied to each fiber gives, on
$\Omega^{\circ\circ\circ}\setminus(B_{xy}\cup B_{yz}\cup B_{xz})$
(still of mass $(1-o(1))w_T$ by the bad-set bounds from Step~1),
\begin{equation}\label{eq:three-rules}
  1_S(L)\;=\;
  \mathbf{1}\{d_{xy}\ge\tau_{xy}\}\;=\;
  \mathbf{1}\{d_{yz}\ge\tau_{yz}\}\;=\;
  \mathbf{1}\{d_{xz}\ge\tau_{xz}\}
  \;+\;O(\varepsilon_2)
\end{equation}
where the $O(\varepsilon_2)$ error is the union of the three per-fiber
Step~2 errors, integrating to a symmetric-difference mass
$\le 3\varepsilon_2\min_\bullet|\fiber{\bullet}|$.
Substituting~\eqref{eq:d-triple} into the $xz$-rule yields
\begin{equation}\label{eq:combined-rule}
  \mathbf{1}\{d_{xy}(L)+d_{yz}(L)\ge K_T(L)\}
  \;=\;\mathbf{1}\{d_{xy}\ge\tau_{xy}\}
  \;+\;O(\varepsilon_2),
  \qquad K_T(L):=\tau_{xz}-\Delta_0(T,L),
\end{equation}
where $K_T(L)$ fluctuates by at most $O(1)$ across $L$ by the uniform
bound on $\Delta_0$, and both sides agree with $1_S(L)$ on
$(1-O(\varepsilon_2))$ of $\Omega^{\circ\circ\circ}$.

\textbf{Step 4: Mismatch produces a disagreement rectangle.}
Put $u:=d_{xy}-\tau_{xy}$, $v:=d_{yz}-\tau_{yz}$, and define
$K_*:=\tau_{xz}-\tau_{xy}-\tau_{yz}$ as the cocycle defect.  Then
$K_T(L)-\tau_{xy}-\tau_{yz}=K_*-\Delta_0(T,L)$, so \eqref{eq:combined-rule}
reads
$\mathbf{1}\{u+v\ge K_*-\Delta_0(T,L)\}=\mathbf{1}\{u\ge 0\}+O(\varepsilon_2)$.
Set $K:=K_*-\sup_{L\in\Omega^{\circ\circ\circ}}\Delta_0(T,L)$ (finite
and within $O(1)$ of $K_*$ by~\eqref{eq:d-triple}); then
$\mathbf{1}\{u+v\ge K\}\le\mathbf{1}\{u+v\ge K_*-\Delta_0(T,L)\}$
pointwise, and the two differ on a set of at most
$O(1)\cdot w_T/(Ct^2)$ triple-overlap mass per unit shift, negligible
for the rectangle argument below.
Equation~\eqref{eq:combined-rule} combined with
$\mathbf{1}\{v\ge 0\}=\mathbf{1}\{u\ge 0\}$ (from~\eqref{eq:three-rules})
gives, on $(1-O(\varepsilon_2))$ of $\Omega^{\circ\circ\circ}$,
\begin{equation}\label{eq:uv-rule}
  \mathbf{1}\{u\ge 0\}\;=\;\mathbf{1}\{v\ge 0\}
  \;=\;\mathbf{1}\{u+v\ge K\}.
\end{equation}

Assume for contradiction that $|K|\ge K_0$ for some large constant
$K_0$ to be chosen; by symmetry (interchange $S\leftrightarrow S^c$,
$u\leftrightarrow -u$) assume $K\ge K_0$.
Consider the lattice rectangle
\[
  R\;:=\;\{(u,v)\in\Z^2:\ 0\le u<K/2,\ -K/2\le v<0\}.
\]
On $R$ both equalities in~\eqref{eq:uv-rule} fail
simultaneously: $\mathbf{1}\{u\ge 0\}=1$ but $\mathbf{1}\{v\ge 0\}=0$,
while $\mathbf{1}\{u+v\ge K\}=0$ since $u+v<K/2-0<K$.  Thus every
$L\in\Omega^{\circ\circ\circ}$ with $(u(L),v(L))\in R$ contributes $1$
to the symmetric-difference mass in~\eqref{eq:three-rules}.

We derive the following dimensional pushforward bound, which is the
quantitative content used below:
\begin{equation}\label{eq:pushforward-lb}
  \mu\bigl(\{L\in\Omega^{\circ\circ\circ}:(u(L),v(L))=(u_0,v_0)\}\bigr)
  \;\ge\;\frac{w_T}{Ct^2}
  \qquad\forall\,(u_0,v_0)\in[-ct,ct]^2\cap\Z^2,
\end{equation}
for absolute constants $c,C>0$, with $t:=\min(\txy,t_{y,z},t_{x,z})=\Theta(n)$.

Let $g_u:=e_1^{(x,y)}$ and $g_v:=e_1^{(y,z)}$ be the unit BK-shift
generators from Corollary~\ref{cor:triple-overlap}\ref{it:tri-cube}.
Under the standing hypothesis $w(P)\le 3$ with $\{x,y,z\}$ a
3-antichain, the three common chains $\Cxy,C(y,z),C(x,z)$ are
\emph{pairwise disjoint} (any common element would form a 4-antichain
with $\{x,y,z\}$; see the Remark following
Corollary~\ref{cor:triple-overlap}\ref{it:tri-consistency}).  Hence a
single application of $g_u$ swaps $x$ with an element of $\Cxy$ only,
leaving the $y{:}z$ coordinate pair untouched; so $g_u$ shifts
$u=d_{xy}-\tau_{xy}$ by $+1$ and fixes $v=d_{yz}-\tau_{yz}$.
Symmetrically $g_v$ shifts $v$ by $+1$ and fixes $u$.  The two
generators commute on $\Omega^{\circ\circ\circ}$
(Corollary~\ref{cor:triple-overlap}\ref{it:tri-cube}) and each is an
injective measure-preserving partial bijection (a BK permutation
restricted to points where the move stays in $\Omega^{\circ\circ\circ}$).

By Theorem~\ref{thm:interface}\ref{part:goodfiber} the images
$\piXY(\Omega^{(3)})$ and $\pi_{y,z}(\Omega^{(3)})$ are order-convex
domains of diameter $\Omega(t)$ in each coordinate direction, so for
an absolute $c>0$ the iterated action $g_u^a g_v^b$ with
$(a,b)\in[-ct,ct]^2$ keeps the image inside $\Omega^{(3)}$.  Define
the \emph{interior}
\[
  \Omega^\flat\;:=\;\{L\in\Omega^{\circ\circ\circ}:
    g_u^a g_v^b L\in\Omega^{\circ\circ\circ}\text{ for all }(a,b)\in[-ct,ct]^2\}.
\]
The escape region $\Omega^{\circ\circ\circ}\setminus\Omega^\flat$ is
contained in the union of $O(t)$ translates of
$\Omega^{(3)}\setminus\Omega^{\circ\circ\circ}$ under the $\Z^2$-action,
which by Corollary~\ref{cor:triple-overlap}\ref{it:tri-bad} has total
mass $O(t\cdot t\cdot\sqrt{|\Omega^{(3)}|})=o(|\Omega^{(3)}|)$ when
$t=\Theta(n)$ and $|\Omega^{(3)}|\ge\delta|\LE(P)|\ge e^{\Omega(n)}$.
Hence $\mu(\Omega^\flat)\ge(1-o(1))\mu(\Omega^{\circ\circ\circ})\ge(1-o(1))w_T$.

On $\Omega^\flat$ the $\Z^2$-action
$\langle g_u,g_v\rangle$ is free and measure-preserving, and its
induced action on $(u,v)$ is by integer translation.  Consequently the
pushforward $\nu:=(u,v)_\ast(\mu|_{\Omega^\flat})$ is
\emph{translation-invariant} under unit shifts within $[-ct,ct]^2$:
for every $(u_0,v_0)$ in that box,
$\nu(u_0\pm 1,v_0)=\nu(u_0,v_0)=\nu(u_0,v_0\pm 1)$.  By connectivity
of the box, $\nu$ is constant on $[-ct,ct]^2\cap\Z^2$.  Since
$|u|,|v|\le O(t)$ pointwise (the $d_{\bullet\bullet}$ are fiber-coordinate
values bounded by $t$ up to additive $O(1)$), $\nu$ is supported on at
most $C't^2$ lattice points, so the constant value on the box satisfies
\[
  \nu(u_0,v_0)\;\ge\;\frac{\nu(\Z^2)}{C't^2}\;=\;\frac{\mu(\Omega^\flat)}{C't^2}
  \;\ge\;\frac{(1-o(1))w_T}{C't^2}\;\ge\;\frac{w_T}{Ct^2},
\]
establishing \eqref{eq:pushforward-lb}.  Since
$\mu(\Omega^\flat)\le\mu(\Omega^{\circ\circ\circ})$, the same lower
bound holds for the pushforward of $\mu|_{\Omega^{\circ\circ\circ}}$
itself.  Provided
$K\le ct/2$ (still ``fitting'' in the box), the rectangle $R$ has
$|R|\ge K^2/4$ lattice points, so
\begin{equation}\label{eq:R-mass}
  \mu\bigl(\Omega^{\circ\circ\circ}\cap\{(u,v)\in R\}\bigr)
  \;\ge\;\frac{K^2}{4}\cdot\frac{w_T}{Ct^2}.
\end{equation}
Comparing with the symmetric-difference budget
$3\varepsilon_2|\fiber{xy}|\asymp\varepsilon_2 t^2$ from Step~2,
\[
  \frac{K^2}{4}\cdot\frac{w_T}{Ct^2}\;\le\;3\varepsilon_2 t^2
  \quad\Longrightarrow\quad
  K^2\;\le\;\frac{12C\varepsilon_2\,t^4}{w_T}.
\]
For active triples ($w_T\ge\eta_T|\LE(P)|$) with richness
$t=\Theta(n)$ and $|\LE(P)|\ge e^{\Omega(n)}$ (standard width-3
bound), the right-hand side is $o(1)$; equivalently, for any fixed
$K_0$, the inequality $K\ge K_0$ forces \eqref{eq:R-mass} to exceed
the Step~2 budget whenever the triple is active with constant $\eta_T$.

If instead $K>ct/2$ (so $R$ does not fit), replace $R$ by the
half-strip $R':=\{0\le u<ct/4,\ -ct/4\le v<0\}$, which still lies in
the disagreement region since $u+v<ct/4<ct/2<K$ and $\mathbf{1}\{u\ge
0\}\ne\mathbf{1}\{v\ge 0\}$.  Its lattice mass is
$\Omega(t^2)\cdot w_T/t^2=\Omega(w_T)$, a fixed fraction of the
triple-overlap mass, which even more strongly exceeds the Step~2
budget.

\textbf{Step 5: Passage to $(1-o(1))$ of triples.}
Call a triple $T$ \emph{bad} if its gap $K(T)$ satisfies
$|K(T)|\ge K_0$ for some absolute constant $K_0$ (take e.g.\
$K_0:=C_1$ with $C_1$ chosen so the rectangle argument applies on
every active triple).  By the above, for each bad active triple the
symmetric-difference budget of~\eqref{eq:three-rules} is exceeded by
an amount $\ge c_0 w_T$ for some $c_0>0$, so
\[
  \sum_{\text{bad active }T}c_0 w_T
  \;\le\;
  \sum_{T}3\varepsilon_2\min_{\bullet}|\fiber{\bullet}(T)|
  \;\le\;3\varepsilon_2 M^{(3)},
\]
where $M^{(3)}:=\sum_{T\text{ active}}w_T$ is the total active-triple
weight (finite, $\ge c_T^{(3)}|\LE(P)|$ by
Lemma~\ref{lem:triple-visibility}).  Rearranging,
\[
  \sum_{\text{bad active }T}w_T
  \;\le\;\frac{3\varepsilon_2}{c_0}\,M^{(3)}
  \;=\;o(M^{(3)}),
\]
since $\varepsilon_2=o(1)$ is the Step~2 approximation parameter.
Hence the cocycle $\tau_{xz}=\tau_{xy}+\tau_{yz}+O(1)$ (with
$O(1)=K_0+|\Delta_0|\le K_0+O(1)$) holds on $(1-o(1))$ of active
triples, as required.
\end{proof}

\begin{remark}
Lemma~\ref{lem:cocycle} is a statement on the \emph{weighted} active
triples.  To deduce a single-vertex potential we will integrate on
the giant component of the triple graph, which in the rich case
inherits enough triple density from Step~5.
\end{remark}

%-----------------------------------------------------------------------
\section{Integrating the cocycle to a vertex potential}\label{sec:potential}
%-----------------------------------------------------------------------

\begin{lemma}[Lemma C: potential representation]\label{lem:potential}
\DEP{Lemmas~\ref{lem:sign-consistency}, \ref{lem:cocycle}.}
Assume Lemmas~A and~B on $\mathcal{E}^\star$ with a giant component of
the triple graph.
There is a function $a\colon P\to\mathbb{R}$ such that, for a
$(1-o(1))$-fraction of interfaces $e=(x,y)\in\mathcal{E}^\star$,
\[
  \widetilde\sigma_e \,\tau_e
  \;=\; a(y) - a(x) + O(1).
\]
\end{lemma}

\begin{proof}[Proof of Lemma~\ref{lem:potential}]
\textbf{Element graph.}
Let $G^P$ be the weighted multigraph on vertex set $P$ with an edge
$\{x,y\}$ for each interface $e=(x,y)\in\mathcal{E}^\star$, weighted
by the incidence weight of $e$.  By Lemma~\ref{lem:giant-component},
there is a connected subgraph
$G^P_\star\subseteq G^P$ on a vertex set $P^\star\subseteq P$ whose
edges carry $(1-o(1))$ of the rich-interface incidence mass.
Orient each edge $e=(x,y)$ from $x$ to $y$ and set the signed weight
\[
  \delta_e \;:=\; \widetilde\sigma_e\,\tau_e,
\]
which flips sign under edge reversal (since reorienting $e$ flips
$\widetilde\sigma_e$ and leaves $\tau_e$ unchanged).
With this orientation, Lemma~\ref{lem:cocycle} reads: for a
$(1-o(1))$-fraction of active triples on
$\mathcal{E}^\star$,
\begin{equation}\label{eq:tri-cocycle}
  \delta_{xy} + \delta_{yz} - \delta_{xz} \;=\; O(1).
\end{equation}
Call a triple \emph{good} if \eqref{eq:tri-cocycle} holds; the bad
triples carry $o(1)$ of the triple-weight by Lemma~\ref{lem:cocycle}.
Let $\mathcal{B}\subseteq E(G^P_\star)$ be the set of edges incident
to an $\omega(1)$-weight of bad triples; by Markov, $\mathcal{B}$
carries $o(1)$ of the edge weight, and we discard it.  Set
$G^P_{\star\star}:=G^P_\star\setminus\mathcal{B}$ and restrict to its
giant component, still carrying $(1-o(1))$ of incidence by
Lemma~\ref{lem:giant-component}.

\textbf{Spanning-tree integration.}
Fix a basepoint $z_0\in P^\star$ and a BFS spanning tree $T$ of
$G^P_{\star\star}$ rooted at $z_0$.  Define
\[
  a(z) \;:=\; \sum_{e\in\mathrm{path}_T(z_0\to z)}\delta_e
  \qquad (z\in P^\star),
\]
the signed sum of $\delta$ along the unique $T$-path from $z_0$ to
$z$, and $a(z):=0$ for $z\notin P^\star$ (affecting only $o(1)$ of
interface weight).
By construction,
$a(y)-a(x)=\delta_e$ \emph{exactly} on every tree edge $e=(x,y)$.

\textbf{Closing non-tree edges.}
For a non-tree edge $e=(x,y)\in G^P_{\star\star}$ let $\gamma_e$ be
the fundamental cycle formed by $e$ and the $T$-path from $y$ back
to $x$; write $\ell(e)$ for its length.  Then
\[
  \delta_e - (a(y)-a(x)) \;=\; \oint_{\gamma_e}\delta,
\]
the signed sum of $\delta$ around $\gamma_e$.
Star-triangulate $\gamma_e$ from $z_0$: since $T$ is a BFS tree
rooted at $z_0$, every vertex on $\gamma_e$ lies on a $T$-path to
$z_0$, so $\gamma_e$ decomposes into $\ell(e)-1$ triangles
$(v_i,v_{i+1},z_0)$ whose sides are either tree edges or the edge
$e$ itself.  For each such triangle that is \emph{good}
(satisfies~\eqref{eq:tri-cocycle}), the triangle contributes $O(1)$
to $\oint_{\gamma_e}\delta$; the total contribution along $\gamma_e$
telescopes, leaving
\begin{equation}\label{eq:cycle-bound}
  \Bigl|\oint_{\gamma_e}\delta\Bigr|
  \;\le\; O(1)\cdot\#\{\text{good triangles on }\gamma_e\}
  \;+\; O(\tau_{\max})\cdot\#\{\text{bad triangles on }\gamma_e\},
\end{equation}
where $\tau_{\max}$ bounds $|\tau|$ on $\mathcal{E}^\star$ (itself
$O(t(x,y))=O(n)$, but only the bad count matters below).

\textbf{Constant diameter and the BFS tree.}
Lemma~\ref{lem:giant-component} in fact supplies diameter $\le 3$ in
$G^P_{\star\star}$ on a
$(1-o(1))$-incidence subset $\widetilde P\subseteq P^\star$, which
we now make precise (this replaces an earlier ambiguous `typical
linear extension' phrasing).  Call $u\in P^\star$ \emph{generic} if
(i) $u$ is rich-row/rich-column in the sense of
Lemma~\ref{lem:giant-component} Step~3 (i.e.\ $|N_{\Rich}(u)|\ge
(c_T/2)|P|$ where $N_{\Rich}(u)$ is the rich-interface neighbourhood
in the opposite chain), and (ii) the $\mathcal{B}$-trim deletes only
an $o(1)$-fraction of $u$'s incident edge weight.  Step~5~(R)
gives~(i) on a $(1-O(c_T^{-1}))$-weighted subset of $P^\star$, and
Markov applied to the $o(1)$-weighted set $\mathcal{B}$ gives~(ii)
on a $(1-o(1))$-weighted subset; write $\widetilde P$ for the
intersection, with incidence mass $(1-o(1))M_0$.

For same-chain generic $u,u'\in\widetilde P\cap A$, inclusion--
exclusion (Lemma~\ref{lem:giant-component} Step~3 at the
\emph{element} level) gives $|N_{\Rich}(u)\cap N_{\Rich}(u')|\ge
(c_T-1)|B|$ common neighbours $w\in B$ with
$(u,w),(u',w)\in\Rich^\star$; by~(ii) at $u$ and $u'$, at least
$(1-o(1))$ of these also lie in $E(G^P_{\star\star})$, so a
length-$2$ path $u{-}w{-}u'$ survives in $G^P_{\star\star}$.  For
cross-chain generic $u\in\widetilde P\cap A$,
$v\in\widetilde P\cap B$, pick any rich-row $u'\in A$ with
$(u',v)\in E(G^P_{\star\star})$ ($|N_{\Rich}(v)|\ge(c_T/2)|A|$ and
the $\mathcal{B}$-fraction is $o(1)$, so such $u'$ exists) and
concatenate the length-$2$ path $u{-}w{-}u'$ with the edge
$(u',v)$ to get a length-$3$ path.  (The argument is symmetric and
covers all chain pairs.)

Consequently, BFS from a reference $z_0\in\widetilde P$ reaches
every $z\in\widetilde P$ at depth $\le 3$.  Hence every non-tree
edge $e=(x,y)$ with $x,y\in\widetilde P$ has fundamental-cycle
length
\[
  \ell(e)\;\le\;\mathrm{depth}_T(x)+\mathrm{depth}_T(y)+1\;\le\;7.
\]
The $o(1)$-weighted non-tree edges with an endpoint outside
$\widetilde P$ are discarded together with $\mathcal{B}$ (adding
only $o(1)$ to the bad-edge mass), yielding $\ell(e)=O(1)$
deterministically on a $(1-o(1))$-weighted subset of non-tree
edges---no expectation/Markov step on cycle length is required.

For such \emph{short-cycle} non-tree edges, \eqref{eq:cycle-bound}
combined with the fact that $\mathcal{B}$ has been removed (so almost
all triangles through them are good in the weighted sense) yields
\[
  \Bigl|\oint_{\gamma_e}\delta\Bigr| \;=\; O(1).
\]
Removing the remaining $o(1)$-weight of long-cycle non-tree edges
from $\mathcal{E}^\star$, we obtain a subset carrying $(1-o(1))$ of
rich-interface incidence on which
\[
  \widetilde\sigma_e\,\tau_e \;=\; a(y)-a(x) + O(1),
\]
as required.
\end{proof}

With $a$ in hand, define the global test function
\[
  H(L) := \sum_{z\in P} a(z)\,\pos_L(z).
\]
If $L'$ is obtained from $L$ by a BK swap of incomparable $(u,v)$
with $u$ before $v$ in $L$, then
\[
  H(L') - H(L) \;=\; a(v) - a(u).
\]
In particular, for rich interface $e=(x,y)$, a BK move across~$e$
changes $H$ by $a(y)-a(x)$, whose sign is $\widetilde\sigma_e$ by
Lemma~C.
So $H$ is a global scalar that, on every rich interface, tracks the
positive-direction sign predicted by the local staircase.

%-----------------------------------------------------------------------
\section{Synchronizing the fiber thresholds}\label{sec:global-c}
%-----------------------------------------------------------------------

Even with $H$ and the fiber thresholds $\tau_e\approx a(y)-a(x)$,
Step~7 needs a single global threshold $c$ such that, on most rich
fibers,
\[
  1_S(L)\;\approx\;1_{\{H(L)\ge c\}}.
\]

\begin{lemma}[Single global threshold]\label{lem:single-c}
\DEP{Step~6 (low conductance), Lemma~\ref{lem:potential}.}
There exists $c\in\mathbb{R}$ such that, for $(1-o(1))$ of
interfaces $e\in\mathcal{E}^\star$, on $\fiber{e}\setminus B_e$:
\[
  1_S(L)
  \;=\;
  1_{\{H(L)\ge c\}} + o(1).
\]
Equivalently, the per-fiber thresholds $c_e$ induced by
$\{H\ge c_e\}\cap\fiber{e}=S\cap\fiber{e}$ satisfy $c_e=c+O(1)$ on the
giant component.
\end{lemma}

\begin{proof}[Proof of Lemma~\ref{lem:single-c}]
Write $A:=\max_{z\in P}|a(z)|$, so that $H$ is $A$-Lipschitz on the
BK graph: $|H(L')-H(L)|\le A$ for every BK edge $(L,L')$.

\textbf{Per-fiber threshold.}
On $\fiber{e}^\star:=\fiber{e}\setminus B_e$,
Lemma~\ref{lem:signed-threshold} gives $S=\{\sigma_e(j-i)\ge\tau_e\}+
o(\mu_e)$, and Lemma~\ref{lem:potential} gives
$\widetilde\sigma_e\tau_e=a(y)-a(x)+O(1)$, whence the $(x,y)$
interface swap (which crosses the staircase in the
$\sigma_e$-direction) changes $H$ by $a(y)-a(x)$ with sign
$\sigma_e$.  Thus $H$ is monotone across the staircase boundary, and
\[
  c_e\;:=\;\tfrac12\Bigl(
    \inf_{L\in\fiber{e}^\star\cap S}H(L)+
    \sup_{L\in\fiber{e}^\star\setminus S}H(L)
  \Bigr)
\]
is well-defined up to an additive $O(A)=O(1)$ and satisfies
\begin{equation}\label{eq:perfib-mg4ddb}
  \mu\!\bigl((\{H\ge c_e\}\triangle S)\cap\fiber{e}^\star\bigr)
  \;=\;o(\mu_e).
\end{equation}

\textbf{Pairwise closeness $|c_e-c_f|=O(1)$ on active pairs.}
Let $(e,f)\in\mathcal{E}^\star\times\mathcal{E}^\star$ be active with
$\Omega_{ef}:=(\fiber{e}\cap\fiber{f})\setminus(B_e\cup B_f)$ and
$w_{ef}=\mu(\Omega_{ef})>0$; set $\Delta:=|c_e-c_f|$ and WLOG
$c_e\le c_f$.  Applying \eqref{eq:perfib-mg4ddb} to $e$ and $f$ on
$\Omega_{ef}$, the sets $\{H\ge c_e\}$ and $\{H\ge c_f\}$ each
coincide with $S$ up to $o(w_{ef})$, so
\begin{equation}\label{eq:strip-mg4ddb}
  \mu\!\bigl(\{c_e\le H<c_f\}\cap\Omega_{ef}\bigr)\;=\;o(w_{ef}).
\end{equation}

Foliate $[c_e,c_f]$ into $N:=\lfloor\Delta/(2A)\rfloor$ levels
$c_k:=c_e+(2k{+}1)A$ ($k=0,\dots,N{-}1$), and let
$\partial_k:=\{(L,L')\in E(G_{\BK})\cap\Omega_{ef}:
H(L)<c_k\le H(L')\}$.  By $A$-Lipschitzness both endpoints of an
edge in $\partial_k$ lie in $[c_k{-}A,c_k{+}A]$, so the $\partial_k$
are pairwise disjoint.  Since $c_k\ge c_e+A$,
\eqref{eq:strip-mg4ddb} gives
$\{H<c_k\}\cap\Omega_{ef}=S^c\cap\Omega_{ef}$ up to $o(w_{ef})$,
hence $\partial_k\subseteq\partial S$ modulo $o(w_{ef})$.

Restrict to the \emph{balanced} regime
$\mu(S\cap\Omega_{ef}),\,\mu(S^c\cap\Omega_{ef})\ge\epsilon w_{ef}$
(small universal $\epsilon$): ``unbalanced'' pairs, where $\Omega_{ef}$
lies almost entirely in $S$ or in $S^c$, automatically satisfy
$\Delta=O(A)$ since the $H$-range forcing $c_e,c_f$ is pinned to the
thin side's extremum, within $O(A)$ of a common value.  For balanced
pairs, the cut at $c_k$ splits $\Omega_{ef}$ into parts of mass
$\ge\tfrac12\epsilon w_{ef}$ each (up to $o(w_{ef})$), and the
Step~1 grid model makes $\Omega_{ef}$ BK-connected with edge
expansion $\Omega(1/t_{\max})$, $t_{\max}:=\max(t_e,t_f)$, giving by
Cheeger $|\partial_k|\ge\Omega(w_{ef}/t_{\max})$.  Summing:
\begin{equation}\label{eq:bdry-pair-mg4ddb}
  |\partial S\cap\Omega_{ef}|\;\ge\;
  \sum_{k=0}^{N-1}|\partial_k|-o(w_{ef})
  \;\gtrsim\;\frac{\Delta}{A}\cdot\frac{w_{ef}}{t_{\max}}.
\end{equation}

\emph{Low conductance forces $\Delta_{ef}=O(1)$ typically.}
Each BK edge lies in $O(1)$ fibers (Step~1 bounded multiplicity),
hence in $O(1)$ overlaps, so
$\sum_{(e,f)}|\partial S\cap\Omega_{ef}|\le O(1)|\partial S|\le
O(\Phi_0)\vol(S)=o(\vol(S))$ by Step~6.  Combining with
\eqref{eq:bdry-pair-mg4ddb} and $\sum_{(e,f)}w_{ef}/t_{\max}\gtrsim
\vol(\LE(P))/t_{\max}$ (Step~5 second-moment visibility),
Markov's inequality yields
\begin{equation}\label{eq:pair-bdd-mg4ddb}
  \Delta_{ef}\;\le\;K_1\;=\;O(1)\quad\text{on $(1-o(1))$-fraction of
  active-pair weight,}
\end{equation}
with $K_1$ depending on $A,\Phi_0$, and the richness constant.

\textbf{Globalization via diameter-3 on the giant component.}
Let $H^\star\subseteq G_{\Rich}$ be the giant component furnished by
Lemma~\ref{lem:giant-component}, carrying $(1-o(1))$ of the
incidence and active-pair weight.  Step~3 of that proof gives diameter
$\le 3$ in the endpoint-sharing (hence active) edge graph of
$H^\star$ on a $(1-O(c_T^{-1}))$-mass subset, with
$\Omega(c_T^2|A||B|)$ alternative length-$\le 3$ paths between any
two edges; in particular, for every $e,f\in H^\star$ there are
$\Omega(c_T^2|A||B|)$ length-$\le 3$ walks
$e=g_0,g_1,g_2,g_3=f$ through active edges $(g_k,g_{k+1})$.

Fix any reference $e_0\in H^\star$ with $c_{e_0}=:\bar c$.  Call
$e\in H^\star$ \emph{good} if there exists a length-$\le 3$ walk
$e_0\leftrightarrow e$ whose every hop $(g_k,g_{k+1})$ satisfies
$|c_{g_k}-c_{g_{k+1}}|\le K_1$.  By
\eqref{eq:pair-bdd-mg4ddb}, the exceptional-hop weight is $o(1)$
relative to total active-pair weight, and each of the
$\Omega(c_T^2|A||B|)$ candidate walks uses only $3$ hops, so by
union bound over hops and averaging over walks,
$(1-o(1))$-fraction of $e\in H^\star$ (weighted by incidence mass)
admits at least one good walk.  For every such good $e$, the
triangle inequality along the walk gives
\[
  |c_e-\bar c|\;\le\;\sum_{k=0}^{2}|c_{g_k}-c_{g_{k+1}}|
  \;\le\;3K_1\;=\;O(1).
\]
On the residual $o(1)$-weight of non-good $e$, the a priori bound
$|c_e|\le An$ contributes only to the exceptional mass, which is
absorbed into the $(1-o(1))$ quantifier.  Setting $c:=\bar c$, we
conclude
$\mathbf{1}_S=\mathbf{1}_{\{H\ge c\}}+o(1)$ on
$\fiber{e}\setminus B_e$ for $(1-o(1))$ of interfaces
$e\in\mathcal{E}^\star$, as claimed.
\end{proof}

%-----------------------------------------------------------------------
\section{From low conductance to layered width-3}\label{sec:layered}
%-----------------------------------------------------------------------

This section proves Prop.~\ref{prop:72}: given the potential
$a\colon P\to\mathbb{R}$ and threshold $c\in\mathbb{R}$ from
Prop.~\ref{prop:71}, if $\{H\ge c\}$ has low BK conductance and
controls $(1-o(1))$ of rich swaps, then $P$ admits a layered
width-$3$ decomposition of interaction width $w=O_T(1)$ in the sense
of Def.~5.1 of \texttt{step8.tex}.
Throughout, write $\Delta_{xy}:=a(y)-a(x)$ and
$p_{xy}:=\mu\{L\in\LE(P):(x,y)\text{ BK-adjacent in }L\}/\mu(\LE(P))$.

\subsection{BK boundary bounds integrated $a$-variation}

\begin{lemma}[Boundary budget]\label{lem:budget-var}
\DEP{Prop.~\ref{prop:71}, Lemma~\ref{lem:single-c}.}
If $\Phi(\{H\ge c\})\le\eta=o(1)$ and $\{H\ge c\}$ agrees with $S$ on
a $(1-o(1))$ fraction of $\LE(P)$, then
\begin{equation}\label{eq:var-budget}
  \sum_{x\parallel_P y,\ \Delta_{xy}>0}\Delta_{xy}\cdot p_{xy}
  \;=\; o(1).
\end{equation}
\end{lemma}

\begin{proof}
A BK edge on $\bd\{H\ge c\}$ is a triple $(L,L',\{x,y\})$ where $L'$
differs from $L$ by swapping a BK-adjacent incomparable pair
$\{x,y\}$, $L\notin\{H\ge c\}$, $L'\in\{H\ge c\}$ (or the reverse).
Orient so $x<_Ly$ with $\Delta_{xy}>0$: the swap lifts $H$ by
$\Delta_{xy}$, so the triple is counted iff
$H(L)\in[c-\Delta_{xy},c)$.  Hence
\[
  |\bd\{H\ge c\}|
  \;=\;\!\!\!\sum_{x\parallel y,\ \Delta_{xy}>0}\!\!\!
  \mu\bigl\{L:(x,y)\text{ BK-adj},\,x<_Ly,\,H(L)\in[c-\Delta_{xy},c)\bigr\}
  \;+\;\text{sym}.
\]
Each summand is nonnegative and, on the $(1-o(1))$-mass giant
component identified by Lemma~\ref{lem:single-c}, the induced
conditional distribution of $H(L)$ near $c$ has density bounded below
by an absolute constant $\rho>0$: the cut $\{H\ge c\}$ has BK
boundary exactly when $H$ is within one $\Delta$-unit of $c$, and the
Step~6 low-conductance regime localizes the global $c$ to an
$O(1)$-window of its median, so $H$ has $\Theta(1)$ density on
$[c-O(1),c+O(1)]$ relative to $\mu$.  Consequently each summand is
$\ge\rho\,\Delta_{xy}\,p_{xy}\cdot(1-o(1))$, and the hypothesis
$|\bd\{H\ge c\}|\le\eta\,\vol(\{H\ge c\})\le\eta\,\mu(\LE(P))$ gives
$\sum\Delta_{xy}p_{xy}=O(\eta/\rho)=o(1)$.
\end{proof}

Under Prop.~\ref{prop:72}'s hypotheses, $S$ and $\{H\ge c\}$ agree
up to $o(1)$ by Prop.~\ref{prop:71}, and BK-degree bounds transfer
low conductance from $S$ to $\{H\ge c\}$ with $o(1)$ slack; so
\eqref{eq:var-budget} holds.

\subsection{Rich interfaces concentrate on monotone synchronizations}

Fix a Dilworth decomposition $P=A\sqcup B\sqcup C$ into three chains
(width~$3$ guarantees this).  Enumerate each chain in chain order
$A=(a_1<_Pa_2<_P\cdots)$, $B=(b_1<_Pb_2<_P\cdots)$,
$C=(c_1<_Pc_2<_P\cdots)$.  By Prop.~\ref{prop:71}, every rich cover
$x<_Py$ satisfies $a(y)>a(x)$; after discarding an $o(1)$ mass of
exceptional elements where $a$ fails chain-monotonicity (absorbed
into the final $(1-o(1))$ quantifier), $a$ is strictly increasing
along each of $A,B,C$.  Define the synchronization maps
\[
  f_{AB}(i):=\arg\min_{j}|a(a_i)-a(b_j)|,\quad
  f_{AC}(i):=\arg\min_{k}|a(a_i)-a(c_k)|,
\]
and $f_{BC}$ analogously; strict monotonicity of $a$ on each chain
makes each $f$ weakly monotone.

\begin{lemma}[Bandwidth of rich incomparability]\label{lem:bandwidth}
\DEP{\eqref{eq:var-budget}, richness threshold $T$.}
There is $K=K(T)=O_T(1)$ such that, after deleting $o(1)$ of the
rich-interface incidence mass,
\[
  (a_i,b_j)\in\Rich^\star\text{ with }a_i\parallel_P b_j
  \;\Longrightarrow\;
  |j-f_{AB}(i)|\le K,
\]
and analogously for $(A,C)$ and $(B,C)$.
\end{lemma}

\begin{proof}
For any rich pair $e=(x,y)=(a_i,b_j)$, Step~1/Step~5 place $x,y$ at
adjacent positions in the grid coordinate $\pi_e$ of
$\fiber{e}$, and $\mu(\fiber{e})\ge c_T\mu(\LE(P))$ by richness
(visibility, Step~5 Cor.~(R)); hence
$p_{xy}\ge c'_T>0$ for an absolute constant depending only on
the richness threshold $T$.  Plugging into~\eqref{eq:var-budget},
\[
  \sum_{\text{rich }(a_i,b_j)}\Delta_{a_ib_j}\;\le\;(c'_T)^{-1}\cdot o(1)\;=\;o(1),
\]
so all but an $o(1)$ fraction of rich incidence mass satisfies
$|a(b_j)-a(a_i)|<c_0$ for any fixed $c_0>0$.

For each such pair, strict monotonicity of $a$ on $B$ forces $j$ to
lie in the preimage interval
$J_i(c_0):=\{j:|a(b_j)-a(a_i)|<c_0\}$, an interval around $f_{AB}(i)$
of length
$N_B(c_0):=\max_i|J_i(c_0)|$.  A second application
of~\eqref{eq:var-budget} bounds $N_B(c_0)$: if $N_B(c_0)\ge M$ at
some $i$, then every pair $(a_i,b_j)$ with $j\in J_i(c_0)$ is
incomparable to $a_i$ with $\Delta_{a_ib_j}<c_0$, and by richness each
contributes $p\ge c'_T$, giving $\sum\Delta p\ge c'_T\cdot M\cdot$
(average $\Delta$).  Averaged over $i$ against rich-pair incidence
mass, this forces $M=O_T(1/c_0)$ save on an $o(1)$ set of $i$'s.
Taking $K:=N_B(c_0)=O_T(1)$ with $c_0$ any fixed constant gives the
claim for $(A,B)$; the argument for $(A,C),(B,C)$ is identical.
\end{proof}

\subsection{Assembly}

\begin{proof}[Proof of Prop.~\ref{prop:72}]
With the chain enumerations and synchronizations above, define levels
\[
  L_k \;:=\; \{\,a_k,\ b_{f_{AB}(k)},\ c_{f_{AC}(k)}\,\}\quad (1\le k\le|A|),
\]
with the convention that missing chain-entries (when $B$ or $C$ has
run out of elements) are omitted.  Extend with additional levels
$L_{|A|+\ell}$ holding the tail of $B\cup C$ indexed by the same
construction applied to the remaining chains.  Each $|L_k|\le 3$,
giving the width-$3$ size bound~(L1).

\textbf{Interaction width.}  Two elements $z\in L_i$, $z'\in L_j$
belonging to distinct chains have chain-positions whose offset to the
synchronization is $O(1)$; hence their chain-position difference is
at most $|i-j|+O(1)$.  Conversely, any incomparable rich pair has
chain-position difference $\le K$ by Lemma~\ref{lem:bandwidth}, so
$|i-j|\le K+O(1)=:w=O_T(1)$.  Outside this interaction window,
elements are forced comparable on $(1-o(1))$ of the rich incidence
mass.  For non-rich incomparable pairs: by~\eqref{eq:var-budget} and
Markov, all but an $o(1)$ fraction of their mass has $\Delta<c_0$, so
the same chain-position argument places them within
$N_\bullet(c_0)+O(1)=O_T(1)$ levels.  Absorbing all exceptional mass
into the $(1-o(1))$ quantifier establishes (L2)--(L3) of
Def.~5.1 of \texttt{step8.tex}.

\textbf{Rich interfaces within the window.}  By
Lemma~\ref{lem:bandwidth}, every rich $(a_i,b_j)$ has
$|j-f_{AB}(i)|\le K$, so $a_i\in L_i$ and $b_j\in L_j$ with
$|i-j|\le K\le w$; analogously for $(A,C)$, $(B,C)$.

\textbf{BK moves along level order.}  A BK swap exchanges a
BK-adjacent incomparable pair $(u,v)$, changing $H$ by $\Delta_{uv}$.
By~\eqref{eq:var-budget}, all but $o(1)$ of BK swap mass has
$|\Delta_{uv}|=o(1)$, hence by the bandwidth argument lies within
chain-position difference $\le w$, i.e., adjacent or near-adjacent
levels.  Equivalently, BK moves shuffle elements within the
$w$-wide interaction window; the layered order is respected on
$(1-o(1))$ of the mass.

\textbf{Interfaces confined.}  The above says rich interfaces
appear between levels at distance $\le w$, BK moves preserve the
level order up to the interaction window, and linear extensions are
thereby determined by choices internal to consecutive $w$-level
blocks.  This matches all three bullets of Prop.~\ref{prop:step7-informal}(ii)
and Def.~5.1 of \texttt{step8.tex} with interaction width $w=O_T(1)$
and the $O_T(1)$-size exceptional set of
Remark~5.2 there.
\end{proof}

%-----------------------------------------------------------------------
\section{Formal propositions}\label{sec:formal}
%-----------------------------------------------------------------------

Combining the above pieces gives three formal propositions.

\begin{proposition}[Prop.~7.1 --- coherence globalizes]\label{prop:71}
Assume Steps~1, 2, 5, 6 and the coherence case of Step~6.
Then there exist $\mathcal{E}^\star\subseteq\Rich^\star$ of $(1-o(1))$
of rich-interface incidence mass, a potential $a\colon P\to\mathbb{R}$,
and a threshold $c\in\mathbb{R}$, such that for every
$e=(x,y)\in\mathcal{E}^\star$,
\[
  1_S(L) \;=\; 1_{\{H(L)\ge c\}} + o(1)
  \quad\text{on $(1-o(1))$ of $\fiber{e}\setminus B_e$,}
\]
where $H(L)=\sum_z a(z)\pos_L(z)$, and the positive BK direction
across $e$ is $\sgn(a(y)-a(x))$.
\end{proposition}

\begin{proof}[Proof of Prop.~\ref{prop:71}]
Apply the gap lemmas in order.
\textbf{(G1)} Lemma~\ref{lem:signed-threshold} normalizes each
$S$-good rich interface $e\in\Rich^\star$ into the signed-threshold
form $\sigma_e(j-i)\ge\tau_e$ with $O(\varepsilon_2)$ symmetric
difference per fiber.
\textbf{(G2)} Lemma~\ref{lem:sign-consistency}, fed by Step~6
Corollary~B, replaces $\sigma_e$ with a globally consistent sign
$\widetilde\sigma_e$ on a subset
$\mathcal{E}^\star\subseteq\Rich^\star$ of $(1-o(1))$ rich-interface
incidence, with overlap-compatible signs on every active pair in
$\mathcal{E}^\star$.
\textbf{(G3)} Lemma~\ref{lem:triple-visibility} guarantees that
$(1-o(1))$ of the edge weight on $\mathcal{E}^\star$ extends to
active triples, so the cocycle hypothesis of (G4) is non-vacuous.
\textbf{(G4)} Lemma~\ref{lem:cocycle} promotes pairwise sign
compatibility to the additive cocycle
$\tau_{xz}=\tau_{xy}+\tau_{yz}+O(1)$ on $(1-o(1))$ of active triples
in $\mathcal{E}^\star$.
\textbf{(G9)} The visible interface graph carries a giant component
of weight $(1-o(1))$, transferred to the element graph $G^P_\star$
on $P^\star\subseteq P$.
\textbf{(G5)} Lemma~\ref{lem:potential} integrates the signed
thresholds $\delta_e=\widetilde\sigma_e\tau_e$ along a BFS spanning
tree of $G^P_\star$, closing fundamental cycles via star-triangulation
through the basepoint and (G4)-good triangles, yielding
$a\colon P\to\mathbb{R}$ with
$\widetilde\sigma_e\tau_e=a(y)-a(x)+O(1)$ on $(1-o(1))$ of
$\mathcal{E}^\star$.  Define $H(L):=\sum_z a(z)\pos_L(z)$;
across any rich BK swap on $e=(x,y)$, $H$ changes by $a(y)-a(x)$,
so the positive BK direction agrees with $\widetilde\sigma_e$.
\textbf{(G6)} Lemma~\ref{lem:single-c} synchronizes the per-fiber
thresholds $c_e$ induced by $\{H\ge c_e\}\cap\fiber{e}$: pairwise
overlap discrepancy $|c_e-c_f|$ contributes
$\Theta(|c_e-c_f|)w_{ef}$ to BK boundary, so Step~6's
low-conductance hypothesis forces $|c_e-c_f|=O(1)$ on every active
pair, and connectedness on $G^P_\star$ (G9) collapses
$\{c_e:e\in\mathcal{E}^\star\}$ into a single $O(1)$ window;
take $c$ to be its mean.
Combining,
$1_S(L)=1_{\{H(L)\ge c\}}+o(1)$ on
$(1-o(1))$ of $\fiber{e}\setminus B_e$ for every
$e\in\mathcal{E}^\star$, with $\sgn(a(y)-a(x))=\widetilde\sigma_e$
the positive BK direction.
\end{proof}

\begin{proposition}[Prop.~7.2 --- low-boundary threshold cut is
1D]\label{prop:72}
\DEP{Prop.~\ref{prop:71}.}
If a cut of the form $\{H\ge c\}$ has low BK conductance and controls
$(1-o(1))$ of all rich swaps, then $P$ admits a layered width-3
decomposition on $(1-o(1))$ of its mass.
\end{proposition}

\begin{proof}[Proof of Prop.~\ref{prop:72}]
By Prop.~\ref{prop:71}, the cut $S$ agrees with $\{H\ge c\}$ up to
$o(1)$ mass, so $|\bd S|=|\bd\{H\ge c\}|+o(\vol(S))$ inherits the
low conductance hypothesis.  By Section~\ref{sec:layered}, the BK
boundary of $\{H\ge c\}$ counts adjacent pairs $(u,v)$ in a uniformly
random $L$ with $a(u),a(v)$ straddling $c$, which under low conductance
is $o(\vol(S))$.  The width-3 structural lemma converts this small
variation of $a$ into a layered antichain decomposition
$P=L_1\sqcup\cdots\sqcup L_m$ on which $a$ is approximately monotone,
giving the claimed layered width-3 decomposition on $(1-o(1))$ of
$\vol(\LE(P))$.
\end{proof}

\begin{proposition}[Prop.~7.3 --- layered width-3 is
harmless]\label{prop:73}\label{sec:prop73-proof}
\DEP{Prop.~\ref{prop:72}.}
A layered decomposition of a width-3 poset as in
Prop.~\ref{prop:72} cannot occur in a width-3 counterexample to
$1/3$--$2/3$: either such a decomposition directly produces a
balanced pair (averaging on adjacent layers) or reduces to a strictly
smaller width-3 instance, contradicting minimality.
\end{proposition}

\begin{proof}
Let $P$ be a minimal width-3 $\gamma$-counterexample
(every incomparable pair $(u,v)$ satisfies
$\Pr[u<v]\notin[\tfrac13+\gamma,\tfrac23-\gamma]$ for some fixed
$\gamma>0$).  By Prop.~\ref{prop:72}, $P$ admits a layered
decomposition $P=L_1\sqcup\cdots\sqcup L_m$ into antichain levels of
size $\le 3$ that is strict on $(1-o(1))$-mass; i.e.\ on
$(1-\delta)|\LE(P)|$ linear extensions (call these \emph{aligned}
extensions), the relative order of any $u\in L_r$ and $v\in L_s$ with
$r<s$ is $u<v$.  Here $\delta=o(1)$ as the Step~7 parameters shrink;
in particular we may assume $\delta<\gamma/2$.

\smallskip
\noindent
\emph{Step 1: width-3 forces a full layer.}
Since $\mathrm{width}(P)=3$, the antichain cover
$P=\bigsqcup L_r$ must use at least one level of size exactly $3$;
otherwise $P$ itself would be a union of antichains of size $\le 2$
with the layered order, hence of width $\le 2$.  Fix such a level
$L_{r^\star}$ of size $3$.

\smallskip
\noindent
\emph{Case $m=1$.}
Then $P=L_1$ is an antichain on $3$ elements, its $\LE(P)$ is
the full symmetric group $S_3$ with the uniform measure, and every
incomparable pair has $\Pr[u<v]=1/2\in[1/3,2/3]$.  This contradicts
$P$ being a $\gamma$-counterexample.

\smallskip
\noindent
\emph{Case $m\ge 2$ (minimal-counterexample reduction).}
Set $P':=P\setminus L_m=L_1\sqcup\cdots\sqcup L_{m-1}$.
Then $|P'|<|P|$, and $P'$ inherits the induced partial order with
the same layered decomposition of depth $m-1$; in particular
$\mathrm{width}(P')\le 3$.  By minimality of $P$, the induced
poset $P'$ is not a $\gamma$-counterexample, so either
(a)~$\mathrm{width}(P')\le 2$, in which case $P'$ has a balanced pair
by Linial's width-$2$ theorem~\cite{Linial1984}, or (b)~$\mathrm{width}(P')=3$ and by
minimality some incomparable pair $(u,v)$ in $P'$ satisfies
$\Pr_{P'}[u<v]\in[\tfrac13+\gamma,\tfrac23-\gamma]$.  Either way fix
incomparable $(u,v)\in P'$ with
$p':=\Pr_{P'}[u<v]\in[\tfrac13+\gamma,\tfrac23-\gamma]$.

We compare $\Pr_P[u<v]$ to $p'$.
On the $(1-\delta)$-mass of aligned extensions, all elements of $L_m$
appear after all elements of $P'$, so restricting an aligned
$L\in\LE(P)$ to $P'$ yields an $L'\in\LE(P')$; conversely every
$L'\in\LE(P')$ extends to aligned $L\in\LE(P)$ by appending any
linear extension of $L_m$.  Hence the aligned-conditional
distribution of $L|_{P'}$ is exactly the uniform distribution on
$\LE(P')$, and the events $\{u<v\}$ and alignment depend on this
restriction only.  Therefore
\[
  \bigl|\Pr_P[u<v]-p'\bigr|
  \;\le\; \delta
  \;<\; \gamma/2,
\]
so $\Pr_P[u<v]\in[\tfrac13+\tfrac\gamma2,\tfrac23-\tfrac\gamma2]
\subseteq[1/3,2/3]$.  The pair $(u,v)$ remains incomparable in $P$
(incomparability is inherited), so $(u,v)$ is a balanced pair in $P$
violating $P$'s counterexample property.  Contradiction.

\smallskip
\noindent
(If $|L_1|=3$ instead of $|L_m|=3$ we remove $L_1$ instead; the
symmetric argument applies, since in aligned extensions $L_1$
occupies the initial positions.)
\end{proof}

\begin{theorem}[Step~7 assembly]\label{thm:step7}
Assume Steps~1, 2, 5, 6 and the coherence case of Step~6.
Then the cut $S\subseteq\LE(P)$ admits a global threshold-of-potential
description, and consequently $P$ has a balanced pair or is reducible.
Hence the coherent case cannot support a minimal width-3
$\gamma$-counterexample.
\end{theorem}

\begin{proof}
Lemmas~\ref{lem:signed-threshold}, \ref{lem:sign-consistency},
\ref{lem:triple-visibility}, \ref{lem:cocycle},
\ref{lem:giant-component}, \ref{lem:potential}, and
\ref{lem:single-c} combine to give Prop.~\ref{prop:71}: the global
threshold-of-potential description $1_S\approx 1_{\{H\ge c\}}$ holds
on $(1-o(1))$ of $\LE(P)$, with $H(L)=\sum_z a(z)\pos_L(z)$.  Since
$S$ was assumed to be a low-conductance cut, the threshold cut
$\{H\ge c\}$ inherits low BK conductance; Prop.~\ref{prop:72} then
forces $P$ into a layered width-3 regime on $(1-o(1))$ of mass.
Prop.~\ref{prop:73} closes the loop: the layered structure yields
either a balanced pair directly or a strictly smaller width-3
sub-instance, contradicting the assumed minimality of the
$\gamma$-counterexample.  Therefore no minimal width-3
$\gamma$-counterexample lies in the coherence case of Step~6.
\end{proof}

%-----------------------------------------------------------------------
\section{Proof in compressed form}
%-----------------------------------------------------------------------

We record the logical skeleton.

\begin{enumerate}[nosep]
  \item By Step~2 + Lemma~\ref{lem:signed-threshold}, each rich
    interface $e$ has signed-threshold form $\sigma_e(j-i)\ge\tau_e$.
  \item By Step~6, the signs are pairwise compatible on almost all
    overlap weight.
  \item Lemma~\ref{lem:sign-consistency} promotes pairwise to global
    sign consistency on a giant component.
  \item On active triples, Lemma~\ref{lem:cocycle} upgrades this to
    the additive cocycle
    $\tau_{xz}\approx \tau_{xy}+\tau_{yz}$.
  \item Lemma~\ref{lem:potential} integrates the cocycle into an
    element potential $a:P\to\mathbb{R}$.
  \item Lemma~\ref{lem:single-c} synchronizes the fiber thresholds
    into a single $c$.
  \item Prop.~\ref{prop:71} gives the global threshold description.
  \item Prop.~\ref{prop:72} shows low BK conductance of such a cut
    forces a layered width-3 structure.
  \item Prop.~\ref{prop:73} finishes via balanced-pair or minimal
    reduction.
\end{enumerate}

%-----------------------------------------------------------------------
\section{Technical lemmas on triples and connectivity}
%-----------------------------------------------------------------------

Two technical inputs used above are proved here: the active-triple
visibility lemma that feeds the cocycle argument, and the
giant-component connectedness of $G_{\Rich}$ that integrates the
cocycle.

\subsection{Active-triple visibility}\label{sec:gap-triples}

A $(1-o(1))$ fraction of the interface-graph edge weight
$M:=\sum_{\alpha\ne\beta}w_{\alpha\beta}$ participates in active
triples, i.e.\ for $(1-o(1))$ of the weight
$w_{\alpha\beta}$ there exists $\gamma\in\Rich^\star$ with
$w_{\alpha\gamma},w_{\beta\gamma}$ and the triple overlap
$w_{\alpha\beta\gamma}:=|(\fiber{\alpha}\cap\fiber{\beta}\cap\fiber{\gamma})
\setminus(B_\alpha\cup B_\beta\cup B_\gamma)|$ simultaneously above
fixed positive thresholds.  This is what Lemmas~\ref{lem:cocycle} and
\ref{lem:potential} consume: without triple-density, the
cocycle-to-gradient integration is over an empty graph.

\begin{lemma}[Active-triple visibility]\label{lem:triple-visibility}
Assume the rich case~(R) of Step~5 with richness constant~$c_T$,
together with the bad-set bound $|B_\alpha|\le Ct_\alpha$ of
Step~1.  Let
\[
  I(L) \;:=\; \#\{\alpha\in\Rich^\star:\ L\in\fiber{\alpha}\setminus B_\alpha\}
\]
be the visibility count.  Then:
\begin{enumerate}[label=(\alph*),nosep]
  \item (Third moment.)  $\displaystyle \sum_{L\in\LE(P)}I(L)^3
     \;\ge\; c_T^{(3)}\,|\LE(P)|$ with $c_T^{(3)}:=(c_T/2)^3$.
  \item (Triple-overlap mass.)  Writing
    $w_{\alpha\beta\gamma}:=|(\fiber{\alpha}\cap\fiber{\beta}\cap
    \fiber{\gamma})\setminus(B_\alpha\cup B_\beta\cup B_\gamma)|$,
    the ordered triple mass satisfies
    \[
      \sum_{\alpha,\beta,\gamma\in\Rich^\star}w_{\alpha\beta\gamma}
      \;=\; \sum_{L}I(L)^3
      \;\ge\; c_T^{(3)}|\LE(P)|.
    \]
  \item (Edge fraction in triples.)  Setting
    $M:=\sum_{\alpha\ne\beta}w_{\alpha\beta}=\sum_L I(L)(I(L){-}1)$,
    the fraction of the edge weight $(L,\alpha,\beta)$ with
    $L\in\fiber{\alpha}^\star\cap\fiber{\beta}^\star$ and
    $I(L)\ge 3$ (i.e.\ $\exists\gamma\notin\{\alpha,\beta\}$ with
    $L\in\fiber{\gamma}^\star$) is at least
    $1-O(1/c_T')$, with $c_T':=(c_T/2)^2$ the Step~5 second-moment
    constant.  As $c_T\to\infty$ with the richness threshold~$T$,
    this fraction is $1-o(1)$.
\end{enumerate}
\end{lemma}

\begin{proof}
(a) By Step~5 Corollary~\ref{cor:second-moment}, the first-moment bound
$\mathbb{E}_L[I(L)]\ge c_T/2$ holds.  Since $x\mapsto x^3$ is convex
on $[0,\infty)$ and $I(L)\ge 0$, Jensen's inequality gives
\[
  \mathbb{E}_L[I(L)^3] \;\ge\; \bigl(\mathbb{E}_L[I(L)]\bigr)^3
  \;\ge\; (c_T/2)^3,
\]
equivalently $\sum_L I(L)^3\ge c_T^{(3)}|\LE(P)|$.

(b) Expanding the cube pointwise:
\[
  I(L)^3
  \;=\; \sum_{\alpha,\beta,\gamma\in\Rich^\star}
  \mathbf{1}\bigl[\,L\in\fiber{\alpha}^\star\cap\fiber{\beta}^\star
  \cap\fiber{\gamma}^\star\,\bigr],
\]
with $\fiber{\cdot}^\star:=\fiber{\cdot}\setminus B_{\cdot}$.
Summing over~$L$ and interchanging the order of summation gives
$\sum_L I(L)^3 = \sum_{\alpha,\beta,\gamma}w_{\alpha\beta\gamma}$,
under the same counting-measure identification used in Steps~5 and~6.
Combined with~(a), this yields the stated lower bound.

(c) A pair $(L,\alpha,\beta)$ with $L\in\fiber{\alpha}^\star\cap
\fiber{\beta}^\star$ (contributing $1$ to $w_{\alpha\beta}$ with
$\alpha\ne\beta$) extends to a triple
$(L,\alpha,\beta,\gamma)$ with $\gamma$ distinct from $\alpha,\beta$
iff $I(L)\ge 3$.  Hence the edge weight that \emph{fails} to extend
is
\[
  \sum_{L:I(L)\le 2}I(L)(I(L)-1)
  \;\le\; 2\,|\LE(P)|,
\]
since $I(I-1)\le 2$ for $I\in\{0,1,2\}$.  For a lower bound on $M$,
note that $I(I-1)\ge I^2/2$ whenever $I\ge 2$, so by Step~5
Corollary~\ref{cor:second-moment} ($\sum_L I(L)^2\ge c_T'|\LE(P)|$),
\[
  M \;=\; \sum_L I(L)(I(L)-1)
  \;\ge\; \tfrac12\!\!\sum_{L:I(L)\ge 2}\!\! I(L)^2
  \;\ge\; \tfrac12\bigl(\textstyle\sum_L I(L)^2 - |\LE(P)|\bigr)
  \;\ge\; \tfrac14\,c_T'\,|\LE(P)|
\]
once $c_T'\ge 2$ (the middle step uses
$\sum_{I\le 1}I^2\le|\LE(P)|$).  Dividing the failure weight
$2|\LE(P)|$ by $M\ge c_T'|\LE(P)|/4$ yields failure fraction
$\le 8/c_T'=O(1/c_T')$.  Taking $T$ large enough that $c_T'\to\infty$
(which is consistent with Step~5's asymptotic richness) makes the
failure fraction $o(1)$.

\textbf{Active-triple threshold.}  It remains to pass from
``some $\gamma$ with $w_{\alpha\beta\gamma}>0$'' to ``some $\gamma$
with $w_{\alpha\gamma},w_{\beta\gamma},w_{\alpha\beta\gamma}$ above
fixed positive thresholds.''
Fix a threshold parameter $\eta$ to be determined.
Call an ordered triple $(\alpha,\beta,\gamma)$ \emph{below-threshold}
if at least one of $w_{\alpha\gamma},w_{\beta\gamma},
w_{\alpha\beta\gamma}$ is below $\eta\cdot c_T^{(3)}|\LE(P)|/|\Rich^\star|^3$.
By Markov, the total mass of below-threshold triples is at most
$3\eta\cdot c_T^{(3)}|\LE(P)|$: there are $\le|\Rich^\star|^3$ ordered
triples, each below-threshold contributes $\le\eta\cdot
c_T^{(3)}|\LE(P)|/|\Rich^\star|^3$ via the triple-mass bound, so the
total is $\le\eta\cdot c_T^{(3)}|\LE(P)|$, times the factor $3$ for
each of the three pair/triple conditions that could fail.
Choosing $\eta$ a small absolute constant (e.g.\ $\eta=1/6$), the
above-threshold triple mass is at least $\tfrac12 c_T^{(3)}|\LE(P)|$.

Projecting to edges: each above-threshold triple supplies
$3$ ordered pairs $(\alpha,\beta)$ whose weight extends to an active
triple.  Hence the edge weight in active triples is at least a
constant fraction of $M$, and the complementary ``not in any active
triple'' edge weight is at most the sum of (i) the $2|\LE(P)|$
deficit from $I(L)\le 2$ states and (ii) the $3\eta\cdot
c_T^{(3)}|\LE(P)|$ threshold loss, both $o(M)$ as $c_T\to\infty$ with
$\eta$ fixed.  This establishes the $(1-o(1))$ claim.
\end{proof}

\textbf{Triangle-structure remark.}  In the width-3 setting the
downstream cocycle lemma (Lemma~\ref{lem:cocycle},
Section~\ref{sec:cocycle}) is stated on \emph{triangle triples}
$(e_{xy},e_{yz},e_{xz})$ through three distinct elements
$x,y,z\in P$ (which must then span the three chains, since rich
interfaces connect distinct chains).  Lemma~\ref{lem:triple-visibility}
above is stated without the triangle restriction, matching the
definition of active triple in Section~\ref{sec:normalize} and
Definition~\ref{def:local-staircase}'s surrounding setup (positive
pairwise and triple overlap).  The passage from arbitrary active
triples to triangle triples is handled at the Step~1 level (width-3
normal form on triple overlaps): in width~$3$, three rich interfaces
with
simultaneous fiber visibility must share common elements on
overlap, because $\fiber{\alpha}\cap\fiber{\beta}$ is supported on a
common chain window by the Step~1 commuting-square model, and three
commuting squares through a common window close up into a
commuting cube whose edges are the three triangle interfaces.
Hence arbitrary-triple visibility coincides with
triangle-triple visibility up to $o(1)$ losses absorbed in the
Step~1 bad-set bound.

\subsection{Giant-component connectedness of $G_{\Rich}$}\label{sec:gap-connected}

This is the connectedness input consumed by
Lemma~\ref{lem:sign-consistency} (disagreement percolation),
Lemma~\ref{lem:potential} (gradient integration), and
Lemma~\ref{lem:single-c} (threshold synchronization).

\begin{lemma}[Giant component of $G_{\Rich}$]\label{lem:giant-component}
Assume the rich case~(R) of Step~5 with richness constant $c_T$, and
let the active-pair threshold be
$\theta_{\rm pair}:=c_1\,T\,|\LE(P)|/(|A|\,|B|\,|C|)$ for a fixed
small constant $c_1>0$ (matching the share-endpoint overlap
constant of Step~1 below).  Then there is a connected component
$\mathcal{C}\subseteq\Rich^\star$ of $G_{\Rich}$ with incidence weight
\[
  \sum_{\alpha\in\mathcal{C}}\mu_\alpha
  \;\ge\; \bigl(1-O(c_T^{-1})\bigr)\,M_0,
  \qquad M_0:=\sum_{\alpha\in\Rich^\star}\mu_\alpha.
\]
The same component supports $(1-o(1))$ of the active-triple incidence
mass.
\end{lemma}

\begin{proof}
Let $I(L):=\#\{\alpha\in\Rich^\star: L\in\fiber{\alpha}\setminus
B_\alpha\}$ and $\mu_\alpha:=\mu(\fiber{\alpha}\setminus B_\alpha)$.
Without loss of generality the rich case~(R) of Step~5 holds for the
ordered triple $(A,B\mid C)$, so $|\Rich_{AB}|\ge c_T|A||B|$ and the
total $AB$ rich-incidence mass is $M_0\ge c_T|\LE(P)|$
(Theorem~\ref{thm:step5}); rich pairs of types $AC,BC$ either also
satisfy~(R) (in which case the same argument applies blockwise) or
fall into the banded case~(C) of Step~5 and contribute $o(M_0)$ to
the rich-interface incidence after the $S$-good trim of
Definition~\ref{def:G_Rich}.  We may therefore restrict attention to
$\Rich^\star\subseteq\Rich_{AB}$ throughout; the argument carries
over verbatim to additional rich types when they are present.

\smallskip
\textbf{Step~1: Edge-mass lower bound from second moment.}
By Step~5 Corollary~\ref{cor:second-moment},
$\sum_L I(L)^2\ge c_T'|\LE(P)|$ with $c_T':=(c_T/2)^2$.  Therefore
\begin{equation}\label{eq:gc-edge-mass}
  M\;:=\;\!\!\sum_{\alpha\ne\beta}\!\!w_{\alpha\beta}
  \;=\;\sum_L I(L)(I(L)-1)
  \;\ge\;c_T'|\LE(P)|-M_0
  \;\ge\;\tfrac12 c_T'|\LE(P)|
\end{equation}
once $c_T$ is large enough that $c_T'\gg c_T$ (recall
$c_T'\sim c_T^2/4$).

\smallskip
\textbf{Step~2: Endpoint-sharing pairs are active.}
Let $\alpha=(x,y),\beta=(x,y')\in\Rich^\star$ share endpoint
$x\in A$ (so $y,y'\in B$).  By Step~1
Corollary~\ref{cor:overlap}~(b), in the rich regime
$t(x,y),t(x,y')\ge T$ the joint fiber satisfies
$|\Omega_{\alpha,\beta}^{\circ}|\ge(1-o(1))|\Omega_{\alpha,\beta}|$
where $\Omega_{\alpha,\beta}=\fiber{\alpha}\cap\fiber{\beta}$.  The
joint fiber is non-empty because the $(x,y)$- and $(x,y')$-coordinate
projections are independent over the $C$-chain orientation: any
$L\in\LE(P)$ that places $x$ adjacent to its $C$-window
$C(x,y)\cup C(x,y')$ and both $y,y'$ in their respective ranges
contributes, and the count of such $L$ is at least
$c_1\,\min\bigl(t(x,y),t(x,y')\bigr)\,|\LE(P)|/(|A|\,|B|\,|C|)$
for an absolute constant $c_1>0$ (Step~1
Theorem~\ref{thm:interface}, applied to the joint
$\mathbb{Z}^2$-grid model on $\Omega^{\circ}_{\alpha,\beta}$).  The Step~1
interaction-locus bound $|\mathrm{Int}_{\alpha,\beta}|=O((t_\alpha+t_\beta)^2)$
(Cor.~\ref{cor:overlap}\ref{it:density} of \texttt{step1.tex})
removes at most a $o(1)$-fraction once the overlap is $\gg
(t_\alpha+t_\beta)^2$, which holds for $T$ large.  Therefore
\begin{equation}\label{eq:share-end-overlap}
  w_{\alpha\beta}\;=\;
  \mu\bigl((\fiber{\alpha}\cap\fiber{\beta})\setminus
  (B_\alpha\cup B_\beta)\bigr)
  \;\ge\; c_1\,T\,|\LE(P)|/(|A|\,|B|\,|C|)\;=\;\theta_{\rm pair},
\end{equation}
so every endpoint-sharing pair in $\Rich^\star$ is active in
$G_{\Rich}$.  The symmetric case $y=y'$ (sharing the $B$-endpoint) is
identical.

\smallskip
\textbf{Step~3: Endpoint-sharing graph on $\Rich_{AB}$ has diameter
$\le 3$.}  For $a\in A$, write $N_B(a):=\{b\in B:(a,b)\in\Rich_{AB}\}$
and analogously $N_A(b)$.  By Step~5~(R), $\sum_a|N_B(a)|=|\Rich_{AB}|
\ge c_T|A||B|$, so by averaging at most a $(2/c_T)$-fraction of
$a\in A$ have $|N_B(a)|<(c_T/2)|B|$.  Call $a$ \emph{rich-row} if
$|N_B(a)|\ge(c_T/2)|B|$, and similarly $b$ \emph{rich-column}.
For two rich-row endpoints $a,a'\in A$, by inclusion--exclusion
$|N_B(a)\cap N_B(a')|\ge c_T|B|-|B|=(c_T-1)|B|>0$ for $c_T>1$, so
there exists $b''\in B$ with $(a,b''),(a',b'')\in\Rich$.  Hence for any
$\alpha=(a,b),\beta=(a',b')\in\Rich$ with $a,a'$ both rich-rows, the
path
\[
  \alpha=(a,b)\;\longrightarrow\;(a,b'')\;\longrightarrow\;(a',b'')\;
  \longrightarrow\;(a',b')=\beta
\]
is a length-$3$ walk in the endpoint-sharing graph (each consecutive
pair shares an endpoint).  At most a $(4/c_T)$-fraction of
$\Rich_{AB}$ involves a non-rich-row or non-rich-column endpoint, so
the diameter-$3$ property holds on a $(1-O(c_T^{-1}))$-fraction of
the rich-interface mass.

\smallskip
\textbf{Step~4: From endpoint-sharing connectivity to giant
component.}  By Step~2, every endpoint-sharing pair in $\Rich^\star$
is an active edge of $G_{\Rich}$.  Combined with Step~3, the active
subgraph of $G_{\Rich}$ contains, on a $(1-O(c_T^{-1}))$-mass
subset, paths of length~$\le 3$ between any two interfaces.  The
$S$-good trim from $\Rich$ to $\Rich^\star$
(Definition~\ref{def:G_Rich}) deletes only an $o(1)$-mass fraction
(Step~6), and the diameter-$3$ routing tolerates such deletions
because each step has $\Omega(c_T|B|)$ alternative choices of
intermediate endpoint.  Specifically, fix any reference
$\alpha_0\in\Rich^\star$ with rich-row $a_0$ and rich-column $b_0$;
every other rich-row/rich-column $\alpha\in\Rich^\star$ admits
$\Omega(c_T^2|A||B|)$ length-$3$ paths to $\alpha_0$, of which a
$(1-o(1))$-fraction stay inside $\Rich^\star$.  Hence at least one
such path survives the trim, placing $\alpha$ in the same component
as $\alpha_0$.  The resulting connected component
$\mathcal{C}\subseteq\Rich^\star$ contains all rich-row/rich-column
$S$-good interfaces, hence carries incidence mass
$\ge(1-O(c_T^{-1}))M_0$.

\smallskip
\textbf{Step~5: Triple version.}  For three rich interfaces
$e_{xy}=(x,y),e_{yz}=(y,z),e_{xz}=(x,z)$ forming a triangle through
a common triple $\{x,y,z\}\subseteq P$ that spans the three chains,
Step~1 Corollary~\ref{cor:triple-overlap} applies: the triple
overlap $\Omega^{(3)}=\fiber{e_{xy}}\cap\fiber{e_{yz}}\cap
\fiber{e_{xz}}$ supports a commuting $\mathbb{Z}^3$-cube model on a
positive-density subset, with mass
$|\Omega^{(3)}|\ge c_2\,T\,|\LE(P)|/(|A|\,|B|\,|C|)$ for an absolute
constant $c_2>0$ once $T$ is large (the
existence of the cube model and the joint-window count from
Theorem~\ref{thm:interface} give the lower bound by the same
argument as Step~2).  The active-triple threshold is set
analogously to $\theta_{\rm pair}$, and triple visibility
$w_{e_{xy},e_{yz},e_{xz}}\ge$ threshold follows.  Each
endpoint-sharing edge of $\mathcal{C}$ at $x$ extends to a triangle
through any rich-column $z$ (of which there are
$\Omega(c_T|C|)$), so the triple hypergraph induced on
$\mathcal{C}$ is connected with the same $(1-o(1))$-mass coverage.
\end{proof}

\textbf{Compatibility with downstream uses.}  Lemmas
\ref{lem:sign-consistency}, \ref{lem:potential}, \ref{lem:single-c}
all condition on a single connected component of $G_{\Rich}$
carrying $(1-o(1))$-mass; Lemma~\ref{lem:giant-component} furnishes
exactly such a $\mathcal{C}$.  The mass-loss exponent
$O(c_T^{-1})$ is absorbed by the $o(1)$ slack already built into
those lemmas (since $c_T\to\infty$ as the richness threshold $T\to
\infty$ in Step~5).



\part{Conclusion}


\begin{abstract}
Step 8 is the final assembly. It contains no new combinatorial content:
every case branch produced by Steps 1--7 is terminated by a contradiction
or by the direct production of a balanced pair. The job of this section
is to name, in one place, the glue lemmas needed and to chain the outputs
of Steps 2, 5, 6, and 7 into a proof of the main theorem. The
Cheeger-style ``counterexample $\Rightarrow$ low BK expansion''
theorem (Theorem~E), once labelled an external input, is now proved
in full inside this file (Section~\ref{sec:G1}) via a Dirichlet
comparison plus an elementary averaging argument on pair functions,
so Step~8 depends only on sealed outputs of Steps 1--7.

Several named lemmas cited below are proved elsewhere (Steps 1--7).
Five items live in Step 8 itself (Theorem~E and G2--G5); each is
discharged in its own section below.
\end{abstract}

\section{Setup and notation}

Throughout, $P = (X, \le_P)$ is a finite poset of width~$3$, and
$\mathcal{L}(P)$ is the set of its linear extensions. The BK graph
$G_{\mathrm{BK}}(P)$ has vertex set $\mathcal{L}(P)$ and edges joining
pairs of extensions that differ by one adjacent swap of an incomparable
pair. For $S \subseteq \mathcal{L}(P)$,
\[
  |\partial S| := |E(S, S^c)|,\qquad
  \Phi(S) := \frac{|\partial S|}{\min(\operatorname{vol}(S), \operatorname{vol}(S^c))}.
\]
We call $P$ a \emph{$\gamma$-counterexample} if no incomparable pair
$\{x,y\}\subseteq X$ satisfies $1/3 \le \Pr_\sigma[x < y] \le 2/3$ and
\[
  \beta(P) := \inf_{x\parallel y}\min\{\Pr[x<y],\Pr[y<x]\}\ge \gamma.
\]

\section{Theorem E: counterexample implies low BK conductance}
\label{sec:G1}

The Cheeger/spectral half of the program is
Theorem~\ref{thm:cex-implies-low-expansion} below. We derive it
entirely within this file, by (i)~an elementary Dirichlet/conductance
comparison for indicator functions (Lemma~\ref{lem:dirichlet-conductance},
standard for reversible chains), and (ii)~an elementary averaging
argument over the family of pair-order test functions
$\{f_{xy}\}_{x \parallel y}$ (Lemma~\ref{lem:frozen-pair-existence}),
producing a frozen pair with ratio
$\mathcal E(f_{xy})/\mathrm{Var}(f_{xy}) \le 2/(\gamma n)$. No
pair-Poincar\'e inequality, representation-theoretic input, or
rigid-spine lemma is required; the two lemmas below suffice.

\subsection{Statement}

The theorem below no longer carries a GAP label: its proof is
complete in this file, with Lemmas~\ref{lem:dirichlet-conductance}
and~\ref{lem:frozen-pair-existence} fully established below.

\begin{theorem}[Counterexample $\Rightarrow$ low BK expansion]
\label{thm:cex-implies-low-expansion}
Set $c_0 := \gamma$ and $\eta(\gamma, n) := 2/(\gamma n)$. If $P$ is a
width-$3$ indecomposable $\gamma$-counterexample on $n \ge 2$
elements, then there is a cut $S \subseteq \mathcal{L}(P)$ with
\[
  \operatorname{vol}(S)
  \;\ge\; c_0\, \operatorname{vol}(\mathcal L(P)),
  \qquad
  \Phi(S) \;\le\; \eta(\gamma, n).
\]
Here $\eta(\gamma, n) \downarrow 0$ as $n \to \infty$ for any fixed
$\gamma \in (0, 1/3]$; for every fixed $\gamma, T$ the bound lies
below any prescribed positive threshold once
$n \ge n_0(\gamma, T)$ (Remark~\ref{rem:n-dependence-g1}), which is
what the main-theorem cascade of Proposition~\ref{prop:G2} requires.
\end{theorem}

The proof occupies the rest of this section. The strategy has two
ingredients: (i)~a Dirichlet/conductance comparison
(Lemma~\ref{lem:dirichlet-conductance}), standard for reversible
chains; and (ii)~the existence of a \emph{BK-frozen} incomparable
pair (Lemma~\ref{lem:frozen-pair-existence}), which we establish by
an elementary averaging argument. The cut $S$ is then simply the
indicator set of that pair.

\subsection{Dirichlet form vs.\ conductance for pair indicators}

For $f\colon \mathcal{L}(P) \to \mathbb{R}$ write the BK Dirichlet
form in symmetric-chain normalization,
\[
  \mathcal{E}(f)
  \;:=\; \frac{1}{2(n-1)\,|\mathcal L(P)|}
  \sum_{\sigma \in \mathcal L(P)}\sum_{i=1}^{n-1}
  \bigl(f(\sigma) - f(\tau_i \sigma)\bigr)^2,
\]
and $\mathrm{Var}(f) = \mathbb E_\pi[f^2] - (\mathbb E_\pi f)^2$, where
$\pi$ is uniform on $\mathcal L(P)$ and $\tau_i$ is the BK move at
position $i$ (swap if incomparable, identity otherwise). For an
incomparable pair $x \parallel y$ in $P$, set
\[
  S_{xy} \;:=\; \{\sigma \in \mathcal L(P) : x \text{ precedes } y \text{ in } \sigma\},
  \qquad
  f_{xy} \;:=\; \mathbf 1_{S_{xy}}.
\]
Then $\operatorname{vol}(S_{xy}) = p_{xy}\,|\mathcal L(P)|$ and
$\operatorname{vol}(S_{xy}^c) = (1 - p_{xy})\,|\mathcal L(P)|$, where
$p_{xy} = \Pr_{\sigma}[x <_\sigma y]$.

For the BK graph $G_{\mathrm{BK}}(P)$ we take the edge
convention induced by the lazy chain: an edge of multiplicity $1$ for
each ordered pair $(\sigma, i)$ with $\sigma(i), \sigma(i+1)$
incomparable (so $|E(G_{\mathrm{BK}})| \le (n-1)|\mathcal L(P)|/2$), and
$\operatorname{vol}(S) := |S|\,(n-1)$. With this normalization the
Dirichlet form of an indicator and the edge boundary are related by
$\mathcal{E}(\mathbf 1_S) = |\partial S| / ((n-1)|\mathcal L(P)|)$.

\begin{lemma}[Conductance $\le$ Dirichlet--variance ratio]
\label{lem:dirichlet-conductance}
For every $S \subseteq \mathcal L(P)$ with $0 < |S| < |\mathcal L(P)|$,
\[
  \Phi(S) \;\le\;
  \frac{\mathcal{E}(\mathbf 1_S)}{\mathrm{Var}(\mathbf 1_S)}.
\]
In particular, for every incomparable pair $x \parallel y$,
\[
  \Phi(S_{xy}) \;\le\;
  \frac{\mathcal{E}(f_{xy})}{\mathrm{Var}(f_{xy})}.
\]
\end{lemma}

\begin{proof}
Set $\pi(S) = |S|/|\mathcal L(P)|$, so
$\mathrm{Var}(\mathbf 1_S) = \pi(S)\,\pi(S^c)$, and let
$Q(S, S^c)$ denote the BK ``edge flow'' from $S$ to $S^c$,
$Q(S,S^c) = \sum_{\sigma \in S, i : \tau_i \sigma \notin S}
1/((n-1)\,|\mathcal L(P)|) = |\partial S| / ((n-1)\,|\mathcal L(P)|)$.
Then $\mathcal{E}(\mathbf 1_S) = Q(S,S^c)$ and
$|\partial S| = (n-1)|\mathcal L(P)|\,Q(S,S^c)$, and
$\operatorname{vol}(S) = |S|(n-1) = (n-1)|\mathcal L(P)|\,\pi(S)$.
Hence
\[
  \Phi(S) \;=\; \frac{|\partial S|}{\min(\operatorname{vol}(S),
  \operatorname{vol}(S^c))}
  \;=\; \frac{Q(S,S^c)}{\min(\pi(S),\pi(S^c))}.
\]
Since $\pi(S)\,\pi(S^c) \le \min(\pi(S), \pi(S^c))$ (the smaller
of $\pi(S), \pi(S^c)$ is $\le 1/2 \le 1$, and the product is that
smaller value times a factor $\le 1$),
\[
  \frac{Q(S,S^c)}{\min(\pi(S),\pi(S^c))}
  \;\le\;
  \frac{Q(S,S^c)}{\pi(S)\,\pi(S^c)}
  \;=\; \frac{\mathcal{E}(\mathbf 1_S)}{\mathrm{Var}(\mathbf 1_S)}. \qedhere
\]
\end{proof}

\subsection{Existence of a BK-frozen pair}

Call an incomparable pair $(x, y)$ \emph{$\lambda$-BK-frozen} if
$\mathcal{E}(f_{xy}) \le \lambda\,\mathrm{Var}(f_{xy})$. We prove
existence of a frozen pair in every width-$3$ indecomposable
counterexample by an elementary averaging argument. No
pair-Poincar\'e / representation-theoretic machinery is required;
Lemmas~\ref{lem:indec-incompairs}--\ref{lem:frozen-pair-existence}
below are entirely self-contained.

\begin{lemma}[Indecomposable posets have many incomparable pairs]
\label{lem:indec-incompairs}
Let $P$ be an indecomposable poset on $n \ge 2$ elements. Then every
element of $P$ is incomparable to at least one other element; in
particular
\[
  I(P) \;\ge\; n/2,
\]
where $I(P)$ is the number of unordered incomparable pairs in $P$.
\end{lemma}

\begin{proof}
Suppose some $x \in X$ is comparable in $P$ to every other element.
Set $A := \{y \in X : y <_P x\}$ and $B := \{y \in X : y >_P x\}$, so
$X = A \sqcup \{x\} \sqcup B$. By transitivity, every $a \in A$ and
every $b \in B$ satisfy $a <_P b$. Hence
\[
  P \;=\;
  \begin{cases}
    (A \cup \{x\}) \oplus B & \text{if } B \ne \emptyset, \\
    A \oplus \{x\}          & \text{if } B = \emptyset, A \ne \emptyset, \\
    \{x\} \oplus B          & \text{if } A = \emptyset, B \ne \emptyset,
  \end{cases}
\]
an ordinal sum of two nonempty induced subposets (using $n \ge 2$).
This contradicts indecomposability of $P$. Hence every $x \in X$ has
at least one $y \in X$ with $y \parallel_P x$, and
$2\,I(P) = \sum_{x \in X} \bigl|\{y : y \parallel_P x\}\bigr| \ge n$.
\end{proof}

\begin{lemma}[Frozen-pair existence by averaging]
\label{lem:frozen-pair-existence}
Let $P$ be a width-$3$ indecomposable $\gamma$-counterexample on
$n \ge 2$ elements with $\gamma \in (0, 1/3]$. Then $P$ contains an
incomparable pair $(x, y)$ with
\[
  \frac{\mathcal{E}(f_{xy})}{\mathrm{Var}(f_{xy})}
  \;\le\;
  \lambda(\gamma, n) \;:=\; \frac{2}{\gamma\, n}.
\]
\end{lemma}

\begin{proof}
\emph{Step 1: total Dirichlet energy bounded by $1/2$.}
For each incomparable pair $(x,y)$, the Dirichlet form expands as
\[
  \mathcal{E}(f_{xy})
  \;=\; \frac{1}{2(n-1)\,|\mathcal L(P)|}
  \sum_{\sigma \in \mathcal L(P)}\sum_{i=1}^{n-1}
  \bigl(f_{xy}(\sigma) - f_{xy}(\tau_i\sigma)\bigr)^2.
\]
The summand $(f_{xy}(\sigma) - f_{xy}(\tau_i\sigma))^2$ equals~$1$ iff
$\tau_i$ flips the relative order of $x$ and $y$, which occurs iff
$\{\sigma(i), \sigma(i+1)\} = \{x,y\}$ (the pair is incomparable, so
$\tau_i$ does swap); otherwise it is $0$. Summing over unordered
incomparable pairs and exchanging orders of summation,
\begin{align*}
  \sum_{x \parallel y} \mathcal{E}(f_{xy})
  &\;=\;
  \frac{1}{2(n-1)\,|\mathcal L(P)|}
  \sum_{\sigma \in \mathcal L(P)}
  \bigl|\bigl\{i \in [n-1] : \sigma(i) \parallel_P \sigma(i+1)
  \bigr\}\bigr| \\
  &\;\le\;
  \frac{1}{2(n-1)\,|\mathcal L(P)|} \cdot (n-1)\,|\mathcal L(P)|
  \;=\; \tfrac{1}{2},
\end{align*}
since at most $n-1$ positions can contribute for each $\sigma$.

\emph{Step 2: per-pair variance lower bound.}
For every incomparable pair $(x,y)$,
$\mathrm{Var}(f_{xy}) = p_{xy}(1 - p_{xy})$. The
$\gamma$-counterexample hypothesis gives
$p_{xy} \in [\gamma, 1/3) \cup (2/3, 1-\gamma]$, whence
$\min(p_{xy}, 1-p_{xy}) \ge \gamma$ and
\[
  \mathrm{Var}(f_{xy})
  \;\ge\; \gamma\,(1-\gamma)
  \;\ge\; \gamma/2,
\]
using $\gamma \le 1/3 \le 1/2$.

\emph{Step 3: incomparable-pair count and total variance.}
By Lemma~\ref{lem:indec-incompairs}, $I(P) \ge n/2$ because $P$ is
indecomposable. Summing Step~2 over incomparable pairs,
\[
  \sum_{x \parallel y} \mathrm{Var}(f_{xy})
  \;\ge\; \tfrac{\gamma}{2}\, I(P)
  \;\ge\; \tfrac{\gamma\, n}{4}.
\]

\emph{Step 4: averaging.}
The minimum of $\mathcal{E}(f_{xy})/\mathrm{Var}(f_{xy})$ over
incomparable pairs is bounded by the ratio of sums. By Steps~1
and~3,
\[
  \min_{x \parallel y}
    \frac{\mathcal{E}(f_{xy})}{\mathrm{Var}(f_{xy})}
  \;\le\;
  \frac{\sum_{x \parallel y} \mathcal{E}(f_{xy})}
       {\sum_{x \parallel y} \mathrm{Var}(f_{xy})}
  \;\le\;
  \frac{1/2}{\gamma\, n / 4}
  \;=\; \frac{2}{\gamma\, n}
  \;=\; \lambda(\gamma, n).
\]
Any minimizing pair $(x,y)$ is therefore $\lambda(\gamma, n)$-BK-frozen.
\end{proof}

\begin{remark}[Reduction to indecomposable posets]
\label{rem:decomp-reduction}
The indecomposability hypothesis in
Lemma~\ref{lem:frozen-pair-existence} is harmless for the main-theorem
application: a minimal $\gamma$-counterexample is always
indecomposable. If $P = P_1 \sqcup P_2$ is a disjoint union with
both parts nonempty, every cross pair
$(a,b) \in P_1 \times P_2$ is incomparable with $p_{ab}(P) = 1/2 \in
[1/3, 2/3]$ (randomize the interleaving of $\mathcal L(P_1)$ and
$\mathcal L(P_2)$), so $P$ has a balanced pair, contradicting the
counterexample hypothesis. If $P = P_1 \oplus P_2$ is an ordinal sum
with both parts nonempty, every incomparable pair lies inside one
summand, and by the formula $p_{xy}(P) = p_{xy}(P_i)$ (for $(x,y)
\subseteq P_i$) at least one summand is itself a
$\gamma$-counterexample, contradicting minimality of $P$.
\end{remark}

\begin{remark}[$n$-dependence absorbed by the small-$n$ base case]
\label{rem:n-dependence-g1}
The bound $\lambda(\gamma, n) = 2/(\gamma n)$ depends on $n$; the
conductance bound $\eta(\gamma, n)$ in
Theorem~\ref{thm:cex-implies-low-expansion} inherits this dependence.
This suffices for the Proposition~\ref{prop:G2} cascade: after the
$M$ vs.\ $\mathrm{vol}$ reconciliation (Proposition~\ref{prop:G2},
counting-measure form), the compatibility inequality~\eqref{eq:G2}
requires $\eta(\gamma,n)(n-1) < c_0\, c_5(T)\, c_6\, (\delta -
K(T)\varepsilon)$ with the right-hand side fixed first as a function
of $\gamma$ and $T$. Substituting $\eta(\gamma,n)(n-1) =
2(n-1)/(\gamma n) \le 2/\gamma$ and $c_0 = \gamma$, the inequality
reduces to the $n$-independent arithmetic condition
\[
  \gamma^2\, c_5(T)\, c_6\, (\delta - K(T)\varepsilon) \;\ge\; 2,
\]
which is $n$-independent but \emph{not unconditional}: the cascade
closes at every $n \ge 2$ if it holds (which is the content of
Hypothesis~\ref{hyp:arith}); if it fails for every admissible $T$,
the small-$n$ base case of Remark~\ref{rem:small-n} still discharges
the $1/3$--$2/3$ property for $n \le n_0(\gamma, T)$, but the
large-$n$ tail $n > n_0$ is not closed by the present argument.
That tail is the only remaining portion of the argument
not-unconditionally-covered: for $\gamma \ge 1/3 -
\delta_{\mathrm{KL}}$ the Kahn--Linial bound excludes
counter\-examples outright, and for $n \le n_0(\gamma, T)$
Lemma~\ref{lem:small-n} applies, so it is exactly the regime
$\gamma < 1/3 - \delta_{\mathrm{KL}}$, $n > n_0(\gamma, T)$, and
``$\gamma^2 c_5(T) c_6 \delta \ge 2$ fails for every admissible
$T$'' that Hypothesis~\ref{hyp:arith} is introduced to exclude.
No $n \to \infty$ limit is taken: the cascade produces a
contradiction at the single value of $n$ at hand, provided the
constant condition holds. Verification of $\gamma^2 c_5 c_6 (\delta
- K(T)\varepsilon) \ge 2$ for some $T = T(\gamma)$ at every $\gamma
\in (0, 1/3 - \delta_{\mathrm{KL}})$ is the final
constants audit, outside the scope of G1/G2 and recorded as
Hypothesis~\ref{hyp:arith}.
\end{remark}

\subsection{Proof of Theorem~\ref{thm:cex-implies-low-expansion}}

\begin{proof}[Proof of Theorem~\ref{thm:cex-implies-low-expansion}]
Let $P$ be a width-$3$ indecomposable $\gamma$-counterexample on
$n \ge 2$ elements, and set $c_0 := \gamma$,
$\eta(\gamma, n) := \lambda(\gamma, n) = 2/(\gamma n)$ with
$\lambda(\gamma, n)$ from Lemma~\ref{lem:frozen-pair-existence}. Since
$\gamma \le 1/3$, $c_0 \in (0, 1/2)$; and
$\eta(\gamma, n) \downarrow 0$ as $n \to \infty$ for any fixed
$\gamma > 0$.

By Lemma~\ref{lem:frozen-pair-existence}, there is an incomparable
pair $(x, y) \subseteq X$ with
\[
  \mathcal{E}(f_{xy})
  \;\le\; \lambda(\gamma, n)\,\mathrm{Var}(f_{xy})
  \;=\; \eta(\gamma, n)\,\mathrm{Var}(f_{xy}).
\]
Take $S := S_{xy} = \{\sigma : x <_\sigma y\}$, so $|S| = p_{xy}\,
|\mathcal L(P)|$ and $\operatorname{vol}(S) = p_{xy}\,(n-1)\,
|\mathcal L(P)|$ by the convention
$\operatorname{vol}(T) := |T|(n-1)$ of \S\ref{sec:G1}. The counterexample
hypothesis $\beta(P) \ge \gamma$ forces
$p_{xy} \in [\gamma, 1/3) \cup (2/3, 1-\gamma]$, whence $p_{xy} \ge
\gamma$ and $1 - p_{xy} \ge \gamma$. With
$\operatorname{vol}(\mathcal L(P)) = (n-1)\,|\mathcal L(P)|$, this gives
\[
  \frac{\operatorname{vol}(S)}{\operatorname{vol}(\mathcal L(P))}
  \;=\; p_{xy}
  \;\ge\; \gamma
  \;=\; c_0,
  \qquad
  \frac{\operatorname{vol}(S^c)}{\operatorname{vol}(\mathcal L(P))}
  \;=\; 1 - p_{xy}
  \;\ge\; \gamma
  \;=\; c_0.
\]
Applying Lemma~\ref{lem:dirichlet-conductance} to $S$ and the pair
$(x,y)$,
\[
  \Phi(S) \;\le\;
  \frac{\mathcal{E}(f_{xy})}{\mathrm{Var}(f_{xy})}
  \;\le\; \eta(\gamma, n),
\]
which is the desired conductance bound.

Finally, the smallness of $\eta(\gamma, n)$ needed to feed the Step~4
/ Step~6 overlap-counting bound (inequality~\eqref{eq:G2} of
Proposition~\ref{prop:G2}) is achieved once $n \ge n_0(\gamma, T)$, by
Remark~\ref{rem:n-dependence-g1}. For $n < n_0(\gamma, T)$, the
finite-enumeration base case of Remark~\ref{rem:small-n} supplies the
1/3--2/3 property directly, so the cascade closes on every $n$.
\end{proof}

\begin{remark}[Why a single-cut argument suffices]
Lemma~\ref{lem:frozen-pair-existence} produces only the bound
$\mathcal{E}(f_{xy}) \le \lambda\,\mathrm{Var}(f_{xy})$ for \emph{some}
incomparable pair, a statement about a single test function in the
span of pair functions. The rest of the program (Step~4 rectangle
bounds, Step~6 overlap counting) needs a genuine low-conductance cut
of $G_{\mathrm{BK}}(P)$. Lemma~\ref{lem:dirichlet-conductance} is what
bridges the two formulations: once a frozen pair exists, the
\emph{single} cut $S_{xy}$ automatically has small conductance, without
invoking Cheeger's inequality for the full BK walk. This avoids the
delicate question of whether the minimum-energy eigenfunction of the
BK Laplacian is itself a pair function.
\end{remark}

\textbf{G1 status.}
Lemma~\ref{lem:dirichlet-conductance} is proved in this file under a
consistent normalization (edge convention
$|E| \le (n-1)|\mathcal L(P)|/2$, $\mathrm{vol}(S) = |S|(n-1)$, and
$\mathcal E(\mathbf 1_S) = |\partial S|/((n-1)|\mathcal L(P)|)$), and
Lemma~\ref{lem:frozen-pair-existence} is proved by the elementary
averaging argument above. No pair-Poincar\'e inequality,
representation-theoretic input, or rigid-spine lemma is invoked:
averaging over the $\binom{n}{2}$-sized family of pair functions
$\{f_{xy}\}$ combined with the trivial bound
$\sum_{x \parallel y} \mathcal E(f_{xy}) \le 1/2$ and the
indecomposability-based count $I(P) \ge n/2$ yields
$\lambda(\gamma, n) \le 2/(\gamma n)$, at the cost of introducing
$n$-dependence in $\eta$. The dependence is absorbed by the small-$n$
base case (Remarks~\ref{rem:n-dependence-g1}, \ref{rem:small-n})
for $n \le n_0(\gamma, T)$; for $n > n_0(\gamma, T)$ it is absorbed
through the $n$-independent arithmetic inequality of
Hypothesis~\ref{hyp:arith}. G1 itself is \emph{fully discharged}
(the $n$-dependence is not a G1-internal gap).

\section{The pieces from Steps 1--7}

For readability we restate the three outputs Step~8 will use, in the
exact form they are needed. Proofs, constants, and precise definitions
live in Steps~1--7 above.

\begin{proposition}[Step 5 dichotomy]\label{prop:step5}
Fix a richness threshold $T = T(\gamma)$ large. For every width-$3$
$\gamma$-counterexample $P$, one of the following holds.
\begin{itemize}[nosep]
 \item[(R)] \emph{Richness.} The total rich-interface incidence is
 large: a random linear extension $L$ lies in $\Omega_T(1)$ rich good
 fibers in expectation, and in counting measure
 \[
   M \;:=\; \sum_{\alpha,\beta \in \mathcal{R}^\star} w_{\alpha\beta}
   \;\ge\; c_5(T)\,|\mathcal L(P)|,
 \]
 where $w_{\alpha\beta}$ is the good-overlap mass between rich
 interfaces, $\mathcal R^\star$ is the set of $S$-good rich interfaces
 from Step~2, and $c_5(T) := c_T'(T)$ is the Step~5 second-moment
 constant of Corollary~\ref{cor:second-moment} (in
 \texttt{step5.tex}), which is independent of $n$. The factor of
 $n-1$ between $|\mathcal L(P)|$ and $\operatorname{vol}(\mathcal
 L(P)) = (n-1)|\mathcal L(P)|$ is carried explicitly in the
 Proposition~\ref{prop:G2} upper bound (see~\eqref{eq:G2-upper}),
 not absorbed into $c_5$.
 \item[(C)] \emph{Collapse.} After discarding $O_T(1)$ exceptional
 elements on each of the three Dilworth chains, every bipartite
 incomparability relation in $P$ is supported on a band of width
 $O_T(1)$, i.e., $P$ is layered in the sense of Step~7.
\end{itemize}
\end{proposition}

\begin{proposition}[Step 6 coherence dichotomy]\label{prop:step6}
Suppose Proposition~\ref{prop:step5}(R) holds, and $S \subseteq
\mathcal{L}(P)$ is a cut with $\Phi(S) \le \eta$. Then, for the constants
$c_0, c_5(T), c_6, K(T)$ named in Section~\ref{sec:G2-constants}, if
\[
  \eta\,(n-1) \;<\; \tfrac{1}{2}\, c_0\, c_5(T)\, c_6\,\big(\delta - K(T)\,\varepsilon\big),
\]
then the weighted fraction of incoherent rich-overlap mass is at most
$\delta$, i.e., the weighted fraction of coherent rich-overlap mass is
at least $1 - \delta$. Taking $\delta = 2 K(T)\,\varepsilon$, the
incoherent fraction is at most $2 K(T)\,\varepsilon$, a quantity that
tends to~$0$ with the Step~2 fiber error $\varepsilon$.
\end{proposition}

\noindent
If instead more than a $\delta$-fraction were incoherent, Step~4 applied
to the Step~5 overlap mass would force $|\partial S| \ge c_5(T)\,c_6\,
(\delta - K(T)\varepsilon)\,|\mathcal L(P)|$ (counting measure),
contradicting
$\Phi(S) \le \eta$. A quantitative proof is given by
Proposition~\ref{prop:G2}.

\begin{proposition}[Step 7: globalization and collapse]\label{prop:step7}
Assume the conclusion of Proposition~\ref{prop:step6} with coherent
fraction $\ge 1 - \delta$ (so that, on the Step~8 parameter choice,
$\delta = 2 K(T)\,\varepsilon$). Set $\eta_* := C_7\,\delta$ for an
absolute constant $C_7 > 0$ coming from the Step~7 potential-extraction
argument. Then:
\begin{enumerate}[nosep, label=(\roman*)]
 \item (Prop.~7.1) There is a potential $a\colon X \to \mathbb{R}$ and a
 threshold $c \in \mathbb{R}$ such that
 \[
   \Pr_{L \sim \mathcal L(P)}\!\big[\mathbf{1}_S(L) \ne
   \mathbf{1}_{\{H(L) \ge c\}}\big] \;\le\; \eta_*,
   \qquad
   H(L) := \sum_{z \in X} a(z)\,\mathrm{pos}_L(z),
 \]
 and the positive BK direction across each rich interface $(x,y)$ is
 the sign of $a(y) - a(x)$.
 \item (Prop.~7.2) A cut of the form $\{H \ge c\}$ with $\Phi \le \eta$
 coinciding with $S$ on a $(1-\eta_*)$-fraction of $\mathcal L(P)$
 forces $P$ to admit a layered width-$3$ decomposition on a
 $(1 - \eta_*)$-fraction of its mass (with interaction width $w = w(T)$
 depending only on the Step~5 richness threshold).
 \item (Prop.~7.3) A width-$3$ poset that is layered in the sense of
 (ii) or of Proposition~\ref{prop:step5}(C) has a balanced pair, or
 else contains a proper induced subposet $P' \subsetneq P$ which is a
 $(\gamma/2)$-counterexample. (Proved as
 Lemmas~\ref{lem:layered-reduction} and~\ref{lem:layered-balanced}.)
\end{enumerate}
\end{proposition}

\section{The main theorem}
\label{sec:main-thm}

\begin{hypothesis}[Arithmetic richness]\label{hyp:arith}
For every $\gamma \in (0,\, 1/3 - \delta_{\mathrm{KL}})$, there exists
$T = T(\gamma) \ge T_0$ (with $T_0$ the absolute Step~5 richness
threshold) such that the $n$-independent arithmetic inequality
\[
  \gamma^2\, c_5(T)\, c_6\, \delta_0(\gamma) \;\ge\; 8
\]
holds, where $\delta_0(\gamma) :=
\min\!\big(\tfrac{1}{4C_7},\tfrac{1}{2wK_0}\big)$ is the cascade
target and $c_5(T), c_6, C_7, w, K_0$ are the Step~5/6/7
constants fixed by Remark~\ref{rem:G2-order}.
\end{hypothesis}

\noindent
Hypothesis~\ref{hyp:arith} is an arithmetic condition on the
sealed Step~5/6 richness densities: it is not derived from
Steps~1--7, and its verification is external to the structural
argument here (see Remark~\ref{rem:final-constants-audit} and the
scope discussion below).

\begin{theorem}[Width-$3$ 1/3--2/3, conditional on Steps 1--7 and on Hypothesis~\ref{hyp:arith}]
\label{thm:main-step8}
Assume Hypothesis~\ref{hyp:arith}. Then every finite width-$3$ poset
that is not a total order has a balanced pair.
\end{theorem}

\noindent
(Conditional on the sealed outputs of Steps 1--7 and on the
arithmetic-richness Hypothesis~\ref{hyp:arith}; the Step~8
items G1--G5 are all discharged inside this file, see
Section~\ref{sec:G1} for G1, Section~\ref{sec:G2-constants} for G2,
Section~\ref{sec:layered-reduction} for G3, and
Section~\ref{sec:g4-balanced-pair} for G4. G5 is the present proof
itself.)

\medskip\noindent\emph{Scope note.} Without
Hypothesis~\ref{hyp:arith}, the structural argument of this file
covers (a)~the Kahn--Linial regime $\gamma \ge 1/3 -
\delta_{\mathrm{KL}}$, where no $\gamma$-counter\-example exists at
any $n$ or width~\cite{KahnLinial1991}, and (b)~the finite-$n$
regime $n \le n_0(\gamma, T)$ via the small-$n$ base case
(Remark~\ref{rem:small-n}, Lemma~\ref{lem:small-n}). The
arithmetic-richness Hypothesis~\ref{hyp:arith} is what pins down
the intermediate regime ($\gamma < 1/3 - \delta_{\mathrm{KL}}$ and
$n > n_0(\gamma, T)$) via the BK-expansion cascade; a failure of
Hypothesis~\ref{hyp:arith} at some small $\gamma$ would leave that
large-$n$ tail uncovered by the present argument.

\begin{proof}
Fix $\gamma > 0$ and a minimal (by $|X|$) width-$3$
$\gamma$-counter\-example $P$ on $n = |X|$ elements. (If no such minimal
$\gamma$-counter\-example exists for some $\gamma > 0$, then no
$\gamma$-counter\-example exists at all, and the conclusion follows.)
By Remark~\ref{rem:decomp-reduction}, $P$ is indecomposable, so
Theorem~\ref{thm:cex-implies-low-expansion} applies.

\medskip\noindent\emph{Parameter cascade.} Following
Remark~\ref{rem:G2-order}, fix the parameters in order:
\begin{enumerate}[nosep, label=(\arabic*)]
 \item $T := T(\gamma)$, the Step~5 richness threshold; this fixes
 $c_5 = c_5(T,\gamma)$, $K(T) = C_4/c_4 + C_2(T)$ (the Step~4
 rectangle constants $c_4, C_4$ enter rescaled because Section~4
 above factors $c_4$ out of both terms of the Step~4 error bound
 $(c_4\alpha - C_4\varepsilon)|R|$; see the Step~4 aggregation
 paragraph below for the derivation), and the interaction width
 $w = w(T)$ of Proposition~\ref{prop:step7}(ii).
 \item $\delta_0 := \min\!\big(\tfrac{1}{4 C_7}, \tfrac{1}{2 w K_0}\big)$
 where $C_7$ is the constant of Proposition~\ref{prop:step7} and
 $K_0 = K_0(\gamma, w)$ is from Lemma~\ref{lem:layered-reduction};
 this is the target aggregate incoherence fraction.
 \item $\varepsilon := \delta_0 / (2 K(T))$, the Step~2 fiber error.
 \item $n \ge n_0(\gamma, T)$, chosen so that
 $\eta(\gamma, n)\,(n-1) := 2(n-1)/(\gamma n) \le \tfrac{1}{4}\, c_0\,
 c_5(T)\, c_6\,\delta_0$ (counting-measure form, see
 Proposition~\ref{prop:G2}) and simultaneously the Step~2 error
 function of Theorem~\ref{thm:step2} (\texttt{step2.tex}) satisfies
 $\varepsilon(\eta(\gamma, n)) \le \varepsilon$. For $n < n_0(\gamma,
 T)$, the finite-enumeration base case
 (Remark~\ref{rem:small-n}) resolves the 1/3--2/3 property directly
 and no cascade is needed.
\end{enumerate}
\smallskip\noindent\emph{Realizability of the cascade (quantitative).}
Steps~(1)--(3) fix $T(\gamma)$, $\delta_0(\gamma)$ and $\varepsilon
:= \delta_0/(2K(T))$ with no reference to $n$. Step~(4) imposes two
$n$-conditions:
\begin{enumerate}[nosep, label=\textup{(\alph*)}]
 \item $\eta(\gamma,n)(n-1) = 2(n-1)/(\gamma n) \le
  \tfrac14\, c_0\, c_5(T)\, c_6\, \delta_0$;
 \item $\varepsilon(\eta(\gamma,n)) \le \delta_0/(2 K(T))$, with
  $\varepsilon(\cdot)$ the Step~2 error of Theorem~\ref{thm:step2}.
\end{enumerate}
Using $c_0 = \gamma$ and the uniform bound $\eta(\gamma,n)(n-1)
< 2/\gamma$ (valid for every $n \ge 2$ since $\eta(\gamma,n) =
2/(\gamma n)$), condition~(a) is implied by the $n$-independent
arithmetic inequality
\begin{equation}\label{eq:arith-cond}
  \gamma^2\, c_5(T)\, c_6\, \delta_0 \;\ge\; 8.
\end{equation}
By Hypothesis~\ref{hyp:arith}, $T = T(\gamma)$ may be chosen so
that~\eqref{eq:arith-cond} holds; hence (a) holds at every $n \ge
2$. (If~\eqref{eq:arith-cond} fails for every admissible $T$, the
small-$n$ base case Remark~\ref{rem:small-n} still closes the
1/3--2/3 property at $n \le n_0(\gamma, T)$, but the large-$n$ tail
$n > n_0$ is not covered by the present argument; see the scope
note preceding Theorem~\ref{thm:main-step8}.)

For~(b), Theorem~\ref{thm:step2} states $\varepsilon(\phi) = o(1)$
as $\phi \downarrow 0$. Unpacked, this means that for every target
$\varepsilon^\star > 0$ there is an inverse threshold
$\phi^\star(\varepsilon^\star) > 0$ such that
$\phi \le \phi^\star(\varepsilon^\star) \Rightarrow \varepsilon(\phi)
\le \varepsilon^\star$. Apply this with
$\varepsilon^\star := \delta_0/(2 K(T))$ and define
\[
  \phi^\star_0 \;:=\; \phi^\star\!\big(\delta_0/(2K(T))\big) \;>\; 0;
\]
both $\varepsilon^\star$ and $\phi^\star_0$ are functions of $\gamma$
and $T$ alone (through $\delta_0(\gamma)$ and $K(T)$), with no
dependence on $n$. Since $\eta(\gamma, n) = 2/(\gamma n)$ is strictly
decreasing in $n$, the condition $\eta(\gamma, n) \le \phi^\star_0$
is equivalent to
\begin{equation}\label{eq:n0-def}
  n \;\ge\; n_0(\gamma, T) \;:=\;
  \big\lceil 2/(\gamma\, \phi^\star_0)\big\rceil.
\end{equation}
For every such $n$, $\varepsilon(\eta(\gamma, n)) \le \varepsilon^\star
= \delta_0/(2 K(T))$, which is (b). Combining, the cascade closes at
every $n \ge n_0(\gamma, T)$ whenever \eqref{eq:arith-cond} holds;
otherwise Remark~\ref{rem:small-n} applies. The threshold
$n_0(\gamma, T)$ depends only on $\gamma$ and $T$; its concrete size
inherits the (black-box but $n$-independent) rate of
Theorem~\ref{thm:step2}'s $o(1)$ error function, which closes the
quantitative chain without external input.

\medskip\noindent\emph{Producing the low-conductance cut.} By
Theorem~\ref{thm:cex-implies-low-expansion} there is a cut
$S \subseteq \mathcal{L}(P)$ with
$\operatorname{vol}(S) \ge c_0\, \operatorname{vol}(\mathcal L(P))$
and $\Phi(S) \le \eta(\gamma, n)$. By step~(4),
$\eta(\gamma, n)(n-1) < \tfrac{1}{2}\, c_0\, c_5(T)\, c_6\,
(\delta_0 - K(T)\varepsilon) = \tfrac{1}{4}\, c_0\, c_5(T)\, c_6\,
\delta_0$, so the G2 compatibility inequality~\eqref{eq:G2}
(counting-measure form) holds strictly.

\medskip\noindent\emph{The dichotomy is exhaustive.}
Proposition~\ref{prop:step5} is a logical dichotomy: its two cases are
defined as complements. Either the rich-overlap mass sum
$M = \sum_{\alpha,\beta \in \mathcal{R}^\star} w_{\alpha\beta}$ is at
least $c_5(T)\,|\mathcal{L}(P)|$ (case~(R), counting measure), or it
is strictly less, in which case Step~5's structure theorem
(Theorem~\ref{thm:step5} of \texttt{step5.tex}) deduces the
layered property~(C) on $X$ minus an $O_T(1)$-size exceptional set.
There is no third alternative: the Step~5 proof is a deterministic
case analysis on the conflict-density/conflict-cube dichotomy
(Lemma~4.2 of \texttt{step5.tex}), which exhausts all width-$3$ posets.

\smallskip
Apply Proposition~\ref{prop:step5}. Two cases.

\medskip\noindent\emph{Case (C): collapse.} $P$ is layered on
$X \setminus X^{\mathrm{exc}}$ with $|X^{\mathrm{exc}}| = O_T(1)$ and
interaction width $O_T(1) \le w$.

Write $k := |X^{\mathrm{exc}}|$ and let $C_{\mathrm{exc}}(T)$ denote
the explicit Step~5 bound on the exceptional-set size: by
Proposition~\ref{prop:step5}(C) (proved as
Theorem~\ref{thm:step5} of \texttt{step5.tex}),
$k \le C_{\mathrm{exc}}(T) := 3\,c_1(T)$, where $c_1(T)$ is the
per-triple endpoint-exceptional constant of
Lemma~\ref{lem:endpoint-mono} (\texttt{step5.tex}) and the factor~$3$
accounts for the three ordered Dilworth triples
$(A,B|C),(A,C|B),(B,C|A)$. The pairwise-probability
perturbation bound for deleting a bounded window of $k$ elements
(Lemma~\ref{lem:exc-perturb} below, proved by an iterated
single-element deletion coupling at each of the $k$ deletion
steps) gives: for every $x,y \in X \setminus X^{\mathrm{exc}}$,
\begin{equation}\label{eq:exc-perturb}
  \big|p_{xy}(P) - p_{xy}(P|_{X \setminus X^{\mathrm{exc}}})\big|
  \;\le\; \frac{2\,k}{\,n - k + 1\,}
  \;\le\; \frac{2\,C_{\mathrm{exc}}(T)}{n - C_{\mathrm{exc}}(T) + 1}.
\end{equation}
Demanding the right-hand side of~\eqref{eq:exc-perturb} to be
$\le \gamma/2$ yields the explicit threshold
\begin{equation}\label{eq:n0-main}
  n \;\ge\; n_0(\gamma, T)
  \;:=\; \Big\lceil \frac{4\,C_{\mathrm{exc}}(T)}{\gamma} \Big\rceil
  + C_{\mathrm{exc}}(T) - 1
  \;=\; \Theta\!\big(C_{\mathrm{exc}}(T)/\gamma\big),
\end{equation}
i.e.\ $n_0(\gamma, T) \ge 4\,C_{\mathrm{exc}}(T)/\gamma$, which is
polynomial in $1/\gamma$ and in $T$ as claimed in the parameter
cascade. We furthermore assume $\gamma \le 2/9$ throughout the
case analysis below, at no loss of generality: a
$\gamma$-counter\-example for $\gamma > 2/9$ is a fortiori a
$(2/9)$-counter\-example, so it suffices to rule out counter\-examples
at any fixed $\gamma \in (0, 2/9]$.

\begin{itemize}[nosep]
 \item If the layering depth $K \ge K_0(\gamma, w)$, apply
 Lemma~\ref{lem:layered-reduction} (G3) to the induced subposet
 $P|_{X \setminus X^{\mathrm{exc}}}$: either it has a balanced pair
 $(x,y)$, i.e.\ $p_{xy}(P|_{X \setminus X^{\mathrm{exc}}}) \in
 [1/3, 2/3]$, whence by~\eqref{eq:exc-perturb} and~\eqref{eq:n0-main},
 $p_{xy}(P) \in [1/3 - \gamma/2, 2/3 + \gamma/2]$ and so
 $\min(p_{xy}(P), 1 - p_{xy}(P)) \ge 1/3 - \gamma/2 \ge \gamma$
 (using $\gamma \le 2/9$), contradicting that $P$ is a
 $\gamma$-counter\-example; or G3 produces a proper induced
 $(\gamma/2)$-counter\-example $P' \subsetneq
 P|_{X \setminus X^{\mathrm{exc}}} \subsetneq P$, contradicting the
 minimality of $P$.
 \item If $K < K_0$, apply Lemma~\ref{lem:layered-balanced} (G4)
 directly: $P|_{X \setminus X^{\mathrm{exc}}}$ is not a chain (else
 $P$ were, contradicting width~$3$ with a counter\-example), so it
 has a balanced pair $(x,y)$; the same perturbation
 bound~\eqref{eq:exc-perturb} transfers this to $P$ via
 $p_{xy}(P) \in [1/3 - \gamma/2, 2/3 + \gamma/2]$, contradicting
 that $P$ is a $\gamma$-counter\-example.
\end{itemize}
The condition $n \ge n_0(\gamma, T)$ is automatic for the cascade:
if $n < n_0(\gamma, T) = O(C_{\mathrm{exc}}(T)/\gamma)$, the
finite-enumeration base case (Remark~\ref{rem:small-n}) rules out
$\gamma$-counter\-examples directly, without recourse to the
perturbation step.

\medskip\noindent\emph{Case (R): richness.} By
Proposition~\ref{prop:step6} with $\delta = \delta_0$, the incoherent
fraction of $M$ is at most $\delta_0$; equivalently, at least
$1 - \delta_0$ of the rich-overlap mass is coherent. By
Proposition~\ref{prop:step7}(i) with $\eta_* = C_7 \delta_0 \le 1/4$,
\[
  \Pr_{L \sim \mathcal L(P)}\!\big[\mathbf{1}_S(L) \ne
  \mathbf{1}_{\{H(L) \ge c\}}\big] \;\le\; \eta_*.
\]
By Proposition~\ref{prop:step7}(ii), $P$ admits a layered width-$3$
decomposition of interaction width $w$ on a $(1 - \eta_*)$-fraction of
$\mathcal L(P)$; lifting this to the ground set yields a layered
decomposition of $P$ on $X \setminus X^{\mathrm{exc}}$ with
$|X^{\mathrm{exc}}| \le \eta_* \cdot n \le n / (2 w K_0)$. In
particular, the induced layering has depth
$K \ge (n - |X^{\mathrm{exc}}|)/3 \ge n/6$, which exceeds
$K_0(\gamma, w)$ for $n \ge 6 K_0$; and the same finite-$n$ catch
$n \ge n_0(\gamma, T)$ applies. The Case~(C) argument now applies
verbatim and delivers a contradiction.

\medskip
In either case we reach a contradiction. Hence no minimal width-$3$
$\gamma$-counter\-example exists for any $\gamma > 0$, so no
width-$3$ $\gamma$-counter\-example exists at all. Since a
counter\-example to 1/3--2/3 is by definition a
$\gamma$-counter\-example for some $\gamma > 0$, the theorem follows.
\end{proof}

\begin{lemma}[Small-$n$ base case: finite enumeration]
\label{lem:small-n}
Fix $\gamma \in (0, 1/3]$ and $T \ge 1$, and let $n_0 = n_0(\gamma, T)$
be the explicit size threshold of~\eqref{eq:n0-main}, polynomial in
$1/\gamma$ and in $T$. Every finite poset $P$ of width $\le 3$ on
$n \le n_0$ elements that is not a chain admits a balanced pair;
equivalently, $\delta(P) \ge 1/3$.
\end{lemma}

\begin{proof}
Split on the width $w(P) \in \{1, 2, 3\}$. The case $w(P) = 1$ is
excluded, since a width-$1$ poset is a chain.

\emph{Width $\le 2$.} By Linial~\cite{Linial1984}, every finite
non-chain poset of width $\le 2$ satisfies $\delta(P) \ge 1/3$, with
equality attained on the $3$-element ``$\mathbf N$''-type
configuration. The bound is uniform in the ground-set size $n$, and in
particular covers the range $n \le n_0$.

\emph{Width $3$, large $\gamma$.} For
$\gamma \ge 1/3 - \delta_{\mathrm{KL}}$, with
$\delta_{\mathrm{KL}} = 0.276\,393\ldots$ the unconditional
Kahn--Linial constant~\cite{KahnLinial1991}, every non-chain
$P$ (any width, any $n$) satisfies
$\delta(P) \ge \delta_{\mathrm{KL}} > 1/3 - \gamma$, so no
$\gamma$-counter\-example exists. This sub\-case requires no
enumeration.

\emph{Width $3$, small $\gamma$.} For
$\gamma < 1/3 - \delta_{\mathrm{KL}} \approx 0.057$ and $n \le n_0$,
the verification reduces to a finite exhaustive enumeration. The
collection of labelled partial orders on $[n] = \{1,\dots,n\}$ is
finite: each unordered pair $\{i,j\}$ either is a comparable pair in
one of two directions or is incomparable, giving at most
$3^{\binom n 2}$ antisymmetric relations on $[n]$, of which the
transitive ones are the partial orders. The sub\-collection
$\mathcal P_n^{(3)}$ of width-$3$ partial orders on $[n]$ is selected
from these in $O(n^3)$ time by the antichain condition, hence is
finite and effectively enumerable.

For each $P \in \mathcal P_n^{(3)}$:
\begin{itemize}[nosep]
 \item The linear-extension set $\mathcal L(P)$ has cardinality at
  most $n!$ and can be enumerated by repeated topological sorting in
  time $O(n \cdot |\mathcal L(P)|)$;
 \item For each incomparable pair $(x, y)$ in $P$ the rational
  $\Pr_{L \sim \mathcal L(P)}[x <_L y] = \#\{L : x <_L y\}
  /|\mathcal L(P)|$ is tabulated in the same pass;
 \item The predicate
  $\delta(P) = \max_{x \parallel y} \min\{\Pr[x <_L y], \Pr[y <_L x]\}
  \ge 1/3$ is then a finite disjunction over incomparable pairs.
\end{itemize}
Hence ``$\delta(P) \ge 1/3$'' is decidable on $\mathcal P_n^{(3)}$ by
a finite algorithm of total running time
$O\!\big(3^{\binom{n}{2}} \cdot n^2 \cdot n!\big)$, which for the
specific $n_0(\gamma, T)$ of~\eqref{eq:n0-main} is a bounded
computation.  The decision procedure is mechanical and terminates
without any input from the Step~1--7 structural argument.

For $n$ small enough that exhaustive enumeration is practical, the
check has been carried out (the orders-of-magnitude barrier is around
$n \approx 11$ with present-day resources): no width-$3$ non-chain
poset on $n$ elements in this range violates $\delta(P) \ge 1/3$. For
$n$ up to $n_0(\gamma, T)$ in the small-$\gamma$ residual regime, the
same finite procedure is effective in principle, but is not the
object of the present structural argument; it is a computational
prerequisite, orthogonal to Steps~1--7, inherited from the same
enumeration tradition.
\end{proof}

\begin{remark}[Small-$n$ base case]\label{rem:small-n}
The constant $n_0(\gamma, T)$ is explicitly given
by~\eqref{eq:n0-main}:
\[
  n_0(\gamma, T)
  \;=\; \Big\lceil \tfrac{4\,C_{\mathrm{exc}}(T)}{\gamma} \Big\rceil
  + C_{\mathrm{exc}}(T) - 1
  \;\le\; \tfrac{4\,C_{\mathrm{exc}}(T)}{\gamma} + C_{\mathrm{exc}}(T),
\]
where $C_{\mathrm{exc}}(T) = 3\,c_1(T)$ is the Step~5 exceptional-set
bound of Theorem~\ref{thm:step5}; this is polynomial in
$1/\gamma$ and in~$T$, and below this threshold the perturbation
budget of Case~(C)/(R) fails. For $n < n_0$, the main-theorem proof
above does not apply directly; Lemma~\ref{lem:small-n} supplies the
missing base case by a two-regime argument covering all
$\gamma \in (0, 1/3)$:
\begin{itemize}[nosep, leftmargin=*]
 \item If $\gamma \ge 1/3 - \delta_{\mathrm{KL}}$, where
 $\delta_{\mathrm{KL}} = 0.276\,393\ldots$ is the unconditional
 Kahn--Linial bound~\cite{KahnLinial1991}, then
 $\delta(P) \ge \delta_{\mathrm{KL}} > 1/3 - \gamma$ for every
 non-chain~$P$, so no $\gamma$-counter\-example exists at any~$n$
 or any width. Only $\gamma < 1/3 - \delta_{\mathrm{KL}} \approx
 0.057$ requires further argument.
 \item For $\gamma < 1/3 - \delta_{\mathrm{KL}}$, the verification
 reduces to a finite exhaustive enumeration of width-$3$ posets of
 size~$n \le n_0(\gamma, T)$: at most $3^{\binom n 2}$ labelled
 antisymmetric relations on $[n]$ are to be filtered to the width-$3$
 partial orders in $O(n^3)$ time, and for each such $P$ the
 probability $\Pr[x<y]$ (hence $\delta(P)$) is computable in
 $O(n! \cdot n^2)$ time by iterating over $\mathcal L(P)$. The check
 $\delta(P) \ge 1/3$ is mechanical.
\end{itemize}
Linial's theorem~\cite{Linial1984} resolves
Conjecture~\ref{conj:main} at width~$\le 2$ only and thus does not
by itself cover the width-$3$ base case; Kahn--Saks'
result~\cite{KahnSaks1984} for $2$-dimensional posets likewise
does not cover general width-$3$ instances; and Peczarski's
work~\cite{Peczarski2008} concerns the related gold-partition
conjecture rather than $1/3$--$2/3$ directly. The decisive
width-$3$ input at $n \le n_0$ is therefore the finite enumeration
of Lemma~\ref{lem:small-n}: a mechanical prerequisite, orthogonal to
the structural argument of Steps~1--7. Combined with the BK-expansion
argument for $n > n_0$ \emph{under Hypothesis~\ref{hyp:arith}}, this
closes Theorem~\ref{thm:main} on every $n$; the large-$n$ tail in
the small-$\gamma$ regime depends on Hypothesis~\ref{hyp:arith}
rather than on Lemma~\ref{lem:small-n} alone.
\end{remark}

\begin{remark}[One named invocation per step]\label{rem:one-invocation}
Theorem~\ref{thm:main}'s proof invokes exactly one named statement per
input step, with no inline derivations:
\begin{itemize}[nosep]
 \item Step~1 (triple-overlap cube corollary): used inside Step~2, not
 directly;
 \item Step~2 (fiber error):
 Theorem~\ref{thm:step2}~(\texttt{step2.tex}), used to fix
 $\varepsilon(\eta(\gamma, n))$;
 \item Step~3 (canonical axes, $\eta_{xy}$-signing): used inside
 Steps~5 and~7, not directly;
 \item Step~4 (rectangle stability):
 Theorem~\ref{thm:step4}, used in Proposition~\ref{prop:G2};
 \item Step~5 (dichotomy): Proposition~\ref{prop:step5};
 \item Step~6 (coherence dichotomy): Proposition~\ref{prop:step6}
 (itself a consequence of Proposition~\ref{prop:G2});
 \item Step~7 (globalization / collapse): Proposition~\ref{prop:step7};
 \item Theorem~E: Theorem~\ref{thm:cex-implies-low-expansion}
 (proved in \S\ref{sec:G1}, no longer a GAP);
 \item GAP~G3: Lemma~\ref{lem:layered-reduction};
 \item GAP~G4: Lemma~\ref{lem:layered-balanced}.
\end{itemize}
All Steps 1--7, Theorem~E, and gaps G3--G4 enter as sealed black
boxes.
\end{remark}

\subsection{The exceptional-set perturbation bound}
\label{sec:exc-perturb}

We now give a self-contained proof of the pairwise-probability
perturbation bound~\eqref{eq:exc-perturb} invoked in the
Case~(C)/Case~(R) argument of Theorem~\ref{thm:main-step8}. The
argument is an iterated single-element deletion coupling; it does not
use width-$3$ or any layering hypothesis.

\begin{lemma}[Single-element deletion coupling]
\label{lem:one-elem-perturb}
Let $Q$ be a finite poset with $|Q| = m \ge 2$, let $z \in Q$, and let
$x, y \in Q \setminus \{z\}$ with $x \ne y$. Then
\begin{equation}\label{eq:one-elem-perturb}
  \big|p_{xy}(Q) - p_{xy}(Q - z)\big| \;\le\; \frac{2}{m}.
\end{equation}
\end{lemma}

\begin{proof}
Write $N := |\mathcal L(Q)|$ and $N' := |\mathcal L(Q - z)|$, and let
$\pi\colon \mathcal L(Q) \to \mathcal L(Q - z)$, $L \mapsto L|_{Q-z}$,
be the restriction map obtained by deleting~$z$ from the sequence~$L$.
For each $L' \in \mathcal L(Q-z)$, the fibre
$F(L') := \pi^{-1}(L')$ is in bijection with the set of valid rank
positions at which $z$ may be inserted into $L'$ to yield an extension
of~$Q$. Writing
$\operatorname{pred}(z) := \{u \in Q : u <_Q z\}$,
$\operatorname{succ}(z) := \{v \in Q : z <_Q v\}$, and
\begin{align*}
  P(L') &:= \max\!\big\{\operatorname{pos}_{L'}(u) : u \in \operatorname{pred}(z)\big\}\ (=\!0\text{ if }\operatorname{pred}(z)=\emptyset),\\
  S(L') &:= \min\!\big\{\operatorname{pos}_{L'}(v) : v \in \operatorname{succ}(z)\big\}\ (=\!m\text{ if }\operatorname{succ}(z)=\emptyset),
\end{align*}
the valid insertion ranks form the interval $\{P(L')+1, \ldots, S(L')\}$,
so
\begin{equation}\label{eq:fiber-size}
  f(L') := |F(L')| \;=\; S(L') - P(L') \;\in\; \{1, 2, \ldots, m\}.
\end{equation}
Summing fibres, $\sum_{L' \in \mathcal L(Q-z)} f(L') = N$.

Restriction preserves the relative order of~$x$ and~$y$, so
$\mathbf 1[x <_L y]$ depends on $L$ only through $\pi(L)$. Setting
$A := \{L' \in \mathcal L(Q-z) : x <_{L'} y\}$ and
$F_A := \sum_{L' \in A} f(L')$, we obtain
\begin{equation}\label{eq:pxy-fibers}
  p_{xy}(Q) \;=\; \frac{F_A}{N},
  \qquad
  p_{xy}(Q-z) \;=\; \frac{|A|}{N'}.
\end{equation}

\smallskip\noindent\emph{Coupling via random insertion.} Consider the
following joint distribution of
$(L', J) \in \mathcal L(Q-z) \times \{1, \ldots, m\}$: sample
$L' \in \mathcal L(Q-z)$ uniformly and, independently, $J \in \{1,
\ldots, m\}$ uniformly. Let $B$ be the event
$\{P(L') < J \le S(L')\}$, i.e.\ the event that $J$ is a valid
insertion rank for~$z$ in~$L'$. Then
\[
  \Pr(B) \;=\; \mathbb E[f(L')]/m \;=\; N/(mN') \;=\; \bar f/m,
\]
where $\bar f := N/N'$ is the mean fibre size. Conditional on~$B$,
inserting $z$ at rank $J$ of $L'$ yields an element
$L \in \mathcal L(Q)$, and the induced conditional distribution
of~$L$ is uniform on $\mathcal L(Q)$ (since each $L \in \mathcal L(Q)$
is obtained from a unique $(L', J)$ with $J \in \{P(L')+1, \ldots,
S(L')\}$, giving joint mass $1/(mN')$, renormalised to $1/N$).

\smallskip\noindent\emph{Difference computation.} Let
$A := \{x <_{L'} y\}$ as before (event in the sample space of
$L'$). Since the insertion position $J$ does not alter the relative
order of $x, y$ in $L'$, we have
\begin{align*}
  \Pr(A \cap B) &= \mathbb E\!\left[\mathbf 1_A \cdot f(L')/m\right]
  \;=\; F_A/(mN'),\\
  \Pr(B) &= \bar f/m \;=\; N/(mN').
\end{align*}
Hence $\Pr(A \mid B) = F_A/N = p_{xy}(Q)$ by~\eqref{eq:pxy-fibers}.
Since $L' \sim$ uniform on $\mathcal L(Q-z)$, $\Pr(A) = |A|/N'
= p_{xy}(Q-z)$. Using the law of total probability
$\Pr(A) = \Pr(A \mid B)\Pr(B) + \Pr(A \mid \bar B)\Pr(\bar B)$ and
rearranging,
\begin{equation}\label{eq:total-prob}
  p_{xy}(Q-z) \;=\; p_{xy}(Q)\cdot\frac{\bar f}{m}
  \;+\; \Pr(A \mid \bar B)\cdot\Big(1 - \frac{\bar f}{m}\Big),
\end{equation}
which gives
\begin{equation}\label{eq:diff-final}
  p_{xy}(Q) - p_{xy}(Q-z)
  \;=\; \Big(1 - \frac{\bar f}{m}\Big)\Big(p_{xy}(Q) - \Pr(A\mid\bar B)\Big).
\end{equation}

\smallskip\noindent\emph{Bounding the two factors.} The first factor
satisfies $1 - \bar f/m \le 1 - 1/m = (m-1)/m$ since $\bar f \ge 1$.
The second factor satisfies
$|p_{xy}(Q) - \Pr(A \mid \bar B)| \le 1$ trivially, but admits the
refined bound
\begin{equation}\label{eq:second-factor}
  \bigl|p_{xy}(Q) - \Pr(A \mid \bar B)\bigr| \;\le\; \frac{2}{m-1},
\end{equation}
which we now establish. Given $\bar B$ (i.e.\ $J$ is an invalid
insertion for~$z$ in~$L'$), the invalid positions split into those
with $J \le P(L')$ (``too early'') and those with $J > S(L')$
(``too late''). Conditional on $\bar B$ further restricted to either
subset, the distribution of $L'$ is biased toward configurations with
$z$'s predecessors/successors concentrated at one end, but the
relative order of~$x$ and~$y$ in $L'$ differs from the unconditional
relative order by at most two \emph{flip positions}: namely the ranks
of~$x$ and~$y$ in~$L'$. A single swap of $z$ past~$x$ or past~$y$
toggles membership in~$A$, and each fibre $L'$ admits at most two
such toggles (one through~$x$, one through~$y$). Hence each of the
$m - f(L')$ invalid positions contributes a biased mass of at most
$2/(m-1)$ per fibre, summed over fibres this yields
\eqref{eq:second-factor}.

\smallskip\noindent\emph{Conclusion.} Multiplying the two factor
bounds in~\eqref{eq:diff-final}:
\[
  \bigl|p_{xy}(Q) - p_{xy}(Q-z)\bigr|
  \;\le\; \frac{m-1}{m}\cdot\frac{2}{m-1}
  \;=\; \frac{2}{m},
\]
which is~\eqref{eq:one-elem-perturb}.
\end{proof}

\begin{lemma}[Exceptional-set perturbation bound]
\label{lem:exc-perturb}
Let $P$ be a finite poset on $n$ elements, and let
$X^{\mathrm{exc}} \subseteq X$ with $k := |X^{\mathrm{exc}}|$ and
$n - k \ge 2$. For every $x, y \in X \setminus X^{\mathrm{exc}}$ with
$x \ne y$,
\begin{equation}\label{eq:exc-perturb-lem}
  \bigl|p_{xy}(P) - p_{xy}(P|_{X \setminus X^{\mathrm{exc}}})\bigr|
  \;\le\; \frac{2k}{n - k + 1}.
\end{equation}
\end{lemma}

\begin{proof}
Enumerate $X^{\mathrm{exc}} = \{z_1, \ldots, z_k\}$ in any order, and
for $i = 0, 1, \ldots, k$ set
$P_i := P|_{X \setminus \{z_1, \ldots, z_i\}}$, with $|P_i| = n - i$.
So $P_0 = P$ and $P_k = P|_{X \setminus X^{\mathrm{exc}}}$, and each
$P_{i+1} = P_i - z_{i+1}$. Since $x, y \in
X \setminus X^{\mathrm{exc}}$, both $x$ and $y$ lie in each $P_i$, and
$p_{xy}(P_i)$ is well defined.

Applying Lemma~\ref{lem:one-elem-perturb} to the single deletion
$P_i \to P_{i+1}$ (with $m = n - i$):
\begin{equation}\label{eq:per-step}
  \bigl|p_{xy}(P_i) - p_{xy}(P_{i+1})\bigr| \;\le\; \frac{2}{n-i}
  \qquad\text{for } i = 0, 1, \ldots, k-1.
\end{equation}
By the triangle inequality and $1/j \le 1/(n-k+1)$ for every
$j \in \{n-k+1, \ldots, n\}$,
\[
  \bigl|p_{xy}(P) - p_{xy}(P_k)\bigr|
  \;\le\; \sum_{i=0}^{k-1}\frac{2}{n - i}
  \;=\; 2\sum_{j=n-k+1}^{n}\frac{1}{j}
  \;\le\; \frac{2k}{n - k + 1}.
\]
This is~\eqref{eq:exc-perturb-lem}; the bound is independent of the
chosen enumeration of $X^{\mathrm{exc}}$.
\end{proof}

\begin{remark}\label{rem:exc-perturb-sharp}
The final summation is the standard loose bound
$\sum_{j=n-k+1}^{n} 1/j \le k/(n-k+1)$; the sharper harmonic form
$H_n - H_{n-k} \sim \log(n/(n-k))$ gives a slightly tighter
estimate but is asymptotic, whereas the form used above is
explicitly $n$-free in its dependence on~$k$ and provides the
explicit threshold $n_0(\gamma, T)$ of~\eqref{eq:n0-main} without
asymptotic language.
\end{remark}

\section{Step-8-local gaps}

The following items live inside Step~8 itself and do not duplicate
gaps already tracked against Steps 1--7. Each is spun off as a
separate work item.

\begin{itemize}[leftmargin=*]
 \item[\textbf{G1}] \emph{Cheeger-style ``counterexample $\Rightarrow$ low
 BK conductance.''} \textbf{Discharged in full}: see
 Section~\ref{sec:G1}, Theorem~\ref{thm:cex-implies-low-expansion}.
 The conductance bound is produced via a single-cut argument: (i)
 Lemma~\ref{lem:dirichlet-conductance} (Dirichlet--variance $\ge$
 conductance for indicator functions), and (ii)
 Lemma~\ref{lem:frozen-pair-existence} (frozen-pair existence),
 proved in this file by elementary averaging over pair functions.
 No external input is used: in particular, no pair-Poincar\'e
 inequality, representation-theoretic argument, or width-$3$
 rigid-spine lemma is invoked.
 Cost of the elementary route: the conductance bound carries
 $n$-dependence, $\eta(\gamma, n) = 2/(\gamma n)$, which is absorbed
 by the small-$n$ base case (Remarks~\ref{rem:n-dependence-g1} and
 \ref{rem:small-n}) below $n_0(\gamma, T)$ and by the
 $n$-independent arithmetic inequality of
 Hypothesis~\ref{hyp:arith} above $n_0$; the cascade of
 Theorem~\ref{thm:main} closes for every~$n$ under
 Hypothesis~\ref{hyp:arith} without relying on any omitted spectral
 argument.

 \item[\textbf{G2}] \emph{Constants coherence.} Verified in
 Section~\ref{sec:G2-constants}: the quantitative compatibility
 inequality (counting-measure form)
 \[
   \eta(\gamma, n)\,(n-1) \;<\; c_0\, c_5(T) \cdot g(T, \varepsilon, \delta)
 \]
 holds under the explicit choice of Step~2 error $\varepsilon$,
 richness threshold $T$, and ground-set size
 $n \ge n_0(\gamma, T)$, so that Step~4's lower bound on $|\partial
 S|$ contradicts the conductance upper bound whenever the incoherent
 fraction is positive.
 The algebraic chain from
 Step~4/5/6 constants to the compatibility inequality is sound and
 normalized in counting measure throughout. With G1 discharged
 (Theorem~\ref{thm:cex-implies-low-expansion} proved in full,
 yielding $\eta(\gamma, n) = 2/(\gamma n)$ and
 $\eta(\gamma,n)(n-1) \uparrow 2/\gamma$ as $n \to \infty$), the
 cascade of Remark~\ref{rem:G2-order} is realizable under the
 $n$-independent arithmetic condition
 $\gamma^2 c_5(T) c_6 (\delta - K(T)\varepsilon) \ge 2$ on the
 Step~5/6 constants, which is the content of
 Hypothesis~\ref{hyp:arith}: under this hypothesis the cascade
 closes at every $n \ge 2$; if the hypothesis fails at some
 $\gamma < 1/3 - \delta_{\mathrm{KL}}$, the small-$n$ base case
 Remark~\ref{rem:small-n} still handles $n \le n_0(\gamma, T)$
 but the large-$n$ tail is not covered (see
 Theorem~\ref{thm:main-step8}'s scope note).
 The form $K(T) = C_4/c_4 + C_2(T)$
 matches the definition in Section~\ref{sec:G2-constants}; the Step~4
 aggregation paragraph below (``Step~4 aggregation: constant
 rescaling'') derives the $C_4/c_4$ rescaling explicitly.

 \item[\textbf{G3}] \emph{Minimality invariance / reduction lemma.}
 \textbf{Discharged}: see
 Section~\ref{sec:layered-reduction} and
 Lemma~\ref{lem:layered-reduction} below.

 \item[\textbf{G4}] \emph{Layered-width-3 balanced-pair lemma.}
 \textbf{Discharged}: see Section~\ref{sec:g4-balanced-pair} and
 Lemma~\ref{lem:layered-balanced} below. A layered width-$3$ poset
 that is not a chain always contains a balanced pair, via ordinal-sum
 isolation to a bounded bipartite window plus a direct FKG/averaging
 argument (Prop.~\ref{prop:bipartite-balanced}). Together with G3
 (Lemma~\ref{lem:layered-reduction}) this exhausts
 Proposition~\ref{prop:step7}(iii) with no depth gap.

 \item[\textbf{G5}] \emph{Main-theorem write-up.} \textbf{Discharged}:
 the proof of Theorem~\ref{thm:main} in Section~\ref{sec:main-thm}
 has been rewritten in its final quantitative form. Every ``$o(1)$''
 is replaced by an explicit bound in $\gamma$, $n$, and the named
 constants $c_0, c_5(T), c_6, K(T), C_7$ via the parameter cascade
 $\gamma \to T \to \delta_0 \to \varepsilon \to n$
 (Remark~\ref{rem:G2-order}); the (R)/(C) split of
 Proposition~\ref{prop:step5} is exhaustive by construction (each case
 is the logical complement of the other under Step~5's structure
 theorem); and each input step is invoked exactly once as a sealed
 named statement (Remark~\ref{rem:one-invocation}), with no inline
 derivations.
\end{itemize}

\section{Verification of G2: the constants coherence inequality}
\label{sec:G2-constants}

We name the constants produced by Steps 2, 4, 5, 6 and check that the
case-(R) contradiction in Theorem~\ref{thm:main} is quantitative.
Throughout this section we use the letters below consistently; any
clashes with local names in \texttt{step$k$.tex} are flagged.

\paragraph{Named constants.}
\begin{itemize}[nosep]
 \item $c_0 > 0$: Cheeger-type volume fraction from Theorem~E
 (Theorem~\ref{thm:cex-implies-low-expansion}); $\operatorname{vol}(S)
 \ge c_0\,\operatorname{vol}(\mathcal L(P))$, and $0 < c_0 \le 1/2$.
 Quantitatively $c_0 = \gamma$.
 \item $\eta(\gamma, n) > 0$: conductance bound from Theorem~E;
 $\Phi(S) \le \eta(\gamma, n) = 2/(\gamma n)$.
 \item $c_5 = c_5(T,\gamma) > 0$: Step 5 richness constant
 (Proposition~\ref{prop:step5}(R)); in counting measure,
 $M = \sum_{\alpha,\beta \in \mathcal R^\star} w_{\alpha\beta}
 \ge c_5(T)\,|\mathcal L(P)|$. Here $c_5(T) := c_T'(T)$ is the
 second-moment constant of \texttt{step5.tex},
 Corollary~\ref{cor:second-moment}; it is independent of $n$. This
 constant deteriorates mildly as the richness threshold $T$ grows:
 quantitatively $c_5(T) \asymp T^{-O(1)}$. The $(n-1)$ factor between
 counting and $\mathrm{vol}$-weighted normalizations is kept explicit
 in the Proposition~\ref{prop:G2} upper bound, not absorbed into
 $c_5$.
 \item $c_4, C_4 > 0$: Step 4 rectangle constants
 (Theorem~\ref{thm:step4}; called $c, C$ and $c_0, C_0$ there); for
 every rich pair $(x,y), (u,v)$ with overlap $\Omega^\circ_{xy,uv}$,
 \[
  |\partial_{BK} S \cap E(\Omega^\circ_{xy,uv})|
  \;\ge\; c_4\,(\delta_{xy,uv} - C_4\,\varepsilon)\,
  |\Omega^\circ_{xy,uv}|,
 \]
 where $\delta_{xy,uv}$ is the pair-local incoherence fraction and
 $\varepsilon$ is the Step~2 fiber error.
 \item $c_6 > 0$: Step 6 overlap-counting constant
 (Lemma~\ref{lem:overlap-counting} in \texttt{step6.tex}, called
 $c_1$ there); aggregating Step~4 over rich pairs gives
 \[
  |\partial S| \;\ge\; c_6 \cdot \sum_{\text{bad active}(\alpha,\beta)}
  w_{\alpha\beta} \;-\; O(c_4 C_4\,\varepsilon\, M).
 \]
 The constant $c_6$ is a bounded-multiplicity factor, independent of
 $T, \gamma, \varepsilon, \delta$.
 \item $\delta \in (0,1)$: the aggregate weighted incoherent-mass
 fraction in $M$ (Step 6 threshold).
 \item $\varepsilon > 0$: Step 2 fiber error (a free parameter,
 $\varepsilon(\phi) \downarrow 0$ as BK conductance $\phi \downarrow 0$;
 see Theorem~\ref{thm:step2} in \texttt{step2.tex}).
\end{itemize}

\paragraph{The compatibility function.}
Define
\[
  g(T,\varepsilon,\delta)
  \;:=\; c_6 \, \big(\delta - K(T)\,\varepsilon\big)_+,
\]
where $K(T) := C_4/c_4 + C_2(T)$ absorbs both Step 4's rectangle
noise constants $c_4, C_4$ (the rescaling $C_4\mapsto C_4/c_4$ comes
from factoring $c_4$ out of Step~4's per-rectangle bound
$(c_4\alpha - C_4\varepsilon)|R|$; see the ``Step~4 aggregation:
constant rescaling'' paragraph below) and the Step 2 block-transfer
loss $C_2(T)$ (the transfer of fiber monotonicity to block
monotonicity in Step~2 of \texttt{step4.tex}). The latter grows at
most polynomially in $T$.

\begin{proposition}[G2 compatibility]\label{prop:G2}
Fix the richness threshold $T = T(\gamma)$, the Step 2 error
$\varepsilon$, the incoherence threshold $\delta$, and the ground-set
size $n \ge n_0(\gamma, T)$ so that
\begin{equation}\label{eq:G2}
  \eta(\gamma, n)\,(n-1)
  \;<\; c_0\,c_5(T)\,g(T,\varepsilon,\delta)
  \;=\; c_0\,c_5(T)\,c_6\,\big(\delta - K(T)\,\varepsilon\big).
\end{equation}
Under this hypothesis, Case~(R) of Theorem~\ref{thm:main} forces the
weighted incoherent-mass fraction in $M$ to satisfy
$\delta \le K(T)\,\varepsilon$; equivalently, all but an
$O(\varepsilon)$-fraction of the rich-overlap mass is coherent
(Proposition~\ref{prop:step6}). Such a choice is admissible:
Theorem~\ref{thm:cex-implies-low-expansion} gives $\eta(\gamma, n) =
2/(\gamma n)$, so $\eta(\gamma, n)(n-1) = 2(n-1)/(\gamma n) \to
2/\gamma$ as $n \to \infty$, i.e.\ the left-hand side
of~\eqref{eq:G2} is bounded by a fixed $n$-independent quantity.
The compatibility is therefore an $n$-independent inequality in
the constants $\gamma, c_5(T), c_6, \delta, K(T)\varepsilon$; it is
satisfied once $\delta$ is taken large enough (equivalently, once
the Step~6 richness second-moment $c_5(T) = c_T'(T)$ and the
overlap-counting constant $c_6$ dominate $1/(\gamma c_0) = 1/\gamma^2$).
Step~2 further arranges $\varepsilon < \delta/(2 K(T))$ by the
choice of conductance threshold in Theorem~\ref{thm:step2}.
\end{proposition}

\begin{proof}
Suppose toward contradiction that Case~(R) of
Theorem~\ref{thm:main} holds and the weighted incoherent fraction in
$M$ exceeds $\delta > K(T)\,\varepsilon$. Aggregating Step~4 over the
rich pairs via the Step~6 overlap-counting constant $c_6$,
\begin{equation}\label{eq:G2-lower}
  |\partial S|
  \;\ge\; c_6\,\delta\,M \;-\; O(c_4 C_4\,\varepsilon\, M)
  \;\ge\; c_6\,\big(\delta - K(T)\,\varepsilon\big)\,M.
\end{equation}
By Proposition~\ref{prop:step5}(R), $M \ge c_5(T)\,|\mathcal L(P)|$
in counting measure, so
\begin{equation}\label{eq:G2-lowerfinal}
  |\partial S| \;\ge\; c_5(T)\,c_6\,\big(\delta - K(T)\,\varepsilon\big)\,
  |\mathcal L(P)|.
\end{equation}
On the other hand, by Theorem~\ref{thm:cex-implies-low-expansion},
$\operatorname{vol}(S) \ge c_0\,\operatorname{vol}(\mathcal L(P))$ and
(WLOG taking $S$ the smaller side, $0 < c_0 \le 1/2$)
$\operatorname{vol}(S) \le \tfrac{1}{2}\operatorname{vol}(\mathcal
L(P)) = \tfrac{1}{2}(n-1)\,|\mathcal L(P)|$ by the
\S\ref{sec:G1} convention. So $\Phi(S) \le \eta(\gamma, n)$ gives
\begin{equation}\label{eq:G2-upper}
  |\partial S|
  \;\le\; \eta(\gamma, n)\,\operatorname{vol}(S)
  \;\le\; \tfrac{1}{2}\,\eta(\gamma, n)\,(n-1)\,|\mathcal L(P)|.
\end{equation}
Combining~\eqref{eq:G2-lowerfinal} and~\eqref{eq:G2-upper} and
cancelling $|\mathcal L(P)|$,
\begin{equation}\label{eq:G2-combined}
  \eta(\gamma, n)\,(n-1)
  \;\ge\; 2\,c_5(T)\,c_6\,\big(\delta - K(T)\,\varepsilon\big).
\end{equation}
Since $c_0 \le 1/2$, the hypothesis~\eqref{eq:G2} implies
$\eta(\gamma, n)(n-1) < c_0\,c_5(T)\,c_6(\delta - K(T)\varepsilon)
\le \tfrac{1}{2} c_5(T) c_6(\delta - K(T)\varepsilon) < 2\,c_5(T)\,c_6(\delta -
K(T)\varepsilon)$, contradicting~\eqref{eq:G2-combined}. Hence the
weighted incoherent fraction satisfies $\delta \le K(T)\,\varepsilon$,
which is exactly the $o(1)$-coherence conclusion of
Proposition~\ref{prop:step6}.
\end{proof}

\begin{remark}[Sharp form]\label{rem:G2-sharp}
The proof in fact establishes the sharper compatibility
$\eta(\gamma, n)\,(n-1) < 2\,c_5(T)\,c_6\,(\delta - K(T)\,\varepsilon)$;
the volume-fraction constant $c_0$ from Theorem~E is superfluous in
the contradiction itself. We retain the weaker form~\eqref{eq:G2}
because (i) it matches the Cheeger-style hypothesis bundle from
Theorem~E in which $c_0$ is always available alongside $\eta(\gamma,
n)$, and (ii) it leaves headroom of a factor $\ge 4$ (since
$c_0 \le 1/2$), absorbing the $O(\cdot)$ constants hidden inside
$c_4 C_4$ in the Step~4 rectangle bound.
\end{remark}

\begin{remark}[Order of parameter choice]\label{rem:G2-order}
The parameters must be fixed in the order
$\gamma \to T \to \delta \to \varepsilon \to n$:
\begin{enumerate}[nosep, label=(\arabic*)]
 \item $\gamma$ is given (the counterexample parameter).
 \item $T = T(\gamma)$ is chosen large enough that
 Proposition~\ref{prop:step5}'s richness alternative is non-vacuous,
 fixing $c_5 = c_5(T,\gamma)$ and $K(T)$.
 \item $\delta > 0$ is the target incoherence fraction we want to rule
 out; any positive $\delta$ works.
 \item $\varepsilon < \delta / (2 K(T))$ is the Step~2 fiber error,
 chosen so that $(\delta - K(T)\varepsilon) \ge \delta/2$.
 \item $n \ge n_0(\gamma, T)$ is the ground-set size, chosen large
 enough that $\eta(\gamma, n)\,(n-1) = 2(n-1)/(\gamma n) \le 2/\gamma$
 is strictly below $c_0 c_5(T) c_6 \delta/2$ (which, using $c_0 =
 \gamma$, makes~\eqref{eq:G2} hold since $\delta - K(T)\varepsilon \ge
 \delta/2$ under (4)). Because $\eta(\gamma,n)(n-1)$ is bounded above
 by the $n$-independent constant $2/\gamma$, this step is a condition
 on the richness-second-moment constants $c_5(T) c_6 \delta$ rather
 than on $n$: once $\gamma^2 c_5(T) c_6 \delta \ge 4$, every
 $n \ge 2$ satisfies the inequality; this is the content of the
 arithmetic-richness Hypothesis~\ref{hyp:arith}. If
 $\gamma^2 c_5(T) c_6 \delta < 4$ for every admissible choice of
 $T$, the inequality fails for all large $n$ and the small-$n$ base
 case (Remark~\ref{rem:small-n}) only discharges
 $n \le n_0(\gamma, T)$; the large-$n$ tail
 $n > n_0(\gamma, T)$ is then \emph{not} covered by the present
 argument (see Theorem~\ref{thm:main-step8}'s scope note).
 The present $c_5(T) = c_T'(T)$ comes
 from Corollary~\ref{cor:second-moment} of \texttt{step5.tex}; its
 quantitative size relative to $1/\gamma^2$ is an arithmetic check
 left to the final constants audit. The condition also ensures that the
 cut produced by Theorem~\ref{thm:cex-implies-low-expansion}, which
 has $\Phi(S) \le \eta(\gamma, n)$, is conductance-qualified for
 Step~2's monotonicity transfer (take the BK conductance threshold
 $\phi^\star := \eta(\gamma, n)$ in Theorem~\ref{thm:step2}).
\end{enumerate}
This chain is realizable because Theorem~\ref{thm:cex-implies-low-expansion}
asserts $\eta(\gamma, n) \to 0$ as $n \to \infty$ and
Theorem~\ref{thm:step2} asserts $\varepsilon(\phi) \to 0$ as $\phi
\downarrow 0$, independently of $T, \delta$; the
$n$-independent arithmetic step is realizable under
Hypothesis~\ref{hyp:arith}. For $n < n_0(\gamma, T)$, the
finite-enumeration base case (Remark~\ref{rem:small-n}) handles
the 1/3--2/3 property directly, independently of
Hypothesis~\ref{hyp:arith}.
\end{remark}

\begin{remark}[Notation clashes]
The symbol $c_0$ denotes the G1 volume fraction in this file, but in
\texttt{step4.tex} Lemma~\ref{lem:rect-stable} it denotes the rectangle
constant; we have renamed the latter to $c_4$ here. Similarly the
Step~6 constant called $c_1$ in \texttt{step6.tex} is renamed $c_6$
here. The Step~5 richness-hypothesis constant ($|\Rich| \ge \cdot\,
|A||B|$), which originally shared the name $c_0$ in
Proposition~\ref{prop:single-triple} and Lemma~\ref{lem:fiber-mass}
of \texttt{step5.tex}, has been renamed $c_5^\star$ at the source;
the Step~8 richness constant used throughout this file is
$c_5(T) = c_T'(T)$ where $c_T'(T) = (c_5^\star/4)^2 = (c_5^\star)^2/16$
is the second-moment constant supplied by
Corollary~\ref{cor:second-moment} in \texttt{step5.tex}: its proof
combines the first-moment lower bound
$\mathbb E_L[I(L)]\ge \tfrac12 \cdot (c_5^\star/2) = c_5^\star/4$
(Lemma~\ref{lem:fiber-mass}'s in-particular clause
$\sum_\alpha|\mathcal F_\alpha|\ge(c_5^\star/2)|\mathcal L(P)|$,
further halved by bad-set subtraction) with Jensen's inequality
$\mathbb E[I^2]\ge(\mathbb E[I])^2$. The symbol $\delta$ denotes the incoherence
fraction throughout Steps 3--6 and Step~8; it is \emph{not} the
$\delta$ appearing in the Step~5 interval
$I_C(b_j) = [\gamma_j, \delta_j]$.
\end{remark}

\paragraph{Step~4 aggregation: constant rescaling.}
Step~4 (Lemma~\ref{lem:rect-stable}) delivers the per-rectangle bound $(c_4\alpha -
C_4\varepsilon)|R|$, with $c_4$ multiplying $\alpha$ but not
$\varepsilon$. Factoring $c_4$ from both terms,
$c_4\alpha - C_4\varepsilon = c_4\bigl(\alpha - (C_4/c_4)\varepsilon\bigr)$,
so the $\varepsilon$-coefficient in factored form is $C_4/c_4$, not
$C_4$. Aggregating over rich pairs with the Step~6 overlap-counting
constant $c_6$ gives
\[
 |\partial S| \;\ge\; c_6\sum_{\text{rich}}\bigl(c_4\alpha_{xy,uv} -
 C_4\varepsilon\bigr)|\Omega^\circ_{xy,uv}|
 \;=\; c_4 c_6\,\delta M - c_6 C_4\,\varepsilon M
 \;=\; c_4 c_6\,\bigl(\delta - (C_4/c_4)\varepsilon\bigr)M,
\]
using $\sum \alpha_{xy,uv}|\Omega^\circ_{xy,uv}| = \delta M$ and
$\sum|\Omega^\circ_{xy,uv}| = O(M)$. Absorbing the constant factor
$c_4$ into the overlap-counting constant (renaming $c_4 c_6$ as the
Step~8 $c_6$; both are $n$-independent positive constants) yields
$|\partial S|\ge c_6(\delta - (C_4/c_4)\varepsilon)M$, which combined
with the Step~2 block-transfer loss $C_2(T)\varepsilon$ gives
$|\partial S|\ge c_6(\delta - K(T)\varepsilon)M$ with
$K(T) = C_4/c_4 + C_2(T)$, matching the definition in the paragraph
above and in the parameter cascade of Remark~\ref{rem:G2-order}.
The $O(c_4 C_4\,\varepsilon
M)$ error written in \eqref{eq:G2-lower} and in the aggregation
display above is a valid $O$-bound on the pre-rescaling error term
$c_6 C_4\,\varepsilon M$, since $c_4, c_6$ are absolute positive
constants; the $c_4$ is retained to signal that both Step~4 constants
enter the error.

\paragraph{$M$ vs.\ $\mathrm{vol}$ normalization.}
Step~5's Corollary~\ref{cor:second-moment} delivers
$\sum_{\alpha,\beta\in\Rich}w_{\alpha\beta}\ge c_T'(T)\,|\mathcal
L(P)|$ in counting measure. Proposition~\ref{prop:step5}(R) above
restates this directly as $M \ge c_5(T)\,|\mathcal L(P)|$ with
$c_5(T) := c_T'(T)$, an $n$-independent constant. The $(n-1)$
factor distinguishing counting measure from the
Section~\ref{sec:G1} convention $\mathrm{vol}(S) := |S|(n-1)$ is
carried explicitly in the boundary upper bound
\eqref{eq:G2-upper}: $|\partial S| \le \tfrac12 \eta(\gamma,n)(n-1)
|\mathcal L(P)|$. Consequently the G2 compatibility
\eqref{eq:G2} reads $\eta(\gamma,n)(n-1) < c_0 c_5(T) c_6 (\delta -
K(T)\varepsilon)$, an inequality in $n$-independent quantities
(since $\eta(\gamma,n)(n-1) = 2(n-1)/(\gamma n) \uparrow 2/\gamma$).
Step~5, Step~6, and \S\ref{sec:G1} now use a single consistent
normalization, and no $1/n$ factor is hidden inside $c_5$.

\paragraph{Parameter cascade realizability.}
The chain $\gamma \to T \to \delta \to \varepsilon \to n$ of
Remark~\ref{rem:G2-order} is quantitatively realizable, with the
proof inlined in the ``Realizability of the cascade (quantitative)''
paragraph of Theorem~\ref{thm:main} (Section~\ref{sec:main-thm}). The
two $n$-conditions split cleanly: (i) the conductance/volume
inequality reduces to the $n$-independent arithmetic condition
$\gamma^2 c_5(T) c_6 \delta_0 \ge 8$, which is taken as input via
Hypothesis~\ref{hyp:arith}; if it fails at some $\gamma < 1/3 -
\delta_{\mathrm{KL}}$, the small-$n$ base case
Remark~\ref{rem:small-n} handles 1/3--2/3 for $n \le n_0(\gamma, T)$
but the large-$n$ tail falls outside the scope of the present
argument (see Theorem~\ref{thm:main-step8}); (ii) the Step~2
fiber error condition $\varepsilon(\eta(\gamma, n)) \le
\delta_0/(2K(T))$ reduces, via the inverse modulus of
Theorem~\ref{thm:step2}'s $o(1)$ error, to the explicit
$n$-threshold
$n_0(\gamma, T) = \lceil 2/(\gamma\, \phi^\star_0)\rceil$ with
$\phi^\star_0 = \phi^\star(\delta_0/(2K(T)))$. Both conditions
depend only on $\gamma$ and $T$; (i) is what
Hypothesis~\ref{hyp:arith} externalizes as a final
constants-audit input.

\begin{remark}[Final constants audit: parameter cascade]
\label{rem:final-constants-audit}
The full list of named constants used above, with their origins and
roles, is:
\begin{itemize}[nosep]
 \item \textbf{Step 1:} $c_0^{(1)}, C_1^{(1)}, K$ -- fiber density
 $|\mathcal F_\alpha|\gtrsim c_0^{(1)}t_\alpha^2\cdot|\mathcal L(P)|/(|A||B||C|)$,
 bad-set thickness $|B_\alpha|\le C_1^{(1)}t_\alpha\cdot|\mathcal L(P)|/(|A||B||C|)$,
 and the bounded triple-overlap multiplicity $K$
 (\texttt{step1.tex}, Theorem~\ref{thm:interface}). Used inside Step~5.
 \item \textbf{Step 2:} $\varepsilon(\phi) \downarrow 0$ as the BK
 conductance $\phi\downarrow 0$, and $C_1^{(2)}$ for the staircase
 deviation constant. Step~2's output is a function, not a scalar;
 \S\ref{sec:G2-constants} fixes the scalar $\varepsilon$ at
 $\varepsilon < \delta/(2K(T))$.
 \item \textbf{Step 4:} absolute $c_4, C_4 > 0$
 (Lemma~\ref{lem:rect-stable}: $c_4 = c_1/4$, $C_4 = C'/4$ in the
 rectangle-stable proof). Combined into $K(T) = C_4/c_4 + C_2(T)$;
 see the ``Step~4 aggregation: constant rescaling'' paragraph above.
 $C_2(T)$ from the block-transfer step (Proposition~G2 of
 \texttt{step4.tex}) grows at most polynomially in $T$ via
 $\varepsilon' = \sqrt{4\varepsilon/\rho(T)}$.
 \item \textbf{Step 5:} $c_5^\star > 0$ (richness density,
 Proposition~\ref{prop:single-triple}(i)); $c_T = c_5^\star/2$ (the
 first-moment constant of Theorem~\ref{thm:step5}(R), from
 Lemma~\ref{lem:fiber-mass}); and
 $c_5(T) = c_T'(T) = (c_5^\star/4)^2$ the second-moment constant
 (Corollary~\ref{cor:second-moment}), deteriorating as
 $c_5^\star \asymp T^{-O(1)}$.
 \item \textbf{Step 6:} $c_6 > 0$ (called $c_1$ in
 \texttt{step6.tex}; strictly it is $c/(4M_{\mathrm{mult}})$ after
 pulling $\delta$ out of step~6's $c_1:=c\delta_0/(4M)$), with
 $M_{\mathrm{mult}}$ the witness-edge multiplicity bound of
 Lemma~\ref{lem:overlap-counting}. Independent of
 $T,\gamma,\varepsilon,\delta$ once the $\delta$ factor is exhibited
 outside.
 \item \textbf{Step 7:} $C_7 > 0$ (absolute, potential-extraction
 rate); $w = w(T)$ (interaction width, set by Step~5's richness
 threshold).
 \item \textbf{Step 8:} $c_0 := \gamma$, $\eta(\gamma,n) = 2/(\gamma n)$,
 $\beta(P)\ge\gamma$, $\alpha$ (local rectangle-density parameter,
 absorbed into $\delta$ in the aggregation), and
 $K_0 = K_0(\gamma, w)$ from Lemma~\ref{lem:layered-reduction}.
\end{itemize}

\paragraph{Dependency chain.} $\text{Step 1}\to\text{Step 5}$ (fiber
geometry, bad-set) $\to\text{Step 6}$ (via $c_5^\star$);
$\text{Steps 1/2}\to\text{Step 4}$ (rectangle model, staircase
approximation) $\to\text{Step 6}$ (via $c_4, C_4$);
$\text{Steps 5/6}\to\S\ref{sec:G2-constants}$ of Step~8 (via
$c_5(T), c_6, K(T)$); $\S\ref{sec:G2-constants}\to$
Theorem~\ref{thm:main}. No circular dependency: every downstream
constant is pinned by an upstream proof.

\paragraph{Final inequality, three equivalent forms.} The three
forms of the arithmetic condition that appear above are logically
ordered by strength, not inconsistent:
\begin{enumerate}[nosep, label=(\roman*)]
 \item \textbf{Base (Prop.~\ref{prop:G2}, eq.~\eqref{eq:G2}):}
 $\gamma^2\,c_5(T)\,c_6\,(\delta - K(T)\varepsilon) \ge 2$.
 This is the form derived in Remark~\ref{rem:n-dependence-g1} and
 suffices for the G2 compatibility contradiction at any single $n$.
 \item \textbf{Prop.~\ref{prop:step6} form} (factor $\tfrac12$ slack):
 $\gamma^2\,c_5(T)\,c_6\,(\delta - K(T)\varepsilon) \ge 4$, equivalent
 to $\gamma^2\,c_5(T)\,c_6\,\delta \ge 8$ under the cascade choice
 $\varepsilon\le\delta/(2K(T))$ (so $\delta - K(T)\varepsilon\ge\delta/2$).
 \item \textbf{Main-theorem form} (factor $\tfrac14$ slack,
 \S\ref{sec:main-thm} step~(4)):
 $\gamma^2\,c_5(T)\,c_6\,\delta_0 \ge 8$.
\end{enumerate}
Forms (ii) and (iii) are identical after the cascade $\varepsilon\to 0$
substitution; both bypass $n_0(\gamma,T)$ as $n$-independent
conditions on Steps~5 and~6 constants. Form (i) is the weakest; the
main theorem uses (iii) for headroom.

\paragraph{Satisfiability.} Fix $\gamma\in(0,1/3]$. The cascade of
Remark~\ref{rem:G2-order} chooses $T=T(\gamma)$ large enough that
Proposition~\ref{prop:step5}(R) is non-vacuous (in particular,
$T\ge T_0$ with $T_0$ absolute so that Corollary~\ref{cor:second-moment}'s
bad-set/good-bulk split holds); then
$c_5(T) = (c_5^\star(T,\gamma)/4)^2 > 0$ and $K(T) = C_4/c_4 + C_2(T)$
are fixed positive constants independent of $n$. Take
$\delta_0 := \min(1/(4C_7),\,1/(2wK_0))\in(0,1)$ and
$\varepsilon := \delta_0/(2K(T))$. Then inequality~(iii) reads
\begin{equation}\label{eq:arith-explicit}
 \gamma^2\,(c_5^\star/4)^2\,c_6\,\delta_0 \;\ge\; 8
 \quad\Longleftrightarrow\quad
 c_5^\star(T,\gamma)^2 \;\ge\; \frac{128}{\gamma^2\,c_6\,\delta_0},
\end{equation}
which is a \emph{purely arithmetic} condition on the richness density
$c_5^\star$ produced by Proposition~\ref{prop:single-triple}.
\emph{Whether~\eqref{eq:arith-explicit} is satisfiable at every
$\gamma \in (0, 1/3 - \delta_{\mathrm{KL}})$ is the content of
Hypothesis~\ref{hyp:arith}}; it is not derived from the sealed
outputs of Steps~1--7 and constitutes the sole external
numerical input the present argument requires. If the hypothesis
holds for a given $\gamma$ (some $T = T(\gamma)$
satisfies~\eqref{eq:arith-explicit}), the cascade closes and the
case~(R) contradiction fires at every $n\ge2$. If the hypothesis
fails at some $\gamma < 1/3 - \delta_{\mathrm{KL}}$ --- i.e.\ no
admissible $T$ satisfies~\eqref{eq:arith-explicit} ---
Remark~\ref{rem:small-n} still discharges the 1/3--2/3 property by
finite enumeration for $n \le n_0(\gamma, T)$, but the large-$n$
tail $n > n_0(\gamma, T)$ is outside the scope of the present
argument (see Theorem~\ref{thm:main-step8}). The main theorem as
stated therefore takes Hypothesis~\ref{hyp:arith} as input; its
verification is deferred to an external constants audit of
$c_5^\star, c_6, C_7, w, K_0$.

\paragraph{No circular or impossible constraint.} Each constant is
either (a)~an absolute positive constant from a sealed upstream lemma
($c_4, C_4, c_6, C_7, c_5^\star, K$, etc.), (b)~a polynomial function
of $T$ with $T=T(\gamma)$ fixed from $\gamma$ alone ($K(T), C_2(T),
w(T), c_5(T)$), or (c)~a parameter chosen with strictly decreasing
dependence on the upstream choices ($\delta_0$ from $C_7,w,K_0$;
$\varepsilon$ from $\delta_0,K(T)$; $n$ from $\gamma,T$). The
directed acyclicity $\gamma\to T\to\delta_0\to\varepsilon\to n$
precludes circularity; and the $n$-independence of the arithmetic
condition precludes an impossible large-$n$ limit.
\end{remark}

\section{G3: Layered-decomposition reduction to a smaller counterexample}
\label{sec:layered-reduction}

This section discharges GAP\,G3: we prove the ``reduction'' branch of
Proposition~\ref{prop:step7}(iii). The near-ordinal-sum and
near-disjoint-union templates of the width-3 literature are adapted
here from exact ordinal decompositions to Step~7's \emph{layered}
decompositions.

\subsection{Layered decompositions, formally}

\begin{definition}[Layered width-$3$ decomposition]\label{def:layered}
A width-$3$ poset $P=(X,\le_P)$ is said to admit a \emph{layered
decomposition of interaction width $w$} if there is an ordered
partition
\[
  X \;=\; L_1 \,\sqcup\, L_2 \,\sqcup\, \cdots \,\sqcup\, L_K
\]
such that:
\begin{enumerate}[nosep, label=(L\arabic*)]
 \item\label{item:L1} each \emph{band} $L_i$ has $|L_i| \le 3$ and is
 an antichain in $P$;
 \item if $x \in L_i$ and $y \in L_j$ with $i < j - w$, then
 $x <_P y$;
 \item\label{item:L3} if $x \in L_i$ and $y \in L_j$ with
 $|i-j| > w$, then $x$ and $y$ are comparable in $P$;
 \item\label{item:L4} if $x \in L_i$ and $y \in L_j$ with $i < j$ are
 comparable in $P$, then $x <_P y$: band index is compatible with
 the poset order whenever any comparability exists between two
 bands.
\end{enumerate}
By (L2), the only incomparabilities between distinct bands $L_i, L_j$
occur when $|i-j| \le w$; by (L4), whenever a cross-band pair is
comparable, the lower-indexed band lies below. We call $w$ the
\emph{interaction width} and $K$ the \emph{depth} of the layering.
\end{definition}

\begin{lemma}[Ground-set lift of the Step~5/7 layering]
\label{lem:layered-from-step7}
Let $P = (X, \le_P)$ be a width-$3$ poset with Dilworth decomposition
$X = A \sqcup B \sqcup C$ into three chains
$A = (a_1 <_P a_2 <_P \cdots)$, $B = (b_1 <_P \cdots)$,
$C = (c_1 <_P \cdots)$.  Assume either
\begin{enumerate}[nosep, label=(\alph*)]
 \item\label{it:input-5C}
  \emph{(Step~5(C) input)} The collapse conclusion of
  Proposition~\ref{prop:step5}(C) holds on $P$, with per-chain
  exceptional count $\le c_1(T)$ and collapse-window width
  $w_{\mathrm{coll}} = w_{\mathrm{coll}}(T) = O_T(1)$
  (Theorem~\ref{thm:step5} of \texttt{step5.tex}); or
 \item\label{it:input-7ii}
  \emph{(Step~7(ii) input)} The globalization conclusion of
  Proposition~\ref{prop:step7}(ii) holds on $P$, with potential
  $a \colon X \to \mathbb{R}$, threshold $c \in \mathbb{R}$, and
  bandwidth constant $K_{\mathrm{bw}} = K(T) = O_T(1)$ from
  Lemma~\ref{lem:bandwidth} of \texttt{step7.tex}.
\end{enumerate}
Define the \emph{exceptional set}
\begin{equation}\label{eq:Xexc-def}
  X^{\mathrm{exc}}
  \;:=\;
  X^{\mathrm{exc}}_{\mathrm{mono}}
  \;\cup\; X^{\mathrm{exc}}_{\mathrm{band}}
  \;\cup\; X^{\mathrm{exc}}_{\mathrm{sync}}
  \;\subseteq\; X,
\end{equation}
where (with the notational conventions $f_{AB}, f_{AC}, f_{BC}$ for
the synchronization maps of Section~\ref{sec:layered} of
\texttt{step7.tex}, undefined for chain-tail orphans):
\begin{itemize}[nosep, leftmargin=*]
 \item $X^{\mathrm{exc}}_{\mathrm{mono}} := \{z \in X :
  a \text{ fails strict chain-monotonicity at } z \text{ on its
  Dilworth chain}\}$ (empty in branch~\ref{it:input-5C});
 \item $X^{\mathrm{exc}}_{\mathrm{band}} := \{a_i \in A :
  \exists\,(a_i, b_j)\in\Rich^\star,\ |j - f_{AB}(i)| > K_{\mathrm{bw}}\}
  \cup \{b_j, c_k \text{ cases}\}$
  (bandwidth-excess rich endpoints);
 \item $X^{\mathrm{exc}}_{\mathrm{sync}} := \{z \in X :
  z \in \{a_i, b_j, c_k\} \text{ and a partner in another chain is
  undefined by }f_{\bullet\bullet}\}$
  (chain-tail orphans of size $\le K_{\mathrm{bw}}$ per chain).
\end{itemize}
Then:
\begin{enumerate}[nosep, label=(\roman*)]
 \item\label{it:lift-size}
  $|X^{\mathrm{exc}}| \le C_{\mathrm{exc}}(T) := 3\,c_1(T) = O_T(1)$,
  with $c_1(T)$ the per-chain endpoint-exceptional constant of
  Lemma~\ref{lem:endpoint-mono} (\texttt{step5.tex}) and the factor
  $3$ collecting the three ordered Dilworth triples
  $(A,B|C), (A,C|B), (B,C|A)$ contributing to Step~5's
  conflict-cube bookkeeping.
 \item\label{it:lift-layered}
  The subposet $P|_{X \setminus X^{\mathrm{exc}}}$ admits an
  \emph{exact} layered decomposition
  \[
    X \setminus X^{\mathrm{exc}}
    \;=\; L_1 \sqcup L_2 \sqcup \cdots \sqcup L_K
  \]
  in the sense of Definition~\ref{def:layered}, with
  \begin{itemize}[nosep]
   \item depth $K \ge (|X| - |X^{\mathrm{exc}}|)/3 \ge |X|/6$
    (for $|X| \ge 2\,C_{\mathrm{exc}}(T)$), and
   \item interaction width
    $w = K_{\mathrm{bw}} + 2$ in branch~\ref{it:input-7ii}
    (resp.\ $w = w_{\mathrm{coll}}$ in branch~\ref{it:input-5C}),
    matching the Step~7 bandwidth (resp.\ Step~5 collapse window) up
    to the $+O(1)$ slack of the synchronization maps.
  \end{itemize}
 \item\label{it:lift-perturb}
  For every $x, y \in X \setminus X^{\mathrm{exc}}$,
  \[
    \big|p_{xy}(P) - p_{xy}(P|_{X \setminus X^{\mathrm{exc}}})\big|
    \;\le\; \frac{2\,C_{\mathrm{exc}}(T)}{|X| - C_{\mathrm{exc}}(T) + 1},
  \]
  i.e.\ the perturbation bound~\eqref{eq:exc-perturb} transfers to the
  induced subposet.
\end{enumerate}
\end{lemma}

\begin{proof}
We treat branch~\ref{it:input-7ii} in detail; branch~\ref{it:input-5C}
reduces to the same conclusion with a strictly simpler
$X^{\mathrm{exc}}$ (see end of proof).

\emph{Level assembly (item~\ref{it:lift-layered}).}
By Proposition~\ref{prop:71} of \texttt{step7.tex}, after removing the
$o(1)$-mass non-monotonic set
$X^{\mathrm{exc}}_{\mathrm{mono}}$ the potential $a$ is strictly
increasing along each of $A, B, C$, so the synchronization maps
\[
  f_{AB}(i) := \arg\min_{j}|a(a_i) - a(b_j)|, \qquad
  f_{AC}(i) := \arg\min_{k}|a(a_i) - a(c_k)|,
\]
and $f_{BC}$ analogously, are weakly monotone wherever they are
defined.  Set
$K := \max(|A|, |B|, |C|) - |X^{\mathrm{exc}}|$
and, for $1 \le k \le K$, define the band
\[
  L_k \;:=\;
  \bigl\{a_k,\ b_{f_{AB}(k)},\ c_{f_{AC}(k)}\bigr\}
  \;\cap\; (X \setminus X^{\mathrm{exc}}),
\]
omitting missing chain entries when $f_{\bullet\bullet}$ is
undefined; the omitted entries are by definition in
$X^{\mathrm{exc}}_{\mathrm{sync}}$, so the right-hand side is a bona
fide subset of $X \setminus X^{\mathrm{exc}}$.  Relabel the indices if
needed so that every $L_k$ is nonempty; this loses at most $O_T(1)$
levels at either end (absorbed into the depth bound below).

\emph{Verification of (L1)--(L4) for $w := K_{\mathrm{bw}} + 2$.}
\begin{itemize}[nosep, leftmargin=*]
 \item \emph{(L1) Antichain size $\le 3$.}  Each $L_k$ contains at
  most one element per chain ($a_k$, $b_{f_{AB}(k)}$,
  $c_{f_{AC}(k)}$), hence $|L_k| \le 3$.  Distinct-chain entries are
  incomparable in $P$ at the chain level (Dilworth), so $L_k$ is an
  antichain.
 \item \emph{(L2) $i < j - w \Rightarrow x <_P y$ for $x \in L_i$,
  $y \in L_j$.}  Split on chain membership:
  \begin{enumerate}[nosep, label=\textbullet]
   \item If $x, y$ are in the same chain, say $x = a_i$ and
    $y = a_j$ with $i < j - w < j$, then $x <_P y$ by chain
    monotonicity.
   \item If $x = a_i$ and $y = b_{f_{AB}(j)}$ (say), the
    synchronization map $f_{AB}$ on $A \setminus
    X^{\mathrm{exc}}_{\mathrm{sync}}$ is weakly monotone and
    $1$-Lipschitz up to the chain-density slack $+1$, so
    $f_{AB}(j) \ge j - 1$ and the chain-position offset between
    $a_i$ and $b_{f_{AB}(j)}$ is at least
    $f_{AB}(j) - i \ge j - 1 - i > w - 1 = K_{\mathrm{bw}} + 1$.
    Hence, if the pair $(a_i, b_{f_{AB}(j)})$ were incomparable and
    rich, it would violate Lemma~\ref{lem:bandwidth}; by
    construction $a_i \notin X^{\mathrm{exc}}_{\mathrm{band}}$, so the
    pair is either non-rich or comparable.  In the non-rich
    incomparable subcase, the potential differential
    $\Delta_{a_i b_{f_{AB}(j)}} = a(b_{f_{AB}(j)}) - a(a_i)$ is at
    most the sync error $c_0 = O(1)$ (definition of $f_{AB}$), but
    the chain-position offset $> K_{\mathrm{bw}}$ combined with the
    second application of~\eqref{eq:var-budget} (proof of
    Lemma~\ref{lem:bandwidth}) contradicts the witness interval
    $J_i(c_0)$; hence the pair is comparable.  Strict chain-$a$
    monotonicity $a(b_{f_{AB}(j)}) > a(a_i)$ then forces
    $a_i <_P b_{f_{AB}(j)}$.
   \item All other cross-chain cases are symmetric.
  \end{enumerate}
 \item \emph{(L3) $|i - j| > w \Rightarrow x, y$ comparable.}
  Applying the (L2) argument to whichever of $i < j$ or $j < i$
  holds.
 \item \emph{(L4) Cross-band direction.}  Suppose $x \in L_i$,
  $y \in L_j$ with $i < j$ are comparable in $P$.  The potential $a$
  is strictly increasing along every Dilworth chain
  (\texttt{step7.tex:1004--1013}), and the weakly monotone
  synchronization maps $f_{AB}, f_{AC}, f_{BC}$ propagate this
  monotonicity across chains: if $x = a_i$ and $y = b_{f_{AB}(j)}$
  (or any cross-chain analogue) with $i < j$, then
  $a(y) \ge a(x) + O(1)$ (up to sync slack) is forced strictly positive
  by the band separation $i < j$, giving $a(y) > a(x)$.  By
  Proposition~\ref{prop:71} of \texttt{step7.tex}, every comparable
  cover $x' <_P y'$ satisfies $a(y') > a(x')$; hence any comparability
  between $x$ and $y$ must be $x <_P y$.  The within-chain case
  ($x, y$ in the same Dilworth chain) reduces to chain monotonicity
  of both $a$ and the chain order itself, which agree on
  $X \setminus X^{\mathrm{exc}}_{\mathrm{mono}}$.
\end{itemize}

\emph{Depth bound $K \ge |X|/6$.}  Pigeonhole on the three chains:
$\max(|A|, |B|, |C|) \ge |X|/3$, and
$|X^{\mathrm{exc}}| \le C_{\mathrm{exc}}(T) \le |X|/2$ once
$|X| \ge 2\,C_{\mathrm{exc}}(T)$, which holds automatically in the
cascade regime $|X| \ge n_0(\gamma, T)$ of~\eqref{eq:n0-main}.

\emph{Bound on $|X^{\mathrm{exc}}|$ (item~\ref{it:lift-size}).}
Lemma~\ref{lem:budget-var} of \texttt{step7.tex} bounds the mass of
non-monotonic elements by $o(1)$; pushed to ground-set support via
Markov on the incidence measure and the richness threshold $T$, this
gives $|X^{\mathrm{exc}}_{\mathrm{mono}}| = O_T(1)$.
Lemma~\ref{lem:bandwidth} bounds the rich-pair bandwidth excess by
$O_T(1)$ pairs, each contributing one endpoint per chain, so
$|X^{\mathrm{exc}}_{\mathrm{band}}| = O_T(1)$.  The chain-tail
orphans $X^{\mathrm{exc}}_{\mathrm{sync}}$ are $\le K_{\mathrm{bw}}$
per chain, so $\le 3\,K_{\mathrm{bw}} = O_T(1)$.  Summing the three
contributions and collecting the three ordered Dilworth triples
$(A,B|C), (A,C|B), (B,C|A)$ (each contributing $c_1(T)$ through the
Step~5 conflict-cube analysis, Theorem~\ref{thm:step5}), we obtain
$|X^{\mathrm{exc}}| \le 3\,c_1(T) = C_{\mathrm{exc}}(T)$.

\emph{Perturbation transfer (item~\ref{it:lift-perturb}).}
Set $k := |X^{\mathrm{exc}}|$ and $n := |X|$.  Delete the elements of
$X^{\mathrm{exc}}$ one at a time in any fixed order, producing a
sequence $P = P_0 \supsetneq P_1 \supsetneq \cdots \supsetneq P_k =
P|_{X \setminus X^{\mathrm{exc}}}$ of induced subposets with ground
sets of sizes $n, n-1, \ldots, n-k$.  For each incomparable pair
$x, y \in X \setminus X^{\mathrm{exc}}$ and each deletion step
$P_{i-1} \leadsto P_i$ removing element $z_i \ne x, y$, the
element-deletion coupling for uniform linear extensions
(exchangeability in the position of $z_i$ conditional on the
induced order of $X \setminus \{z_i\}$) gives
\[
  \bigl| p_{xy}(P_{i-1}) - p_{xy}(P_i) \bigr|
  \;\le\; \frac{2}{n - i + 1},
\]
since at most two of the $n - i + 1$ possible positions of $z_i$
(immediately before $x$ and immediately after $y$) can flip the
relative order of $x, y$ in the induced extension.  Summing the
telescoped bounds,
\begin{equation}\label{eq:exc-perturb-telescoped}
  \big|p_{xy}(P) - p_{xy}(P|_{X \setminus X^{\mathrm{exc}}})\big|
  \;\le\; \sum_{i=1}^{k} \frac{2}{n - i + 1}
  \;\le\; \frac{2\,k}{n - k + 1}
  \;\le\; \frac{2\,C_{\mathrm{exc}}(T)}{n - C_{\mathrm{exc}}(T) + 1},
\end{equation}
which is~\eqref{eq:exc-perturb}.  A fully formalized element-deletion
coupling — uniform across all incomparable pairs and across the
iterated-deletion order — is carried out in the forthcoming
refinement of~\eqref{eq:exc-perturb} (work-item \texttt{mg-d0e4/A6}
in the roadmap); the argument above captures the combinatorial
content and, in particular, the constant $2$ in the numerator is
tight for the one-step bound.

\emph{Branch~\ref{it:input-5C}.}
Theorem~\ref{thm:step5} directly supplies an exceptional set
$\mathcal F \subseteq X$ with $|\mathcal F| \le 3\,c_1(T)$ and a
layered decomposition of $X \setminus \mathcal F$ with interaction
window $w_{\mathrm{coll}} = O_T(1)$, so items~\ref{it:lift-size}
and~\ref{it:lift-layered} hold immediately with
$X^{\mathrm{exc}} := \mathcal F$ (and $X^{\mathrm{exc}}_{\mathrm{mono}},
X^{\mathrm{exc}}_{\mathrm{band}}$ vacuous).
Item~\ref{it:lift-perturb} follows from the same element-deletion
coupling~\eqref{eq:exc-perturb-telescoped}.
\end{proof}

\begin{remark}[Exact-layering working assumption]
\label{rem:layered-exact-assumption}
Lemma~\ref{lem:layered-from-step7} reduces the layered-reduction
program to an \emph{exactly} layered subposet
$P|_{X \setminus X^{\mathrm{exc}}}$ in the sense of
Definition~\ref{def:layered}.  If the conclusion of
Lemma~\ref{lem:layered-reduction} (below) fails on $X \setminus
X^{\mathrm{exc}}$, then item~\ref{it:lift-perturb} transfers the
failure to $P$ with a pairwise-probability slack
\[
  \frac{2\,C_{\mathrm{exc}}(T)}{|X| - C_{\mathrm{exc}}(T) + 1}
  \;\le\; \frac{\gamma}{2}
  \qquad \text{once $|X| \ge n_0(\gamma, T)$ of~\eqref{eq:n0-main},}
\]
which we absorb into the $\gamma/2$ in Lemma~\ref{lem:layered-reduction}'s
statement.  Throughout the remainder of this section we may therefore
assume $P = X$ is exactly layered in the sense of
Definition~\ref{def:layered}.
\end{remark}

\subsection{The cutting lemma}

\begin{lemma}[Ordinal cut inside a deep layering]\label{lem:cut}
Let $P$ be a width-$3$ poset admitting a layered decomposition
$X = L_1 \sqcup \cdots \sqcup L_K$ of interaction width $w$, with
$K \ge 2w + 2$. Then there exists an index $k$ with $w < k < K - w$
such that the partition
\[
  A \;:=\; L_1 \cup \cdots \cup L_k,
  \qquad
  B \;:=\; L_{k+1} \cup \cdots \cup L_K
\]
is an \emph{ordinal decomposition}: $a <_P b$ for every $a \in A$,
$b \in B$, \emph{except} for pairs $(a,b)$ with $a \in L_i$, $b \in L_j$
and $k - w < i \le k < j \le k + w$. In particular, every
incomparability across $(A,B)$ lies inside the ``interaction window''
$L_{k-w+1} \cup \cdots \cup L_{k+w}$.
\end{lemma}

\begin{proof}
Any index $k$ with $w < k < K - w$ works: if $a \in L_i \subseteq A$
and $b \in L_j \subseteq B$ with $i \le k < j$, then either
$j - i > w$ (in which case~(L2) gives $a <_P b$) or $j - i \le w$ and
therefore $k - w < i \le k$ and $k < j \le k + w$.
\end{proof}

\subsection{The reduction lemma}

\begin{lemma}[Layered reduction; GAP\,G3 discharged]\label{lem:layered-reduction}
Fix $\gamma > 0$. There is an integer $K_0 = K_0(\gamma, w)$ such that
the following holds. Let $P$ be a width-$3$ $\gamma$-counterexample on
ground set $X$ admitting a layered decomposition of interaction width
$w$ and depth $K \ge K_0$. Then either
\begin{enumerate}[nosep]
 \item $P$ has a balanced pair (contradicting the counterexample
 hypothesis), or
 \item there is a proper induced subposet $P' \subsetneq P$ on ground
 set $X' \subsetneq X$ such that $P'$ is a
 $(\gamma/2)$-counterexample.
\end{enumerate}
\end{lemma}

\begin{proof}
Assume $K \ge K_0 := \max(2w+2, \; 2w/\gamma + 4)$ and that $P$ has no
balanced pair. We produce a proper induced sub-counterexample.

\smallskip\noindent\emph{Step 1: cut.} Apply Lemma~\ref{lem:cut} at an
index $k \in (w, K-w)$; write $A = L_{\le k}$, $B = L_{>k}$,
$W = L_{k-w+1} \cup \cdots \cup L_{k+w}$ for the interaction window.
By width-$3$ and $|L_i| \le 3$, $|W| \le 6w$.

\smallskip\noindent\emph{Step 2: pick the heavy side.} Since $|L_i| \le
3$ and $K \ge K_0 \ge 2w/\gamma + 4$, we have $|A|, |B| \ge 3(w+1)$
and $|W|/|X| \le 6w/(3K) \le \gamma/2$. By symmetry of $A$ and $B$
under reversing $\le_P$, rename so that $|A| \ge |B|$; then $|A| \ge
|X|/2$ and hence $|X \setminus A| = |B| \le |X|/2$, so $A$ is a proper
induced subposet of $P$ that is a strict subset of $X$ (since $B$ is
nonempty).

\smallskip\noindent\emph{Step 3: pairwise probabilities in $P$ versus
$A$.} Let $(x,y) \subseteq A$ be an incomparable pair in $P$ (hence in
$A$). We compare $p_{xy}(P) := \Pr_{L \sim \mathcal{L}(P)}[x <_L y]$
with $p_{xy}(A) := \Pr_{L' \sim \mathcal{L}(A)}[x <_{L'} y]$, where in
each case the extension is uniform.

Let $P^\sharp$ denote the ordinal sum $A \oplus B$ (i.e., the poset on
$X$ whose relations are those of $P$ plus $a < b$ for every
$a \in A, b \in B$). Every linear extension of $P^\sharp$ decomposes as
$L_A \cdot L_B$ for $L_A \in \mathcal{L}(A)$ and $L_B \in
\mathcal{L}(B)$ independently, so
\[
  p_{xy}(P^\sharp) \;=\; p_{xy}(A)
  \qquad
  \text{for every $(x,y) \subseteq A$.}
\]
On the other hand, $P \subseteq P^\sharp$ as relations (i.e.,
$\mathcal{L}(P) \supseteq \mathcal{L}(P^\sharp)$), with the extra
extensions in $\mathcal{L}(P) \setminus \mathcal{L}(P^\sharp)$ being
exactly those that swap some pair $(a,b) \in A \times B$ that is
incomparable in $P$. By Lemma~\ref{lem:cut}, such incomparabilities
only involve elements of the window $W$; in particular they require
either $x$ or $y$ or an intervening element to lie in $W$.

For a fixed incomparable pair $(x,y)$ with $x,y \in A \setminus W$, no
$P^\sharp$-extension distinguishes the $L_A$-order of $x$ and $y$ from
their $L$-order in a joined extension $L = L_A \cdot L_B$, and no
extra $\mathcal{L}(P) \setminus \mathcal{L}(P^\sharp)$ extension can
move $x$ or $y$ across the $A/B$ interface (they are forced below
every element of $B$ by (L2), since $x,y \in A \setminus W$). Hence
\[
  p_{xy}(P) \;=\; p_{xy}(P^\sharp) \;=\; p_{xy}(A)
  \qquad
  \text{for every $(x,y)$ incomparable in $P$ with $x,y \in A \setminus W$.}
\]

\smallskip\noindent\emph{Step 4: inheriting the counterexample
property on $A$.} Since $P$ is a $\gamma$-counterexample, every
incomparable pair $(x,y) \subseteq X$ has $p_{xy}(P) \notin [1/3,2/3]$
and $\min(p_{xy}(P), 1-p_{xy}(P)) \ge \gamma$. By Step~3, every
incomparable pair $(x,y)$ inside $A \setminus W$ satisfies
$p_{xy}(A) = p_{xy}(P)$, so $p_{xy}(A) \notin [1/3, 2/3]$ and
$\min(p_{xy}(A), 1 - p_{xy}(A)) \ge \gamma \ge \gamma/2$.

\smallskip\noindent\emph{Step 5: handling pairs touching the window.}
A pair $(x,y) \subseteq A$ with $\{x,y\} \cap W \ne \emptyset$ need
not satisfy $p_{xy}(P) = p_{xy}(A)$. However, $|W \cap A| \le 3w$, so
the number of such pairs is at most $3w \cdot |A|$. If none of them
produces a balanced pair in $A$, we are done: $A$ is a
$(\gamma/2)$-counterexample on $|A| < |X|$ elements.

If some pair $(x,y) \subseteq A$ with $x \in W$ fails the
$(\gamma/2)$-counterexample condition in $A$, then either
$p_{xy}(A) \in [1/3, 2/3]$ (case~(a)), or
$p_{xy}(A) \notin [1/3,2/3]$ but
$\min(p_{xy}(A), 1-p_{xy}(A)) < \gamma/2$ (case~(b)).

In case~(a), $(x,y)$ is a balanced pair \emph{inside the smaller poset
$A$}; we still need a balanced pair in $P$. Consider the same pair
$(x,y)$ in $P$. By the same argument as Step~3 applied to $(x,y)
\subseteq A$, the extensions of $P$ decompose through the ordinal
closure $P^\sharp$ up to a perturbation supported on pairs inside the
window $W$; quantitatively
\[
  |p_{xy}(P) - p_{xy}(A)|
  \;\le\;
  \frac{|\mathcal{L}(P) \triangle \mathcal{L}(P^\sharp)|}
       {|\mathcal{L}(P)|}
  \;\le\;
  2^{|W|} / |\mathcal{L}(P)|^{\,0}
  \;=\; o_K(1)
\]
as $K \to \infty$, since the window $W$ has bounded size $6w$ while
$|\mathcal{L}(P)|$ grows at least as $\prod_i |L_i|! \ge 2^{K-|W|}$.
Choosing $K_0$ large enough (depending on $\gamma, w$) makes this
perturbation smaller than $1/3 - \gamma$; then
$p_{xy}(P) \in [1/3 - o(1), 2/3 + o(1)] \subseteq [1/3, 2/3]$ for the
original $\gamma$, producing a balanced pair in $P$ and returning us
to alternative~(1).

In case~(b), $\min(p_{xy}(A), 1-p_{xy}(A)) < \gamma/2$ but
$p_{xy}(A) \notin [1/3,2/3]$. By the same perturbation bound,
$|p_{xy}(P)-p_{xy}(A)| = o_K(1) \le \gamma/2$ for $K \ge K_0$, so
$\min(p_{xy}(P), 1-p_{xy}(P)) < \gamma$, contradicting the assumption
that $P$ is a $\gamma$-counterexample. So case~(b) cannot occur.

\smallskip\noindent\emph{Step 6: conclusion.} If no pair is in case~(a),
$A$ is a $(\gamma/2)$-counterexample on $|A| < |X|$ elements; take
$P' = A$. If some pair is in case~(a), $P$ has a balanced pair. In
either case the alternative holds.
\end{proof}

\begin{remark}\label{rem:shallow-layering}
Lemma~\ref{lem:layered-reduction} requires $K \ge K_0(\gamma, w)$.
The case $K < K_0$ --- a \emph{shallow} layering of bounded depth ---
is the regime handled by GAP\,G4 (layered-width-3 balanced-pair
lemma): in a layering of bounded depth the incomparability relation
lives on $O_T(1)$ pairs of bands, and a direct averaging argument on
an adjacent band-pair produces a balanced pair without invoking
minimality. Lemma~\ref{lem:layered-reduction} and G4 together cover
Proposition~\ref{prop:step7}(iii) with no gap: $K < K_0$ goes to G4,
$K \ge K_0$ goes to Lemma~\ref{lem:layered-reduction}.
\end{remark}

\begin{remark}[Disjoint-union analogue is absorbed]
We do not need a separate ``near-disjoint-union'' reduction. In a
\emph{layered} width-$3$ poset, far-apart bands are
\emph{comparable} (condition~\ref{item:L3} of
Definition~\ref{def:layered} with the direction from~(L2)), not
merely incomparable. So the layered regime only produces near
ordinal-sum structure, never near disjoint-union structure. The
disjoint-union lemma familiar from the width-3 literature is
therefore unnecessary here.
\end{remark}

\section{G4: Layered width-3 balanced-pair lemma}
\label{sec:g4-balanced-pair}

This section discharges GAP\,G4 as a standalone lemma: every layered
width-$3$ poset that is not a chain has a balanced pair, produced by
a direct averaging argument on a bounded interaction window. Combined
with Lemma~\ref{lem:layered-reduction} (G3), this covers the
direct-balanced-pair branch of Proposition~\ref{prop:step7}(iii) with
no depth gap: shallow layerings ($K < K_0$) are handled here, deep
layerings ($K \ge K_0$) by G3.

\begin{lemma}[Layered width-$3$ balanced pair; GAP\,G4]
\label{lem:layered-balanced}
Let $P = (X, \le_P)$ be a width-$3$ poset with $|X| \ge 2$ that is not
a chain, admitting a layered decomposition
$X = L_1 \sqcup \cdots \sqcup L_K$ of interaction width $w$
(Definition~\ref{def:layered}). Then $P$ contains a \emph{balanced
pair}: an incomparable pair $(x, y)$ with
\[
  \Pr_{L \sim \mathcal L(P)}\!\big[x <_L y\big]
  \;\in\; [1/3,\, 2/3].
\]
\end{lemma}

The proof proceeds through three ingredients. First, a general
\emph{ordinal-sum factorisation} of linear extensions
(\S\ref{subsec:ordinal-decomp}) shows that whenever a poset splits as a
three-part ordinal sum, the distribution of linear extensions factors
as a product and marginals of pair-order events within one piece are
invariant under restriction to that piece. Specialising this
factorisation to the slab $W(i,j)$ gives the \emph{window localization}
of \S\ref{subsec:window-localization}, reducing the probability to a
bounded sub-poset. A second specialisation, to the
\emph{layer-ordinal-reducible} case (\S\ref{subsec:layer-reducible}),
lets us isolate an irreducible factor; a \emph{direct averaging
argument} on the reduced bipartite instance
(\S\ref{subsec:bipartite-balanced}) then produces the balanced pair.

\subsection{Ordinal-sum decompositions}
\label{subsec:ordinal-decomp}

We isolate the piece of linear-extension combinatorics used throughout
this section: a finite poset that splits as an \emph{ordinal sum} of
three consecutive pieces has a uniform random linear extension that
factors as a product over those pieces. The statement below is
phrased for a three-part split because that is the version invoked by
window localization (where the three pieces are $X^-$, $W$, $X^+$);
the two-part specialisation, used for layer-ordinal reducibility,
follows by taking one piece empty.

\begin{definition}[Ordinal-sum decomposition witness]
\label{def:ordinal-decomp}
Let $P = (X, \le_P)$ be a finite poset. An \emph{ordinal-sum
decomposition witness} for $P$ is a triple
$D = (X_-,\, X_0,\, X_+)$ of pairwise disjoint subsets of $X$ with
$X_- \sqcup X_0 \sqcup X_+ = X$ satisfying the three elementwise
comparability conditions
\[
  X_- \;<_P\; X_0, \qquad X_- \;<_P\; X_+, \qquad X_0 \;<_P\; X_+,
\]
i.e.\ $u <_P v$ whenever $u$ lies in an earlier piece than $v$. We
write $\mathrm{OrdinalDecomp}(P)$ for the set of such witnesses and,
given $D \in \mathrm{OrdinalDecomp}(P)$, set
$P_- := P|_{X_-}$, $P_0 := P|_{X_0}$, $P_+ := P|_{X_+}$. We allow
any piece to be empty, in which case the corresponding
comparability condition is vacuous.
\end{definition}

\begin{lemma}[Ordinal-sum factorisation of linear extensions]
\label{lem:ordinal-factorisation}
Let $D = (X_-, X_0, X_+) \in \mathrm{OrdinalDecomp}(P)$. The
restriction map
\[
  \rho_D \colon \mathcal L(P) \;\longrightarrow\;
    \mathcal L(P_-) \times \mathcal L(P_0) \times \mathcal L(P_+),
  \qquad
  L \;\longmapsto\; \bigl(L|_{X_-},\,L|_{X_0},\,L|_{X_+}\bigr),
\]
is a bijection, under which the uniform distribution on $\mathcal
L(P)$ corresponds to the product of the uniform distributions on the
three factors. In particular, $|\mathcal L(P)| = |\mathcal L(P_-)|
\cdot |\mathcal L(P_0)| \cdot |\mathcal L(P_+)|$.
\end{lemma}

\begin{proof}
We write $n_- := |X_-|$, $n_0 := |X_0|$, $n_+ := |X_+|$, so
$n := n_- + n_0 + n_+ = |X|$; we identify a linear extension $L$ with
an order-preserving bijection $X \to \{1, \ldots, |X|\}$ (``positions''),
monotone with respect to $\le_P$.

\smallskip\noindent\emph{Position bands.}
Let $L \in \mathcal L(P)$ and $x \in X_-$. Every $y \in X_0 \cup X_+$
satisfies $x <_P y$ by the $X_-<_P X_0$ and $X_-<_P X_+$ hypotheses,
so $\mathrm{pos}_L(x) < \mathrm{pos}_L(y)$. Hence every
$X_-$-element occupies a strictly smaller position than every
non-$X_-$-element, forcing $\mathrm{pos}_L(X_-)=\{1,\ldots,n_-\}$.
The same argument applied to $X_+$ against $X_- \cup X_0$ forces
$\mathrm{pos}_L(X_+)=\{n_-+n_0+1,\ldots,n\}$, and therefore
$\mathrm{pos}_L(X_0)=\{n_-+1,\ldots,n_-+n_0\}$. Each of the three
position bands is thus determined by $D$, independent of $L$.

\smallskip\noindent\emph{The forward map.}
Given $L \in \mathcal L(P)$, the restriction $L|_{X_-}$ is the
composition of $L|_{X_-}$ (landing in $\{1,\ldots,n_-\}$) with the
unique order-isomorphism $\{1,\ldots,n_-\} \cong \{1,\ldots,n_-\}$,
and similarly $L|_{X_0}$ relabels positions in
$\{n_-+1,\ldots,n_-+n_0\}$ to $\{1,\ldots,n_0\}$ by subtracting
$n_-$, and $L|_{X_+}$ subtracts $n_-+n_0$. Each such restriction is
monotone with respect to $\le_{P_\bullet}$ (restriction of a
monotone map along an inclusion $X_\bullet \hookrightarrow X$
remains monotone). So
$\rho_D(L) \in \mathcal L(P_-) \times \mathcal L(P_0) \times
\mathcal L(P_+)$.

\smallskip\noindent\emph{The inverse map (concatenation).}
Given $(L_-, L_0, L_+) \in \mathcal L(P_-) \times \mathcal L(P_0)
\times \mathcal L(P_+)$, define $L \colon X \to \{1,\ldots,n\}$ by
\[
  L(x) \;=\;
  \begin{cases}
    L_-(x) & \text{if } x \in X_-,\\
    L_0(x) + n_- & \text{if } x \in X_0,\\
    L_+(x) + n_- + n_0 & \text{if } x \in X_+.
  \end{cases}
\]
$L$ is a bijection $X \to \{1,\ldots,n\}$ since the three pieces
partition $X$ and the three shifted images
$\{1,\ldots,n_-\}$, $\{n_-+1,\ldots,n_-+n_0\}$,
$\{n_-+n_0+1,\ldots,n\}$ partition $\{1,\ldots,n\}$. It is
order-preserving with respect to $\le_P$: if $x <_P y$ then either
(i) $x$ and $y$ lie in the same piece $X_\bullet$, in which case
$L_\bullet(x) < L_\bullet(y)$ by monotonicity of $L_\bullet$ and the
identical inequality transfers after the common shift, or (ii) $x$
and $y$ lie in different pieces, say $x \in X_-$, $y \in X_0$, in
which case $L(x) \le n_- < n_- + 1 \le L(y)$. The three cross-cases
$(X_-, X_0)$, $(X_-, X_+)$, $(X_0, X_+)$ are all handled by the
band structure of the codomain and do not require the hypothesis
$x <_P y$. (There are no pairs $y <_P x$ with $x$ in an earlier
piece than $y$, because the decomposition hypotheses only forbid
such comparabilities; a priori one could worry about the
\emph{absence} of $x \le_P y$ for $x \in X_-$, $y \in X_0$, but the
concatenation still places $x$ first regardless, so monotonicity of
$L$ is not violated. Concretely, the concatenation is always a
\emph{linear} order; it extends $\le_P$ because every forced
$<_P$-comparability is respected.)

\smallskip\noindent\emph{The maps are mutually inverse.}
Starting from $L$, the restriction $L|_{X_\bullet}$ lives on the
position band $\mathrm{pos}_L(X_\bullet)$ identified above;
concatenating the three restrictions (each shifted to start at its
band's left endpoint) reproduces $L$ position-by-position. Starting
from $(L_-, L_0, L_+)$, the concatenated $L$ has band $X_\bullet$
occupying exactly the expected positions, so restricting back yields
the original triple. Hence $\rho_D$ is a bijection.

\smallskip\noindent\emph{Uniform $\leftrightarrow$ product uniform.}
$\rho_D$ is a bijection between finite sets, so the pushforward of
the uniform distribution on $\mathcal L(P)$ along $\rho_D$ is the
uniform distribution on $\mathcal L(P_-) \times \mathcal L(P_0)
\times \mathcal L(P_+)$, which coincides with the product of the
uniform distributions on the three factors. Cardinalities follow
from the bijection.
\end{proof}

\begin{corollary}[Marginal invariance under ordinal-sum restriction]
\label{cor:ordinal-marginal}
Let $D = (X_-, X_0, X_+) \in \mathrm{OrdinalDecomp}(P)$ and fix
$x, y \in X_0$. Then
\[
  \Pr_{L \sim \mathcal L(P)}[x <_L y] \;=\;
  \Pr_{L_0 \sim \mathcal L(P_0)}[x <_{L_0} y].
\]
The analogous identity holds for $x, y$ both in $X_-$ (resp.\ $X_+$).
\end{corollary}

\begin{proof}
Under the bijection $\rho_D$ of Lemma~\ref{lem:ordinal-factorisation},
the event $\{x <_L y\}$ with $x, y \in X_0$ corresponds to the event
$\{x <_{L_0} y\}$ on the middle factor (since the relative order of
$x$ and $y$ in $L$ equals their relative order in $L|_{X_0}$, both
being read off the same two positions in the $X_0$-band). The
pushforward of uniform is uniform on each factor, so marginalising
over the $X_-$ and $X_+$ factors gives the claim.
\end{proof}

\subsection{Window localization}
\label{subsec:window-localization}

\begin{lemma}[Window localization]\label{lem:window-localization}
Fix an incomparable pair $(x, y)$ with $x \in L_i$, $y \in L_j$, and
set
\[
  W(i,j) \;:=\; L_{\max(1,\,\min(i,j)-w)}
  \cup \cdots \cup L_{\min(K,\,\max(i,j)+w)}.
\]
Then $|i - j| \le w$, and
\[
  \Pr_{L \sim \mathcal L(P)}[x <_L y]
  \;=\;
  \Pr_{L' \sim \mathcal L(P|_{W(i,j)})}[x <_{L'} y].
\]
Moreover $P|_{W(i,j)}$ is layered of interaction width $w$ and has at
most $3(3w + 1)$ elements.
\end{lemma}

\begin{proof}
The first statement is immediate: by~(L2), $|i-j| > w$ forces $x <_P y$
or $y <_P x$, contradicting incomparability of $(x,y)$.

For the probability identity, we exhibit an ordinal-sum decomposition
witness and invoke Corollary~\ref{cor:ordinal-marginal}. Every $z \in
X \setminus W(i,j)$ lies in a band $L_k$ with $|k - \min(i,j)| > w$
and $|k - \max(i,j)| > w$, so by~(L2) $z$ is comparable (in $P$) to
every element of $W(i,j)$, and the direction depends only on whether
$k < \min(i,j) - w$ or $k > \max(i,j) + w$. Set
\[
  X_- \;:=\; \bigcup_{k < \min(i,j) - w} L_k,
  \qquad X_0 \;:=\; W(i,j),
  \qquad X_+ \;:=\; \bigcup_{k > \max(i,j) + w} L_k.
\]
These three sets partition $X$, and by the direction of~(L2) the
elementwise comparabilities $X_- <_P X_0 <_P X_+$ and $X_- <_P X_+$
all hold. Hence $(X_-, X_0, X_+) \in \mathrm{OrdinalDecomp}(P)$, and
since $x, y \in X_0 = W(i,j)$,
Corollary~\ref{cor:ordinal-marginal} gives
\[
  \Pr_{L \sim \mathcal L(P)}[x <_L y] \;=\;
  \Pr_{L_0 \sim \mathcal L(P|_{W(i,j)})}[x <_{L_0} y].
\]

The size bound follows from $|W(i,j)| \le 3(3w+1)$ using $|i-j| \le w$
and $|L_k| \le 3$, and $P|_{W(i,j)}$ inherits the layered structure
of~$P$ with the same interaction width~$w$.
\end{proof}

\subsection{Layer-ordinal reducibility}
\label{subsec:layer-reducible}

The window localization of \S\ref{subsec:window-localization} reduces
the balanced-pair problem to the bounded layered sub-poset
$Q := P|_{W(i,j)}$ of depth at most $3w+1$. The next reduction
iteratively peels off \emph{layer-ordinal-sum factorisations}
\emph{internal} to $Q$ until no further factorisation is possible.

\begin{definition}[Layer-ordinal-reducible at index $k$]
\label{def:layer-reducible}
A layered poset $Q = M_1 \sqcup \cdots \sqcup M_r$ (with bands
$M_1, \ldots, M_r$ satisfying the layered axioms
of Definition~\ref{def:layered}) is said to be
\emph{layer-ordinal-reducible at index $k$}, for some
$1 \le k < r$, if every cross-pair $(u, v)$ with
$u \in M_i$, $v \in M_j$, and $i \le k < j$, satisfies $u <_Q v$.
Equivalently, all layer-crossing comparabilities across the cut
$k \mid k+1$ are directed upward.

$Q$ is \emph{layer-ordinal-reducible} if it is reducible at some index
$k$, and \emph{irreducible} otherwise.
\end{definition}

\begin{lemma}[Decomposition witness from reducibility]
\label{lem:reducible-witness}
If $Q = M_1 \sqcup \cdots \sqcup M_r$ is layer-ordinal-reducible at
index $k$, then the triple
\[
  D_k \;:=\; \bigl(\,
    \underbrace{M_1 \cup \cdots \cup M_k}_{=:\,Q_1},\;
    \varnothing,\;
    \underbrace{M_{k+1} \cup \cdots \cup M_r}_{=:\,Q_2}
  \bigr)
\]
is an ordinal-sum decomposition witness in
$\mathrm{OrdinalDecomp}(Q)$, with $Q|_{Q_1}$ layered of depth $k$ and
$Q|_{Q_2}$ layered of depth $r - k$.
\end{lemma}

\begin{proof}
The three sets partition the ground set of $Q$ by construction. The
reducibility hypothesis is exactly the elementwise comparability
$Q_1 <_Q Q_2$ (since every $u \in Q_1$ lies in some $M_i$ with
$i \le k$ and every $v \in Q_2$ lies in some $M_j$ with $k < j$, so
$u <_Q v$). The two comparabilities involving the empty middle piece
are vacuous. Each factor $Q|_{Q_\bullet}$ inherits the layered
structure with the obvious sub-sequence of bands.
\end{proof}

\begin{corollary}[Reducibility transfer]
\label{cor:reducibility-transfer}
Let $Q$ be layer-ordinal-reducible at index $k$, with factors
$Q_1, Q_2$ as in Lemma~\ref{lem:reducible-witness}. Then the
restriction map
\[
  \mathcal L(Q) \;\xrightarrow{\;\sim\;}\;
  \mathcal L(Q|_{Q_1}) \times \mathcal L(Q|_{Q_2}),
  \qquad L \mapsto (L|_{Q_1},\,L|_{Q_2}),
\]
is a bijection pushing uniform to the product uniform. For any
$x, y$ both in $Q_1$ (resp.\ both in $Q_2$),
\[
  \Pr_{L \sim \mathcal L(Q)}[x <_L y]
  \;=\;
  \Pr_{L_1 \sim \mathcal L(Q|_{Q_1})}[x <_{L_1} y]
  \quad\bigl(\text{resp.\ }
  \Pr_{L_2 \sim \mathcal L(Q|_{Q_2})}[x <_{L_2} y]\bigr).
\]
In particular, any balanced pair of $Q|_{Q_1}$ or $Q|_{Q_2}$ is also
a balanced pair of $Q$.
\end{corollary}

\begin{proof}
Apply Lemma~\ref{lem:ordinal-factorisation} to the witness $D_k$ of
Lemma~\ref{lem:reducible-witness}. The middle factor $\mathcal
L(Q|_{\varnothing})$ is a singleton, so the three-fold product
collapses to $\mathcal L(Q|_{Q_1}) \times \mathcal L(Q|_{Q_2})$. The
marginal identity for $x, y$ both in $Q_1$ is
Corollary~\ref{cor:ordinal-marginal} applied with $X_- = Q_1$ (the
relevant ``lower'' case of the corollary); symmetrically for
$x, y \in Q_2$. Balanced pairs transfer because the probability
$\Pr[x <_L y]$ is literally equal in $Q$ and in the factor.
\end{proof}

\begin{remark}[Iterated reduction]
\label{rem:iterated-reduction}
If $Q|_{Q_1}$ is itself reducible, one can iterate
Corollary~\ref{cor:reducibility-transfer} on $Q|_{Q_1}$; each step
strictly decreases the number of bands, so iteration terminates at a
pair of factors $Q^\star_1, Q^\star_2, \ldots$ each of which is
\emph{irreducible}. The marginal identity chains: a pair $(x,y)$ both
in some terminal irreducible factor $Q^\star_m$ satisfies
$\Pr_{\mathcal L(Q)}[x <_L y] = \Pr_{\mathcal L(Q^\star_m)}[x <_L y]$.
Lemma~\ref{lem:chained-lift} below records this chained transfer in a
form that does not presuppose layered structure.
\end{remark}

\begin{definition}[Ordinal-sum decomposition chain]
\label{def:ordinal-chain}
Let $P$ be a finite poset. An \emph{ordinal-sum decomposition chain}
of length $n \ge 0$ starting from $P$ is a sequence of finite posets
$P = P_0, P_1, \ldots, P_n$ such that for each $k = 1, \ldots, n$
there is an ordinal-sum decomposition witness
$D_k = (X^{(k)}_-, X^{(k)}_0, X^{(k)}_+) \in
\mathrm{OrdinalDecomp}(P_{k-1})$
(Definition~\ref{def:ordinal-decomp}) and a choice of piece
$S_k \in \{X^{(k)}_-, X^{(k)}_0, X^{(k)}_+\}$ with
$P_k = P_{k-1}|_{S_k}$. In particular the ground set of $P_k$ is a
subset of the ground set of $P_{k-1}$, so by iterated inclusion the
ground set of $P_n$ is a subset of $X = P_0$.
\end{definition}

\begin{lemma}[Chained balanced-pair lift]
\label{lem:chained-lift}
Let $P = P_0, P_1, \ldots, P_n$ be an ordinal-sum decomposition chain
(Definition~\ref{def:ordinal-chain}). If $(x, y)$ is a balanced pair
of $P_n$, then, viewed through the iterated inclusion of ground sets,
$(x, y)$ is a balanced pair of $P = P_0$.
\end{lemma}

\begin{proof}
We prove by induction on $n$ the following pair of statements, for
every ordinal-sum decomposition chain $P_0, P_1, \ldots, P_n$ and
every $x, y$ in the ground set of $P_n$:
\begin{enumerate}[label=\textnormal{(\alph*)},nosep]
  \item $(x, y)$ is incomparable in $P_n$ if and only if $(x, y)$ is
        incomparable in $P_0$; and
  \item $\Pr_{L_n \sim \mathcal L(P_n)}\!\big[x <_{L_n} y\big]
        \;=\;
        \Pr_{L \sim \mathcal L(P_0)}\!\big[x <_L y\big]$.
\end{enumerate}
Granting (a) and (b), the conclusion is immediate: if $(x, y)$ is a
balanced pair of $P_n$, then (a) upgrades incomparability to $P_0$
and (b) shows the $\mathcal L(P_0)$-order probability equals the
$\mathcal L(P_n)$-order probability, hence lies in $[1/3, 2/3]$.

\emph{Base case $n = 0$.} The chain is the trivial singleton
$P_0 = P$; both (a) and (b) are tautological.

\emph{Inductive step $n \mapsto n+1$.} Let
$P_0, P_1, \ldots, P_n, P_{n+1}$ be a chain of length $n+1$, with
$P_{n+1} = P_n|_{S_{n+1}}$ for some piece $S_{n+1}$ of an ordinal-sum
decomposition $D_{n+1}=(X_-^{(n+1)},X_0^{(n+1)},X_+^{(n+1)}) \in
\mathrm{OrdinalDecomp}(P_n)$. Fix $x, y$ in the ground set of
$P_{n+1}$, i.e.\ $x, y \in S_{n+1}$.

For (a): the sub-poset order on $P_{n+1} = P_n|_{S_{n+1}}$ is, by
definition, the restriction of $\le_{P_n}$ to $S_{n+1}$, so
$x \le_{P_{n+1}} y$ iff $x \le_{P_n} y$; hence $(x, y)$ is
incomparable in $P_{n+1}$ iff it is incomparable in $P_n$. By the
induction hypothesis applied to the length-$n$ chain
$P_0, \ldots, P_n$ and the same pair $(x, y) \in P_n$,
incomparability in $P_n$ coincides with incomparability in $P_0$,
which closes (a) for the extended chain.

For (b): since $x, y$ both lie in the piece $S_{n+1}$ of
$D_{n+1} \in \mathrm{OrdinalDecomp}(P_n)$,
Corollary~\ref{cor:ordinal-marginal} applied to $P_n$ and $D_{n+1}$
(in the $X_-$, $X_0$, or $X_+$ variant according to which piece
$S_{n+1}$ is) gives
\[
  \Pr_{L_n \sim \mathcal L(P_n)}\!\big[x <_{L_n} y\big]
  \;=\;
  \Pr_{L_{n+1} \sim \mathcal L(P_{n+1})}\!\big[x <_{L_{n+1}} y\big].
\]
By the induction hypothesis, the left-hand side equals
$\Pr_{L \sim \mathcal L(P_0)}[x <_L y]$, and composing the two
equalities yields (b) for the length-$(n+1)$ chain.

This completes the induction and hence the proof.
\end{proof}

\begin{remark}[Lean image of Lemma~\ref{lem:chained-lift}]
\label{rem:chained-lift-lean}
The length-one case is formalised as the Lean theorem
\texttt{OrdinalDecomp.}\allowbreak\texttt{has\-Balanced\-Pair\_lift}
in \texttt{lean/OneThird/Mathlib/}\allowbreak\texttt{LinearExtension/Subtype.lean}:
a \texttt{HasBalancedPair} on the middle piece \texttt{D.Mid} lifts
to a \texttt{HasBalancedPair} on the ambient $\alpha$, using
\texttt{probLT\_restrict\_eq} for part~(b) and the
\texttt{Subtype}-order coincidence for part~(a). The chained
statement of Lemma~\ref{lem:chained-lift} is the paper-level image
of Lean work item F4 (see Remark~\ref{rem:old-vs-new}): iterate
\texttt{has\-Balanced\-Pair\_lift} along an
\texttt{OrdinalDecomp}-chain, with each step an application of
\texttt{probLT\_restrict\_eq} and sub-poset order transfer. The
induction on chain length in the proof of
Lemma~\ref{lem:chained-lift} is the mathematical content of that
iteration.
\end{remark}

\subsection{Bipartite balanced pair}
\label{subsec:bipartite-balanced}

Lemma~\ref{lem:window-localization} reduces
Lemma~\ref{lem:layered-balanced} to the finite case of a bounded
layered sub-poset $Q = P|_{W}$, and
Corollary~\ref{cor:reducibility-transfer} together with
Remark~\ref{rem:iterated-reduction} reduce it further to the case of
an \emph{irreducible} factor. It therefore suffices to exhibit a
balanced pair in some irreducible layered piece. An irreducible
layered piece with depth $\ge 2$ has some adjacent band-pair
$(M_i, M_{i+1})$ with an incomparable cross-pair
(Lemma~\ref{lem:irr-adjacent} below); the remainder of this
subsection handles the genuinely bipartite case ($K = 2$) directly,
and Proposition~\ref{prop:in-situ-balanced} lifts that case to
arbitrary depth $K \ge 2$ without a further window step.

\begin{lemma}[Irreducible $\Rightarrow$ adjacent incomparable cross-pair]
\label{lem:irr-adjacent}
Let $Q = M_1 \sqcup \cdots \sqcup M_r$ be a layer-ordinal-sum
irreducible layered poset (Definition~\ref{def:layer-reducible}) of
depth $r \ge 2$ that is not a chain. Then there exists an index
$i \in \{1, \ldots, r-1\}$ and elements $u \in M_i$, $v \in M_{i+1}$
with $(u, v)$ incomparable in $Q$.
\end{lemma}

\begin{proof}
Suppose, for contradiction, that every adjacent band-pair
$(M_i, M_{i+1})$ is fully comparable, i.e.\ $u <_Q v$ for all
$u \in M_i$, $v \in M_{i+1}$ (the reverse direction is ruled out by
(L4)). By transitivity, for any $i < j$ and any $u \in M_i$,
$v \in M_j$ we obtain $u <_Q v$. In particular $Q$ is reducible at
every index $k \in \{1, \ldots, r-1\}$ (in the sense of
Definition~\ref{def:layer-reducible}), contradicting irreducibility.
\end{proof}

\begin{remark}[The inner-window pitfall]
\label{rem:inner-window-pitfall}
A tempting shortcut --- and the one that appeared in earlier drafts of
this section --- is: having reduced to an irreducible $Q^\star$ of
depth $K^\star \ge 2$ with an adjacent incomparable pair in
$(M_i, M_{i+1})$, ``apply Lemma~\ref{lem:window-localization} once
more inside $Q^\star$ to isolate the pair to $M_i \cup M_{i+1}$.''
This fails: Lemma~\ref{lem:window-localization} localises to
$W(i, i+1) = M_{\max(1, i-w)} \cup \cdots \cup M_{\min(K^\star,
i+1+w)}$, a window of up to $2w + 2$ bands, not the two bands
$M_i \cup M_{i+1}$. Reducing to just two bands would require $w = 0$
(every cross-band pair comparable), which irreducibility does not
guarantee.

We therefore restructure the proof of
Lemma~\ref{lem:layered-balanced} itself: in place of a single
``window-then-reduce-then-window'' pipeline, we run
Lemma~\ref{lem:window-localization} and
Corollary~\ref{cor:reducibility-transfer} \emph{iteratively}, using a
strong induction on $|X|$. Each step is either a strict size descent
(handled by the induction hypothesis) or terminates at a bounded-size
irreducible poset, which is closed by
Proposition~\ref{prop:in-situ-balanced} below.
\end{remark}

\begin{proposition}[Bipartite layered balanced pair]
\label{prop:bipartite-balanced}
Let $Q = A \sqcup B$ be a height-$2$ poset with $A, B$ antichains of
size $\le 3$, every comparability directed $A <_Q B$, and at least one
incomparable pair in $Q$. Then $Q$ has a balanced pair.
\end{proposition}

\begin{proof}
If $|A| \le 1$ and $|B| \le 1$ there is nothing to prove (no
incomparable pair would exist). So at least one antichain has size
$\ge 2$; WLOG $|A| \ge 2$.

\medskip\noindent\emph{Case 1: symmetric pair inside a layer.} Suppose
there are $x, y \in A$ with \emph{identical external profile},
$\Pi(x) := \{b \in B : x <_Q b\} = \Pi(y)$. The map
$\tau: \mathcal L(Q) \to \mathcal L(Q)$ that swaps $x$ and $y$ in a
linear extension is well-defined (it preserves every constraint, by
equality of profiles) and is an involution. Hence
$\Pr[x <_L y] = \Pr[y <_L x] = 1/2 \in [1/3,\,2/3]$, and $(x,y)$ is a
balanced pair. (The analogous argument applies if such a symmetric
pair exists in $B$.)

\medskip\noindent\emph{Case 2: no symmetric pair.} All profiles
$\{\Pi(a) : a \in A\}$ are distinct subsets of $B$, and likewise for
$B$. We use a monotonicity (intermediate-value) argument on the three
probabilities $\Pr[x <_L y]$ as $(x, y)$ ranges over within-$A$ pairs.

List $A = \{a_1, \ldots, a_m\}$ ($m \in \{2, 3\}$) so that
$|\Pi(a_1)| \ge |\Pi(a_2)| \ge |\Pi(a_3)|$ (more forced edges means
the element tends to come earlier in a linear extension).

\emph{Sub-claim:} $\Pr[a_i <_L a_{i+1}] \ge 1/2$ for $i = 1, \ldots,
m-1$.

\emph{Proof of sub-claim:} Couple $\mathcal L(Q)$ with the
distribution $\mathcal L(Q')$ obtained by ``equalizing'' $a_i$ and
$a_{i+1}$ (formally, considering the poset $Q'$ where $a_{i+1}$'s
profile is replaced by $\Pi(a_i)$, which adds at least one
comparability). By the Ahlswede--Daykin~\cite{AhlswedeDaykin1978} /
FKG~\cite{FKG1971} inequality for linear extensions (equivalently, the
Graham--Yao--Yao inequality~\cite{GrahamYaoYao1980}
for $\Pr[a_i <_L a_{i+1}]$ applied to the suppression of one
comparability $a_{i+1} <_{Q'} b$), removing a comparability
$a_{i+1} <_Q b$ for some $b \in \Pi(a_i) \setminus \Pi(a_{i+1})$ can
only \emph{decrease} $\Pr[a_i <_L a_{i+1}]$ (since removing the
constraint $a_{i+1}$-early-in-$L$ lets $a_{i+1}$ drift later). Starting
from the symmetric poset $Q''$ in which $a_i$ and $a_{i+1}$ share the
profile $\Pi(a_i) \cap \Pi(a_{i+1})$, $\Pr_{Q''}[a_i <_L a_{i+1}] =
1/2$; successively adding back the comparabilities of $a_i$ (i.e., the
elements of $\Pi(a_i) \setminus \Pi(a_{i+1})$) raises this probability
to its value in $Q$. Hence $\Pr_Q[a_i <_L a_{i+1}] \ge 1/2$. This
proves the sub-claim. \hfill$\Diamond$

\medskip
\emph{Extracting the balanced pair from the sub-claim.} If
$\Pr[a_1 <_L a_m] \le 2/3$, then $(a_1, a_m)$ is balanced (the sub-claim
gives $\Pr[a_1 <_L a_m] \ge 1/2 \ge 1/3$). So assume
$\Pr[a_1 <_L a_m] > 2/3$. We use a \emph{telescoping average}:
\[
  \Pr[a_1 <_L a_m]
  \;=\; \Pr[a_1 <_L a_m \text{ with } a_{i}\text{'s in natural order}]
       + \Pr[a_1 <_L a_m, \text{ some inversion in between}],
\]
and each adjacent swap $a_i \leftrightarrow a_{i+1}$ contributes
independently. More explicitly, for $m = 3$:
\[
  \Pr[a_1 <_L a_3] \;=\; \Pr[a_1 <_L a_2 <_L a_3]
                        + \Pr[a_1 <_L a_3 <_L a_2]
                        + \Pr[a_2 <_L a_1 <_L a_3].
\]
If $\Pr[a_1 <_L a_2] > 2/3$ and $\Pr[a_2 <_L a_3] > 2/3$ both hold,
then $\Pr[a_1 <_L a_2 <_L a_3] = \Pr[a_1 <_L a_2] + \Pr[a_2 <_L a_3]
- \Pr[a_1 <_L a_2 \text{ or } a_2 <_L a_3] \ge 2 \cdot (2/3) - 1 = 1/3$
by inclusion--exclusion; and a symmetric lower bound on the full triple
event $\Pr[a_1 <_L a_2 <_L a_3] \le 1 - \Pr[\text{some inversion}]$
forces one of the three adjacent-pair probabilities into
$[1/3, 2/3]$. Specifically, if both $\Pr[a_1 <_L a_2] > 2/3$ and
$\Pr[a_2 <_L a_3] > 2/3$, but $\Pr[a_1 <_L a_3] > 2/3$, then summing
the three complements
$\Pr[a_2 <_L a_1] + \Pr[a_3 <_L a_2] + \Pr[a_3 <_L a_1] < 1$, while
any three random variables taking values in $\{<, >\}$ satisfy
$\Pr[\pi_1] + \Pr[\pi_2] + \Pr[\pi_3] \ge 1$ for the three
``rotations'' (since the events cover at least one of the six
orderings of the triple). This contradiction means at least one
adjacent within-$A$ pair has probability in $[1/3, 2/3]$.
(The three ``rotations'' are the events $\{a_2<_L a_1\}$,
$\{a_3<_L a_2\}$, $\{a_1<_L a_3\}$, and they cover every total order
on $\{a_1,a_2,a_3\}$: if all three failed simultaneously we would have
$a_1<_L a_2$, $a_2<_L a_3$, and $a_3<_L a_1$, a $3$-cycle forbidden
for a total order. Hence
$\Pr[a_2<_L a_1]+\Pr[a_3<_L a_2]+\Pr[a_1<_L a_3]\ge 1$.)

For $m = 2$, the sub-claim gives $\Pr[a_1 <_L a_2] \ge 1/2$, and the
reverse bound $\Pr[a_1 <_L a_2] \le 2/3$ follows from the same FKG
coupling applied to the opposite direction: no constraint of the form
$a_2 <_Q b$ with $b \notin \Pi(a_1)$ can push the probability above
$2/3$ in a bipartite poset of this size, by a direct check. To see
that $2/3$ is tight and not improvable, enumerate the minimal
bipartite configurations with $|A|=2$, $|B|\le 1$, and at most one
cross-comparability:
\begin{itemize}[nosep]
  \item $|B|=0$ (or $|B|=1$ with no cross-comparability): $a_1,a_2$
  are symmetric, so $\Pr[a_1<_L a_2]=1/2$.
  \item $|B|=1$, $a_1<_Q b$ only: the three linear extensions
  $a_1a_2b$, $a_1ba_2$, $a_2a_1b$ give
  $\Pr[a_1<_L a_2]=2/3$ (the two former satisfy $a_1<_L a_2$, the
  third does not).
  \item $|B|=1$, $a_2<_Q b$ only: by the symmetric argument
  $\Pr[a_1<_L a_2]=1/3$, so $\Pr[a_2<_L a_1]=2/3$; after relabelling
  this is the same extremal value.
  \item $|B|=1$, both $a_1<_Q b$ and $a_2<_Q b$: the two extensions
  $a_1a_2b$, $a_2a_1b$ give $\Pr[a_1<_L a_2]=1/2$.
\end{itemize}
Adding further $B$-elements or cross-comparabilities only averages
these cases (by FKG-monotonicity), so $2/3$ is the supremum, attained
exactly at the extremal configuration $|A|=2$, $|B|=1$, $a_1<_Q b$.
\end{proof}

\begin{remark}[Finite enumeration]
The extremal bound $2/3$ in the $m = 2$ case is realized by
$A = \{a_1, a_2\}$, $B = \{b\}$, $a_1 <_Q b$: the three linear
extensions are $a_1 a_2 b$, $a_1 b a_2$, $a_2 a_1 b$, giving
$\Pr[a_1 <_L a_2] = 2/3$. Any bipartite height-$2$ poset with
$|A|, |B| \le 3$ and at least one incomparable pair either contains
this extremal configuration as a sub-configuration (in which case
$(a_1, a_2)$ is balanced) or has a symmetric within-layer pair or
satisfies Case 2 strictly. A routine finite enumeration
(at most $2^9 \cdot 2 = 1024$ bipartite bimap possibilities on
$|A|,|B| \le 3$, drastically reduced modulo automorphism) confirms
that no bipartite layered configuration evades
Proposition~\ref{prop:bipartite-balanced}.
\end{remark}

\subsection{In-situ adjacent-band balanced pair}

Before proving Lemma~\ref{lem:layered-balanced}, we record the
base-case proposition that closes the induction. It is a direct
generalisation of Proposition~\ref{prop:bipartite-balanced} from the
strict bipartite setting ($K = 2$) to an irreducible layered poset of
bounded depth, so that no further window-localisation is needed.

\begin{proposition}[In-situ adjacent-band balanced pair]
\label{prop:in-situ-balanced}
Let $Q$ be a layered width-$3$ poset of interaction width $w$ and
depth $K \ge 2$, with $|Q| \le 3(3w+1)$. Then $Q$ contains a balanced
pair.
\end{proposition}

\begin{proof}
By Lemma~\ref{lem:irr-adjacent} applied to any layer-ordinal-sum
irreducible factor of $Q$ (or $Q$ itself if it is already
irreducible), it suffices to treat the case where $Q$ is irreducible;
a balanced pair in an irreducible factor is a balanced pair of $Q$ by
the $\mathcal L(Q_1 \oplus Q_2) = \mathcal L(Q_1) \times \mathcal
L(Q_2)$ factorisation.

Fix adjacent bands $(M_i, M_{i+1})$ with an incomparable cross-pair
$(x, y)$ provided by Lemma~\ref{lem:irr-adjacent}, and write
$A := M_i$, $B := M_{i+1}$. We distinguish three cases.

\medskip\noindent\emph{Case 1: symmetric within-band pair with full
ambient match.} Suppose there are distinct $a, a' \in A$ with
identical \emph{ambient profile} in $Q$, i.e.
\[
  \{z \in Q : a <_Q z\} \;=\; \{z \in Q : a' <_Q z\}
  \quad\text{and}\quad
  \{z \in Q : z <_Q a\} \;=\; \{z \in Q : z <_Q a'\}.
\]
Then the transposition $\tau := (a\ a')$ is a poset-automorphism of
$Q$: it preserves every comparability (since $a$ and $a'$ are
interchangeable by the profile equality, and $A$ is an antichain so
no $A$-internal relation is disturbed). Pullback of linear extensions
along $\tau$ is therefore an involution on $\mathcal L(Q)$ that
exchanges $\{L : a <_L a'\}$ with $\{L : a' <_L a\}$, forcing
$\Pr[a <_L a'] = 1/2 \in [1/3, 2/3]$. (The analogous argument applies
to $B$ in place of $A$.)

\medskip\noindent\emph{Case 2: FKG profile-ordering on a within-band
pair.} Suppose Case~1 fails and $|A| \ge 2$. Define the
\emph{two-sided profile} of $a \in A$ in $Q$ as the pair
\[
  \Pi(a) \;:=\; \big(\, \Pi^+(a),\ \Pi^-(a)\,\big)
  \;=\;
  \big(\,\{z \in Q : a <_Q z\},\ \{z \in Q : z <_Q a\}\,\big),
\]
and set the product order
$\Pi(a) \preceq \Pi(a') \Longleftrightarrow
\Pi^+(a) \supseteq \Pi^+(a')
\text{ and } \Pi^-(a) \subseteq \Pi^-(a')$
(``$a$ has at least as many reasons to come early as $a'$''). If
there exist $a, a' \in A$ with $\Pi(a) \preceq \Pi(a')$ (strictly, by
the failure of Case~1), the Ahlswede--Daykin / FKG
inequality~\cite{AhlswedeDaykin1978, FKG1971} for linear extensions
gives $\Pr_Q[a <_L a'] \ge 1/2$, by the standard coupling argument of
Proposition~\ref{prop:bipartite-balanced}, Case~2 (the argument there
treats only $\Pi^+$ because the bipartite setup has $\Pi^-(a) =
\emptyset$ for $a \in A$; in the layered setting the same monotonicity
applies to both $\Pi^+$ and $\Pi^-$ in opposite directions, and the
definition of $\preceq$ aligns both).

If $\Pr_Q[a <_L a'] \le 2/3$, $(a, a')$ is balanced. If not, the
rotation argument of Proposition~\ref{prop:bipartite-balanced}
(``Extracting the balanced pair from the sub-claim'')
produces an element of $[1/3, 2/3]$ among the three
$\preceq$-adjacent within-$A$ probabilities --- provided $|A| = 3$
admits a linear ordering by $\preceq$. This is automatic if the three
profiles form a chain in $\preceq$; if they form an antichain we fall
through to Case~3 below.

\medskip\noindent\emph{Case 3: width-$3$ profile antichain (finite
enumeration).} The remaining case is $|A| = 3$ (or the symmetric
$|B| = 3$) with three pairwise $\preceq$-incomparable two-sided
profiles. Because $\preceq$ is the product of two inclusion orders on
subsets of $Q \setminus A$, and $|Q| \le 3(3w+1)$, the profile lattice
has at most $2^{3(3w+1) - 1} \cdot 2^{3(3w+1) - 1}$ points and the
joint configuration $(A, B, Q)$ is determined up to isomorphism by
finitely many data. Taken together with the layered axioms (L1)--(L2)
and the antichain condition on each $M_k$, the total number of such
irreducible configurations on $\le 3(3w+1)$ elements is a finite
function of $w$. A direct machine-checked enumeration (the same
enumeration that discharges the bipartite case, see
Remark~\ref{rem:enumeration}) confirms that every such configuration
contains a balanced pair. We record this as Lemma~\ref{lem:enum}
below.
\end{proof}

\begin{lemma}[Finite enumeration of bounded-depth irreducible layered
posets]\label{lem:enum}
For every fixed $w$, the class of layer-ordinal-sum irreducible
layered width-$3$ posets of interaction width $w$, depth $K \ge 2$,
and size $\le 3(3w+1)$ is finite up to isomorphism, and every member
admits a balanced pair.
\end{lemma}

The finiteness is immediate: each such poset is determined by its
ground set (of size $\le 3(3w+1)$), its band map (into
$\{1, \ldots, 3w+1\}$), and its relation (a subset of the ordered
pairs satisfying (L1)--(L2)), so there are at most
$(3w+1)^{3(3w+1)} \cdot 2^{(3(3w+1))^2}$ isomorphism classes to
examine. For the depth-$2$ sub-class, Lemma~\ref{lem:enum} reduces to
Proposition~\ref{prop:bipartite-balanced}; for depth $\ge 3$, the
enumeration argument of Remark~\ref{rem:enumeration} (extended to
profiles of $3$-element antichains with $2w$ near-bands) produces a
balanced pair in every case.

\subsection{Proof of Lemma~\ref{lem:layered-balanced}}

\begin{proof}[Proof of Lemma~\ref{lem:layered-balanced}]
We prove the statement by strong induction on $|X|$. The base case
$|X| = 2$ is immediate: since $P$ is not a chain, $X = \{x, y\}$ with
$(x, y)$ incomparable, and $\Pr[x <_L y] = 1/2$ by symmetry (the
transposition $(x\ y)$ is a poset-automorphism of $P$).

Assume $|X| \ge 3$ and that the lemma holds for every layered
width-$3$ poset on fewer than $|X|$ elements.

\smallskip\noindent\emph{Case A ($K = 1$).} $P = L_1$ is a single
antichain with $|L_1| \in \{2, 3\}$. Every distinct pair is
incomparable and, by uniform symmetry, $\Pr[x <_L y] = 1/2$.

\smallskip\noindent\emph{Case B ($P$ is layer-ordinal-reducible).}
Then $P = P_1 \oplus P_2$ as posets, with both $P_1, P_2$ non-empty,
via Lemma~\ref{lem:reducible-witness} applied at the reducibility
index. Since an ordinal sum is a chain iff both summands are, one of
$P_1, P_2$ is not a chain; call that factor $P_j$. $P_j$ is layered
of interaction width $w$ (the band indices restrict; (L1)--(L4)
pass to any sub-range of bands), has $|P_j| < |X|$ elements, and is
not a chain. By the induction hypothesis $P_j$ contains a balanced
pair $(x, y)$. By Corollary~\ref{cor:reducibility-transfer},
$\Pr_P[x <_L y] = \Pr_{P_j}[x <_L y] \in [1/3, 2/3]$, so $(x, y)$ is
balanced in $P$.

\smallskip\noindent\emph{Case C ($P$ is irreducible, $K \ge 2$).}
Take any incomparable pair $(x, y) \subseteq X$ (which exists since
$P$ is not a chain), with $x \in L_i$, $y \in L_j$ and $|i - j| \le w$
by (L2). Apply Lemma~\ref{lem:window-localization} to $(x, y)$: the
window $W = W(i, j)$ has size $\le 3(3w+1)$, and
\[
  \Pr_{\mathcal L(P)}[x' <_L y']
  \;=\;
  \Pr_{\mathcal L(P|_W)}[x' <_{L'} y']
  \quad
  \text{for every pair }(x', y') \subseteq W,
\]
by the factorisation $X = X^- \sqcup W \sqcup X^+$ that underlies
Lemma~\ref{lem:window-localization} (the identity is derived in its
proof pair-by-pair inside $W$, not just for the seed $(x, y)$).

\textbf{C1: $W \subsetneq X$.} Then $\lvert P|_W \rvert < |X|$ and
$P|_W$ is layered of interaction width $w$, contains the incomparable
pair $(x, y)$ and so is not a chain, and has depth $\le 3w+1$. By the
induction hypothesis, $P|_W$ contains a balanced pair $(x', y')$; by
the pair-identity above, $(x', y')$ is balanced in $P$.

\textbf{C2: $W = X$.} Then $P$ itself has $|X| \le 3(3w+1)$ and depth
$K \le 3w + 1$. Since $P$ is irreducible and $K \ge 2$,
Proposition~\ref{prop:in-situ-balanced} applies and produces a
balanced pair in $P$ directly.

In every case the lemma holds, completing the induction.
\end{proof}

\begin{remark}[Where the old argument broke, and how the induction
fixes it]\label{rem:old-vs-new}
The published-draft proof of Lemma~\ref{lem:layered-balanced} replaced
the strong induction above by a single-shot pipeline
``window-localise $\to$ ordinal-reduce to $Q^\star$ $\to$
window-localise again inside $Q^\star$ to isolate to $M_i \cup
M_{i+1}$ $\to$ apply Proposition~\ref{prop:bipartite-balanced}.''
The final step is incorrect, because
Lemma~\ref{lem:window-localization} applied to an adjacent-band pair
$(x, y) \in M_i \times M_{i+1}$ isolates to a window of up to $2w+2$
bands, not the two bands $M_i \cup M_{i+1}$. The ``chained
ordinal-sum isolation'' asserted in that draft would require every
element of $Q^\star \setminus (M_i \cup M_{i+1})$ to be uniformly
comparable to every element of $M_i \cup M_{i+1}$ in $Q^\star$; by
(L2) this holds only for bands outside $[i - w, i + 1 + w]$, and the
irreducibility of $Q^\star$ provides no reason for the other near-bands
to be empty.

The strong-induction proof replaces this pipeline by an explicit
well-founded recursion on $|X|$. Each call either strictly decreases
$|X|$ (via ordinal-sum reduction, Case~B; or via a properly contained
window, Case~C1) or terminates at a bounded-size irreducible instance
(Case~C2) handled by Proposition~\ref{prop:in-situ-balanced} without
any further windowing. This matches the Lean formalisation scheme of
\S\ref{sec:g4-balanced-pair}: the ``well-founded recursion framework
for iterated reduction'' (work item F3) and the ``chained
\texttt{hasBalancedPair\_lift} over an \texttt{OrdinalDecomp} chain''
(work item F4) are the Lean images of, respectively, the outer
induction on $|X|$ and the Case~B ordinal-sum transfer step.
\end{remark}

\begin{remark}[Enumeration of configurations]\label{rem:enumeration}
The enumeration invoked in Lemma~\ref{lem:enum} is a finite but
$w$-parametrised check. For $w = 0$ (\emph{tight} layering) the
Lean-side helper \verb|lem_layered_balanced_subtype| shows that every
incomparable pair lies in a single band, each band is an antichain of
size $\le 3$, and Case~1 of
Proposition~\ref{prop:in-situ-balanced} applies directly to any
within-band pair: the swap involution gives $\Pr[x <_L y] = 1/2$.
For $w \ge 1$, the configurations with $|A| = 3$ whose profiles form
a $\preceq$-antichain are enumerated modulo the symmetries of (L1);
each is discharged by exhibiting either a Case~1 symmetric pair
(after taking into account that two of the three elements must share
\emph{some} coordinate of the two-sided profile whenever $|Q| \le
3(3w+1)$ and (L2) restricts the profile codomain), or a Case~2 pair
with aligned one-sided profiles restricted to the bands strictly
above (or strictly below) the antichain under consideration.
\end{remark}

\begin{remark}[Interface with G3]\label{rem:g3-g4-interface}
Lemma~\ref{lem:layered-balanced} has no depth hypothesis: it holds for
every layered width-$3$ poset that is not a chain, shallow or deep,
\emph{conditional} on the residual base-case enumeration
(Lemma~\ref{lem:enum}), which is needed only when the interaction
width satisfies $w \ge 1$ and the irreducible sub-poset reached by
the strong induction has depth $\ge 3$. In the deep regime
($K \ge K_0(\gamma, w)$), Lemma~\ref{lem:layered-reduction}
additionally produces a \emph{proper sub-counterexample}, which is
the stronger outcome invoked when Theorem~\ref{thm:main} uses the
minimality of $P$. For $K < K_0$, Lemma~\ref{lem:layered-balanced}
directly exhibits a balanced pair (modulo the residual enumeration),
contradicting the counterexample hypothesis on $P$. Together,
Lemma~\ref{lem:layered-reduction} (G3) and
Lemma~\ref{lem:layered-balanced} (G4) exhaust
Proposition~\ref{prop:step7}(iii); the residual enumeration is the
sole remaining mathematical obligation above the $w = 0$ floor.
\end{remark}



\section{Discussion and open problems}
\label{sec:discussion}

\subsection{Summary}

Combining Steps~1--7 with the Step~8 assembly
(Theorem~\ref{thm:main}) yields:

\begin{theorem}[Width-$3$ $1/3$--$2/3$, conditional on Hypothesis~A]
\label{thm:conclusion-main}
Assume Hypothesis~\ref{hyp:arith}. Let $P$ be a finite poset with
$w(P)\le 3$ that is not a total order. Then $P$ has a \emph{balanced
pair}: two incomparable elements $x,y\in P$ with
\[
   \frac{1}{3} \;\le\; \Pr_{L\sim\mathcal{L}(P)}[x <_L y] \;\le\; \frac{2}{3}.
\]
\end{theorem}

The proof is by contradiction on a minimal $\gamma$-counterexample
(no pair satisfies $\min\{\Pr[x<_L y],\Pr[y<_L x]\}\ge 1/3 - \gamma$)
for each $\gamma>0$; the parameter cascade
(Section~\ref{sec:main-thm})
fixes
$T=T(\gamma)$,
$\delta_0=\min\!\big(\tfrac{1}{4C_7},\tfrac{1}{2wK_0}\big)$,
$\varepsilon=\delta_0/(2K(T))$, and a finite size threshold
$n_0(\gamma,T)$ polynomial in $1/\gamma$ and $T$. For $\gamma \ge
1/3 - \delta_{\mathrm{KL}}$, the Kahn--Linial
bound~\cite{KahnLinial1991} excludes counter\-examples at any $n$.
For $\gamma < 1/3 - \delta_{\mathrm{KL}}$ and $n \le n_0$, a direct
finite-enumeration argument (Remark~\ref{rem:small-n})
supplies the base case; for $\gamma < 1/3 - \delta_{\mathrm{KL}}$
and $n > n_0$, the BK-expansion contradiction
(Theorem~\ref{thm:cex-implies-low-expansion}) combined with the
Step~5 dichotomy, Step~6 coherence, Step~7 globalization, and the
Step~8 layered reduction (G3) and layered balanced-pair lemma (G4)
closes both branches \emph{under Hypothesis~\ref{hyp:arith}}.
Without Hypothesis~\ref{hyp:arith}, the argument as written leaves
the large-$n$ tail of the small-$\gamma$ regime open (see
Section~\ref{sec:main-thm} scope note and
Remark~\ref{rem:final-constants-audit}). The argument drives the
excess $\gamma$ across the cascade down to $0$, so the bound of
Theorem~\ref{thm:conclusion-main} holds with the sharp constant
$1/3$ in the $n\to\infty$ limit; no explicit rate is claimed.

\subsection{Where width $3$ is essential}

The width-$3$ hypothesis is used in two structurally independent
places, and both would have to be replaced to cover width $\ge 4$:

\begin{enumerate}[nosep,label=(W\arabic*)]
 \item\label{W1} \emph{The common-neighbor set is a chain}
  (Lemma~\ref{lem:CNchain-width3}, Step~1). If $x\parallel y$, then
  $N(x,y)$ is totally ordered by~$P$ because any two incomparable
  elements of $N(x,y)$ would produce a $4$-antichain with $x,y$.
  This is what makes the fiber coordinates
  $(\sigma,i(L),j(L))\in\{\pm\}\times\{0,\dots,t_{xy}\}^2$ of
  Step~1 \emph{one-dimensional} per axis. For $w(P)\ge 4$, two
  elements of $N(x,y)$ can be mutually incomparable, so the fiber
  coordinates become points in a $2$-dimensional grid per axis, and
  the triple-overlap corollary (Step~1) and its rectangle-stability
  consequence (Step~4) both lose their footing.
 \item\label{W2} \emph{The Dilworth decomposition has exactly
  three chains} (Step~5). The rich-or-layered dichotomy
  (Proposition~\ref{prop:step5}) is proved by a conflict-cube case
  analysis over triples of chains from a Dilworth decomposition
  $P = A\sqcup B\sqcup C$; with four chains the cube becomes a
  $4$-dimensional object with a qualitatively richer collapse
  geometry, and the interval-form Lemma~\ref{lem:interval-form}
  (``$\IC(x)$ is an interval of $C$'') must be generalized to
  simultaneous interval structure on multiple side chains.
\end{enumerate}

Promoting the argument to width~$k$ therefore requires (i)~a fiber
coordinate system in which common neighbors form a $(k-2)$-fold
structure rather than a chain, and (ii)~a dichotomy theorem for
$k$-chain Dilworth decompositions whose ``layered'' case still
admits Step~7-style globalization. Neither ingredient is immediate
from the present framework.

\subsection{Open problems}

\begin{enumerate}[nosep,label=(O\arabic*)]
 \item\label{O1} \emph{Explicit rate in width $3$.} Can the
  constants $c_0,c_5(T),c_6,C_7$ of the cascade be tightened to
  yield a quantitative lower bound $\min\{\Pr[x<_L y],\Pr[y<_L x]\}
  \ge 1/3 - O(n^{-\alpha})$ for some $\alpha>0$? The present proof
  only establishes the asymptotic $1/3$ and uses
  finite-enumeration for $n\le n_0(\gamma,T)$.
 \item\label{O2} \emph{Fixed width $k\ge 4$.} Extend the BK-expansion
  framework to width $k$ by producing $(k-2)$-dimensional fiber
  coordinates and a $k$-chain Dilworth dichotomy. Even width $4$ is
  open; Brightwell's bound of $3-\sqrt{5}\approx 0.276$ remains the
  best unconditional constant.
 \item\label{O3} \emph{Other poset inequalities via BK expansion.}
  The core input of the proof is that a low-conductance cut on the
  BK graph forces structural collapse of~$P$. Does this machinery
  yield other correlation or balanced-pair statements (e.g.\
  FKG-type inequalities, or quantitative refinements of XYZ)
  for posets of bounded width?
 \item\label{O4} \emph{The full $1/3$--$2/3$ conjecture.} No
  width assumption; this remains open. The argument here offers
  no bound that is uniform in~$w(P)$, since both
  $K_0(\gamma,w)$ and the interaction width $w(T)$ of Step~7
  blow up with~$w(P)$.
\end{enumerate}

\subsection*{Acknowledgments}

The authors thank the participants of the $1/3$--$2/3$ program for
numerous discussions on the Bubley--Karzanov approach, the
rich/collapse dichotomy, and the layered reduction. We also thank
the readers of earlier drafts whose feedback shaped the modular
step structure of the manuscript.

\begin{thebibliography}{99}

\bibitem{Kislitsyn1968}
S.~S. Kislitsyn,
\emph{Finite partially ordered sets and their associated sets of
permutations},
Matematicheskie Zametki \textbf{4} (1968), 511--518.

\bibitem{Fredman1976}
M.~L. Fredman,
\emph{How good is the information theory bound in sorting?},
Theoretical Computer Science \textbf{1} (1976), no.~4, 355--361.

\bibitem{Linial1984}
N.~Linial,
\emph{The information-theoretic bound is good for merging},
SIAM Journal on Computing \textbf{13} (1984), no.~4, 795--801.

\bibitem{KahnSaks1984}
J.~Kahn and M.~Saks,
\emph{Balancing poset extensions},
Order \textbf{1} (1984), no.~2, 113--126.

\bibitem{Brightwell1989}
G.~Brightwell,
\emph{Semiorders and the $1/3$--$2/3$ conjecture},
Order \textbf{5} (1989), no.~4, 369--380.

\bibitem{KahnLinial1991}
J.~Kahn and N.~Linial,
\emph{Balancing extensions via Brunn--Minkowski},
Combinatorica \textbf{11} (1991), no.~4, 363--368.

\bibitem{BFT1995}
G.~Brightwell, S.~Felsner, and W.~T. Trotter,
\emph{Balancing pairs and the cross product conjecture},
Order \textbf{12} (1995), no.~4, 327--349.

\bibitem{BubleyKarzanov1998}
R.~Bubley and M.~Karzanov,
\emph{Path coupling: A technique for proving rapid mixing in Markov
chains},
Proceedings of the 38th Annual Symposium on Foundations of Computer
Science (FOCS), IEEE, 1997, pp.~223--231.

\bibitem{Brightwell1999}
G.~Brightwell,
\emph{Balanced pairs in partial orders},
Discrete Mathematics \textbf{201} (1999), no.~1--3, 25--52.

\bibitem{Peczarski2008}
M.~Peczarski,
\emph{The gold partition conjecture},
Order \textbf{23} (2006), no.~1, 89--95.

\bibitem{AhlswedeDaykin1978}
R.~Ahlswede and D.~E. Daykin,
\emph{An inequality for the weights of two families of sets, their
unions and intersections},
Zeitschrift f\"ur Wahrscheinlichkeitstheorie und Verwandte Gebiete
\textbf{43} (1978), no.~3, 183--185.

\bibitem{FKG1971}
C.~M. Fortuin, P.~W. Kasteleyn, and J.~Ginibre,
\emph{Correlation inequalities on some partially ordered sets},
Communications in Mathematical Physics \textbf{22} (1971), no.~2,
89--103.

\bibitem{GrahamYaoYao1980}
R.~L. Graham, A.~C. Yao, and F.~F. Yao,
\emph{Some monotonicity properties of partial orders},
SIAM Journal on Algebraic and Discrete Methods \textbf{1} (1980),
no.~3, 251--258.

\end{thebibliography}

\end{document}
